\part{DFT, FPGA, HLS and analog}\label{tm:part:dft-fpga-hls-and-analog}

\chapter{DFT Architecture: scan + MBIST + JTAG for the Sky130 SoC (Stage
0)}\label{tm:ch:dft-architecture}

This chapter describes the DFT architecture. It covers scan, JTAG and memory self-test.

\begin{relocated}{docs/design/DFT_ARCHITECTURE.md}
% Converted from docs/design/DFT_ARCHITECTURE.md by tools/docs/md2tex.py; edit this file, not the original.

Bead \code{j41m.1} (epic \code{j41m}), \ghref{244}. Scope: Sky130 only
(\src{pnr/sky130/soc/}, \src{pnr/sky130/cpu/}). No RTL, flow, macro or
LibreLane-install change was made. Experiments live in
\code{/nobackup/claude_sim_build/dft_stage0/}; the reproducible scripts
are in \src{tools/dft/} (see its README).

\textbf{Evidence labels.} Every statement carries one: \textbf{{[}M{]}}
measured (command and output named), \textbf{{[}R{]}} read from source
(\code{file:line}), \textbf{{[}A{]}} assumed or estimated (reasoning
given). Where a source said something that a measurement contradicted,
the measurement is reported and the conflict is called out.

\section{Decisions already taken (2026-10-10)}\label{tm:ch:dft-architecture:0-decisions-already-taken-2026-10-10}\origsec{0}

These are \textbf{DECIDED} by the project owner and are not open
questions. The rest of the document is written against them.

{\footnotesize\setlength{\tabcolsep}{3pt}

{
\begin{longtable}[]{@{}
  >{\raggedright\arraybackslash}p{\dimexpr(\linewidth-4\tabcolsep)*49/1000\relax}
  >{\raggedright\arraybackslash}p{\dimexpr(\linewidth-4\tabcolsep)*466/1000\relax}
  >{\raggedright\arraybackslash}p{\dimexpr(\linewidth-4\tabcolsep)*466/1000\relax}@{}}
\toprule\noalign{}
\begin{minipage}[b]{\linewidth}\raggedright
\#
\end{minipage} & \begin{minipage}[b]{\linewidth}\raggedright
Decision
\end{minipage} & \begin{minipage}[b]{\linewidth}\raggedright
Consequence in this document
\end{minipage} \\
\midrule\noalign{}
\endhead
\bottomrule\noalign{}
\endlastfoot
1 & \textbf{Goal is real tape-out intent on Sky130.} DFT must be
production-usable. & A stuck-at coverage target is stated and justified
(section 2.1). Patterns must reach a tester (section 2.3). Section 2.4
lists what the open-source tools cannot deliver. \\
2 & ATPG tools may be added to \textbf{this repo\textquotesingle s
\src{flake.nix} devshell only}. Non-commercial-licensed tools (Atalanta)
are acceptable. & Licence text quoted in section 2.2. A tape-out with
commercial intent changes this and is flagged there. \\
3 & \textbf{Crypto key registers are EXCLUDED from scan.} & Quantified
and enforced in section 7, with a check script that fails on violation.
Exclusion alone leaks the key; masking is required (section 7). \\
4 & \textbf{MBIST is tester-only.} No APB register, no power-on
auto-run, no firmware visibility. & The APB-slot analysis was dropped.
MBIST is started and read only through the TAP (decision 5). \\
5 & \textbf{Test access is an IEEE 1149.1 JTAG TAP} controlling both
scan and MBIST. & The TAP is a Stage 1 deliverable, not an optional last
stage (sections 9, 10). \\
6 & \textbf{Pad ring / chip top is NOT in scope now.} & The core macro
exposes the TAP and any unavoidable raw ports. Section 9.5 lists what a
future chip-top epic needs. \\
7 & \textbf{Scan access = JTAG TAP for control and MBIST, PLUS parallel
scan ports} (direct scan-in / scan-out / scan-enable / test clock for
several chains). \emph{(2026-10-10, follow-up)} & Resolves the old N3 in
favour of parallel ports. Pin plan, chain proposal, test time and the
at-speed statement are rewritten in sections 8, 9.5, 9.6, 10. Zero
delay-fault coverage is \textbf{not} accepted as a goal; section 2.4
says what would be needed and that it is out of reach of the current
tools. \\
8 & \textbf{99 \% test / 98 \% fault coverage is a GOAL, not a
commitment.} \emph{(2026-10-10, follow-up)} & Stage 3 begins with a
time-boxed ATPG tool spike on the real netlist. The sign-off number is
fixed only after a tool is shown to complete at scale (section 2.1, 10).
\textbf{2026-10-30 (bead \bead{j41m.23}): adopted as the Stage 3 target,
provisionally} (patched Quaigh completed the fabric; sign-off waits on
the cell-pin re-grade). Section 2.1, 18.8. \\
9 & \textbf{Sky130 CPU macro Synlig exposure is bead \bead{dud4} (P1):
test first, then the user decides on re-hardening.} \emph{(2026-10-10,
follow-up)} & The test is the \code{ma7} branch-dependent differential
on the Sky130 CPU macro gate netlist against RTL, \textbf{no PD run}.
Stage 2\textquotesingle s CPU re-harden depends on its outcome (sections
4, 5, 10). \\
\end{longtable}
}

}

\textbf{Schedule decisions (2026-10-10).} Stage 1 is approved to start
now that \prref{250} has merged. FPGA Stage 0 (\ghref{245}) follows Stage 1,
not before it.

\section{1a. Summary}\label{tm:ch:dft-architecture:1a-summary}

{\footnotesize\setlength{\tabcolsep}{3pt}

{
\begin{longtable}[]{@{}
  >{\raggedright\arraybackslash}p{\dimexpr(\linewidth-4\tabcolsep)*158/1000\relax}
  >{\raggedright\arraybackslash}p{\dimexpr(\linewidth-4\tabcolsep)*653/1000\relax}
  >{\raggedright\arraybackslash}p{\dimexpr(\linewidth-4\tabcolsep)*169/1000\relax}@{}}
\toprule\noalign{}
\begin{minipage}[b]{\linewidth}\raggedright
Q
\end{minipage} & \begin{minipage}[b]{\linewidth}\raggedright
Headline
\end{minipage} & \begin{minipage}[b]{\linewidth}\raggedright
Basis
\end{minipage} \\
\midrule\noalign{}
\endhead
\bottomrule\noalign{}
\endlastfoot
1 Scan tool & The Sky130 flow\textquotesingle s OpenROAD \textbf{does}
have a \code{dft} module and stitched a legal chain on a real block
(\code{timer}, 178 flops) and on the whole synthesised SoC fabric
(18,752 flops, 5 s). It honours \code{set_dont_touch} (needed for
decision 3). It has real defects (below). No lock-up latches, no
clock-gate/\allowbreak{}test-enable handling, no macro scan pins.
\textbf{Usable for stitching, with workarounds.} & {[}M{]} \\
2 ATPG & Nothing was installed. Quaigh 0.0.6 (MIT/\allowbreak{}Apache,
Rust, SAT) built and ran: \textbf{92.41 \% raw (4336/\allowbreak{}4692),
358 patterns, 1.5 s} on the full-scan cut of \code{timer} (generic-gate
model, not Sky130 cells). No tool is proven at SoC scale, and
\textbf{none of them delivers production-grade ATPG} (section 2.4). &
{[}M{]} / {[}A{]} \\
3 RTL audit & 18,753 sequential cells in the fabric + 4,829 inside the
CPU macro. One live clock gate (no test enable), one latch, no RTL
negedge flops, no tri-states, no comb loops, no ring oscillator in the
build. \textasciitilde20 Stage 1 changes (section 4). &
{[}M{]}/{[}R{]} \\
4 Hard macros & The CPU is a hard macro with \textbf{no scan or test
pins}; 4,829 flops are invisible to SoC-level insertion.
\textasciitilde6-7 h per CPU+SoC re-harden round. & {[}R{]}/{[}A{]} \\
5 MBIST & 12 SRAM macros (2 x 4 KB in the fabric, 10 x 1 KB inside the
CPU macro), 147,456 bits. March C- (10N): 10,240 cycles fabric / 2,560
cycles CPU in parallel = 256 us at 40 MHz. Two controllers, driven and
read through the TAP. & {[}R{]} + {[}A{]} \\
6 Crypto & DECIDED: exclude. 256 flops (1.37 \% of flops).
\textbf{Unmasked, the exclusion costs \textasciitilde11 \% of the fault
universe and still leaks the key}; masking brings it to
\textasciitilde0.7 \%. & {[}M{]} \\
7 Cost & Scan flops +117,350 um2 fabric (+10.2 \% of synthesised
stdcell, \textasciitilde+0.6 pp utilisation), +30,209 um2 CPU macro.
D-pin setup +0.07 to +0.45 ns by corner. & {[}M{]}
area/\allowbreak{}setup \\
8 Access & DECIDED: JTAG TAP (control + MBIST) \textbf{plus parallel
scan ports} (decision 7): TAP 5-6 pins + \code{scan_clk_i},
\code{scan_en_i}, 8 x \code{scan_in_i}/\code{scan_out_o} = 23-24 raw
pins (proposal, section 8). & decision / {[}A{]} \\
\end{longtable}
}

}

\textbf{Go/no-go:} Stage 1 (TAP + parallel ports + scan-ready RTL) GO
and approved to start (\prref{250} merged). Stage 2 CONDITIONAL GO. Stage 3
CONDITIONAL: GO for a time-boxed tool spike, NO-GO for a promised
coverage number until a tool is shown to complete at scale (decision 8).
Stage 2\textquotesingle s CPU re-harden additionally waits on
\code{dud4} (decision 9). Stage 4 (MBIST) GO. Details in section 11.

\section{Scan insertion tool}\label{tm:ch:dft-architecture:1-scan-insertion-tool}\origsec{1}

\subsection{Which OpenROAD}\label{tm:ch:dft-architecture:11-which-openroad}\origsec{1.1}

The task brief expected \code{openroad-2026-02-17}. \textbf{{[}M{]}} The
Sky130 flow does not use it. \code{nix-shell} in the LibreLane checkout
(\code{shell_probe.log}) puts \code{openroad} at
\code{/nix/store/784j83yn...-openroad-python3-3.11.9-env/bin/openroad},
reporting revision \code{edf00dff99f6c40d67a30c0e22a8191c5d2ed9d6}. This
matches \src{pnr/Makefile:112-136} {[}R{]}, which says Sky130 and
FreePDK45 are deliberately \textbf{not} overridden and only ASAP7 uses
26Q2. Three builds are on this host:

{\small\setlength{\tabcolsep}{4pt}

{
\begin{longtable}[]{@{}
  >{\raggedright\arraybackslash}p{\dimexpr(\linewidth-4\tabcolsep)*298/1000\relax}
  >{\raggedright\arraybackslash}p{\dimexpr(\linewidth-4\tabcolsep)*140/1000\relax}
  >{\raggedright\arraybackslash}p{\dimexpr(\linewidth-4\tabcolsep)*542/1000\relax}@{}}
\toprule\noalign{}
\begin{minipage}[b]{\linewidth}\raggedright
Build
\end{minipage} & \begin{minipage}[b]{\linewidth}\raggedright
Revision
\end{minipage} & \begin{minipage}[b]{\linewidth}\raggedright
DFT commands {[}M: \src{tools/dft/help_probe.tcl}{]}
\end{minipage} \\
\midrule\noalign{}
\endhead
\bottomrule\noalign{}
\endlastfoot
LibreLane devshell (\textbf{Sky130 flow}) & \code{edf00dff...} &
\code{set_dft_config}, \code{report_dft_config}, \code{preview_dft},
\code{insert_dft}, \code{scan_replace} \\
\code{openroad-2026-02-17} & \code{dcf36133...} & \code{set_dft_config},
\code{report_dft_config}, \code{report_dft_plan},
\code{execute_dft_plan}, \code{scan_replace}, \code{scan_opt} \\
\code{openroad-26Q2} & \code{26Q2} & same as 2026-02-17 \\
\end{longtable}
}

}

\code{set_dft_config} options: \code{edf00dff}:
\texttt{-\allowbreak{}max\_\allowbreak{}length\ -\allowbreak{}max\_\allowbreak{}chains\ -\allowbreak{}clock\_\allowbreak{}mixing}.
26Q2 adds
\texttt{-\allowbreak{}scan\_\allowbreak{}enable\_\allowbreak{}name\_\allowbreak{}pattern\ -\allowbreak{}scan\_\allowbreak{}in\_\allowbreak{}name\_\allowbreak{}pattern\ -\allowbreak{}scan\_\allowbreak{}out\_\allowbreak{}name\_\allowbreak{}pattern}.
Valid \code{-clock_mixing} values are \code{no_mix} and
\code{clock_mix}; \code{mix} raises DFT-0006 and \textbf{aborts the
process with a stack trace} {[}M: \code{toy_edf.log}{]}.

\subsection{Real-block experiment {[}M{]}}\label{tm:ch:dft-architecture:12-real-block-experiment-m}\origsec{1.2}

\code{timer} (+ \code{apb4_register_bank}) -\textgreater{} sv2v
-\textgreater{} Yosys 0.46
\texttt{synth\ -\allowbreak{}flatten;\ dfflibmap;\ abc} against
\code{sky130_fd_sc_hd__tt_025C_1v80.lib} (\src{tools/dft/synth.ys})
-\textgreater{} 891 cells, 178 \code{dfxtp_1}. Then
\texttt{scan\_\allowbreak{}replace;\ preview\_\allowbreak{}dft;\ insert\_\allowbreak{}dft}
(\src{tools/dft/dft_timer.tcl}):

\begin{itemize}
\tightlist
\item
  Result: 178 \code{dfxtp_1} became 178 \code{sdfxtp_1}, one chain of
  178, new ports \code{scan_enable_1}, \code{scan_in_1},
  \code{scan_out_1}.
\item
  \src{tools/dft/chain_check.py} walked the netlist: chain length 178, 0
  unreached flops, 0 non-scan flops left, one SCE net, one clock net.
\item
  The scan netlist re-reads and links in OpenROAD (891 cells). It also
  reads and times in OpenSTA.
\item
  \textbf{Not done:} functional-mode equivalence and a
  shift/\allowbreak{}flush simulation. The chain was checked
  structurally only.
\item
  Whole fabric,
  \code{RUN_2026-10-06_22-43-34/04-yosys-synthesis/soc_top.nl.v} (91,424
  cells) with the real CPU and SRAM LEFs/libs
  (\src{tools/dft/soc_dft.tcl}): \code{scan_replace} plus
  \code{insert_dft} took \textbf{5 s}. 17,903 \code{dfxtp_2}
  -\textgreater{} \code{sdfxtp_2}, 843 \code{dfrtp_2} -\textgreater{}
  \code{sdfrtp_2}, 6 \code{dfstp_2} -\textgreater{} \code{sdfsbp_2}, the
  1 latch (\code{dlxtn_1}) left alone. With
  \texttt{-\allowbreak{}max\_\allowbreak{}chains\ 4\ -\allowbreak{}clock\_\allowbreak{}mixing\ no\_\allowbreak{}mix}
  it made \textbf{8 chains}: four of 235-236 flops on \code{cpu_clk_i}
  (944 flops) and four of 4,450-4,453 on \code{clk_i} (17,809). So
  \code{-max_chains} is per clock domain, and the CPU-domain flops are
  only the 173 + 771 outside the macro.
\end{itemize}

\subsection{What the module cannot do, and the defects found}\label{tm:ch:dft-architecture:13-what-the-module-cannot-do-and-the-defects-found}\origsec{1.3}

{\footnotesize\setlength{\tabcolsep}{3pt}

{
\begin{longtable}[]{@{}
  >{\raggedright\arraybackslash}p{\dimexpr(\linewidth-2\tabcolsep)*312/1000\relax}
  >{\raggedright\arraybackslash}p{\dimexpr(\linewidth-2\tabcolsep)*668/1000\relax}@{}}
\toprule\noalign{}
\begin{minipage}[b]{\linewidth}\raggedright
Capability
\end{minipage} & \begin{minipage}[b]{\linewidth}\raggedright
Result
\end{minipage} \\
\midrule\noalign{}
\endhead
\bottomrule\noalign{}
\endlastfoot
Multiple clock domains & Groups by (clock, edge) in \code{no_mix}. In
\code{clock_mix} it stitches rising/\allowbreak{}falling and different
clocks into \textbf{one chain with no lock-up latch} (toy: \code{fn0}
falling clka -\textgreater{} \code{fb1} rising clkb -\textgreater{}
\code{fg0} -\textgreater{} ...). {[}M \code{toy_edf.log}{]} Use
\code{no_mix} only. \\
Lock-up latches & None inserted (no \code{dlx*} cell added in the toy;
the only latch in the output is the one that was in the input).
{[}M{]} \\
Clock gates & A gated flop\textquotesingle s clock is traced through the
gate to the root clock (\code{fg0} joined clka\textquotesingle s chain).
The gate itself is \textbf{not} scan-replaced (\code{DFT-0002}) and no
test-enable is driven. A library \code{dlclkp} has an \code{sdlclkp}
equivalent in the Liberty but the tool does not use it. {[}M{]} \\
Latches & Skipped with a warning, left unscanned (\code{DFT-0002/0007}).
{[}M{]} \\
Negedge flops & Own chain in \code{no_mix}. \textbf{Defect, edf00dff
only:} \code{dfrtn_1} was replaced by \code{sdfbbn_1} (set+reset flop)
and its \code{SET_B} pin is \textbf{left unconnected}; 26Q2 picks
\code{sdfrtn_1} correctly. {[}M \code{toy_scan_edf.v},
\code{toy_scan_26q2.v}{]} \\
\code{report_dft_config} & Fails with a Tcl error
(\texttt{can't\ read\ "args"}) on edf00dff. {[}M{]} \\
Chain count/\allowbreak{}balancing & Chain count follows the
clock-domain count; \code{-max_chains} below the domain count is not
honoured (toy: asked 2, got 3). The 8-chain fabric result is 236 vs
4,453, unbalanced. \code{-max_length} was not shown to split anything.
{[}M{]} \\
Scan-enable & One \code{scan_enable_N} port, \textbf{one net with 18,752
sinks, unbuffered}. No buffer tree. {[}M{]} \\
Hard-macro scan pins & Macro instances are ignored (not scan cells). A
chain through a macro cannot be built by the tool. {[}M: CPU/SRAM macros
untouched{]} \\
Netlist legality & \code{write_verilog} emits
\texttt{assign\ scan\_\allowbreak{}out\_\allowbreak{}N\ =\allowbreak{}\ <internal\ net>}
for each scan-out (8 assigns in the SoC; the input netlist had 0).
\code{Checker.NetlistAssignStatements} runs at flow step 06 and would
object if the insertion ran before it. {[}M{]} \\
DEF \code{SCANCHAINS} & \textbf{Not tested.} {[}A{]} Assume absent: no
command in either build writes it, and the module has no scan-def
API. \\
\code{set_dont_touch} & \textbf{Honoured} by \code{scan_replace} and the
stitcher (256 marked flops stayed \code{dfxtp_2} and off the chains).
Needed for decision 3. {[}M \code{soc_excl.log}{]} \\
Reset/set handling in shift & Not handled. Async resets stay functional,
so any internally generated reset can corrupt a chain in shift (section
4). \\
\end{longtable}
}

}

\subsection{Alternatives (not exercised)}\label{tm:ch:dft-architecture:14-alternatives-not-exercised}\origsec{1.4}

\begin{itemize}
\tightlist
\item
  \textbf{Yosys-level replacement}: Yosys \code{dfflibmap} warned
  \texttt{Found\ unsupported\ expression\ 'D\&!SCE|SCD\&SCE'} on every
  Sky130 \code{sdf*} cell {[}M{]}, so it cannot map to scan flops from
  Liberty; you would hand-write a techmap. More work than the OpenROAD
  route for no gain.
\item
  \textbf{Fault\textquotesingle s chain stitcher}
  (\texttt{fault\ chain}): Swift, not built {[}A{]}. It also produces
  boundary-scan registers. Held in reserve.
\item
  \textbf{Recommendation:} OpenROAD \code{dft} with \code{no_mix}, run
  from the \textbf{26Q2 binary on the Verilog netlist} (bug-free
  \code{sdfrtn_1} choice, name patterns), with the chain count set
  explicitly. Netlist-level avoids reading an ODB written by a different
  OpenROAD build {[}A: ODB schema compatibility across builds is not
  guaranteed; not tested{]}.
\end{itemize}

\subsection{Slotting into LibreLane as a project-local plugin}\label{tm:ch:dft-architecture:15-slotting-into-librelane-as-a-project-local-plugin}\origsec{1.5}

The existing plugin
\src{pnr/sky130/soc/plugin/librelane_plugin_cvt_sky130/__init__.py}
registers \code{Step} subclasses and activates them through
\code{meta.substituting_steps} {[}R{]}. Classic\textquotesingle s step
list is \code{librelane/flows/classic.py} (\code{Yosys.Synthesis},
\code{Checker.YosysUnmappedCells}, ...,
\code{Checker.NetlistAssignStatements}, \code{OpenROAD.CheckSDCFiles},
..., \code{OpenROAD.Floorplan} ...) {[}R{]}; the Sky130
run\textquotesingle s own order (steps 04-21) is \code{Yosys.Synthesis}
04, \code{Checker.*} 05-06, \code{OpenROAD.STAPrePNR} 09,
\code{OpenROAD.Floorplan} 10, tap/PDN 14-18, global placement skip-IO
20, \code{OpenROAD.IOPlacement} 21 {[}M: step directory names in
\code{RUN_2026-10-09_06-18-39}{]}. \code{substituting_steps} only
replaces a step, so adding one needs a plugin-registered \code{Flow}
subclass of \code{Classic} with \code{Steps} edited {[}A: not tried; the
\code{SequentialFlow} class exposes \code{Steps} as a plain list{]}.

{\footnotesize\setlength{\tabcolsep}{3pt}

{
\begin{longtable}[]{@{}
  >{\raggedright\arraybackslash}p{\dimexpr(\linewidth-6\tabcolsep)*139/1000\relax}
  >{\raggedright\arraybackslash}p{\dimexpr(\linewidth-6\tabcolsep)*288/1000\relax}
  >{\raggedright\arraybackslash}p{\dimexpr(\linewidth-6\tabcolsep)*288/1000\relax}
  >{\raggedright\arraybackslash}p{\dimexpr(\linewidth-6\tabcolsep)*265/1000\relax}@{}}
\toprule\noalign{}
\begin{minipage}[b]{\linewidth}\raggedright
Option
\end{minipage} & \begin{minipage}[b]{\linewidth}\raggedright
Where
\end{minipage} & \begin{minipage}[b]{\linewidth}\raggedright
For
\end{minipage} & \begin{minipage}[b]{\linewidth}\raggedright
Against
\end{minipage} \\
\midrule\noalign{}
\endhead
\bottomrule\noalign{}
\endlastfoot
\textbf{N (recommended)} netlist step & after
\code{Checker.NetlistAssignStatements} (06), before \code{STAPrePNR} &
any OpenROAD build; ports exist before IO placement; no ODB version
coupling & must remove or accept the \code{assign} lines \\
O ODB step & after \code{OpenROAD.Floorplan} (10), before IO placement
(21) & native ODB, no netlist re-read & pinned \code{edf00dff} only
(SET\_B defect, no name patterns) \\
\end{longtable}
}

}

What breaks downstream {[}A, except where noted{]}:

\begin{itemize}
\tightlist
\item
  \textbf{SDC}: the single functional SDC has no
  \code{scan_enable}/\code{scan_mode} constraint. Needs
  \texttt{set\_\allowbreak{}case\_\allowbreak{}analysis\ 0} for
  functional timing (or a false path), plus a separate shift-mode check.
  \code{cpu_clk_i} has \textbf{no \code{create_clock}} in
  \code{sky130_soc.sdc} today {[}R: audit; the file defines only
  \code{core_clk} at \code{:237}{]}, so CPU-domain flops are
  unconstrained in the current SoC signoff. This must be fixed before
  scan timing means anything.
\item
  \textbf{Hold}: every chain link is Q -\textgreater{} SCD with no
  logic. Sky130 \code{GRT_RESIZER_HOLD_SLACK_MARGIN} is 0.3 ns {[}R: PD
  history{]}, so expect a hold-buffer count proportional to chain links
  (18.7 k). Not measured; this is the largest unknown in Stage 2.
\item
  \textbf{CTS}: unchanged sink count; the gated-clock structure is the
  same. The scan-enable net needs a buffer tree
  (resizer/\code{repair_design}).
\item
  \textbf{Antenna and \code{e45j}}: one new net with \textasciitilde18.8
  k pins plus 8-12 long scan-out/\allowbreak{}in nets. The repair-design
  steps (\code{CVT.RepairDesignPostGRT}) and the diode-on-port step will
  see new nets; slew/\allowbreak{}cap/\allowbreak{}antenna counts must
  be re-compared against the Gate B and \code{rccal} baselines.
\item
  \textbf{Cell policy}: the flow\textquotesingle s resizer must be
  allowed to keep \code{sdf*} cells. Whether the PDK
  \code{DONT_USE_CELLS} list contains them was \textbf{not verified} (I
  could not locate the list in \code{resolved.json}).
\end{itemize}

\section{ATPG tool, coverage target, pattern export}\label{tm:ch:dft-architecture:2-atpg-tool-coverage-target-pattern-export}\origsec{2}

\subsection{Coverage goal (decision 8: a GOAL, not a commitment)
{[}A{]}}\label{tm:ch:dft-architecture:21-coverage-goal-decision-8-a-goal-not-a-commitment-a}\origsec{2.1}

\textbf{Decision 8 (2026-10-10):} the numbers below are a goal. Nothing
is promised until a tool has been shown to complete on the real netlist.
Stage 3 therefore starts with a \textbf{time-boxed ATPG tool spike}, and
the sign-off target is fixed only afterwards (section 10, Stage 3). If
the spike fails, the per-block fallback of section 2.4 applies and the
goal is restated per block.

\textbf{Goal: stuck-at test coverage \textgreater= 99.0 \% of the
scanned-domain fault list, and stuck-at fault coverage \textgreater=
98.0 \% of the whole-SoC fault list, with every exclusion itemised.}
Test coverage = detected / (total - proven untestable); fault coverage =
detected / total.

\textbf{Adopted 2026-10-30 (owner decision, bead \bead{j41m.23}),
PROVISIONAL.} After the patched Quaigh completed the fabric rung (18.9),
the 99.0 \% / 98.0 \% numbers above are the Stage 3 \textbf{target}, not
only a goal. Status: \textbf{met provisionally} on this
flow\textquotesingle s universe (generic-gate faults, collapsed list, no
X): fabric test coverage 99.95 \%, fault coverage 99.16 \% on the
excluded list, about 98.26 \% counting every excluded fault as
undetected (18.9). Scope: the "whole-SoC fault list" is the
\textbf{fabric} list. The CPU macro (5,299 scan flops, a black box in
the fabric cut) gets \textbf{its own rung with the same 99.0 / 98.0
target}, reported next to the fabric and never merged into one number.
The target counts as met for sign-off only after (1) the cell-pin
re-grade (bead \bead{j41m.20}) still shows \textgreater= 99.0 \% /
\textgreater= 98.0 \%, (2) the CPU-macro rung meets it, and (3) the
coverage around the black boxes (CPU macro, SRAMs), an upper bound here
because the flow has no X, is restated with X-aware handling or
shadowing (\code{j41m.21}). The per-block fallback of 2.4 / 18.8 item 2
is \textbf{dropped}.

Why 99.0 \%: the DFT orchestrator\textquotesingle s own sign-off
criterion in this repo is
\texttt{saf\_\allowbreak{}coverage\_\allowbreak{}pct\ >=\allowbreak{}\ 99.\allowbreak{}0}
{[}R: dft-orchestrator contract{]}, and the standard Williams-Brown
estimate shows what lower numbers cost. Defect level DL = 1 - Y\^{}(1-T)
for yield Y and coverage T (arithmetic computed here, but \textbf{Y is
assumed}; no Sky130 yield data for this die exists in the repo):

{
\begin{longtable}[]{@{}lll@{}}
\toprule\noalign{}
T (stuck-at) & DL at Y=0.8 (ppm) & DL at Y=0.5 (ppm) \\
\midrule\noalign{}
\endhead
\bottomrule\noalign{}
\endlastfoot
95 \% & 11,095 & 34,064 \\
98 \% & 4,453 & 13,767 \\
99 \% & 2,229 & 6,908 \\
99.5 \% & 1,115 & 3,460 \\
99.9 \% & 223 & 693 \\
\end{longtable}
}

Reading: even 99 \% is thousands of ppm, which is \textbf{not}
production-grade for a volume product (that needs \textasciitilde99.5
\%+ plus transition and functional tests). It is a reasonable floor for
a low-volume MPW-class die, where every part is also functionally
tested. 98 \% is the minimum I would accept, and only with itemised
exclusions. The model is stuck-at only: \textbf{delay defects are not
covered} (section 2.4).

Budget: the 1 \% gap between 100 \% and 99 \% is nearly consumed by the
already-decided exclusions: crypto keys \textasciitilde0.7 \% with
masking {[}M proxy, section 7{]}, TAP flops (\textasciitilde100-300
flops plus their cone, not on any chain) \textasciitilde0.2-0.4 \%
{[}A{]}. Therefore \textbf{SRAM-boundary shadowing (section 6.3) and key
masking are both mandatory to hit the target}; without them the number
is \textasciitilde89 \%.

\subsection{Tools, licences, packaging}\label{tm:ch:dft-architecture:22-tools-licences-packaging}\origsec{2.2}

{\footnotesize\setlength{\tabcolsep}{3pt}

{
\begin{longtable}[]{@{}
  >{\raggedright\arraybackslash}p{\dimexpr(\linewidth-8\tabcolsep)*162/1000\relax}
  >{\raggedright\arraybackslash}p{\dimexpr(\linewidth-8\tabcolsep)*205/1000\relax}
  >{\raggedright\arraybackslash}p{\dimexpr(\linewidth-8\tabcolsep)*205/1000\relax}
  >{\raggedright\arraybackslash}p{\dimexpr(\linewidth-8\tabcolsep)*205/1000\relax}
  >{\raggedright\arraybackslash}p{\dimexpr(\linewidth-8\tabcolsep)*205/1000\relax}@{}}
\toprule\noalign{}
\begin{minipage}[b]{\linewidth}\raggedright
Tool
\end{minipage} & \begin{minipage}[b]{\linewidth}\raggedright
Licence
\end{minipage} & \begin{minipage}[b]{\linewidth}\raggedright
Packaging
\end{minipage} & \begin{minipage}[b]{\linewidth}\raggedright
Maintenance
\end{minipage} & \begin{minipage}[b]{\linewidth}\raggedright
Scale evidence
\end{minipage} \\
\midrule\noalign{}
\endhead
\bottomrule\noalign{}
\endlastfoot
\textbf{Quaigh} 0.0.6 (Coloquinte/\allowbreak{}quaigh) & MIT OR
Apache-2.0 {[}M: \texttt{gh\ api}{]} & crates.io.
\texttt{cargo\ install} needs OpenSSL headers because
\code{rustsat-kissat} fetches Kissat at build time. Built here with a
throwaway \code{CARGO_HOME} (\src{tools/dft/build_quaigh.sh}). & last
push 2026-02-02, 53 stars {[}M{]} & 2.9 k gates, 1.5 s {[}M{]}. Larger
\textbf{unmeasured}. \\
\textbf{Fault} (AUCOHL) & Apache-2.0 for Fault itself {[}M: Readme{]} &
Swift 5.6; ships \src{flake.nix}, \code{default.nix},
\code{nix/atalanta.nix}, \code{nix/podem.nix} {[}M: repo tree{]}.
\textbf{Not built here.} & last push 2025-08-30, release 0.9.4 on
2025-05-28, 206 stars {[}M{]} & Own engine = pseudo-random patterns +
fault simulation through iverilog (\code{atpg.swift} references
\code{iverilogExecutable}) {[}R{]}; algorithmic back ends are
Atalanta/\allowbreak{}PODEM. SoC scale unknown. \\
\textbf{Atalanta} (hsluoyz/\allowbreak{}atalanta, from Virginia Tech) &
See quote below. GitHub reports no licence file
(\texttt{license:\allowbreak{}\ null}) {[}M{]}. Fault\textquotesingle s
\code{nix/atalanta.nix} marks it \code{licenses.unfree} {[}M{]}. & Plain
C, built from rev \code{a8e07fe} by Fault\textquotesingle s nix file &
last push 2024-05-07 {[}M{]} & ISCAS-class circuits {[}A{]} \\
\textbf{PODEM} (via Fault) & Fault\textquotesingle s Readme, below & via
Fault\textquotesingle s nix file & not checked & ISCAS-class {[}A{]} \\
\end{longtable}
}

}

\textbf{Licence text, verbatim.} Atalanta README: "The source code is
released for teaching and research use only. Any publication in which
ATALANTA was used to obtain the results should cite the reference given
below. ... This program, or any derivative thereof, may not be
reproduced nor used for any commercial product without a written
permission form from Prof. Dong S. Ha. For commercial use of ATALANTA
... please contact to Prof. Dong S. Ha" {[}M:
\texttt{gh\ api\ repos/\allowbreak{}hsluoyz/\allowbreak{}atalanta/\allowbreak{}readme}{]}.
Fault Readme: "SOFTWARE INCLUDED WITH SOME FAULT DISTRIBUTIONS, I.E.
ATALANTA AND PODEM, WHILE FREE TO DISTRIBUTE, ARE PROPRIETARY, AND MAY
NOT BE USED FOR COMMERCIAL PURPOSES." {[}M{]}.

\textbf{Flag (decision 2 and decision 1 interact).} The owner accepted
non-commercial tools. Atalanta\textquotesingle s grant is "teaching and
research use only"; patterns it generates for a die that is sold, or
taped out with commercial intent, are arguably a use "for a commercial
product" and need written permission. This is a legal reading I cannot
settle. If the tape-out becomes commercial, Atalanta/\allowbreak{}PODEM
must be dropped or licensed; \textbf{keep
Quaigh/\allowbreak{}Fault\textquotesingle s own engine as the
licence-clean path} so the flow does not depend on Atalanta. Installing:
add the tools to a new devshell in this repo\textquotesingle s
\src{flake.nix} (existing shells: \code{default}, \code{bench},
\code{hls} at \code{flake.nix:161,229,250} {[}R{]}); nothing
system-wide. Not done in Stage 0 (no flake edit was made).

\textbf{Measured run (Quaigh).} \src{tools/dft/cut.ys} makes a full-scan
combinational cut of \code{timer}
(\texttt{dffunmap;\ expose\ -\allowbreak{}evert-\allowbreak{}dff;\ abc}
to
AND/\allowbreak{}NAND/\allowbreak{}OR/\allowbreak{}NOR/\allowbreak{}XOR/\allowbreak{}XNOR/\allowbreak{}MUX,
written as BLIF): 554 inputs, 976 outputs, 2,925 gates.
\texttt{quaigh\ atpg\ timer\_\allowbreak{}cut.\allowbreak{}blif\ -\allowbreak{}o\ timer.\allowbreak{}test}:

\begin{quote}
Generated 128 random patterns, detecting 2547/4692 faults (54.28 \%
coverage) ... Generated 6528 patterns total, detecting 4336/4692 faults
(92.41 \% coverage) ... Kept 358 patterns, detecting 4336/4692 faults
(92.41 \% coverage)
\end{quote}

It also reported \code{unobservable=351}. \textbf{{[}A{]}} If those are
proven-redundant faults, test coverage would be
4336/\allowbreak{}(4692-351) = 99.88 \%; I did not verify that reading.
Caveats: the fault universe is generic gates after \code{abc}, not
Sky130 cell pins; the cut ignores the scan path (no
shift/\allowbreak{}flush test, no SCE/SCD faults); clocks and resets are
plain inputs; no cell-library awareness, no fault dictionary, no tester
format. Quaigh offers \code{--num-cycles} for sequential random patterns
only.

\subsection{Pattern export to a tester (decision 1)}\label{tm:ch:dft-architecture:23-pattern-export-to-a-tester-decision-1}\origsec{2.3}

Patterns are driven through the TAP (decision 5), so the natural
exchange format is \textbf{SVF} (Serial Vector Format:
\code{SIR}/\code{SDR}/\code{RUNTEST}), which JTAG controllers and, I
believe, OpenOCD can play {[}A: not verified here{]}. Plan: (1) ATPG
gives pseudo-primary-input/\allowbreak{}output vectors for the full-scan
cut; (2) a \textbf{project script} maps them onto chain order (the
\code{chain_check.py} walk gives the order) and emits SVF plus a
cocotb/\allowbreak{}Verilator replay of the same TAP sequence on the
gate-level netlist; (3) the test house or a project converter produces
STIL (IEEE 1450) / WGL for a production ATE. \textbf{No open-source tool
here does (2) or (3)}; both are glue to be written. {[}A{]}

\subsection{What the open-source tools cannot deliver (strict list,
tape-out context) {[}A unless noted{]}}\label{tm:ch:dft-architecture:24-what-the-open-source-tools-cannot-deliver-strict-list-tape-out-context-a-unless-noted}\origsec{2.4}

\begin{itemize}
\tightlist
\item
  \textbf{No scan DRC tool.} Nothing checks, in the commercial sense,
  clock-as-data, reset controllability, X sources, bus contention, latch
  transparency. Stage 1 must carry its own netlist checks
  (\src{tools/dft/chain_check.py} and \code{check_scan_exclusions.py}
  are the start).
\item
  \textbf{No transition / at-speed ATPG, no path-delay.} Quaigh and
  Fault target stuck-at {[}R: Fault Readme topics and \code{atpg.swift};
  \texttt{quaigh\ atpg\ -\allowbreak{}-\allowbreak{}help} offers only
  stuck-at generation plus random sequential patterns via
  \code{--num-cycles}{]}. This survey is not exhaustive.
  \textbf{Decision 7 does not accept zero delay-fault coverage as a
  goal}, so this is an unmet requirement, not an accepted loss. What
  parallel scan ports do and do not change is in section 9.6. What
  delay-fault coverage needs, and the current tools lack: (1) a
  transition-delay fault model (and path-delay for chosen critical
  paths); (2) two-time-frame test generation (launch-on-capture, or
  launch-on-shift), because a single-frame stuck-at
  pattern\textquotesingle s expected response is \textbf{wrong} if the
  capture is applied with two pulses; (3) a transition-fault simulator
  to grade coverage; (4) timing-derived targeting from STA/SDF for
  small-delay defects; (5) capture-clock control: ATE-supplied bursts
  are enough here (section 9.6), an on-chip clock controller is not
  required; (6) masking of the asynchronous-domain crossings, the clock
  gate and false/\allowbreak{}multicycle paths in the capture sequence.
  Commercial ATPG (TetraMAX, Modus, Tessent) provides these; nothing
  available here does. It is \textbf{out of reach of the current tools}.
  Route options are open item N10.
\item
  \textbf{No pattern compression (EDT-style), no diagnosis, no fault
  dictionary.} Test time grows linearly with pattern count (section 8).
\item
  \textbf{No clock-domain-aware capture sequencing} for two asynchronous
  domains with a clock gate; the capture protocol (which domain pulses
  when) is hand-written.
\item
  \textbf{No clock-gate / latch modelling.} \code{en_latch} and the gate
  are untested structures (section 4 item 5).
\item
  \textbf{No X-masking engine} beyond what the cut makes explicit; SRAM
  outputs and the unscanned TAP/key flops must be made X-free by design
  (sections 6.3, 7).
\item
  \textbf{No guarantee of correlation between gate-level fault coverage
  and the Sky130 cell netlist.} Measured coverage on a generic-gate cut
  is an approximation; the final number must be re-simulated on the real
  netlist with Sky130 models.
\item
  \textbf{Scale unproven.} Whether a SAT-based generator finishes on a
  \textasciitilde20 k-flop cut (a few hundred thousand gates {[}A{]}) is
  unmeasured.
\end{itemize}

Plainly: \textbf{no tool available here is credible at SoC scale today
for a signed production coverage number.} The honest fallback if the
Stage 3 spike fails is scan with chain-integrity tests plus partial
coverage on blocks that do complete, reported per block, not as one SoC
number.

\section{Memory inventory}\label{tm:ch:dft-architecture:3-memory-inventory}\origsec{3}

\subsection{Hard macros (12 instances, 147,456 bits = 18 KiB)}\label{tm:ch:dft-architecture:31-hard-macros-12-instances-147456-bits-18-kib}\origsec{3.1}

{\small\setlength{\tabcolsep}{4pt}

{
\begin{longtable}[]{@{}
  >{\raggedright\arraybackslash}p{\dimexpr(\linewidth-10\tabcolsep)*259/1000\relax}
  >{\raggedright\arraybackslash}p{\dimexpr(\linewidth-10\tabcolsep)*253/1000\relax}
  >{\raggedright\arraybackslash}p{\dimexpr(\linewidth-10\tabcolsep)*107/1000\relax}
  >{\raggedright\arraybackslash}p{\dimexpr(\linewidth-10\tabcolsep)*152/1000\relax}
  >{\raggedright\arraybackslash}p{\dimexpr(\linewidth-10\tabcolsep)*107/1000\relax}
  >{\raggedright\arraybackslash}p{\dimexpr(\linewidth-10\tabcolsep)*101/1000\relax}@{}}
\toprule\noalign{}
\begin{minipage}[b]{\linewidth}\raggedright
Instance
\end{minipage} & \begin{minipage}[b]{\linewidth}\raggedright
Macro
\end{minipage} & \begin{minipage}[b]{\linewidth}\raggedright
Words x bits
\end{minipage} & \begin{minipage}[b]{\linewidth}\raggedright
Ports
\end{minipage} & \begin{minipage}[b]{\linewidth}\raggedright
Clock
\end{minipage} & \begin{minipage}[b]{\linewidth}\raggedright
Where
\end{minipage} \\
\midrule\noalign{}
\endhead
\bottomrule\noalign{}
\endlastfoot
\code{u_sram.u_sram_macro} & \code{sky130_sram_4kbyte_1rw1r_32x1024_8} &
1024 x 32 & 1RW + 1R & \code{core_clk} & SoC fabric \\
\code{u_npu.g_on.u_wmem.u_sram_macro} & same & 1024 x 32 & 1RW + 1R,
port 1 tied off & \code{core_clk} & SoC fabric \\
icache \code{u_tag_sram} + 4 x \code{gen_data_sram} &
\code{sky130_sram_1kbyte_1rw1r_32x256_8} & 256 x 32 each (5) & 1RW + 1R,
port 1 tied off & \code{cpu_gated_clk} & \textbf{inside CPU macro} \\
dcache \code{u_tag_sram} + 4 x \code{gen_data_sram} & same & 256 x 32
each (5) & same & \code{cpu_gated_clk} & \textbf{inside CPU macro} \\
\end{longtable}
}

}

Sources: macro instances {[}M: \code{grep} on
\code{soc_top.nl.v:659127,659179,659192};
\texttt{design\_\allowbreak{}\_\allowbreak{}instance\_\allowbreak{}\_\allowbreak{}count\_\allowbreak{}\_\allowbreak{}macros\ =\allowbreak{}\ 3}
in \code{62-cvt-maxcapviolations/state_out.json}{]}, CPU-internal SRAMs
{[}R: \src{pnr/sky130/cpu/macro_placement.cfg}, 10 instances{]}, wrapper
\code{ifdef} arms {[}R{]}:

{\footnotesize\setlength{\tabcolsep}{3pt}

{
\begin{longtable}[]{@{}
  >{\raggedright\arraybackslash}p{\dimexpr(\linewidth-4\tabcolsep)*297/1000\relax}
  >{\raggedright\arraybackslash}p{\dimexpr(\linewidth-4\tabcolsep)*336/1000\relax}
  >{\raggedright\arraybackslash}p{\dimexpr(\linewidth-4\tabcolsep)*347/1000\relax}@{}}
\toprule\noalign{}
\begin{minipage}[b]{\linewidth}\raggedright
Wrapper
\end{minipage} & \begin{minipage}[b]{\linewidth}\raggedright
\code{SRAM_SKY130} arm
\end{minipage} & \begin{minipage}[b]{\linewidth}\raggedright
Other arms
\end{minipage} \\
\midrule\noalign{}
\endhead
\bottomrule\noalign{}
\endlastfoot
\src{rtl/soc/sram_controller.sv} & macro \code{:856-881}, 2-entry skid
read FSM \code{:663-854} & default: flat 32 K-flop array \code{:146-148}
(no ASAP7 macro arm) \\
\src{rtl/npu/npu_weight_mem.sv} & one 4 KB macro \code{:85-104} &
\code{SRAM_ASAP7} 4 x 256x32 behind \code{rv32i_clock_gate}
\code{:119-129}; default FreePDK45 \code{:130-139} \\
\src{rtl/mem/rv32i_icache.sv} / \code{rv32i_dcache.sv} & tag
\code{:82-95} / \code{:105-118}; data x4 \code{:130-143} /
\code{:153-166} & \texttt{elsif\ SRAM\_\allowbreak{}ASAP7}, \code{else}
FreePDK45 \\
\end{longtable}
}

}

Macro interface {[}M: \texttt{grep\ PIN} on the LEF, 127 pins{]}: port 0
\texttt{clk0\ csb0\ web0\ wmask0[3:\allowbreak{}0]\ addr0[9:\allowbreak{}0]\ din0[31:\allowbreak{}0]\ dout0[31:\allowbreak{}0]},
port 1
\texttt{clk1\ csb1\ addr1[9:\allowbreak{}0]\ dout1[31:\allowbreak{}0]}.
Both \code{dout} are negedge-launched (Liberty \code{falling_edge},
clk-to-q 0.336 / 0.365 / 0.481 ns) {[}R{]}. The sim model used by cocotb
(\src{sim/sky130_sram_4kbyte_1rw1r_32x1024_8.sv}) captures inputs on
posedge and acts on negedge {[}R{]}. \textbf{No spare
rows/\allowbreak{}columns}: pass/fail only. Functional port use: the SoC
SRAM writes on port 0 and reads on \textbf{port 1} (\code{dout0} unused,
\code{sram_controller.sv:856-881}); NPU and caches read on port 0.

\subsection{Not hard macros but memory-like (flops)}\label{tm:ch:dft-architecture:32-not-hard-macros-but-memory-like-flops}\origsec{3.2}

8,192-flop \code{dma_engine.linebuf} (\code{:318}), 944-bit
\code{cdc_gray_fifo.mem_q},
UART/\allowbreak{}SPI/\allowbreak{}I2C/\allowbreak{}TRNG FIFOs, SHA
\code{w_q}, NPU \code{ain_q/aout_q}, regfile 992 flops inside the CPU
macro, cache valid/\allowbreak{}dirty arrays (256 flops each) {[}M for
the totals, R for locations{]}. These are scanned as ordinary flops;
they do not need MBIST. \code{boot_rom.mem} folds to constant zero under
\code{__pnr__} {[}R \code{boot_rom.sv:105-122}; not present in the
netlist{]}.

\section{Scan-readiness audit (source reading plus netlist counts)}\label{tm:ch:dft-architecture:4-scan-readiness-audit-source-reading-plus-netlist-counts}\origsec{4}

Measured flop census,
\code{RUN_2026-10-09_06-18-39/04-yosys-synthesis/reports/stat.rpt}
{[}M{]}: \code{dfxtp_2} 17,903 (no reset), \code{dfrtp_2} 843 (async
reset to 0), \code{dfstp_2} 6 (async set), \code{dlxtn_1} 1;
\textbf{18,753 sequential cells, 855 with async reset/set}. CPU macro
(\code{rv32i_cpu_top.nl.v.gz}): \textbf{4,829 \code{dfxtp}, 0 reset
flops, 0 latches} {[}M{]}. By group (Q-net-name attribution; totals
exact, split approximate): DMA \textasciitilde9,012 (8,192 linebuf),
CRYPTO \textasciitilde3,008, NPU \textasciitilde1,169,
\code{u_cpu_axi_cdc} \textasciitilde1,044, I2C \textasciitilde621, TRNG
\textasciitilde515, bus \textasciitilde507, GPIO \textasciitilde352.
Sky130 build facts {[}R{]}: CPU is a blackbox
(\code{rv32i_cpu_top_stub.sv}), GPU is a tie-off, \code{pll_clkgen_stub}
is a pure passthrough
(\texttt{assign\ out\_\allowbreak{}clk\_\allowbreak{}o\ =\allowbreak{}\ ref\_\allowbreak{}clk\_\allowbreak{}i}),
\code{USE_ICG_CELL} is not set, \code{trng_ro_sky130.sv} is not in the
build (the TRNG is an LFSR, \code{trng_lfsr_entropy.sv}).

{\footnotesize\setlength{\tabcolsep}{3pt}

{
\begin{longtable}[]{@{}
  >{\raggedright\arraybackslash}p{\dimexpr(\linewidth-8\tabcolsep)*48/1000\relax}
  >{\raggedright\arraybackslash}p{\dimexpr(\linewidth-8\tabcolsep)*165/1000\relax}
  >{\raggedright\arraybackslash}p{\dimexpr(\linewidth-8\tabcolsep)*256/1000\relax}
  >{\raggedright\arraybackslash}p{\dimexpr(\linewidth-8\tabcolsep)*256/1000\relax}
  >{\raggedright\arraybackslash}p{\dimexpr(\linewidth-8\tabcolsep)*256/1000\relax}@{}}
\toprule\noalign{}
\begin{minipage}[b]{\linewidth}\raggedright
\#
\end{minipage} & \begin{minipage}[b]{\linewidth}\raggedright
Item
\end{minipage} & \begin{minipage}[b]{\linewidth}\raggedright
Where
\end{minipage} & \begin{minipage}[b]{\linewidth}\raggedright
What it is
\end{minipage} & \begin{minipage}[b]{\linewidth}\raggedright
Stage 1 must
\end{minipage} \\
\midrule\noalign{}
\endhead
\bottomrule\noalign{}
\endlastfoot
1 & Scan ports & \code{soc_top.sv:196-261} & none exist (35 ports / 266
bits {[}M{]}) & add \code{scan_mode_i}, \code{scan_en_i},
\code{scan_rst_ni}, \code{scan_in_i[C]}, \code{scan_out_o[C]}; update
\code{soc_top_ports.v} and IO config \\
2 & Tie-offs & \code{soc_top.sv:288-289} & \code{SCAN_MODE_TIE_OFF=0},
\code{SCAN_RST_TIE_OFF=1}, 6 sites / 12 connections & replace with the
ports (sites: \code{:647}, \code{:707}, \code{:1021}, \code{:1248},
\code{:2076}, \code{:2182}) \\
3 & \code{cdc_reset_sync} hook & \code{cdc_reset_sync.sv:82,128} & mux
is on the \textbf{input} of the sync chain only;
\texttt{rst\_\allowbreak{}n\_\allowbreak{}o\ =\allowbreak{}\ sync\_\allowbreak{}q[.\allowbreak{}.\allowbreak{}.\allowbreak{}]}
is a scan-chain flop that toggles during shift and would pulse every
downstream async reset (262+83+133+56+... flops) & move/add the mux to
the \textbf{output}; 10 instances \\
4 & \code{cdc_sync} & \code{cdc_sync.sv:151-177} & async reset from raw
\code{rst_n_i}, no hook; \textasciitilde60 instances (GPIO 32, FIFO
pointers, 4 in \code{soc_top}) & add \code{scanmode_i/scan_rst_ni} \\
5 & Clock gate & \code{rv32i_clock_gate.sv:13-27}; users
\code{soc_top.sv:2320} (live), \code{:2321} (GPU, dead),
caches/\allowbreak{}NPU (ASAP7 arm only) & latch + AND, no test enable;
\code{ICGx1} arm hard-ties \texttt{SE=\allowbreak{}1'b0} & add
\code{test_en}, \code{en_latch} input =
\texttt{en\ |\ test\_\allowbreak{}en}; tie \code{.SE(test_en)} in the
ICG arm. \textbf{One} live site on Sky130 \\
6 & Latch & \code{rv32i_clock_gate.sv:24} & the only \code{always_latch}
in \src{rtl/}; \code{dlxtn_1} in netlist & none beyond item 5; ATPG must
treat it as a clock-gate latch \\
7 & Resets & \code{pll_subsystem.sv:97,100}, \code{soc_top.sv:826},
\code{async_axi_fifo.sv:167},
\texttt{apb\_\allowbreak{}cdc\_\allowbreak{}bridge.\allowbreak{}sv:\allowbreak{}\textasciitilde{}325}
&
\texttt{core\_\allowbreak{}rst\_\allowbreak{}n\ =\allowbreak{}\ rst\_\allowbreak{}n\_\allowbreak{}i\ \&\ pll\_\allowbreak{}locked},
\texttt{pll\_\allowbreak{}rst\_\allowbreak{}n\ =\allowbreak{}\ rst\_\allowbreak{}n\_\allowbreak{}i\ \&\ pll\_\allowbreak{}enable}
(an APB register), \code{cpu_domain_rst_n} = AND of 3 flop-derived
resets, \code{wdt_cpu_rst_req_q}, PMU resets. Only
\code{rst_n_i}/\code{cpu_rst_n_i} are real pins; all others are flop- or
register-driven & in scan mode force every internal async reset inactive
(mux with \code{scan_rst_ni}) at the \emph{net} feeding the flops, not
at a sync input \\
8 & PLL stub & \code{pll_clkgen_stub.sv:55,69} & passthrough clock,
async-reset lock counter & make \code{locked_q} controllable; a real PLL
(Phase 7) will need a test-clock bypass. \textbf{Not needed while it is
a stub} \\
9 & Two domains & see below & \code{clk_i} 17,809 flops,
\code{cpu_clk_i} 173, \code{cpu_gated_clk} 771 + 4,829 macro {[}M{]} &
separate chains per domain, no \code{clock_mix}; add
\texttt{create\_\allowbreak{}clock\ cpu\_\allowbreak{}clk\_\allowbreak{}i}
to the SDC \\
10 & \code{cdc_gray_fifo} & \code{cdc_gray_fifo.sv:191-199} & no-reset
\code{mem_q} (944 bits) read combinationally across domains & ATPG
capture constraint or a lock-up element at the read path \\
11 & \code{apb_cdc_bridge} & \code{:549,649,716} & async-reset toggle
handshakes, 3 instances & reset override; chains per side \\
12 & CPU macro & \code{rv32i_cpu_top_stub.sv} & 4,829 flops, no scan
pins; \textbf{built with
\texttt{USE\_\allowbreak{}SYNLIG:\allowbreak{}\ true}}
(\src{pnr/sky130/cpu/config.json:88}) & add scan ports to
\code{rv32i_cpu_top}, pins in \code{pin_order.cfg}, re-harden (section
5) \\
13 & SRAM macros & 2 in SoC, 10 in CPU & opaque, negedge-launched
outputs, X to ATPG & collar with bypass (section 6) \\
14 & Large no-reset arrays & \code{dma_engine.sv:318-327},
\code{cdc_gray_fifo}, FIFOs & \textasciitilde11 k flops; X until written
& none for scannability; they dominate chain length \\
15 & Key material & \code{crypto_accel.sv:441},
\code{aes128_core.sv:442} & scan-reachable & policy (section 7) \\
16 & TRNG & \code{trng.sv:107-111} & LFSR only in this build; the real
ring oscillator (a comb loop) is blackboxed and not in any file list &
none now; a future RO needs a scan bypass \\
17 & Boot ROM & \code{boot_rom.sv:105-122} & constant 0 & none (coverage
accounting only) \\
18 & \code{initial}/\allowbreak{}tri-state/\allowbreak{}neg-edge RTL
flops & grep over \src{rtl/} & \textbf{none found}; the negedge
behaviour lives in the SRAM macro model only {[}R{]} & none \\
19 & Combinational loops & grep & \textbf{none found} & none \\
20 & SDC & \code{sky130_soc.sdc:237} & only \code{core_clk};
\code{cpu_clk_i} flops unconstrained & add clock and
\texttt{set\_\allowbreak{}clock\_\allowbreak{}groups\ -\allowbreak{}asynchronous}
before Stage 2 \\
\end{longtable}
}

}

\textbf{Stage 1a status (bead \bead{j41m.2}, section 14):} items 1, 2,
3, 5, 6, 7, 8 and 15 are implemented in the fabric; item 4 was resolved
without a \code{cdc_sync} hook (section 14.2); items 9 and 11 are
covered by the reset/\allowbreak{}clock overrides but the per-domain
chains are Stage 2; items 12 (CPU macro scan pins) and 13 (SRAM collars)
are \textbf{not} done (section 14.6, Stage 4); item 20
(\texttt{create\_\allowbreak{}clock\ cpu\_\allowbreak{}clk\_\allowbreak{}i})
is deferred (section 14.5).

Existing hook in detail (bead \bead{07n}) {[}R{]}: declared in
\code{cdc_reset_sync.sv:61-62}, forwarded by
\code{async_axi_fifo.sv:108-109,194-195,204-205} and
\code{apb_cdc_bridge.sv:304-305,353-354,363-364}. Behaviour
\texttt{assign\ rst\_\allowbreak{}n\_\allowbreak{}async\ =\allowbreak{}\ scanmode\_\allowbreak{}i\ ?\ scan\_\allowbreak{}rst\_\allowbreak{}ni\ :\allowbreak{}\ rst\_\allowbreak{}n\_\allowbreak{}i;}
(\code{:82}). The comment at \code{soc_top.sv:279-283} warns that new
top ports would "perturb the pinned CPU macro"; that note is about
ASAP7\textquotesingle s \code{pin_order.cfg}. The Sky130 SoC has no
\code{FP_PIN_ORDER_CFG}; the Sky130 \emph{CPU} macro does.

Domain crossings that need chain separation or lock-up handling {[}R{]}:
\code{u_cpu_axi_cdc} (\code{soc_top.sv:1247}, 5 gray FIFOs),
\code{u_apb_dbg_cdc} (\code{:1018}), \code{u_apb_pll_cdc}
(\code{:2073}), \code{u_apb_pll2_cdc} (\code{:2179}), \code{cdc_sync} x2
for PMU clock/\allowbreak{}isolate (\code{:583,611}), IRQ syncs
(\code{:771,780}), \code{cdc_reset_sync} x2 (\code{:642,702}).
\code{cpu_gated_clk} and \code{cpu_clk_i} are the same root and phase,
so they can share a chain \emph{if} the gate is forced on in scan.

\textbf{Side finding, now bead \bead{dud4} (P1; decision 9).}
\src{CLAUDE.md} says Sky130 is unaffected by the \code{ma7} Synlig
miscompile because "both use the sv2v frontend".
\src{pnr/sky130/soc/config.json:41} has
\texttt{USE\_\allowbreak{}SYNLIG:\allowbreak{}\ false}, but
\src{pnr/sky130/cpu/config.json:88} has
\texttt{USE\_\allowbreak{}SYNLIG:\allowbreak{}\ true}, and the CPU macro
contains \code{rv32i_regfile.sv}\textquotesingle s runtime-indexed read
idiom {[}R{]}. I did not test whether the Sky130 CPU macro is affected
{[}A: the defect was \code{OPT_MUXTREE} on a Synlig-elaborated design,
PDK-independent in principle{]}. \textbf{Decided:} test first, with the
\code{ma7} branch-dependent differential (the
\code{tb_cpu_macro_check.v} method) on the Sky130 CPU macro gate netlist
against RTL, \textbf{no PD run}; then the user decides on re-hardening.
\textbf{Stage 2\textquotesingle s CPU-macro re-harden depends on that
outcome} (sections 5, 10), because a frontend switch changes the netlist
the scan chains are inserted into.

\section{Hard macros}\label{tm:ch:dft-architecture:5-hard-macros}\origsec{5}

\begin{itemize}
\tightlist
\item
  The Sky130 SoC consumes the CPU as a hard macro
  (\code{MACROS.rv32i_cpu_top}, 9 corner libs, LEF, GDS, spice from
  \code{../cpu/macro/}) {[}R \code{soc/config.json}{]}; 3600 x 1800 um,
  403 pins {[}R LEF{]}; produced by \src{pnr/sky130/cpu/config.json}
  (period 13.333 ns, die \code{[0,0,3600,1800]},
  \texttt{FP\_\allowbreak{}PIN\_\allowbreak{}ORDER\_\allowbreak{}CFG\ pin\_\allowbreak{}order.\allowbreak{}cfg},
  \texttt{ERRORS\_\allowbreak{}ON\_\allowbreak{}UNMATCHED\_\allowbreak{}IO\ both})
  {[}R{]}.
\item
  \code{rv32i_cpu_top.sv} contains no \code{scan} identifier {[}M:
  grep{]}. Required for scan through it: (1) scan ports on
  \code{rv32i_cpu_top}; (2) every new pin added to \code{pin_order.cfg},
  because unmatched IO is an error, and earlier pin-order attempts
  crashed at \code{Odb.CustomIOPlacement} and caused GRT-0118 congestion
  (\code{pin_order.README.md}) {[}R{]}; (3) insertion \textbf{inside}
  the macro flow; (4) regenerate LEF, 9 Liberty views, GDS, spice,
  netlist; (5) re-run the SoC.
\item
  The 10 cache SRAMs are inside it, so MBIST ports (\code{mbist_en},
  start, done, fail, per-memory fail optional) are macro pins too.
\item
  Chain architecture: the macro gets its own chains (e.g. 4 x
  \textasciitilde1,207 flops), wired in \code{soc_top} RTL straight to
  top-level \code{scan_in/scan_out}. The SoC-level OpenROAD pass never
  sees those flops, so they are stitched in the macro run. Fabric chains
  are made at SoC level.
\item
  \textbf{Dependency on \code{dud4} (decision 9).} The CPU-macro
  re-harden for scan should be the same re-harden that fixes the Synlig
  exposure if the user chooses to re-harden for \code{dud4}; doing them
  as two rounds would cost another \textasciitilde6-7 h per round
  (below). So Stage 2\textquotesingle s CPU-macro half waits for the
  \code{dud4} differential result and the user\textquotesingle s
  decision. The SoC-fabric half does not depend on it.
\item
  Re-harden cost {[}R from run history; the sum is {[}A{]}{]}: CPU full
  flow \textasciitilde2h17m for 74 steps
  (\code{RUN_2026-07-21_18-02-11}), adopted run estimated 2.5-3 h; CPU
  harden had 8 runs and one pin-order attempt crashed twice before
  succeeding; SoC harden 3.4-3.6 h (\code{runtime.txt} sums {[}M{]}:
  12,328 s and the prior accepted run 3.62 h); SoC Magic DRC peak
  \textasciitilde15.4 GiB RSS on a 15 GiB host; host reboots every 2-8
  h. So **\textasciitilde6-7 h per CPU+SoC round, realistically 2-3
  rounds: 12-21 h of wall time**, with the KLayout DRC still skipped.
\item
  Timing headroom is thin: CPU macro standalone nom\_tt setup was +0.174
  ns at 13.333 ns {[}R
  \src{memory/pd/run_state.archived_pd_20260918_repo_root.md}{]}. A scan
  mux adds +0.066/\allowbreak{}+0.085 ns (rise/\allowbreak{}fall,
  nom\_tt) to every \code{dfxtp_2} D input at a mid-table point {[}M,
  section 8{]}, so the CPU macro may need a period relaxation or a
  retime. This is the single most likely Stage 2 blocker.
\end{itemize}

\section{MBIST architecture (tester-only, TAP-driven; decisions 4 and
5)}\label{tm:ch:dft-architecture:6-mbist-architecture-tester-only-tap-driven-decisions-4-and-5}\origsec{6}

\begin{quote}
\textbf{Superseded for Stage 4 by section 16 (2026-10-17): algorithm
set, TAP registers, the control-word crossing (depth-2
\code{cdc_gray_fifo}) and the CPU-macro proposal are specified there.}
The text below is the Stage-0 plan, kept as history.

\textbf{Status (2026-10-13): Stage 4 plan, NOT built.} No MBIST RTL
exists; \code{MBIST_CTRL}/\code{MBIST_STATUS} (IR 6/7) decode to BYPASS
(\code{jtag_tap.sv:17}). The "2-FF sync" crossings drawn and described
in this section (and the "write, then wait \textgreater= 3 destination
cycles" rule in section 9.3) apply \textbf{only to the future MBIST
control/\allowbreak{}status words}. They do \textbf{not} describe the
as-built Stage 1b scan controls, which are deliberately unsynchronised
(section 15.5, crossing 2). The crossing mechanism for the MBIST control
word is not settled; see "Deferred: MBIST control-word crossing" at the
end of section 15.5.
\end{quote}

\begin{lstlisting}[language={},caption={Diagram source (Mermaid)}]
flowchart LR
  TAP[JTAG TAP, TCK domain] -->|MBIST_CTRL TDR, 2-FF sync| C1[mbist_ctrl core_clk]
  TAP -->|same, 2-FF sync| C2[mbist_ctrl cpu_gated_clk, inside CPU macro]
  C1 --> K1[collar: SoC SRAM]
  C1 --> K2[collar: NPU wmem]
  C2 --> K3[10 collars: I$/D$ tag+data]
  K1 <--> M1[4KB macro]
  K2 <--> M2[4KB macro]
  K3 <--> M3[10 x 1KB macros]
  C1 -->|done/fail sticky, 2-FF sync| TAP
  C2 -->|done/fail sticky, 2-FF sync| TAP
\end{lstlisting}

\subsection{Algorithm and coverage}\label{tm:ch:dft-architecture:61-algorithm-and-coverage}\origsec{6.1}

\begin{itemize}
\tightlist
\item
  \textbf{March C- (10N)}:
  \texttt{up(w0)\ up(r0,\allowbreak{}w1)\ up(r1,\allowbreak{}w0)\ down(r0,\allowbreak{}w1)\ down(r1,\allowbreak{}w0)\ down(r0)}.
  Covers stuck-at, transition, address-decoder faults, unlinked
  inversion/\allowbreak{}idempotent/\allowbreak{}state coupling faults.
  \textbf{Does not cover} neighbourhood-pattern-sensitive faults,
  read-disturb and retention faults (need a pause/\allowbreak{}hammer
  element), write-mask (\code{wmask0}) lane faults unless each byte lane
  is exercised separately, linked coupling faults, and timing faults
  beyond what a 40 MHz functional clock implies. Add: \textbf{two data
  backgrounds} (solid, checkerboard), a \textbf{byte-lane pass}, and for
  the 1R port a \textbf{port-1 read pass} (write via port 0, read
  everything via port 1, plus a simultaneous port-0 write / port-1 read
  of different addresses). Port 1 is the SoC SRAM\textquotesingle s
  functional read path (\code{sram_controller.sv:856-881} {[}R{]}), so
  it must be tested, not assumed. No spare rows/\allowbreak{}columns
  exist: \textbf{pass/fail only}.
\item
  \textbf{Test time {[}A, arithmetic{]}}: N=1024 -\textgreater{} 10,240
  cycles; N=256 -\textgreater{} 2,560. Both fabric macros in parallel
  (shared address/\allowbreak{}data generator, two comparators) =
  \textbf{256 us at 40 MHz}; the 10 CPU macros in parallel = 2,560
  cycles (\textless= 34 us at 75 MHz). Serial worst case 46,080 cycles =
  1.15 ms. With 2 backgrounds + lane + port-1 passes,
  \textasciitilde3-4x: under 5 ms. Test time is not a constraint, so
  parallel-within-domain is chosen for cheap control.
\end{itemize}

\subsection{Control and observation through the TAP only (decision
4)}\label{tm:ch:dft-architecture:62-control-and-observation-through-the-tap-only-decision-4}\origsec{6.2}

\begin{itemize}
\tightlist
\item
  \textbf{Start:} the tester loads \code{MBIST_RUN}, shifts a control
  TDR (domain select, algorithm/\allowbreak{}background select,
  \code{go}), and \code{Update-DR} latches it in the TCK domain. A
  \code{go} level crosses to \code{core_clk} and to the CPU
  macro\textquotesingle s clock through a \code{cdc_sync} per control
  bit; the control word is quasi-static (written, then held) so a 2-FF
  synchroniser is sufficient {[}A{]} (\textbf{Stage 4 plan only; the
  mechanism is a PROPOSAL, see the deferred note at the end of section
  15.5}). The "quasi-static, written then held" contract is the one the
  repo already states for \code{scanmode_i} {[}R
  \code{cdc_reset_sync.sv:48-49}{]}, but \code{scanmode_i} and the other
  as-built TAP scan controls are \textbf{not} synchronised (section
  15.5, crossing 2).
\item
  \textbf{Run:} the memory test runs on the \textbf{functional clock}
  (\code{core_clk} / \code{cpu_gated_clk}), i.e. at speed for the
  memories. The tester idles in \code{Run-Test/Idle} and polls.
\item
  \textbf{Observe:} sticky \code{done} and \code{fail} flops (one
  \code{done} and a \code{fail} bit per memory, 12 + 2) sit in the
  functional domain and are captured into \code{MBIST_STATUS} at
  \code{Capture-DR} through a 2-FF synchroniser (\code{done}/\code{fail}
  are stable once \code{done} is high). The tester reads them from TDO.
\item
  \textbf{What is deliberately absent:} no APB slot (the
  \code{0x2001_1000} slot analysis in the previous draft is withdrawn),
  no IRQ bit, no power-on auto-run, no firmware-readable status.
  Enforcement: \code{mbist_go} is generated only from the TAP control
  register, reset by \code{TRST_N}/\allowbreak{}Test-Logic-Reset, and
  the controller has no bus slave port. A check in Stage 4:
  grep/\allowbreak{}elaboration test that no \code{soc_top} bus slave
  decodes to the MBIST block, and a cocotb test that firmware-visible
  address space is bit-identical with and without MBIST compiled in.
\item
  The controller and the TAP control register must come out of a
  \textbf{functional reset in a known idle state}; \code{rst_n_i} must
  not start MBIST.
\end{itemize}

\subsection{Collar placement, scan interaction, SRAM-boundary
shadowing}\label{tm:ch:dft-architecture:63-collar-placement-scan-interaction-sram-boundary-shadowing}\origsec{6.3}

\begin{itemize}
\tightlist
\item
  \textbf{Collar placement}: the mux sits on the \textbf{functional side
  of the \code{ifdef} arms}, between the
  cache/\allowbreak{}controller/\allowbreak{}NPU logic and the
  arm\textquotesingle s macro or model instance, so the FreePDK45 and
  ASAP7 arms and the cocotb sim model sit behind the same collar and the
  sim exercises it {[}R: arms listed in section 3{]}. For
  \code{sram_controller} the default arm is a flat array; BIST there
  tests nothing physical but must still elaborate. Collar per macro:
  \textasciitilde48 muxes on port 0 (\code{addr0} 10, \code{din0} 32,
  \code{wmask0} 4, \code{csb0}, \code{web0}), \textasciitilde11 on port
  1, a 32-bit compare per read port {[}A{]}.
\item
  \textbf{Scan interaction}: in scan mode an SRAM macro is a black box,
  so (a) its outputs must be controllable: a 32-bit shadow flop per
  \code{dout} (on a chain) drives the downstream logic in scan mode; (b)
  its inputs must be observable: an XOR-reduction of
  \code{addr/din/wmask/csb/web} into a few scan flops per macro
  (aliasing accepted) {[}A: standard technique, not built{]}. The MBIST
  controller\textquotesingle s own flops (\textasciitilde150-300 per
  domain {[}A{]}) go on the chains. Without (a)/(b) the SRAM
  fan-in/\allowbreak{}fan-out cones are X-blocked and the coverage
  target of section 2.1 is not reachable.
\item
  \textbf{Timing risk}: the collar mux lands on exactly the pins that
  \code{e45j} lists as max-slew residue. In
  \code{RUN_2026-10-09_06-18-39}, 18 SRAM-input pins violate max slew at
  max\_ss (limit 0.5 ns): worst \code{u_npu...addr0[0]} 0.7009 ns,
  \code{u_sram...addr1[5]} 0.6553, \code{wmask0[3]} 0.6260 {[}M:
  \code{51-openroad-stapostpnr/max_ss_100C_1v60/checks.rpt}{]}. A
  mux2\_1 driving the pin adds a drive stage (could help) or load (could
  hurt) {[}A: untested{]}; those pins get re-closed either way. Setup is
  not the concern: no SRAM path in the top 1,000 setup paths at max\_ss,
  worst setup +3.3286 ns (max\_ss) {[}M{]}. The NPU MAC path was +1.84
  ns at Gate A but +3.56 ns at max\_ss post-route in Gate B {[}R
  \texttt{Doc:\allowbreak{}\ ds:\allowbreak{}ch:\allowbreak{}phase6-\allowbreak{}ip-\allowbreak{}expansion-\allowbreak{}plan\ (line\ 896\ of\ the\ original\ file)}{]};
  the collar does not touch it.
\item
  \textbf{Repair}: none. \textbf{Area {[}A{]}}: 12 collars x
  \textasciitilde1.2 k um2 + two controllers \textasciitilde12 k +
  shadow flops \textasciitilde20 k = \textasciitilde45-60 k um2,
  \textasciitilde0.3 pp of the 20.36 M um2 core.
\end{itemize}

\section{Crypto key registers: EXCLUDED from scan (decision 3)}\label{tm:ch:dft-architecture:7-crypto-key-registers-excluded-from-scan-decision-3}\origsec{7}

\subsection{What is excluded {[}M{]}}\label{tm:ch:dft-architecture:71-what-is-excluded-m}\origsec{7.1}

\code{rk_q} (AES round key, 128 flops, invertible to the master key) and
the key shadow register (128 flops; the netlist Q-net name is
\code{u_crypto.g_aes.u_aes.key_i}, the AES core\textquotesingle s key
input, which is \code{crypto_accel.sv}\textquotesingle s \code{key_q}
{[}A: inferred from the name and the 128 count; \code{key_q} itself does
not appear as a net name{]}). Total \textbf{256 flops = 1.37 \% of the
18,752 fabric flops} {[}M: \src{tools/dft/crypto_cone.py} on
\code{RUN_2026-10-06_22-43-34/04-yosys-synthesis/soc_top.nl.v}{]}. They
are \code{dfxtp_2} (no reset).

\subsection{Coverage lost, quantified {[}M, graph analysis, not
ATPG{]}}\label{tm:ch:dft-architecture:72-coverage-lost-quantified-m-graph-analysis-not-atpg}\origsec{7.2}

\code{crypto_cone.py} builds the connectivity graph of the synthesised
netlist (91,424 cells). Fault counts use an \textbf{uncollapsed
pin-based proxy} (2 faults per cell pin; 778,780 total), not an ATPG
number.

{
\begin{longtable}[]{@{}
  >{\raggedright\arraybackslash}p{\dimexpr(\linewidth-6\tabcolsep)*632/1000\relax}
  >{\raggedright\arraybackslash}p{\dimexpr(\linewidth-6\tabcolsep)*116/1000\relax}
  >{\raggedright\arraybackslash}p{\dimexpr(\linewidth-6\tabcolsep)*137/1000\relax}
  >{\raggedright\arraybackslash}p{\dimexpr(\linewidth-6\tabcolsep)*95/1000\relax}@{}}
\toprule\noalign{}
\begin{minipage}[b]{\linewidth}\raggedright
Quantity
\end{minipage} & \begin{minipage}[b]{\linewidth}\raggedright
Cells
\end{minipage} & \begin{minipage}[b]{\linewidth}\raggedright
Fault proxy
\end{minipage} & \begin{minipage}[b]{\linewidth}\raggedright
Share
\end{minipage} \\
\midrule\noalign{}
\endhead
\bottomrule\noalign{}
\endlastfoot
Excluded flops\textquotesingle{} own D/CLK/Q pins & 256 flops & 1,536 &
0.20 \% \\
Fan-in cone observable \textbf{only} through the excluded flops & 386 &
4,018 & 0.52 \% \\
Fan-out cone of the excluded flops (affected, upper bound) & 9,142 &
87,388 & \textbf{11.22 \%} \\
Cells with no scan-controllable input at all & 0 & 0 & 0 \% \\
\end{longtable}
}

Interpretation {[}A{]}: the fan-out cone is large because the round key
feeds the AES datapath. Those cells are not dead (the scanned \code{s_q}
state controls their other inputs), but the excluded flops hold values
the tester cannot set; ATPG must treat them as X, and X propagates
through every XOR/AND on the way to a scanned capture flop. \textbf{Left
unmasked, up to \textasciitilde11 \% of the fault universe is X-blocked
and the 99 \% target is unreachable.}

\subsection{Exclusion alone leaks the key; masking is mandatory
{[}A, reasoned{]}}\label{tm:ch:dft-architecture:73-exclusion-alone-leaks-the-key-masking-is-mandatory-a-reasoned}\origsec{7.3}

Not scanning a flop does not stop its value reaching scanned flops. In
scan mode a tester can shift a chosen \code{s_q} in, pulse one capture,
and shift out
\texttt{f(s\_\allowbreak{}q,\allowbreak{}\ rk\_\allowbreak{}q)} (the AES
round output). That is key-dependent data on TDO, from which the round
key is recoverable. So the exclusion is only secure together with a
\textbf{scan-mode mask on the excluded flops\textquotesingle{} outputs}:
\texttt{key\_\allowbreak{}eff\ =\allowbreak{}\ key\_\allowbreak{}q\ \&\ \{128\{\textasciitilde{}scan\_\allowbreak{}mode\}\}},
\texttt{rk\_\allowbreak{}eff\ =\allowbreak{}\ rk\_\allowbreak{}q\ \&\ \{128\{\textasciitilde{}scan\_\allowbreak{}mode\}\}}
(256 AND2 gates, \textasciitilde1.6 k um2 {[}A{]}). The mask also makes
the cone controllable (constant 0), which is what shrinks the loss from
\textasciitilde11 \% to the first two rows plus the 256 masked
nets\textquotesingle{} faults: **\textasciitilde0.7 \% {[}M proxy{]} +
the AND-gate key-side faults**. Mask placement is an RTL change in
\code{crypto_accel.sv}/\code{aes128_core.sv} (Stage 1) and must not be
optimised away (hierarchy \code{keep}, or a dedicated cell). The key
\emph{load} path (APB write-snoop into the flops) is not observable
through scan because the excluded flops\textquotesingle{} D pins are not
captured into any chain.

\subsection{How the exclusion is enforced in the insertion flow
{[}M{]}}\label{tm:ch:dft-architecture:74-how-the-exclusion-is-enforced-in-the-insertion-flow-m}\origsec{7.4}

\code{set_dont_touch} on the excluded instances is honoured by the
Sky130 flow\textquotesingle s OpenROAD \code{scan_replace} +
\code{insert_dft}: with 256 instances marked, all 256 stayed
\code{dfxtp_2}, the other 18,496 became scan cells, and the chains
totalled 943 + 17,553 = 18,496 stitched flops
(\src{tools/dft/soc_dft_exclude.tcl}, log \code{soc_excl.log}) {[}M{]}.
The exclusion list must be \textbf{derived from the netlist, not typed}:
instance names are anonymous (\code{_147697_}), so the plugin step
derives them from the Q-net names
\texttt{u\_\allowbreak{}crypto.\allowbreak{}g\_\allowbreak{}aes.\allowbreak{}u\_\allowbreak{}aes.\allowbreak{}(rk\_\allowbreak{}q|key\_\allowbreak{}i)}
(regex in \src{tools/dft/check_scan_exclusions.py}) before calling
\code{set_dont_touch}.

\subsection{How it is checked {[}M{]}}\label{tm:ch:dft-architecture:75-how-it-is-checked-m}\origsec{7.5}

\texttt{tools/\allowbreak{}dft/\allowbreak{}check\_\allowbreak{}scan\_\allowbreak{}exclusions.\allowbreak{}py\ scan\_\allowbreak{}netlist.\allowbreak{}v\ -\allowbreak{}-\allowbreak{}expect\ 256}
exits non-zero unless (a) exactly 256 flops match the regex, (b) none is
a scan cell or sits on a chain, (c) every other flop is a scan cell.
Results: the \code{dont_touch} netlist \textbf{passes} (256 excluded,
18,496 other, 0 problems); the full-scan netlist (negative control)
\textbf{fails} with 256 problems; a wrong regex \textbf{fails} with a
count mismatch, so a renamed or merged register cannot silently drop the
exclusion. A second check, to add in Stage 2, must prove the mask
exists: every fan-out sink of an excluded Q net is an AND gate whose
other input is \texttt{\textasciitilde{}scan\_\allowbreak{}mode}
(structural check on the netlist). The key is also still readable by
anyone with functional access (no privilege split exists, per
\src{CLAUDE.md} item 14); scan exclusion does not change that.

\subsection{Residual}\label{tm:ch:dft-architecture:76-residual}\origsec{7.6}

The excluded flops get no structural test. A functional known-answer
test (FIPS-197 vector, already in \code{test_crypto}) run by boot
firmware would exercise them; that is a production-test item, outside
the TAP flow {[}A{]}. The same applies to the TAP\textquotesingle s own
flops (section 9.4).

\section{Cost estimate}\label{tm:ch:dft-architecture:8-cost-estimate}\origsec{8}

\textbf{Area.} Cell areas from the tt Liberty {[}M:
\code{sky130_fd_sc_hd__tt_025C_1v80.lib}{]}: \code{dfxtp_2} 21.2704
-\textgreater{} \code{sdfxtp_2} 27.5264 (+6.256); \code{dfrtp_2} 26.2752
-\textgreater{} \code{sdfrtp_2} 32.5312 (+6.256); \code{dfstp_2} 26.2752
-\textgreater{} \code{sdfsbp_2} 38.7872 (+12.512); \code{dfxtp_1/_4}
also +6.256.

{\small\setlength{\tabcolsep}{4pt}

{
\begin{longtable}[]{@{}
  >{\raggedright\arraybackslash}p{\dimexpr(\linewidth-6\tabcolsep)*206/1000\relax}
  >{\raggedright\arraybackslash}p{\dimexpr(\linewidth-6\tabcolsep)*426/1000\relax}
  >{\raggedright\arraybackslash}p{\dimexpr(\linewidth-6\tabcolsep)*142/1000\relax}
  >{\raggedright\arraybackslash}p{\dimexpr(\linewidth-6\tabcolsep)*206/1000\relax}@{}}
\toprule\noalign{}
\begin{minipage}[b]{\linewidth}\raggedright
Item
\end{minipage} & \begin{minipage}[b]{\linewidth}\raggedright
Calculation
\end{minipage} & \begin{minipage}[b]{\linewidth}\raggedright
um2
\end{minipage} & \begin{minipage}[b]{\linewidth}\raggedright
Label
\end{minipage} \\
\midrule\noalign{}
\endhead
\bottomrule\noalign{}
\endlastfoot
Fabric scan-flop swap & 17,903 x 6.256 + 843 x 6.256 + 6 x 12.512 &
\textbf{117,350} & {[}M{]} counts x {[}M{]} areas \\
as fraction & of synthesised stdcell area 1,150,630 / of core 20,359,700
& +10.2 \% / +0.58 pp & {[}M{]} \\
CPU macro & 4,829 x 6.256 & \textbf{30,209} & {[}M{]}; macro die fixed
6.48 M \\
SCE buffer tree & \textasciitilde18.8 k sinks / \textasciitilde20 per
buffer & \textasciitilde4 k & {[}A{]} \\
Hold buffers on chain links & unknown & ? & {[}A{]} largest unknown \\
MBIST + collars & section 6 & 45-60 k & {[}A{]} \\
\end{longtable}
}

}

Current utilisation 45.21 \% {[}M:
\code{design__instance__utilization}{]}. Scan alone moves it to
\textasciitilde45.8 \%, with MBIST \textasciitilde46.1 \%; plus hold
buffers {[}A{]}. Headroom is not the constraint.

\textbf{Timing.} Setup of the D pin, mid-table point,
\src{tools/dft/lib_setup.py} {[}M, read from the Liberty, not STA{]}:

{
\begin{longtable}[]{@{}
  >{\raggedright\arraybackslash}p{\dimexpr(\linewidth-4\tabcolsep)*253/1000\relax}
  >{\raggedright\arraybackslash}p{\dimexpr(\linewidth-4\tabcolsep)*283/1000\relax}
  >{\raggedright\arraybackslash}p{\dimexpr(\linewidth-4\tabcolsep)*444/1000\relax}@{}}
\toprule\noalign{}
\begin{minipage}[b]{\linewidth}\raggedright
Cell
\end{minipage} & \begin{minipage}[b]{\linewidth}\raggedright
tt D rise/fall (ns)
\end{minipage} & \begin{minipage}[b]{\linewidth}\raggedright
ss\_\allowbreak{}100C\_\allowbreak{}1v60 D rise/fall (ns)
\end{minipage} \\
\midrule\noalign{}
\endhead
\bottomrule\noalign{}
\endlastfoot
\code{dfxtp_2} & 0.086 / 0.201 & 0.282 / 0.590 \\
\code{sdfxtp_2} (D pin) & 0.152 / 0.286 & 0.469 / 0.709 \\
delta & +0.066 / +0.085 & +0.187 / +0.119 \\
\code{dfrtp_2} -\textgreater{} \code{sdfrtp_2} &
0.090/\allowbreak{}0.205 -\textgreater{} 0.205/\allowbreak{}0.405 &
0.260/\allowbreak{}0.599 -\textgreater{} 0.625/\allowbreak{}1.049
(\textbf{+0.365/\allowbreak{}+0.450}) \\
\end{longtable}
}

Clock-to-Q is unchanged (0.370 vs 0.369 ns for \code{dfxtp_2}). Against
the latest accepted run\textquotesingle s worst setup slack
\textbf{+3.3286 ns (max\_ss), +10.05 ns (nom\_tt)} {[}M{]}, the fabric
absorbs this easily. Thin spots: the CPU macro (+0.174 ns, section 5)
and the 843 reset flops (+0.45 ns at ss). Hold: worst +0.2822 ns
(min\_ff) {[}M{]}; the new Q-\textgreater SCD links add the hold risk of
section 1.5. \textbf{Pins} (core macro, no pad ring {[}R{]}; decisions 5
and 7). Raw ports to add to \code{soc_top}:

{\small\setlength{\tabcolsep}{4pt}

{
\begin{longtable}[]{@{}
  >{\raggedright\arraybackslash}p{\dimexpr(\linewidth-4\tabcolsep)*182/1000\relax}
  >{\raggedright\arraybackslash}p{\dimexpr(\linewidth-4\tabcolsep)*490/1000\relax}
  >{\raggedright\arraybackslash}p{\dimexpr(\linewidth-4\tabcolsep)*308/1000\relax}@{}}
\toprule\noalign{}
\begin{minipage}[b]{\linewidth}\raggedright
Group
\end{minipage} & \begin{minipage}[b]{\linewidth}\raggedright
Ports
\end{minipage} & \begin{minipage}[b]{\linewidth}\raggedright
Count
\end{minipage} \\
\midrule\noalign{}
\endhead
\bottomrule\noalign{}
\endlastfoot
TAP &
\texttt{tck\_\allowbreak{}i\ tms\_\allowbreak{}i\ tdi\_\allowbreak{}i\ tdo\_\allowbreak{}o\ tdo\_\allowbreak{}oe\_\allowbreak{}o}
(+ optional \code{trst_ni}) & 5-6 \\
Parallel scan control & \code{scan_clk_i} (test clock), \code{scan_en_i}
& 2 \\
Parallel scan data & \code{scan_in_i[7:0]}, \code{scan_out_o[7:0]} &
16 \\
\textbf{Total} & & \textbf{23-24} (proposal {[}A{]}; N=8 chain pairs) \\
\end{longtable}
}

}

\code{scan_mode}, \code{scan_rst_n}, chain select and all MBIST control
stay internal, set through the TAP. The direct pins are \textbf{OR-ed}
with the TAP-held test controls
(\texttt{scan\_\allowbreak{}mode\ =\allowbreak{}\ scan\_\allowbreak{}mode\_\allowbreak{}i\ |\ tap\_\allowbreak{}scan},
section 15.4): the 1a pins work on their own and the TAP never gates
them (user decision 2026-10-11, Q-PIN; an earlier draft here said "ANDed
with the TAP-held \code{scan_mode}"). Protection of a floating or noisy
\code{scan_mode_i}/\code{scan_clk_i} is therefore \textbf{not} in the
RTL: it is a pad pull-down on those pins, to be added at pad-ring time
(section 9.5). Pin count is 2N+2 for N chain pairs: N=4 gives 10, N=16
gives 34 (plus the TAP). At chip top the pairs are candidates for
GPIO-pad sharing (section 9.5). The earlier 15-31 dedicated-pin estimate
and the TAP-only 5-6-pin plan are both superseded.

\textbf{Chain proposal (decision 7) {[}A, from measured flop counts{]}.}
Scanned flops: core domain 17,553, \code{cpu_clk_i} fabric domain 943
{[}M: \code{soc_excl.log}, 256 key flops excluded{]}, CPU macro 4,829
{[}M: macro netlist{]}; total 23,325. Chains are single-domain
(\code{no_mix}) and the macro chains are stitched inside the macro run
(section 5):

{
\begin{longtable}[]{@{}
  >{\raggedright\arraybackslash}p{\dimexpr(\linewidth-6\tabcolsep)*88/1000\relax}
  >{\raggedright\arraybackslash}p{\dimexpr(\linewidth-6\tabcolsep)*683/1000\relax}
  >{\raggedright\arraybackslash}p{\dimexpr(\linewidth-6\tabcolsep)*132/1000\relax}
  >{\raggedright\arraybackslash}p{\dimexpr(\linewidth-6\tabcolsep)*77/1000\relax}@{}}
\toprule\noalign{}
\begin{minipage}[b]{\linewidth}\raggedright
Chains
\end{minipage} & \begin{minipage}[b]{\linewidth}\raggedright
Domain
\end{minipage} & \begin{minipage}[b]{\linewidth}\raggedright
Flops each
\end{minipage} & \begin{minipage}[b]{\linewidth}\raggedright
Pairs
\end{minipage} \\
\midrule\noalign{}
\endhead
\bottomrule\noalign{}
\endlastfoot
5 & \code{clk_i} fabric & \textasciitilde3,511 & 5 \\
1 & \code{cpu_clk_i} fabric (incl. \code{cpu_gated_clk} flops outside
the macro) & 943 & 1 \\
2 & CPU macro & \textasciitilde2,415 & 2 \\
\end{longtable}
}

Longest chain 3,511, so the 943-flop chain is under-used. The OpenROAD
stitcher does not balance (section 1.3), so the split is made by the
plugin (it picks flop subsets per chain), which is Stage 2 work. N=16
would give a longest chain of \textasciitilde1,755 (10 core chains) at
34 raw pins.

\textbf{Test time {[}A{]}.} Shift cycles per pattern = longest chain +
1. At an assumed 1,000 patterns (no pattern count exists until Stage 3):
N=8 gives 3.51 M cycles = \textbf{0.14 s at 25 MHz, 0.35 s at 10 MHz}.
Serial TAP scan (23,325 bits) is 23.3 M cycles = 0.93 s at 25 MHz, 2.3 s
at 10 MHz, so parallel ports are \textasciitilde6.6x faster at N=8 (the
imbalance limits it below 8x). Parallel shift is clocked by
\code{scan_clk_i}, which is not limited by JTAG pad timing; the
achievable rate is unknown until a pad ring exists.

\section{Test access: IEEE 1149.1 TAP (decisions 5 and 6)}\label{tm:ch:dft-architecture:9-test-access-ieee-11491-tap-decisions-5-and-6}\origsec{9}

\subsection{Architecture}\label{tm:ch:dft-architecture:91-architecture}\origsec{9.1}

\begin{lstlisting}[language={},caption={Diagram source (Mermaid)}]
flowchart LR
  P[TCK TMS TDI TDO TRST_N] --> T[TAP FSM + IR, TCK domain]
  PP[scan_clk_i scan_en_i scan_in_i/scan_out_o x N] --> CH
  T --> D[TDR bank: IDCODE, BYPASS, SCAN_SEL, SCAN, MBIST_CTRL, MBIST_STATUS]
  D --> TC[dft_ctrl: scan_mode, scan_en, scan_rst_n, chain select]
  TC -->|glitch-free clock mux per root| CK[clk_i / cpu_clk_i roots]
  TC --> CH[fabric chains + CPU-macro chains]
  D -->|MBIST only, Stage 4 plan: 2-FF sync| MB[mbist_ctrl x2]
\end{lstlisting}

TAP, IR, TDRs and \code{dft_ctrl} are new RTL (Stage 1). \textbf{None of
this exists in the repo today} {[}R:
\texttt{grep\ -\allowbreak{}ri\ jtag\ rtl/\allowbreak{}} finds no TAP;
not exhaustively checked, {[}A{]}{]}.

\subsection{Instruction set (private codes to be fixed in the RTL
spec)}\label{tm:ch:dft-architecture:92-instruction-set-private-codes-to-be-fixed-in-the-rtl-spec}\origsec{9.2}

{\footnotesize\setlength{\tabcolsep}{3pt}

{
\begin{longtable}[]{@{}
  >{\raggedright\arraybackslash}p{\dimexpr(\linewidth-4\tabcolsep)*191/1000\relax}
  >{\raggedright\arraybackslash}p{\dimexpr(\linewidth-4\tabcolsep)*394/1000\relax}
  >{\raggedright\arraybackslash}p{\dimexpr(\linewidth-4\tabcolsep)*394/1000\relax}@{}}
\toprule\noalign{}
\begin{minipage}[b]{\linewidth}\raggedright
Instruction
\end{minipage} & \begin{minipage}[b]{\linewidth}\raggedright
DR
\end{minipage} & \begin{minipage}[b]{\linewidth}\raggedright
Purpose
\end{minipage} \\
\midrule\noalign{}
\endhead
\bottomrule\noalign{}
\endlastfoot
\code{BYPASS} (all ones) & 1 bit & mandatory \\
\code{IDCODE} (reset default) & 32 bit & mandatory in practice.
\textbf{Manufacturer ID field needs a decision} (section 12, N1) \\
\code{SAMPLE/PRELOAD}, \code{EXTEST} & none until a pad ring exists &
see 9.5 and decision N2: they cannot be implemented without a boundary
register \\
\code{SCAN_SEL} & chain-group select (fabric \code{clk_i}, fabric
\code{cpu_clk_i}, CPU macro, all) & selects which chain(s) \code{SCAN}
connects to TDI/TDO \\
\code{SCAN} & selected chain(s); a \code{PARALLEL} bit in
\code{SCAN_SEL} hands shift/\allowbreak{}capture control to
\code{scan_clk_i}/\code{scan_en_i}/\code{scan_in_i}/\code{scan_out_o}
(decision 7) & enters scan mode (\code{scan_mode}=1, internal resets
forced inactive, clock muxes to TCK). \code{Shift-DR} shifts with
\code{scan_en}=1; \code{Capture-DR} performs one capture pulse with
\code{scan_en}=0 \\
\code{MBIST_CTRL} & control word &
domain/\allowbreak{}algorithm/\allowbreak{}background select,
\code{go} \\
\code{MBIST_STATUS} & status word & 12 \code{done} + 12 \code{fail} +
controller state, read-only, captured at \code{Capture-DR} \\
\code{SCAN_RST} (optional) & 1 bit & pulse the functional reset network
under test control \\
\end{longtable}
}

}

Tester-only MBIST (decision 4) is exactly \code{MBIST_CTRL} then idle
then \code{MBIST_STATUS}; see section 6.2.

\subsection{Clocking, domains, reset}\label{tm:ch:dft-architecture:93-clocking-domains-reset}\origsec{9.3}

\begin{itemize}
\tightlist
\item
  \textbf{TCK is a third asynchronous domain}, unrelated to \code{clk_i}
  and \code{cpu_clk_i} (and these two are already asynchronous to each
  other {[}R{]}).
\item
  \textbf{Scan shift clock = TCK in serial mode, \code{scan_clk_i} in
  parallel mode} (decision 7). One
  \texttt{test\_\allowbreak{}clk\ =\allowbreak{}\ parallel\ ?\ scan\_\allowbreak{}clk\_\allowbreak{}i\ :\allowbreak{}\ TCK}
  is formed first (the select comes from the TAP and is quasi-static),
  so each domain\textquotesingle s clock root still has a single 2-input
  \textbf{glitch-free clock mux}
  (\texttt{scan\_\allowbreak{}mode\ ?\ test\_\allowbreak{}clk\ :\allowbreak{}\ functional\ clock},
  select synchronised into each source domain, enable-before-switch)
  placed at the port side, ahead of the \code{u_cpu_cg} gate and ahead
  of CTS. Consequences: (1) no crossing exists \emph{during shift},
  because every scanned flop sees the same test clock; (2) the mux is on
  the clock path, so the clock tree insertion delay and the
  TCK-to-each-tree skew change, and \textbf{chains must stay
  single-domain} (\code{no_mix}, as planned) with a TDR-level junction
  (negative-edge retiming flop or lock-up latch) wherever chains are
  concatenated {[}A{]}; (3) CTS must treat the mux inputs; this changes
  \code{cts} results and needs a re-closure of hold on both trees {[}A:
  not measured{]}.
\item
  \textbf{Capture} uses a TCK pulse as well (a slow capture): fine for
  stuck-at, \textbf{not} at-speed. See 9.6.
\item
  \textbf{Control crossings TCK -\textgreater{} \code{core_clk} /
  \code{cpu_clk}:} \textbf{Stage 4 plan, applies ONLY to the future
  MBIST control word (not built); the mechanism is a PROPOSAL, see
  "Deferred: MBIST control-word crossing" at the end of section 15.5.}
  The quasi-static MBIST bits (\code{go}, algorithm selects) were
  planned to go through \code{cdc_sync}, with the protocol "write, then
  wait \textgreater= 3 destination cycles in
  Run-Test/\allowbreak{}Idle". \textbf{As built (Stage 1b), the scan
  controls (\code{scan_mode}, \code{test_en}, \code{scan_rst},
  \code{scan_en}, \code{serial}, \code{parallel}, \code{chain_grp},
  \code{shift_clk_en}) are NOT synchronised and there is no such wait
  rule}; they cross unsynchronised by design under the quasi-static
  tester contract, with the reasons in section 15.5 crossing 2 (the
  original plan listed \code{scan_mode} for the reset overrides here;
  that was superseded by Stage 1b). \textbf{Status crossings back}
  (\code{done}, \code{fail}) are sticky and sampled at \code{Capture-DR}
  through 2-FF synchronisers (also Stage 4 plan). The
  repo\textquotesingle s existing primitives (\code{cdc_sync},
  \code{cdc_reset_sync} {[}R{]}) are reused; the SDC needs
  \texttt{set\_\allowbreak{}clock\_\allowbreak{}groups\ -\allowbreak{}asynchronous}
  for TCK and
  \texttt{set\_\allowbreak{}max\_\allowbreak{}delay\ -\allowbreak{}datapath\_\allowbreak{}only}
  on these crossings, following \code{phase5_soc_multiclock.sdc} {[}R:
  item 12 of CLAUDE.md{]}.
\item
  \textbf{TRST\_N / reset:} TAP FSM, IR and TDRs are reset by
  \code{TRST_N} (async, optional pin) \textbf{or} by five TCK cycles
  with TMS high (mandatory). The TAP must \textbf{not} be reset by
  \code{rst_n_i} or the PLL-derived resets (it must keep working while
  the core is in reset and while scan has forced the reset network).
  Entering Test-Logic-Reset clears \code{scan_mode}, \code{mbist_go} and
  the clock-mux selects. \code{TRST_N} and the TAP flops need a defined
  power-up state; they are not on any scan chain (9.4).
\item
  \textbf{\code{scan_rst_ni} becomes an internal signal} from
  \code{dft_ctrl}, not a pin: the \code{07n} hooks
  (\texttt{scanmode\_\allowbreak{}i\ ?\ scan\_\allowbreak{}rst\_\allowbreak{}ni\ :\allowbreak{}\ rst\_\allowbreak{}n\_\allowbreak{}i})
  are wired to it {[}R{]}, with the \code{cdc_reset_sync} output-side
  mux fix of section 4 item 3.
\end{itemize}

\subsection{The TAP\textquotesingle s own flops}\label{tm:ch:dft-architecture:94-the-taps-own-flops}\origsec{9.4}

The TAP FSM, IR and TDRs (\textasciitilde100-300 flops {[}A{]}) cannot
be on the chains they control. They are an additional \textbf{unscanned
set} (like the keys): tested by \code{IDCODE}/\code{BYPASS} readback and
TDR write/\allowbreak{}readback, not by ATPG. Their coverage cost
(\textasciitilde0.2-0.4 \% {[}A{]}) is in the budget of section 2.1.

\subsection{Raw ports on the core macro (decisions 6 and 7)}\label{tm:ch:dft-architecture:95-raw-ports-on-the-core-macro-decisions-6-and-7}\origsec{9.5}

The macro boundary exposes the TAP ports \code{tck_i}, \code{tms_i},
\code{tdi_i}, \code{tdo_o}, \textbf{\code{tdo_oe_o}} (a separate
output-enable, matching this tree\textquotesingle s no-\code{inout}
style {[}R{]}; the chip top will build the tristate), \code{trst_ni}
(optional), and, from decision 7, the parallel scan ports
\code{scan_clk_i}, \code{scan_en_i}, \code{scan_in_i[N-1:0]},
\code{scan_out_o[N-1:0]} (N=8 proposed, section 8). \code{clk_i},
\code{cpu_clk_i}, \code{rst_n_i}, \code{cpu_rst_n_i} already exist
{[}R{]}. \code{scan_mode}, \code{scan_rst_n} and the \code{mbist_*}
signals stay internal. The CPU macro boundary (403 pins today {[}R{]})
gains \code{scan_in[2]/scan_out[2]}, \code{scan_en}, \code{scan_mode},
MBIST \code{go/done/fail[10]} and the shadow/\allowbreak{}observation
ports, all through \code{pin_order.cfg} {[}R{]}.

\textbf{What a future chip-top epic will need from DFT (not done now):}
pad pull-downs on \code{scan_mode_i} and \code{scan_clk_i} and a pull-up
on \code{scan_rst_ni} (their inactive levels) (the RTL OR-s the direct
pins with the TAP and does not protect a floating pin, section 15.6 item
6); boundary-scan register and cells on every pad, real
\code{EXTEST}/\code{SAMPLE}/\code{PRELOAD}/\code{CLAMP}/\code{HIGHZ}
(making the TAP 1149.1-compliant); a registered JEDEC manufacturer ID in
\code{IDCODE}; the TDO tristate and pad enables; the pad-level pin
budget and whether the 18 parallel scan pins (\code{scan_clk_i},
\code{scan_en_i}, 8 pairs) can share GPIO pads; pad-ring test (the pad
cells and ESD are outside scan); package pinout and ATE pin map;
TCK/\allowbreak{}TMS/\allowbreak{}TDI input timing and TDO output timing
constraints; the SDC \code{set_input_delay}/\code{set_output_delay} for
the JTAG pins; and a probe/\allowbreak{}bond test plan including
\code{TRST_N} handling.

\subsection{Throughput and at-speed with parallel ports (decision 7)
{[}A unless noted{]}}\label{tm:ch:dft-architecture:96-throughput-and-at-speed-with-parallel-ports-decision-7-a-unless-noted}\origsec{9.6}

\begin{itemize}
\tightlist
\item
  \textbf{Throughput.} Serial TAP scan is \textasciitilde23.3 k bits per
  pattern {[}M counts{]}; parallel ports at N=8 cut the longest chain to
  \textasciitilde3.5 k (section 8). Serial TAP scan stays available for
  bring-up and as the path that needs only 5-6 pins.
\item
  \textbf{What parallel ports buy.} (a) Shift throughput, as above. (b)
  A shift clock decoupled from TCK. (c) A tester-controlled
  \code{scan_en_i} and test clock, so the \textbf{capture sequence is
  under ATE control}: shift with \code{scan_en_i} high, drop
  \code{scan_en_i} and let it settle, then apply the capture pulses on
  the test clock (or, with a select, on the functional clock pins; the
  PLL in this build is a passthrough, \code{pll_clkgen_stub.sv:69}
  {[}R{]}, so the functional clocks are ATE-supplied pins anyway). A
  two-pulse launch-on-capture burst at core-clock rate is therefore
  \textbf{physically possible}, with \code{scan_en_i} held static low
  across both pulses; no on-chip clock controller and no pipelined
  \code{scan_en} are needed for launch-on-capture. Launch-on-shift (a
  \code{scan_en_i} transition at speed) is \textbf{not} supported: an
  external pin and its buffer tree cannot switch inside one core cycle.
\item
  \textbf{What they do not buy with stuck-at-only open-source ATPG.} (1)
  The tools generate \textbf{single-frame stuck-at patterns}. Applied
  with two pulses their expected responses are wrong, so each pattern
  must be re-simulated through two frames on the gate netlist
  (cocotb/\allowbreak{}Verilator) to obtain expected data; that gives
  pass/fail data, not a coverage figure. (2) No transition-fault list,
  fault simulation or targeting exists, so \textbf{delay-fault coverage
  remains unmeasured and unclaimed}; at-speed application of stuck-at
  patterns detects delay defects only incidentally. (3) Nothing masks
  the \code{clk_i}/\code{cpu_clk_i} asynchronous crossings,
  \code{en_latch} or false/\allowbreak{}multicycle paths in a two-pulse
  capture; the capture protocol per domain is hand-written. (4) The root
  clock mux must pass an ATE burst glitch-free and its insertion delay
  changes the capture timing seen by each domain (section 9.3). What
  delay coverage would need is in section 2.4 and open item N10.
\item
  \textbf{TCK ceiling.} TCK and \code{scan_clk_i} rates are limited by
  pad timing, not logic; with no pad ring the real limit is not yet
  known.
\item
  \textbf{Clock-mux cost.} The mux in each clock root adds insertion
  delay on the functional clock path in \emph{every} mode and perturbs
  CTS and the macro/SoC clock skew; the CPU macro\textquotesingle s
  +0.174 ns margin (section 5) makes this the most likely Stage 2
  blocker.
\end{itemize}

\section{Proposed architecture and staged plan (restaged for
decision 5)}\label{tm:ch:dft-architecture:10-proposed-architecture-and-staged-plan-restaged-for-decision-5}\origsec{10}

\begin{itemize}
\tightlist
\item
  Test-mode controls are \textbf{internal}, from \code{dft_ctrl} inside
  the TAP: \code{scan_mode}, \code{scan_en}, \code{scan_rst_n};
  \texttt{test\_\allowbreak{}en\ =\allowbreak{}\ scan\_\allowbreak{}mode\ |\ mbist\_\allowbreak{}en}
  OR-ed into every clock-gate enable.
\item
  Chains (decision 7): 8 pairs on parallel ports: 5 fabric \code{clk_i},
  1 fabric \code{cpu_clk_i}, 2 CPU macro, all \code{no_mix}; the same
  chains are also reachable serially through the TAP; key flops and TAP
  flops unscanned; SRAM macros shadowed (section 8).
\item
  MBIST: two controllers, collars on the functional side of the
  \code{ifdef} arms, TAP-only start/\allowbreak{}status.
\item
  Raw ports: TCK, TMS, TDI, TDO, TDO\_OE, TRST\_N, \code{scan_clk_i},
  \code{scan_en_i}, 8 x \code{scan_in_i}/\code{scan_out_o} (section
  9.5).
\end{itemize}

{\footnotesize\setlength{\tabcolsep}{3pt}

{
\begin{longtable}[]{@{}
  >{\raggedright\arraybackslash}p{\dimexpr(\linewidth-4\tabcolsep)*327/1000\relax}
  >{\raggedright\arraybackslash}p{\dimexpr(\linewidth-4\tabcolsep)*327/1000\relax}
  >{\raggedright\arraybackslash}p{\dimexpr(\linewidth-4\tabcolsep)*327/1000\relax}@{}}
\toprule\noalign{}
\begin{minipage}[b]{\linewidth}\raggedright
Stage
\end{minipage} & \begin{minipage}[b]{\linewidth}\raggedright
Content
\end{minipage} & \begin{minipage}[b]{\linewidth}\raggedright
Concrete changes
\end{minipage} \\
\midrule\noalign{}
\endhead
\bottomrule\noalign{}
\endlastfoot
\textbf{1 TAP + parallel ports + scan-ready RTL} (\textbf{approved to
start}, \prref{250} merged; FPGA Stage 0 \ghref{245} follows it) & no
behaviour change with the TAP idle and the scan pins inert & (a) TAP
FSM, IR, TDR bank, \code{dft_ctrl}, \code{test_clk} select,
parallel-port gating by \code{scan_mode}, glitch-free clock muxes, TCK
crossings, standalone cocotb suite (TAP state walk, IDCODE, BYPASS,
reset). (b) Items 1-9, 11-13, 20 of section 4: un-tie the \code{07n}
hooks to \code{dft_ctrl}; \code{cdc_reset_sync} output-side mux; scan
hook in \code{cdc_sync}; \code{test_en} on \code{rv32i_clock_gate} (+
\code{.SE}); internal-reset override mux; scan/\allowbreak{}MBIST ports
on \code{rv32i_cpu_top} + \code{pin_order.cfg};
\texttt{create\_\allowbreak{}clock\ cpu\_\allowbreak{}clk\_\allowbreak{}i}
and TCK clocks in \code{sky130_soc.sdc}; key/\code{rk_q} scan-mode masks
(section 7). Gate: \code{soc_all} unchanged with TAP idle; Verilator
lint; netlist X/loop check; \code{check_scan_exclusions.py}. \\
\textbf{2 Insertion} & plugin step (option N, section 1.5), CPU macro
first; \textbf{the CPU-macro re-harden waits on \code{dud4} (decision
9)} & Plugin \code{Flow} subclass; \code{set_dont_touch} list derived
from net names; CPU macro re-harden with chains; SoC re-harden; compare
LVS/\allowbreak{}DRC/\allowbreak{}setup/\allowbreak{}hold/\allowbreak{}antenna/\code{e45j}
against \code{RUN_2026-10-09_06-18-39}; SET\_B workaround (use 26Q2 on
the netlist); scan/TCK SDC. \\
\textbf{3 ATPG} & time-boxed tool spike first (decision 8); sign-off
number fixed only after it & Spike {[}A, proposed box: about 5 working
days, 12 GB RAM cap, 24 h per tool run{]}: Quaigh on the fabric-sized
full-scan cut (\textasciitilde18.5 k flops) via a ladder \code{npu}
(\textasciitilde1.2 k flops), \code{crypto} (\textasciitilde3 k),
\code{dma} (\textasciitilde9 k), fabric, recording completion, wall
time, peak RAM, coverage. Then: add the tools to a new \src{flake.nix}
devshell; Quaigh on \code{dma}/\code{npu}-sized cuts, then the fabric;
Sky130-cell glue; gate-level replay of patterns through the TAP incl. a
chain flush; SVF/STIL export (2.3). \\
\textbf{4 MBIST} & via \code{rtl-design-orchestrator}; needs Stage 1 TAP
& \code{mbist_ctrl}, \code{mem_collar} x12, SRAM
shadow/\allowbreak{}observe logic, cocotb with injected
stuck-at/\allowbreak{}coupling/\allowbreak{}decoder faults on the sim
models (a BIST that always passes is the failure to rule out), then
re-harden. \\
\textbf{5 Chip-top hand-off} & notes only now & Boundary scan, IDCODE
ID, TDO tristate, pad-level test; future epic (decision 6). \\
\end{longtable}
}

}

\section{Go / no-go}\label{tm:ch:dft-architecture:11-go-no-go}\origsec{11}

{\footnotesize\setlength{\tabcolsep}{3pt}

{
\begin{longtable}[]{@{}
  >{\raggedright\arraybackslash}p{\dimexpr(\linewidth-4\tabcolsep)*66/1000\relax}
  >{\raggedright\arraybackslash}p{\dimexpr(\linewidth-4\tabcolsep)*457/1000\relax}
  >{\raggedright\arraybackslash}p{\dimexpr(\linewidth-4\tabcolsep)*457/1000\relax}@{}}
\toprule\noalign{}
\begin{minipage}[b]{\linewidth}\raggedright
Stage
\end{minipage} & \begin{minipage}[b]{\linewidth}\raggedright
Call
\end{minipage} & \begin{minipage}[b]{\linewidth}\raggedright
Reason
\end{minipage} \\
\midrule\noalign{}
\endhead
\bottomrule\noalign{}
\endlastfoot
1 & \textbf{GO, approved to start} & RTL-only, reversible,
regression-checkable with the TAP idle; every item read from source.
Stage 1 now includes the TAP, so it is larger than the
issue\textquotesingle s original estimate. \\
2 & \textbf{CONDITIONAL GO} (CPU-macro half also waits on \code{dud4}) &
Tool works at SoC scale in 5 s {[}M{]} and honours \code{dont_touch}
{[}M{]}; Sky130 build has defects (SET\_B, \code{report_dft_config}), no
lock-up/\allowbreak{}clock-gate/\allowbreak{}scan-enable handling; the
CPU macro setup margin is +0.174 ns and now also pays for a clock
mux. \\
3 & \textbf{CONDITIONAL}: GO for a time-boxed spike; NO-GO for a
promised number (99 \% / 98 \% is a goal, decision 8) & Quaigh works on
a 2.9 k-gate cut {[}M{]}; SoC scale unproven; Fault unbuilt;
Atalanta/\allowbreak{}PODEM non-commercial. Fallback is itemised
per-block coverage. \textbf{Spike run 2026-10-28 (section 18): Quaigh
completes npu, crypto and dma (up to 9,088 flops, 11.4 GiB peak) and is
OOM-killed at the 12 GiB cap on the real fabric netlist (412,472 faults,
29,376 patterns, 679 s). The spike did NOT demonstrate the SoC-level
goal; the per-block fallback applies unless the owner chooses a
commercial or patched engine (18.8, not adopted).} \textbf{Update
2026-10-29/\allowbreak{}30 (beads \code{j41m.18}, \code{j41m.23}): the
patched Quaigh completes the fabric (99.95 \% TC / 99.16 \% FC, 18.9).
The 99 / 98 target is adopted provisionally (2.1); the per-block
fallback is dropped.} \\
4 MBIST & \textbf{GO} & independent of ATPG; memory is 88 \% of placed
area; cheap; known SRAM-pin timing risk. Depends on the Stage 1 TAP. \\
5 & \textbf{N/A now} & pad ring out of scope. \\
\end{longtable}
}

}

\section{Open items}\label{tm:ch:dft-architecture:12-open-items}\origsec{12}

IDs are kept stable because other sections cite them. Resolved items
move to the second table.

\textbf{Still open} (each needs an answer before the stage that uses
it):

{\footnotesize\setlength{\tabcolsep}{3pt}

{
\begin{longtable}[]{@{}
  >{\raggedright\arraybackslash}p{\dimexpr(\linewidth-6\tabcolsep)*48/1000\relax}
  >{\raggedright\arraybackslash}p{\dimexpr(\linewidth-6\tabcolsep)*311/1000\relax}
  >{\raggedright\arraybackslash}p{\dimexpr(\linewidth-6\tabcolsep)*311/1000\relax}
  >{\raggedright\arraybackslash}p{\dimexpr(\linewidth-6\tabcolsep)*311/1000\relax}@{}}
\toprule\noalign{}
\begin{minipage}[b]{\linewidth}\raggedright
\#
\end{minipage} & \begin{minipage}[b]{\linewidth}\raggedright
Question
\end{minipage} & \begin{minipage}[b]{\linewidth}\raggedright
Options
\end{minipage} & \begin{minipage}[b]{\linewidth}\raggedright
Recommendation
\end{minipage} \\
\midrule\noalign{}
\endhead
\bottomrule\noalign{}
\endlastfoot
N1 & \code{IDCODE} manufacturer/\allowbreak{}part values & placeholder
(documented non-JEDEC) / registered ID & placeholder until the chip-top
epic \\
N2 & 1149.1 compliance level: mandatory
\code{SAMPLE/PRELOAD}/\code{EXTEST} need a boundary register &
compliant-minus-boundary (documented) / boundary register on macro ports
now & compliant-minus-boundary; revisit at chip top \\
N4 & TAP flops unscanned (\textasciitilde0.2-0.4 \%) & accept with
IDCODE/\allowbreak{}TDR readback / develop a separate TAP scan mode &
accept \\
N7 & CPU macro target period (13.333 ns, +0.174 ns margin) with a clock
mux and scan muxes & keep / relax & relax if Stage 2 shows negative
slack; do not hide it \\
N8 & Atalanta use if the tape-out turns commercial (2.2) & drop /
license & keep the flow licence-clean (Quaigh/\allowbreak{}Fault
engine) \\
N10 & \textbf{Route to delay-fault coverage} (new; follows from decision
7 rejecting zero delay coverage, section 2.4) & commercial ATPG licence
or pattern-generation service for the tape-out / build transition ATPG
on Quaigh/\allowbreak{}Fault / accept incidental detection only & scope
as its own bead after the Stage 3 spike; a commercial service for
pattern generation is the only option that meets the goal at tape-out
scale \\
N11 & Parallel chain count N (section 8: 8 pairs, 23-24 raw pins) & 4
(10+TAP pins) / 8 / 16 (34+TAP) & 8; revisit when the chip-top pad
budget is known \\
\end{longtable}
}

}

\textbf{Resolved} (2026-10-10):

{\footnotesize\setlength{\tabcolsep}{3pt}

{
\begin{longtable}[]{@{}
  >{\raggedright\arraybackslash}p{\dimexpr(\linewidth-4\tabcolsep)*49/1000\relax}
  >{\raggedright\arraybackslash}p{\dimexpr(\linewidth-4\tabcolsep)*454/1000\relax}
  >{\raggedright\arraybackslash}p{\dimexpr(\linewidth-4\tabcolsep)*477/1000\relax}@{}}
\toprule\noalign{}
\begin{minipage}[b]{\linewidth}\raggedright
\#
\end{minipage} & \begin{minipage}[b]{\linewidth}\raggedright
Was
\end{minipage} & \begin{minipage}[b]{\linewidth}\raggedright
Resolution
\end{minipage} \\
\midrule\noalign{}
\endhead
\bottomrule\noalign{}
\endlastfoot
N3 & Serial TAP scan vs parallel scan ports & \textbf{Decision 7:} TAP
for control and MBIST plus parallel scan ports. Sections 8, 9.5, 9.6. \\
N5 & Confirm coverage target 99.0 \% test / 98.0 \% fault &
\textbf{Decision 8:} a goal, not a commitment; fixed after the Stage 3
tool spike. Section 2.1. \textbf{Resolved 2026-10-30 (owner,
\code{j41m.23}): adopted provisionally as the target; fabric and CPU
macro are reported separately; sign-off after the cell-pin re-grade
\code{j41m.20}.} \\
N12 & Spike box for Stage 3 (section 10) & \textbf{Resolved 2026-10-28
(user): box confirmed as proposed}: about 5 working days, 12 GB RAM and
24 h wall per run, Quaigh only (Fault only as a fallback), stop at the
first rung that breaks a limit. The box was used; the fabric rung broke
the 12 GiB limit. Record in section 18. \\
N6 & \code{ma7}: Sky130 CPU macro built with
\texttt{USE\_\allowbreak{}SYNLIG:\allowbreak{}\ true} & \textbf{Decision
9:} bead \bead{dud4}, test first (no PD run), then the user decides on
re-hardening. Section 4. \\
N9 & Delay-fault coverage: accept none from scan & \textbf{Not accepted}
(decision 7). Replaced by N10. \\
\end{longtable}
}

}

\section{Not done or not verified}\label{tm:ch:dft-architecture:13-not-done-or-not-verified}\origsec{13}

\begin{itemize}
\tightlist
\item
  Fault was not built. Quaigh now ran the \code{npu}, \code{crypto} and
  \code{dma} rungs (section 18) but was OOM-killed at 12 GiB on the real
  fabric netlist: \textbf{still no SoC-scale deterministic ATPG number
  exists} (only a 69.59 \% random-pattern floor). A \code{dft} devshell
  was added to \src{flake.nix} (yosys, sv2v, python; Quaigh stays an
  external pinned binary). Section 18.7 lists what the spike did not
  cover (cell-pin faults, pattern-to-chain mapping and SVF replay, SRAM
  shadowing, delay faults, the CPU macro as a rung, the whole-SoC fault
  count including the excluded cones).
\item
  No functional-mode equivalence or shift/\allowbreak{}flush simulation
  of any scan netlist (structural chain walk only). Scan netlists were
  not run through OpenSTA with a real SDC.
\item
  DEF \code{SCANCHAINS}, lock-up latch behaviour beyond the toy, and
  26Q2 \code{-max_length} splitting were not tested. The
  \code{dont_touch} experiment used \code{edf00dff} only.
\item
  Whether \code{sdf*} cells are in the PDK \code{DONT_USE_CELLS} list
  was not verified.
\item
  Hold-buffer impact of chain links, the clock-mux CTS effect, and the
  collar muxes on SRAM pins are estimates.
\item
  The crypto cone numbers are a pin-count proxy on one netlist
  (\code{RUN_2026-10-06_22-43-34}), not an ATPG result; the \code{key_q}
  = \code{key_i} identification is by name and count.
\item
  SVF/\allowbreak{}OpenOCD playback and the TAP RTL itself are not
  built; the TAP section is a design, not a measurement.
\item
  \code{hdl-kgraph} was not used; the RTL audit is grep and file
  reading, so a construct outside the greps could be missed.
\item
  The audit used RTL from a sibling worktree for the synthesis counts
  (same \code{main} at the time; small differences possible). Flop-count
  attribution by group is approximate; totals are exact.
\item
  Follow-up (decisions 7-9): the chain split, pin count and test times
  in section 8 are proposals from measured flop counts, not an insertion
  result; the parallel-port gating by \code{scan_mode} and the two-pulse
  launch-on-capture feasibility are design reasoning, not simulated; the
  \code{dud4} differential has not been run; the spike box is a
  proposal.
\end{itemize}

\section{\texorpdfstring{Stage 1a implementation record (bead
\code{j41m.2}): scan-ready fabric
RTL}{Stage 1a implementation record (bead ): scan-ready fabric RTL}}\label{tm:ch:dft-architecture:14-stage-1a-implementation-record-bead-j41m2-scan-ready-fabric-rtl}\origsec{14}

Added after the 2026-10-10 decisions (parallel scan ports in addition to
the TAP; TAP is Stage 1b, bead \bead{j41m.3}). Scope is the
\textbf{fabric only}: the CPU hard macro boundary is untouched (section
14.6). Labels as above: \textbf{{[}M{]}} measured, \textbf{{[}R{]}}
read, \textbf{{[}A{]}} assumed.

\subsection{\texorpdfstring{Ports added to \code{soc_top} and the
internal seam}{Ports added to  and the internal seam}}\label{tm:ch:dft-architecture:141-ports-added-to-soc-top-and-the-internal-seam}\origsec{14.1}

{\small\setlength{\tabcolsep}{4pt}

{
\begin{longtable}[]{@{}
  >{\raggedright\arraybackslash}p{\dimexpr(\linewidth-8\tabcolsep)*135/1000\relax}
  >{\raggedright\arraybackslash}p{\dimexpr(\linewidth-8\tabcolsep)*51/1000\relax}
  >{\raggedright\arraybackslash}p{\dimexpr(\linewidth-8\tabcolsep)*135/1000\relax}
  >{\raggedright\arraybackslash}p{\dimexpr(\linewidth-8\tabcolsep)*135/1000\relax}
  >{\raggedright\arraybackslash}p{\dimexpr(\linewidth-8\tabcolsep)*524/1000\relax}@{}}
\toprule\noalign{}
\begin{minipage}[b]{\linewidth}\raggedright
Port
\end{minipage} & \begin{minipage}[b]{\linewidth}\raggedright
Dir
\end{minipage} & \begin{minipage}[b]{\linewidth}\raggedright
Width
\end{minipage} & \begin{minipage}[b]{\linewidth}\raggedright
Inactive value
\end{minipage} & \begin{minipage}[b]{\linewidth}\raggedright
Meaning
\end{minipage} \\
\midrule\noalign{}
\endhead
\bottomrule\noalign{}
\endlastfoot
\code{scan_mode_i} & in & 1 & 0 & test mode; quasi-static (set before
the first test clock edge, not toggled while a clock runs) \\
\code{scan_en_i} & in & 1 & 0 & 1 = shift, 0 = capture; no RTL consumer
until Stage 2 inserts scan flops \\
\code{scan_rst_ni} & in & 1 & 1 & scan reset, active low; delivered to
every async reset in scan mode \\
\code{scan_clk_i} & in & 1 & 0 (any) & the single shared
shift/\allowbreak{}capture clock for all scanned domains \\
\code{scan_in_i} & in & \code{SCAN_CHAINS} & 0 & chain inputs,
\textbf{placeholder} until Stage 2 \\
\code{scan_out_o} & out & \code{SCAN_CHAINS} & 0 & chain outputs,
\textbf{tied 0} until Stage 2 \\
\end{longtable}
}

}

\code{SCAN_CHAINS} defaults to \textbf{8}, the figure fixed by the
chain-count decision (\prref{253}): 5 \code{clk_i} fabric chains, 1
\code{cpu_clk_i} fabric chain, 2 CPU-macro chains. Indices \code{[7:6]}
are reserved for the macro; they are wired through to
\code{rv32i_cpu_top} since 14.10 (\code{j41m.4}) and the macro ties its
own \code{scan_out_o} 0 until Stage 2. My own count (17.5 k
\code{clk_i}-domain scannable flops {[}M, section 4{]}, 0.9 k
\code{cpu_clk_i} outside the macro) does not contradict it: 5 chains of
\textasciitilde3.5 k is within 40 \% of 7 of \textasciitilde2.5 k, and
the shift time is set by the longest chain, so the pin saving of two
pairs costs \textasciitilde1.4x shift length.

All consumers take their control from one instance of
\src{rtl/soc/dft/dft_ctrl_ports.sv} (\code{dft_scan_mode},
\code{dft_scan_rst_n}, \code{dft_test_en}, \code{dft_test_clk},
\code{dft_scan_en}). In 1a it is a pass-through of the ports;
\textbf{Stage 1b replaces that one module} with the TAP-driven
\code{dft_ctrl} of identical port list, so nothing in \code{soc_top} is
re-plumbed.
\texttt{test\_\allowbreak{}en\ =\allowbreak{}\ scan\_\allowbreak{}mode}
(OR \code{mbist_en} in Stage 4).

\subsection{What each scope item became}\label{tm:ch:dft-architecture:142-what-each-scope-item-became}\origsec{14.2}

{\footnotesize\setlength{\tabcolsep}{3pt}

{
\begin{longtable}[]{@{}
  >{\raggedright\arraybackslash}p{\dimexpr(\linewidth-4\tabcolsep)*221/1000\relax}
  >{\raggedright\arraybackslash}p{\dimexpr(\linewidth-4\tabcolsep)*457/1000\relax}
  >{\raggedright\arraybackslash}p{\dimexpr(\linewidth-4\tabcolsep)*302/1000\relax}@{}}
\toprule\noalign{}
\begin{minipage}[b]{\linewidth}\raggedright
Item
\end{minipage} & \begin{minipage}[b]{\linewidth}\raggedright
Change
\end{minipage} & \begin{minipage}[b]{\linewidth}\raggedright
Where
\end{minipage} \\
\midrule\noalign{}
\endhead
\bottomrule\noalign{}
\endlastfoot
Test clock & \code{dft_clk_mux} on each PLL reference
(\code{ref_clk_w}), so lock counter, PLL registers and the
stub\textquotesingle s passthrough \code{core_clk} run on
\code{scan_clk_i} in scan mode; a non-stub PLL\textquotesingle s output
gets its own mux (\code{g_core_clk_pll}). The APB CDC
bridges\textquotesingle{} destination faces use \code{ref_clk_o}, not
the raw port & \code{pll_subsystem.sv}, \code{soc_top.sv} \\
Gate test enable & \code{rv32i_clock_gate.test_en}: latch input
\texttt{en\ |\ test\_\allowbreak{}en}; \code{.SE(test_en)} on the
ICG-cell arm. \code{u_cpu_cg}, \code{u_gpu_cg} get \code{dft_test_en}.
icache/\allowbreak{}dcache (x2 each) and \code{npu_weight_mem} tie
\texttt{1'b0}: their gates clock \textbf{SRAM macros only}, no scannable
flop behind them & \code{rv32i_clock_gate.sv}, callers \\
Reset sync & \code{cdc_reset_sync}: output mux
\texttt{rst\_\allowbreak{}n\_\allowbreak{}o\ =\allowbreak{}\ scanmode\_\allowbreak{}i\ ?\ scan\_\allowbreak{}rst\_\allowbreak{}ni\ :\allowbreak{}\ sync\_\allowbreak{}q[last]}
\textbf{plus} the existing input mux (the chain\textquotesingle s own
clear must also stay quiet during shift). Scan-mode release is not
re-timed; the tester owns it & \code{cdc_reset_sync.sv} \\
Internal resets & \code{pll_rst_n}, \code{core_rst_n}, the PLL
register-file reset and \code{gpu_domain_rst_n} get a
\code{dft_rst_mux}. \code{cpu_domain_rst_n} needs none: its three terms
all equal \code{scan_rst_ni} in scan mode. \code{async_axi_fifo} and
\code{apb_cdc_bridge} resets come out of \code{cdc_reset_sync}.
\code{cdc_sync} was \textbf{not} given a hook (section 4 item 4 proposed
one): every one of its instances is reset from a net that is already
scan-controlled at its source, which \code{check_scan_clk_rst.py} proves
on the netlist & \code{soc_top.sv}, \code{pll_subsystem.sv} \\
Crypto &
\texttt{key\_\allowbreak{}eff\_\allowbreak{}w\ =\allowbreak{}\ key\_\allowbreak{}q\ \&\ \textasciitilde{}scan\_\allowbreak{}mode\_\allowbreak{}i}
(128 AND, the AES core\textquotesingle s \code{key_i}) and
\texttt{rk\_\allowbreak{}eff\_\allowbreak{}w\ =\allowbreak{}\ rk\_\allowbreak{}q\ \&\ \textasciitilde{}scan\_\allowbreak{}mode\_\allowbreak{}i}
(128 AND, the \textbf{only} reader of \code{rk_q}). 256 gates {[}A:
area{]}. Nothing else is masked & \code{crypto_accel.sv},
\code{aes128_core.sv} \\
Latch & the only latch in the build is the clock-gate latch;
\code{check_scan_clk_rst.py} requires its D cone to contain
\code{scan_mode} & - \\
\code{cpu_clk_i} \code{create_clock} & \textbf{NOT done} (14.5) & - \\
\end{longtable}
}

}

\subsection{Findings the mask does not close}\label{tm:ch:dft-architecture:143-findings-the-mask-does-not-close}\origsec{14.3}

\begin{enumerate}
\def\labelenumi{\arabic{enumi}.}
\tightlist
\item
  \textbf{Residual state at scan entry. CLOSED in hardware by bead
  \bead{j41m.6} (section 14.10), with one stated limit.} The mask stops
  the \emph{capture-path} leak. \code{s_q} is scanned and holds
  \texttt{block\ \textasciicircum{}\ key} (reloaded every idle cycle)
  and round states from a prior functional operation, and the SHA
  working state and the DOUT/\allowbreak{}DIGEST/\allowbreak{}IV words
  are scanned too; shifting out before any test-clock capture reads
  them. The original answer was a protocol requirement on the TAP/ATE
  (reset the block first). The user decision of
  2026-10-15/\allowbreak{}17 replaced it with a \textbf{hardware zeroize
  on every \code{scan_mode} edge}: \code{crypto_accel} clears its key,
  round-key, AES state, SHA state, message shadow and result/IV words in
  one cycle, on entry and on exit. Limit: the clear is functional D
  logic that a shift edge overrides, so it protects against software and
  a cooperating tester, not an adversarial one (14.10 item 5).
  \textbf{Closed for scan\_en by \code{j41m.17} (14.14): the shift
  enable is now withheld in hardware until the wipe is applied.}
\item
  The unmasked \code{msg_q}/IV/\code{DOUT} are scanned and are not
  secret by themselves.
\item
  A reset is not enough for \emph{DOUT} if a prior ciphertext is
  sensitive; it is ciphertext, which is already an output.
\end{enumerate}

\subsection{Scan port contract for Stage 2 insertion {[}A: the
OpenROAD behaviour is untested{]}}\label{tm:ch:dft-architecture:144-scan-port-contract-for-stage-2-insertion-a-the-openroad-behaviour-is-untested}\origsec{14.4}

\code{scan_in_i}/\code{scan_out_o}/\code{scan_en_i} exist now;
\code{scan_out_o} is one
\texttt{assign\ scan\_\allowbreak{}out\_\allowbreak{}o\ =\allowbreak{}\ '0;\ /\allowbreak{}/\allowbreak{}\ DFT\_\allowbreak{}PLACEHOLDER}.
Insertion (plugin option N, section 1.5) must (a) delete that tie from
the netlist (a driver conflict otherwise), (b) connect chain
\emph{k}\textquotesingle s last flop to \code{scan_out_o[k]} and
\code{scan_in_i[k]} to its first, (c) connect every scan
flop\textquotesingle s \code{SCE} to \code{scan_en_i}. Whether 26Q2
\code{insert_dft} reuses a pre-existing port by
\code{-scan_in_name_pattern}/\code{-scan_out_name_pattern} or always
creates \code{scan_in_N} is \textbf{not verified}; the fallback is to
let it create its ports and drop ours in the same plugin step. Lint
stays clean because the unread inputs are consumed into a waived sink
and the output is driven.

\subsection{\texorpdfstring{Clock/\allowbreak{}reset paths the new
muxes sit on, and
SDC}{Clock/reset paths the new muxes sit on, and SDC}}\label{tm:ch:dft-architecture:145-clockreset-paths-the-new-muxes-sit-on-and-sdc}\origsec{14.5}

\begin{itemize}
\tightlist
\item
  \textbf{Clock path} (matters for CTS): one \code{dft_clk_mux} between
  each of \code{clk_i} and \code{cpu_clk_i} and its whole tree (two
  muxes, every flop behind them); in the stub build
  \code{pll_clkgen_stub}\textquotesingle s passthrough remains. ASAP7
  history records that this passthrough already costs a real buffer
  under \code{deferred_flatten}; the mux is one more cell in front of
  it. \textbf{No PD run was made}, so CTS
  skew/\allowbreak{}insertion-delay change is unmeasured.
\item
  \textbf{Reset paths} (recovery/\allowbreak{}removal):
  \code{dft_rst_mux} adds one mux delay in front of \code{pll_rst_n},
  \code{core_rst_n}, \code{gpu_domain_rst_n} and each
  \code{cdc_reset_sync} output. \code{cdc_reset_sync} now has two muxes
  on its async path (input into the chain clear, output to the
  destination).
\item
  \textbf{SDC} (functional mode only): \code{set_case_analysis} pins
  \code{scan_mode_i}=0, \code{scan_en_i}=0, \code{scan_rst_ni}=1 (so
  \code{scan_clk_i} is blocked from the clock network) and
  \code{set_false_path} removes the test ports from timing, in
  \code{sky130_soc.sdc} and the three ASAP7 SoC SDCs. A
  \textbf{scan-mode SDC is a Stage 2 deliverable}.
  \texttt{create\_\allowbreak{}clock\ cpu\_\allowbreak{}clk\_\allowbreak{}i}
  was \emph{not} added to \code{sky130_soc.sdc}: it newly constrains
  \textasciitilde5.7 k CPU-domain flops and needs the async clock groups
  and CDC budgets of \code{phase5_soc_multiclock.sdc}, which changes
  sign-off timing and cannot be validated without a PD run. Filed as a
  bead.
\end{itemize}

\subsection{\texorpdfstring{CPU macro boundary: deliberately not
changed (SUPERSEDED for the Sky130 RTL by 14.10, bead
\code{j41m.4})}{CPU macro boundary: deliberately not changed (SUPERSEDED for the Sky130 RTL by 14.10, bead )}}\label{tm:ch:dft-architecture:146-cpu-macro-boundary-deliberately-not-changed-superseded-for-the-sky130-rtl-by-1410-bead-j41m4}\origsec{14.6}

\code{rv32i_cpu_top} has no scan pins, and adding them invalidates the
committed LEF/\allowbreak{}Liberty/\allowbreak{}blackbox views and
\code{pin_order.cfg} on both PDKs; the ASAP7 macro cannot be re-hardened
on this host (\code{ma7}/\code{lxv}/\code{2kn}), and \code{dud4} is
currently investigating gate-versus-RTL differences in the Sky130 CPU
macro netlist. So: the macro instantiates \code{rv32i_clock_gate} only
in its ASAP7 cache arms, with \code{test_en} tied \texttt{1'b0} (no
netlist change after constant propagation; the ASAP7 \code{ICG}
\code{.SE} pin was already tied 0), the Sky130 macro is bit-for-bit
unaffected by this change, and the macro\textquotesingle s 4,829 flops,
its internal gates and resets are the named follow-up.

\subsection{Checks added}\label{tm:ch:dft-architecture:147-checks-added}\origsec{14.7}

\begin{itemize}
\tightlist
\item
  \src{tools/dft/check_scan_clk_rst.py}: on the flattened \code{proc}-ed
  netlist, every async reset/set cone ends at a pin, a constant or a
  scan-mode mux (never a flop output), and every clock cone passes a
  scan-mode mux or ends at a gate latch whose enable contains
  \code{scan_mode}. On \code{soc_top_sv2v.v} (Sky130): \textbf{0
  violations, 1,640 sequential cells, 240 async pins, 28 scan-mode
  muxes}; negative controls (clock mux -\textgreater{} wire, reset mux
  -\textgreater{} wire, \code{cdc_reset_sync} output mux -\textgreater{}
  wire) fail with 556, 103 and 117 violations.
  \texttt{make\ -\allowbreak{}C\ pnr\ dft-\allowbreak{}scan-\allowbreak{}check}.
  Re-measured 2026-10-10 on the netlist regenerated from the final
  merged tree (yosys 0.62): unchanged (1,640 / 240 / 28, 0 violations).
  Per-instance negative controls (one instance bypassed):
  \code{u_core_rm} 99 violations, \code{u_pll_rm} 4, \code{u_ref_cm}
  556, \textbf{\code{u_gpu_domain_rm} 0 (not detected)}: the GPU is a
  stub in both PDK builds, so no flop sits behind that mux and the
  netlist check cannot see it (see 14.9, mutant g). The checker was
  hardened the same day: it crashed (\code{KeyError}) on the ASAP7
  netlist and would have passed vacuously on any cell with no
  definition; it now fails on those and takes \code{--models}
  (\src{tools/dft/asap7_cell_models.v}: \code{ICGx1_ASAP7_75t_R},
  \code{sram_1rw_256x32_asap7}). ASAP7 netlist: 0 violations, 3,726
  sequential cells, 240 async pins, 29 muxes
  (\texttt{make\ -\allowbreak{}C\ pnr\ dft-\allowbreak{}scan-\allowbreak{}check-\allowbreak{}asap7}).
  \src{tb/tests/test_dft_scan_check.py} (6 tests, skipped without yosys)
  holds the negative controls the checker\textquotesingle s docstring
  referred to.
\item
  \src{tools/dft/check_scan_exclusions.py} regex follows the renamed Q
  nets (\code{u_crypto.key_q}, \code{u_crypto.g_aes.u_aes.rk_q}) {[}A:
  net names of the next synthesis are not yet seen{]}.
\end{itemize}

\subsection{Stage 1a done / not done}\label{tm:ch:dft-architecture:148-stage-1a-done-not-done}\origsec{14.8}

{\footnotesize\setlength{\tabcolsep}{3pt}

{
\begin{longtable}[]{@{}
  >{\raggedright\arraybackslash}p{\dimexpr(\linewidth-2\tabcolsep)*490/1000\relax}
  >{\raggedright\arraybackslash}p{\dimexpr(\linewidth-2\tabcolsep)*490/1000\relax}@{}}
\toprule\noalign{}
\begin{minipage}[b]{\linewidth}\raggedright
Done in this change
\end{minipage} & \begin{minipage}[b]{\linewidth}\raggedright
Not done (bead)
\end{minipage} \\
\midrule\noalign{}
\endhead
\bottomrule\noalign{}
\endlastfoot
\code{soc_top} test ports + \code{dft_ctrl_ports} seam (TAP-ready);
\code{scan_clk_i} shared test clock & CPU macro scan ports, its internal
clock gates/\allowbreak{}resets, \code{pin_order.cfg}, re-harden
(CPU-boundary bead) \\
\code{rv32i_clock_gate.test_en}, live on
\code{u_cpu_cg}/\code{u_gpu_cg}; SRAM-only gates tied 0 &
\texttt{create\_\allowbreak{}clock\ cpu\_\allowbreak{}clk\_\allowbreak{}i}
in \code{sky130_soc.sdc} (SDC bead) \\
\code{cdc_reset_sync} output-side scan select;
PLL/\allowbreak{}GPU-domain reset overrides; PLL test-clock bypass &
Scan-mode SDC, scan insertion, lock-up latches (Stage 2) \\
Crypto \code{key_q}/\code{rk_q} read masks, tested & \sout{Scan-entry
scrub protocol for crypto residue} done in hardware by \code{j41m.6}
(14.10); scan\_en withholding done by \code{j41m.17} (14.14) \\
\code{check_scan_clk_rst.py} netlist proof + negative controls; both
sv2v netlists regenerate & Re-run
\texttt{make\ soc\_\allowbreak{}coverage} and prune waivers that the new
suite makes stale (coverage bead) \\
\end{longtable}
}

}

\subsection{Gap-closure results (2026-10-10, after the implementing
agent ran out of turns)}\label{tm:ch:dft-architecture:149-gap-closure-results-2026-10-10-after-the-implementing-agent-ran-out-of-turns}\origsec{14.9}

\begin{itemize}
\tightlist
\item
  \textbf{Merge of \code{origin/main}} (\#255-\#257): one conflict,
  \src{memory/rtl-design/experiences.jsonl} (both sides appended; both
  kept). No RTL conflict. \code{PASS_FLOOR} 622 in both
  \code{cocotb.yml} and \code{soc_coverage.yml} (the latter was still
  602): 625 measured pre-merge minus 3 (\textbf{superseded by j41m.9:
  624 = 627 measured minus 3, \code{soc_all_ci} from clean on the merged
  tree, 627/0/17 over 62 summaries}); the merge touched no
  \src{tb/cocotb/soc} file, so the count is carried over, not
  re-measured.
\item
  \textbf{sv2v netlists regenerated on the final tree, both PDKs}
  (\texttt{make\ -\allowbreak{}C\ pnr\ sky130-\allowbreak{}soc-\allowbreak{}sv2v\ asap7-\allowbreak{}soc-\allowbreak{}sv2v};
  neither is tracked). Real yosys 0.62: Sky130
  \texttt{hierarchy\ -\allowbreak{}check\ -\allowbreak{}top\ soc\_\allowbreak{}top}
  passes with no undefined module; ASAP7 leaves exactly
  \code{ICGx1_ASAP7_75t_R} and \code{sram_1rw_256x32_asap7} (library
  cell and hard macro, undefined before this change too).
  \code{dft_clk_mux}, \code{dft_rst_mux}, \code{dft_ctrl_ports} are
  defined in both.
\item
  \textbf{SDC}: all four edited SDCs sourced in OpenSTA against a
  port-only \code{soc_top} netlist built from the generated netlists.
  \code{scan_mode_i}, \code{scan_en_i}, \code{scan_rst_ni},
  \code{scan_clk_i} resolve to 1 port each, \code{scan_in_i} and
  \code{scan_out_o} to 8; sky130 (with the hd tt liberty) and
  \code{phase5_soc.sdc} source to completion with no new message against
  the \code{origin/main} version; the two multiclock SDCs abort at their
  own pre-existing section 10b guard (CPU macro pins are absent from a
  port-only netlist), after their DFT block, with identical messages on
  main and PR.
\item
  \textbf{Formal equivalence (test inputs tied inactive vs
  \code{origin/main}, yosys 0.62, capped at 4 GB, finished well under
  the cap)}: \code{rv32i_clock_gate}, \code{cdc_reset_sync} (earlier),
  \code{aes128_core} (399/399 \texttt{\$equiv} cells proven; the same
  miter with \code{scan_mode_i} tied 1 leaves 252 unproven, so it is not
  vacuous) and \code{pll_subsystem} (140/140, via sv2v because yosys
  rejects \texttt{parameter\ string}). \code{crypto_accel} was not run
  separately: its only change is the \code{key_q} mask, covered by
  mutant a.
\item
  \textbf{Mutation checks} (each applied, suite run, restored; diff
  empty afterwards): a \code{key_q} mask removed -\textgreater{}
  \code{test_crypto} 2 FAIL; b \code{rk_q} mask removed -\textgreater{}
  \code{test_crypto} 2 FAIL; c \code{u_cpu_cg} \code{test_en} tied 0
  -\textgreater{} \code{soc_dft_scan.test_cpu_clock_gate_forced_open}
  FAIL; d \code{cdc_reset_sync} input-side only -\textgreater{}
  \code{cdc_reset_sync_scan} 2 FAIL, \code{soc_dft_scan} 1 FAIL; e
  \code{dft_clk_mux} ignores scan -\textgreater{} \code{dft_muxes} 1,
  \code{soc_dft_scan} 3 FAIL; f1 \code{dft_rst_mux} ignores scan (all
  instances) -\textgreater{} \code{dft_muxes} 1, \code{soc_dft_scan} 5
  FAIL; f2 one instance (\code{pll_subsystem.u_core_rm}) -\textgreater{}
  \code{soc_dft_scan} 2 FAIL. \textbf{g \code{u_gpu_domain_rm} ignores
  scan: first SURVIVED} (\code{soc_dft_scan} 8/8 pass): with
  \code{pmu_gpu_rst_n} high the functional side already equals
  \code{scan_rst_ni}. \textbf{Killed 2026-10-10 (bead \bead{j41m.9})} by
  \code{soc_dft_scan.test_scan_gpu_reset_with_pmu_gpu_off} (9 PASS / 1
  FAIL on the mutant, message
  \texttt{gpu\_\allowbreak{}domain\_\allowbreak{}rst\_\allowbreak{}n\ ignored\ scan\_\allowbreak{}rst\_\allowbreak{}ni=\allowbreak{}1\ while\ the\ PMU\ held\ the\ GPU\ in\ reset};
  10/10 on the real RTL). The test boots a 4-instruction CPU firmware
  (\texttt{sw\ MODE\_\allowbreak{}GPU\_\allowbreak{}OFF} to PMU CTRL
  \code{0x2000_8000}, then spin) so the PMU walks the GPU domain to
  \code{DOM_OFF} itself, no deposit; asserts in functional mode that
  \code{pmu_gpu_rst_n}, \code{gpu_domain_rst_n_func} and
  \code{gpu_domain_rst_n} are all 0; stops the functional clocks and
  enters scan mode with \code{scan_rst_ni} held at 1 (the PMU is reset
  by \code{core_rst_n}, which equals \code{scan_rst_ni} in scan mode, so
  holding it high preserves the power-down state); re-asserts the guard
  (\code{dom_state_q[GPU]==DOM_OFF}, \code{pmu_gpu_rst_n==0},
  \code{gpu_domain_rst_n_func==0}) and only then requires
  \code{gpu_domain_rst_n==1}. The converse leg pulses
  \code{scan_rst_ni}=0, which resets the PMU (\code{pmu_gpu_rst_n} back
  to 1, asserted as a guard) and requires \code{gpu_domain_rst_n==0}.
  \textbf{CPU analogue} (\code{test_scan_cpu_reset_with_pmu_cpu_off},
  firmware -\textgreater{} DMA -\textgreater{} PMU CTRL
  \code{MODE_CPU_OFF} via the existing \code{pmu_cycle.hex}):
  \code{cpu_domain_rst_n} follows \code{scan_rst_ni} with the PMU
  holding the CPU in reset; \textbf{passes on the real RTL, so the
  "needs no mux" claim for \code{cpu_domain_rst_n} holds}. Mutant h
  (\code{u_cpu_pmu_rst_sync.scanmode_i} tied 0) fails it (and also
  \code{test_cpu_clock_gate_forced_open}). A mutant that removed only
  the output-side mux of \code{cdc_reset_sync} while keeping its
  async-input select is equivalent once the 2-stage chain fills (the
  chain input is constant 1 and its async clear is \code{scan_rst_ni}),
  so it is not distinguishable by any test and is not claimed.
\end{itemize}

\subsection{\texorpdfstring{CPU macro scan/test boundary, Sky130
half (bead \code{j41m.4}, RTL step
only)}{CPU macro scan/test boundary, Sky130 half (bead , RTL step only)}}\label{tm:ch:dft-architecture:1410-cpu-macro-scantest-boundary-sky130-half-bead-j41m4-rtl-step-only}\origsec{14.10}

Supersedes 14.6 for the macro boundary. RTL only: \textbf{no PD run, no
macro view regenerated} (the committed \code{pnr/sky130/cpu/macro/*}
views and the ASAP7 views are still the pre-DFT ones). The ASAP7 half is
bead \bead{j41m.13}. Labels as above.

\textbf{Ports added to \code{rv32i_cpu_top}} (inactive values =
today\textquotesingle s behaviour; Verilator leaves an undriven input at
0, so every existing cocotb top that instantiates the macro bare is
unaffected):

{\small\setlength{\tabcolsep}{4pt}

{
\begin{longtable}[]{@{}
  >{\raggedright\arraybackslash}p{\dimexpr(\linewidth-8\tabcolsep)*155/1000\relax}
  >{\raggedright\arraybackslash}p{\dimexpr(\linewidth-8\tabcolsep)*58/1000\relax}
  >{\raggedright\arraybackslash}p{\dimexpr(\linewidth-8\tabcolsep)*68/1000\relax}
  >{\raggedright\arraybackslash}p{\dimexpr(\linewidth-8\tabcolsep)*97/1000\relax}
  >{\raggedright\arraybackslash}p{\dimexpr(\linewidth-8\tabcolsep)*602/1000\relax}@{}}
\toprule\noalign{}
\begin{minipage}[b]{\linewidth}\raggedright
Port
\end{minipage} & \begin{minipage}[b]{\linewidth}\raggedright
Dir
\end{minipage} & \begin{minipage}[b]{\linewidth}\raggedright
Width
\end{minipage} & \begin{minipage}[b]{\linewidth}\raggedright
Inactive
\end{minipage} & \begin{minipage}[b]{\linewidth}\raggedright
Driven in \code{soc_top} from
\end{minipage} \\
\midrule\noalign{}
\endhead
\bottomrule\noalign{}
\endlastfoot
\code{scan_mode_i} & in & 1 & 0 & \code{dft_scan_mode} (pin OR TAP) \\
\code{scan_en_i} & in & 1 & 0 & \code{dft_scan_en} (\code{dft_ctrl}:
TAP-serial ? TAP : pin). \textbf{Not} the raw \code{scan_en_i} port \\
\code{scan_in_i} & in & 2 & 0 & top
\texttt{scan\_\allowbreak{}in\_\allowbreak{}i[SCAN\_\allowbreak{}CHAINS-\allowbreak{}1\ -\allowbreak{}:\allowbreak{}\ 2]},
i.e. indices 7:6 \\
\code{scan_out_o} & out & 2 & tied 0 & top \code{scan_out_o[7:6]};
fabric bits \code{[5:0]} stay the old placeholder tie \\
\end{longtable}
}

}

\textbf{Deliberately no \code{scan_clk} and no \code{scan_rst_ni} port}
{[}R{]}:

\begin{itemize}
\tightlist
\item
  \code{clk_i} of the macro is \code{cpu_gated_clk}, which in scan mode
  is already the test clock (\code{dft_clk_mux} on the \code{cpu_clk_i}
  root, ahead of \code{u_cpu_cg}, 14.2). The macro keeps one clock root
  and gains \textbf{no clock mux, so no insertion delay on its tree};
  the +0.174 ns margin of section 5 is not spent on a clock mux.
\item
  Every reset inside the macro is the synchronous \code{rst_n_i} (no
  \texttt{or\ negedge} flop, no derived reset:
  \code{check_scan_clk_rst.py} block mode reports \textbf{0 async
  reset/set pins} on both macro netlists {[}M{]}). \code{soc_top} drives
  it from \code{cpu_domain_rst_n}, which equals \code{scan_rst_ni} in
  scan mode (14.2). A mux inside the macro would only put a gate in
  front of a \textasciitilde4.8 k-flop reset tree. If a later change
  introduces any async or derived reset,
  \texttt{-\allowbreak{}-\allowbreak{}max-\allowbreak{}async-\allowbreak{}pins\ 0}
  fails and the override becomes mandatory then.
\end{itemize}

\textbf{Clock gates} {[}R{]}: the Sky130 macro has \textbf{none}
(\code{SKY130_CPU_SV_FILES} omits \code{rv32i_clock_gate.sv}, the SRAM
arm instantiates the macro directly). The ASAP7 arm has 10
(\code{u_tag_cg} + 4x \code{u_data_cg}, in each of \code{u_icache} and
\code{u_dcache}); their \code{test_en} was tied \texttt{1'b0} and is now
\texttt{test\_\allowbreak{}en\_\allowbreak{}w\ =\allowbreak{}\ scan\_\allowbreak{}mode\_\allowbreak{}i}
(via new \code{test_en_i} ports on \code{rv32i_core},
\code{rv32i_icache}, \code{rv32i_dcache};
\texttt{|\ mbist\_\allowbreak{}en} is added where Stage 4 lands). They
clock \textbf{SRAM macros only}, so nothing scannable is gained; the
reason to do it is that the clock-gate rule of the checker requires a
test term in every clock-gate latch. The new port is unused in the
Sky130/\allowbreak{}FPGA/\allowbreak{}FreePDK45 arms (waived with a
scoped \texttt{lint\_\allowbreak{}off\ UNUSEDSIGNAL}).

\textbf{Placeholders}: \code{scan_en_i}/\code{scan_in_i} have no RTL
consumer and \code{scan_out_o} is one
\texttt{assign\ scan\_\allowbreak{}out\_\allowbreak{}o\ =\allowbreak{}\ '0;\ /\allowbreak{}/\allowbreak{}\ DFT\_\allowbreak{}PLACEHOLDER}
until Stage 2 stitches the two chains \textbf{inside the macro flow}
(14.4 contract: remove the tie, last flop to \code{scan_out_o[k]},
\code{scan_in_i[k]} to first flop, every \code{SCE} to
\code{scan_en_i}). The two blackbox stubs
(\texttt{pnr/\allowbreak{}\{sky130,\allowbreak{}asap7\}/\allowbreak{}soc/\allowbreak{}rv32i\_\allowbreak{}cpu\_\allowbreak{}top\_\allowbreak{}stub.\allowbreak{}sv})
carry the same four ports. The committed ASAP7 macro LEF/lib lack these
pins: an ASAP7 SoC P\&R against them fails at macro pin match until
\code{j41m.13}.

\textbf{Pins} (\src{pnr/sky130/cpu/pin_order.cfg}): 6 added at the end
of the N edge (401 -\textgreater{} 407 signal bits; every port bit
matches exactly one entry {[}M{]}). \src{pnr/asap7/cpu/pin_order.cfg} is
\textbf{unchanged}: the ASAP7 CPU harden will fail loudly in
\code{io_place.py} (design pin missing from the file) until
\code{j41m.13} adds them; it could not complete on this host anyway
(\code{ma7}/\code{lxv}/\code{2kn}).

\textbf{Checks} {[}M{]}:
\texttt{make\ -\allowbreak{}C\ pnr\ dft-\allowbreak{}scan-\allowbreak{}check-\allowbreak{}cpu-\allowbreak{}macro}
/ \code{dft-scan-check-cpu-macro-asap7} (new;
\code{check_scan_clk_rst.py} BLOCK MODE:
\texttt{-\allowbreak{}-\allowbreak{}clock-\allowbreak{}port\ clk\_\allowbreak{}i\ -\allowbreak{}-\allowbreak{}max-\allowbreak{}async-\allowbreak{}pins\ 0\ -\allowbreak{}-\allowbreak{}expect-\allowbreak{}seq-\allowbreak{}min\ 900\ -\allowbreak{}-\allowbreak{}all-\allowbreak{}latches-\allowbreak{}need-\allowbreak{}scan-\allowbreak{}mode},
ASAP7 also
\texttt{-\allowbreak{}-\allowbreak{}expect-\allowbreak{}latches-\allowbreak{}min\ 10}):
Sky130 macro 944 word-level sequential cells, 0 async pins, 0
violations; ASAP7 macro 954 / 0 / 0 violations, 10 clock-gate latches
all with \code{scan_mode_i} in D. Negative control on the real ASAP7
macro netlist with \texttt{.\allowbreak{}test\_\allowbreak{}en(1'b0)}:
\textbf{10 violations}. The SoC-level checks needed
\code{u_dft_ctrl.tap_scan_en} added to \code{--tap-allowed-nets} (its
first consumer is now the macro pin), as 15.9 / the bead note predicted.
\src{tb/tests/test_dft_scan_check.py} 25 -\textgreater{} 32 tests, CI
floor raised to 32 and the two macro targets added to the CI loop.

\textbf{Risks and Stage 2 notes} {[}A unless noted{]}:

\begin{itemize}
\tightlist
\item
  Fabric insertion (\code{j41m.8}) must take its \code{SCE} from
  \code{dft_scan_en}, not the raw \code{scan_en_i} port, or serial (TAP)
  scan would shift the macro and not the fabric.
\item
  Re-harden of the Sky130 macro will add the scan mux to
  \textasciitilde4.8 k D pins (+0.066/\allowbreak{}+0.085 ns, section 5)
  on a +0.174 ns margin: the relaxation/\allowbreak{}retime question of
  section 5 is still open and is \emph{the} Stage 2 risk. This RTL step
  adds only the 6 pins and one tie cell, no logic on a timing path.
\item
  Whether appending 6 names to the N list moves the other 397 pins is
  \textbf{unmeasured}; compare the regenerated LEF against the committed
  one before accepting (see \code{pin_order.README.md}).
\item
  The macro netlist now has 3 input ports and 2 output ports without
  loads/\allowbreak{}with tie cells;
  \texttt{DIODE\_\allowbreak{}ON\_\allowbreak{}PORTS:\allowbreak{}\ in}
  may place diodes on the unloaded inputs. Harmless, not checked.
\item
  \code{sky130_cpu.sdc} is unchanged (new ports are unconstrained in
  functional STA); a scan-mode macro SDC is a Stage 2 deliverable.
\end{itemize}

\subsection{\texorpdfstring{PD re-harden for the CPU macro scan
boundary and the \code{cpu_clk_i} SDC (beads \code{j41m.4} PD step +
\code{j41m.5},
2026-10-16)}{PD re-harden for the CPU macro scan boundary and the  SDC (beads  PD step + , 2026-10-16)}}\label{tm:ch:dft-architecture:1411-pd-re-harden-for-the-cpu-macro-scan-boundary-and-the-cpu-clk-i-sdc-beads-j41m4-pd-step-j41m5-2026-10-16}\origsec{14.11}

Labels: {[}M{]} measured this session, {[}R{]} read from a file, {[}A{]}
inferred. Runs: CPU macro \code{RUN_2026-10-16_02-27-26}
(\code{/nobackup/sky130_cpu_runs}), SoC \code{RUN_2026-10-16_04-53-49}
(\texttt{make\ -\allowbreak{}C\ pnr\ librelane-\allowbreak{}sky130-\allowbreak{}soc-\allowbreak{}rccal-\allowbreak{}noklayout},
KLayout DRC skipped as always). Baselines: CPU
\code{RUN_2026-10-11_06-00-12}, SoC \code{RUN_2026-10-11_08-39-41} and
\code{RUN_2026-10-09_06-18-39}. \textbf{Not a single-variable comparison
with either baseline} (RTL vintage changed: JTAG TAP, clock switch, scan
boundary; and the SoC SDC now clocks the CPU domain).

\textbf{CPU macro {[}M{]}.} LVS 0, KLayout DRC 0, routing DRC 0, Magic
27,733,913 (bit-identical to the dud4 macro), antenna 94/96 (was 96/96),
4,880 flops (was 4,868; not attributed). Setup: tt/ff corners clean and
slightly better (+3.79 ns nom\_tt, was +3.44); \textbf{ss still fails
exactly as before} (max\_ss -4.009 ns / 972 violators, was -3.704 / 996;
nom\_ss -3.333 / 861, was -3.198 / 854; icache-tag path class, reported
not tuned). Hold: I/\allowbreak{}O-port-only violators remain (nom\_tt
-0.076 ns / 7, was -0.379 / 84; max\_ss -0.467 / 124, was -0.834 / 144);
reg-to-reg hold clean. Gate-vs-RTL proof
(\src{tools/verif/gls/sky130/run_netlist_proof.sh} on the post-fill
netlist): 6/6 programs identical, 124/124 matrix, negative control
detected (4/6 diverge, 90 cells fail). \textbf{Pin result: the 6
appended N-edge names re-spaced the whole N edge: 95 of the 403 existing
pins moved (max 214 um, mean 106 um),
S/\allowbreak{}E/\allowbreak{}W/\allowbreak{}power identical}
(\code{pin_order.README.md}). The Liberty views carry the 6 new pins.

\textbf{SoC {[}M{]}} (vs \code{RUN_2026-10-11_08-39-41}): LVS 0 (same),
routing DRC 0 (same), Magic 8,984 (same), setup and hold 0 violators at
all 9 corners (same), hold worst +0.113 ns (max\_ff; was +0.287),
antenna 197/252 nets/pins (was 187/265), max-slew 304 (was 370), util
45.37 \% (45.22), power 100.3 mW (84.6; the CPU domain\textquotesingle s
945 flops and the new CTS tree are now clocked and counted, {[}A{]}),
433 hold buffers (was 2: the CPU-macro input hold of about 1 ns,
previously unchecked because the macro clock was undefined). CTS now
builds \code{cpu_core_clk} (174 sinks) and \code{cpu_gated_clk_regs}
(771 sinks) trees. The flow\textquotesingle s reported setup worst slack
is a literal 0.000 at every corner: a latch time-borrow artifact of the
TAP\textquotesingle s negative-level latch in the \code{tck} group, not
a real margin. Real per-clock figures (replica STA with the
run\textquotesingle s SPEF, propagated clocks): core\_clk +9.97 ns
nom\_tt / +2.55 max\_ss (baseline +2.48), cpu\_clk +10.1 nom\_tt / +7.09
max\_ss, hold cpu\_clk +0.59 / +1.34.
\textbf{\code{CVT.MaxCapViolations} FAILS (deferred): 2 unwaived nets,
both clock nets}: \code{cpu_gated_clk} (driver \code{and2_2} clock-gate
AND, 0.60-0.68 pF vs limit 0.19-0.40 at 8 corners) and \code{core_clk}
(driver \code{a21bo_2}, the \code{dft_clk_mux} output, 0.29-0.31 pF vs
0.18-0.29 at 4 tt/ss corners). Both old waivers are STALE (the by-name
driver waivers for \code{RUN_2026-10-11_08-39-41} match nothing).
Attribution {[}A{]}: the \code{core_clk} driver is the clock switch from
the JTAG TAP work, absent from the baseline netlist (no \code{tck_i}
there), so that one is not caused by this change; \code{cpu_gated_clk}
was an \code{and2_4} in the baseline and is now a clock net owned by
CTS, which \code{repair_design} does not buffer. No waiver was added or
changed (owner decision).

\textbf{Clock-root drive fix (owner decision: fix, not waive; RTL done,
PD re-harden pending)}: root cause is that CTS owns the clock net, hangs
a long root wire on the synthesised driver (the register-tree root
buffer is placed at the sink centroid: \code{and2_2} at (4533, 2026) um,
\code{clkbuf_regs_0_cpu_clk} at (2276, 193) um, about 4 mm), and
\code{repair_design} skips clock nets, so nothing re-drives it. The
drivers must be strong from synthesis on. Under
\code{+define+USE_SKY130_CLKCELLS} (only in
\code{SKY130_SOC_SV2V_DEFINES}, so the CPU macro is untouched) the two
roots are hand-instantiated library cells, black-box stubs in
\src{pnr/sky130/soc/sky130_clkcells_stub.v}: \code{u_cpu_cg}
(\texttt{rv32i\_\allowbreak{}clock\_\allowbreak{}gate\ \#(.\allowbreak{}HARD\_\allowbreak{}ROOT(1))})
is \code{sky130_fd_sc_hd__sdlclkp_4} (GATE/SCE =
\code{en}/\code{test_en}, same \texttt{en\ |\ test\_\allowbreak{}en}
latched while CLK is low) driving \code{sky130_fd_sc_hd__buf_16}; every
\code{dft_clk_mux} (\code{u_pll_sub.u_ref_cm},
\code{u_cpu_pll_sub.u_ref_cm}) is \code{mux2_4} + \code{buf_16}.
Flattened instance names (stable, checked in the yosys-flattened
netlist): \code{u_cpu_cg.g_cell.u_icg},
\code{u_cpu_cg.g_cell.u_root_buf}, \code{u_pll_sub.u_ref_cm.u_mux},
\code{u_pll_sub.u_ref_cm.u_root_buf}, and the same two under
\code{u_cpu_pll_sub}. \code{u_gpu_cg} stays behavioural so it still
optimises away with the tie-off GPU. \code{RSZ_DONT_TOUCH_RX} in the
three \code{pnr/sky130/soc/config*.json} pins these instances.
\textbf{Why \code{buf_16} and not \code{clkbuf_16} (CTS-0105,
\runid{RUN_2026-10-17_08-53-43})}: the first cell
arm used \code{clkbuf_16} for the three root buffers
(\code{u_cpu_cg.g_cell.u_root_buf},
\code{u_pll_sub.u_ref_cm.u_root_buf},
\code{u_cpu_pll_sub.u_ref_cm.u_root_buf}). \code{clkbuf_16} is in
\code{CTS_CLK_BUFFERS}, so TritonCTS treated each one as a clock buffer
that already existed on the net and logged
\texttt{CTS-\allowbreak{}0105:\allowbreak{}\ Net\ "core\_\allowbreak{}clk"\ already\ has\ clock\ buffer\ u\_\allowbreak{}pll\_\allowbreak{}sub.\allowbreak{}u\_\allowbreak{}ref\_\allowbreak{}cm.\allowbreak{}u\_\allowbreak{}root\_\allowbreak{}buf.\allowbreak{}\ Skipping.\allowbreak{}.\allowbreak{}.\allowbreak{}}
(likewise \code{cpu_core_clk} and \code{cpu_gated_clk}). It built no
clock tree on any of the three roots, and setup ended at WNS -130 ns.
Owner decision: use \code{sky130_fd_sc_hd__buf_16}, a plain buffer that
is not in the CTS buffer list, so CTS sees an ordinary driver and builds
its tree below it. \code{CTS_CLK_BUFFERS} / \code{CTS_ROOT_BUFFER} still
name \code{clkbuf_16} (CTS\textquotesingle s own cell choice,
unchanged). Instance names are unchanged, so \code{RSZ_DONT_TOUCH_RX}
and \code{MANUAL_GLOBAL_PLACEMENTS} need no edit. Liberty check
(MEASURED, \code{max_capacitance} of the output pin): \code{buf_16} 5.0
pF (the library cap) at ss\_\allowbreak{}n40C\_\allowbreak{}1v28,
ss\_\allowbreak{}100C\_\allowbreak{}1v60 and
tt\_\allowbreak{}025C\_\allowbreak{}1v80, 2.846 pF at
ss\_\allowbreak{}n40C\_\allowbreak{}1v40, against \code{clkbuf_16} 0.757
/ 1.017 / 1.532 / 2.826 pF at the same corners, so the drive margin is
not reduced. Input pin cap is 13.3 fF (\code{buf_16}) vs 7.5 fF
(\code{clkbuf_16}), negligible against the 0.33 pF
\code{mux2_4}/\code{sdlclkp_4} limit. Cell area is 27.5 vs 25.0 um\^{}2.
NOT yet measured: whether CTS now builds the trees and closes setup,
which is the next PD run. Load vs limit (measured loads from
\runid{RUN_2026-10-16_04-53-49}; the limit quoted is
the \code{clkbuf_16} Liberty value used at the time, \code{buf_16} is
higher, see above): \code{cpu_gated_clk} 0.60-0.68 pF vs 1.017 pF at ss
(about 67 \% worst case, was 0.19 pF), \code{core_clk} 0.29-0.31 pF vs
1.017 pF at ss (about 31 \%, was 0.18 pF). INFERRED, not measured: the
new 2-pin ICG-to-buffer and mux-to-buffer nets are placed by global
placement and must stay under the \code{sdlclkp_4}/\code{mux2_4} ss
limit (0.33 pF, roughly 3 mm of wire); the post-CTS run decides.
Scan-mode behaviour of the cell arm: \code{dft-scan-check} runs against
\src{tools/dft/sky130_cell_models.v} (0 violations, 1 clock gate checked
for the scan-mode force: \code{u_cpu_cg.g_cell.u_icg}, via the j41m.16
\texttt{-\allowbreak{}-\allowbreak{}expect-\allowbreak{}clock-\allowbreak{}gates\ 1}
proof; negative controls: a model without the \code{SCE} term fails it
with the gate-not-forced-open message, a broken-mux model fails the mux
rule). The proof is on the analysis model, whose \texttt{GATE\ |\ SCE}
latch input is taken from the Liberty statetable
(\texttt{SCE=\allowbreak{}H\ -\allowbreak{}>\ M0=\allowbreak{}H}
regardless of \code{GATE}), not on the library cell itself.

\textbf{CDC budget check {[}M{]}}
(\texttt{make\ -\allowbreak{}C\ pnr\ check-\allowbreak{}cdc-\allowbreak{}timing-\allowbreak{}sky130-\allowbreak{}soc},
real SPEF): PASS at nom\_tt and max\_ss; 11 groups (10
FIFO/\allowbreak{}bridge/\allowbreak{}PMU/\allowbreak{}IRQ + the
\code{u_m_rst_status_sync}) plus the two domain-wide fallbacks all
bound; worst margins at max\_ss: fifo\_\allowbreak{}rd2wr +0.95 ns of
12.5, fifo\_\allowbreak{}wr2rd +1.68. Negative control (all budgets
forced to 1 ns): FAIL with 9 findings. The requested
\texttt{set\_\allowbreak{}max\_\allowbreak{}delay\ -\allowbreak{}datapath\_\allowbreak{}only}
cannot live in \code{sky130_soc.sdc}: OpenSTA has no such flag and drops
a max\_delay across \code{-asynchronous} clocks (bead \bead{k07}), so the
budgets are in \code{sky130_soc_cdc_check.sdc} as for ASAP7.

\textbf{Scan checkers {[}M{]}}: \code{dft-scan-check} (SoC, sv2v netlist
of the same RTL): 1,694 sequential cells, 295 async pins, 28 scan-mode
muxes, 0 violations; \code{dft-scan-check-cpu-macro}: 944 sequential
cells, 0 async pins, 0 violations. They run on the sv2v netlists, not on
the placed netlist (which has Sky130 cells the checker has no models
for).

\subsection{\texorpdfstring{Stage 2 scan insertion: method, port
reuse, CPU macro result (bead \code{j41m.8},
2026-10-26)}{Stage 2 scan insertion: method, port reuse, CPU macro result (bead , 2026-10-26)}}\label{tm:ch:dft-architecture:1412-stage-2-scan-insertion-method-port-reuse-cpu-macro-result-bead-j41m8-2026-10-26}\origsec{14.12}

Labels as above. Code: \src{pnr/sky130/soc/plugin/cvt_scan_step.py}
(step \code{CVT.ScanInsert}, flow \code{CVTClassicScan} = Classic plus
the step right after \code{Checker.NetlistAssignStatements}),
\code{.../tcl/scan_insert.tcl}, \src{tools/dft/check_scan_chains.py},
\src{tools/dft/run_scan_step_isolated.sh}. No shared LibreLane file is
edited.

\textbf{Port reuse {[}M{]}.} OpenROAD 26Q2 \code{insert_dft} does
\textbf{not} reuse existing ports (answers 14.4). It always creates
\code{scan_enable_0}, \code{scan_in_<k>}, \code{scan_out_<k>};
\code{-scan_*_name_pattern} only renames them (no \texttt{\%d}
expansion, a pattern without an index collides,
\texttt{scan\_\allowbreak{}in\_\allowbreak{}i[\%d]} aborts with
\code{stoi}). The step therefore (1) lets the tool insert, (2) merges
its nets into the existing \code{scan_in_i[n]}, the existing enable net
(\code{dft_scan_en} on the SoC, \code{scan_en_i} on the macro, not the
raw \code{scan_en_i} port, per 14.10), (3) destroys the \code{conb}
placeholder tie on \code{scan_out_o[n]} and drives it through a
\code{buf_2} from the chain\textquotesingle s last Q, (4) strips the
tool\textquotesingle s own ports from the written netlist header (they
otherwise reach IO placement as
\texttt{scan\_\allowbreak{}in\_\allowbreak{}0\ not\ found\ in\ config\ but\ found\ in\ design}).
The netlist has no \code{assign}.

\textbf{Chains {[}M, synthesis netlists of
\runid{RUN_2026-10-25_08-28-45} /
\runid{RUN_2026-10-24_19-19-02}{]}.}
\code{-max_length} (not \code{-max_chains}) splits per clock domain.
SoC, 3700: 5 \code{core_clk} chains (3565 x4, 3562) plus 1 CPU-domain
chain (981) = \code{scan_in_i/out_o[5:0]}; 18,803 scan flops, 922 left
unscanned on purpose = 666 flops clocked by \code{tck_i} (TAP, MBIST
TDRs, incl. the negedge-tck ones behind an inverter) or
\code{scan_clk_i} (clock switch) + the 256 crypto key registers
(decision 3). The 3 latches are not scanned. The MBIST TDR flops on
\code{tck} are therefore \textbf{not} scan-tested (no port is left for a
\code{tck} chain): a decision for Stage 3. CPU macro, 2660: 2 chains
(2650, 2649) on \code{scan_in_i/out_o[1:0]}, 5,299 flops (4,880 before
the MBIST controller), 0 unscanned.

\textbf{CPU macro harden \code{RUN_2026-10-26_05-39-04} vs run 6
\code{RUN_2026-10-24_19-19-02} {[}M{]}.} LVS 0 errors, KLayout DRC 0,
routing DRC 0. Magic DRC 27,733,913: the report is
\textbf{byte-identical} to run 6 (\code{cmp}), so the 45a waiver basis
is unchanged. LEF: 503 signal pins, 503 identical, 0 moved
(\texttt{tools/\allowbreak{}verif/\allowbreak{}pin\_\allowbreak{}lock.\allowbreak{}py\ check});
die 3600 x 1800. Setup: all six tt/ff corners clean (nom\_tt +3.008 ns,
was +3.031); ss still fails (max\_ss -4.584 ns, 1,005 violators, was
-3.996 / 1,149): reg-to-reg 910 (worst -4.584), reg-to-out 74,
\textbf{new: 16 \code{apb_paddr_i} input-to-reg} (worst -0.756) at
max\_ss. Hold: every violator is input-port-to-register in both runs, no
reg-to-reg violation (nom\_tt 11, worst -0.091 ns, was 156 / -0.458;
max\_ss 136 / -0.491, was 217 / -1.296). Area: stdcell 775,983 um2
(707,819), util 42.90 \% (41.81 \%), power 115.8 mW (103.4), wirelength
6.23 M (5.59 M), antenna 88 pins / 87 nets (119 / 116). Netlist proof on
the post-fill netlist
(\texttt{run\_\allowbreak{}netlist\_\allowbreak{}proof.\allowbreak{}sh\ .\allowbreak{}.\allowbreak{}.\allowbreak{}\ negctl}):
census 5,299 flops with \code{fwd_b_*} present, 6/6 programs IDENTICAL,
124/124 port-matrix cells, negative control detected (4/6 diverge, 90
cells fail). \code{check_scan_chains.py} on the final netlist: 2 chains
2650/\allowbreak{}2649, 5,299 scan flops, 0 unscanned, every SCE traces
through the repair buffer tree to \code{scan_en_i}.

\textbf{Limits {[}A{]}.} The proof runs the macro in functional mode
(scan ports tied 0): chain shift is verified structurally, not by
simulation. The macro Liberty has \textbf{no timing arcs on
\code{scan_en_i} / \code{scan_in_i}} (unconstrained inputs,
\code{sky130_cpu.sdc} unchanged), so the SoC does not time them: a
scan-mode macro SDC is still a Stage 2/3 deliverable. SoC-level numbers:
see 14.13.

\subsection{\texorpdfstring{Stage 2 SoC harden with scan:
\code{RUN_2026-10-26_09-59-03} vs \code{RUN_2026-10-25_08-28-45} (bead
\code{j41m.8},
2026-10-26)}{Stage 2 SoC harden with scan:  vs  (bead , 2026-10-26)}}\label{tm:ch:dft-architecture:1413-stage-2-soc-harden-with-scan-run-2026-10-26-09-59-03-vs-run-2026-10-25-08-28-45-bead-j41m8-2026-10-26}\origsec{14.13}

Labels as above. Command:
\texttt{make\ -\allowbreak{}C\ pnr\ librelane-\allowbreak{}sky130-\allowbreak{}soc-\allowbreak{}noklayout}
(rccal default, \code{CVT_SYNTH_GATE=warn}, KLayout DRC skipped as in
the baseline), flow \code{CVTClassicScan}, CPU macro = the scan macro of
14.12. Not a single-variable comparison with the baseline: 18,803 SoC
flops became scan flops and the CPU macro was re-hardened with scan.
Runtime 5.21 h (baseline 4.28 h, sum of step runtimes). Final exit:
deferred errors Magic DRC and max-cap only.

\textbf{Scan structure {[}M, \code{check_scan_chains.py} on the
post-fill netlist{]}.} 6 fabric chains (3565 x4, 3562, 981), 18,803 scan
flops, 922 unscanned, every one accounted for: 664 on \code{tck_i} (TAP
and MBIST TDR flops, incl. the negedge ones), 2 on \code{scan_clk_i},
256 crypto key registers. All 18,803 SCE pins trace through the repair
buffer tree to \code{dft_scan_en} (root driver: the TAP/pin
\code{mux2_8}). Chains {[}7:6{]} are inside the CPU macro (14.12). The
Tcl insertion, the strip of the tool\textquotesingle s own ports and the
State output are exercised alone by
\src{tools/dft/run_scan_step_isolated.sh}.

\textbf{Delta table {[}M{]}.}

{\small\setlength{\tabcolsep}{4pt}

{
\begin{longtable}[]{@{}
  >{\raggedright\arraybackslash}p{\dimexpr(\linewidth-4\tabcolsep)*475/1000\relax}
  >{\raggedright\arraybackslash}p{\dimexpr(\linewidth-4\tabcolsep)*253/1000\relax}
  >{\raggedright\arraybackslash}p{\dimexpr(\linewidth-4\tabcolsep)*253/1000\relax}@{}}
\toprule\noalign{}
\begin{minipage}[b]{\linewidth}\raggedright
\end{minipage} & \begin{minipage}[b]{\linewidth}\raggedright
scan run
\end{minipage} & \begin{minipage}[b]{\linewidth}\raggedright
baseline
\end{minipage} \\
\midrule\noalign{}
\endhead
\bottomrule\noalign{}
\endlastfoot
LVS errors / routing DRC & 0 / 0 & 0 / 0 \\
Magic DRC & 11,003 & 8,952 \\
setup / hold violators, 9 corners (flow) & 0 / 0 & 0 / 0 \\
hold worst slack nom\_tt / max\_ss / max\_ff (ns) & +0.415 / +0.223 /
+0.258 & +0.444 / +0.534 / +0.166 \\
setup worst, \code{core_clk} nom\_tt / nom\_ss / max\_ss (ns) & +5.671 /
+0.156 / +0.002 & +5.668 / n.m. / +0.134 \\
setup worst, \code{cpu_clk} nom\_tt / max\_ss (ns) & +15.92 / +7.76 &
+14.73 / +5.50 \\
CPU macro input pins, hold worst nom\_tt / max\_ss / max\_ff (ns) &
+0.568 / +1.457 / +0.274 & +0.338 / +1.108 / n.m. \\
max-slew violators, worst corner (max\_ss) & 525 & 795 \\
max-cap violators, worst corner (max\_ss) & 5 & 5 \\
antenna pins / nets & 368 / 286 & 313 / 225 \\
power (mW) / util (all) / stdcell area (um2) & 119.0 / 46.55 \% /
2,052,900 & 104.8 / 45.82 \% / 1,904,520 \\
instances / timing repair bufs / hold bufs / diodes & 328,829 / 28,045 /
111 / 17,642 & 326,231 / 27,140 / 135 / 16,243 \\
wirelength & 9.21 M & 8.25 M \\
\end{longtable}
}

}

Flow setup slack is a literal 0.000 at every corner (the TAP latch
time-borrow artifact of 14.11); the \code{core_clk}/\code{cpu_clk} rows
are replica STA with the run\textquotesingle s SPEF and propagated
clocks (\src{pnr/scripts/sta_load_sky130_soc.tcl} +
\code{sta_scan_queries_sky130_soc.tcl}; the worst \code{core_clk}
endpoint at max\_ss is \code{_167843_/D}, an \code{sdfrtp_1} in
\code{u_mbist.u_stsync_core}, +2.4 ps, against +134 ps on the same class
in the baseline). Setup per clock at all 9 corners,
nom\_\allowbreak{}tt/\allowbreak{}min\_\allowbreak{}tt/\allowbreak{}max\_\allowbreak{}tt:
core +5.67/\allowbreak{}+5.89/\allowbreak{}+5.47, cpu
+15.9/\allowbreak{}+16.3/\allowbreak{}+15.6; ss: core
+0.156/\allowbreak{}+0.359/\allowbreak{}+0.002, cpu
+8.20/\allowbreak{}+8.68/\allowbreak{}+7.76; ff: core
+7.86/\allowbreak{}+8.03/\allowbreak{}+7.68, cpu
+18.7/\allowbreak{}+19.0/\allowbreak{}+18.4. Hold per clock: core +0.223
(max\_ss) to +0.486 (min\_ss), cpu +0.270 (max\_ff) to +1.126 (max\_ss).

\textbf{Magic DRC +2,051 root cause {[}M{]}.} Two rules only:
\code{nwell.4} ("all nwells must contain metal-connected N+ taps")
10,971 vs 8,920 (+2,051) and \code{met4.4a} 32 vs 32. The
macro-footprint counts are identical to the baseline (84 CPU + 32
NPU-SRAM); \textbf{none} of the 2,051 is inside an SRAM or macro
footprint: all of them are in the standard-cell rows, and the rule is
the same one that already accounts for 8,836 standard-cell rectangles of
the baseline. 2,666 rectangles are at positions not in the baseline (615
baseline rectangles went away): placement changed. Cell attribution
under the rectangles (DEF of the post-fill stage): 743 sit under the new
scan flops (\code{sdfxtp_1} 486, \code{sdfrtp_4} 88, \code{sdfrtp_1} 57,
\code{sdfxtp_4} 46, \code{sdfxtp_2} 38, \code{sdfrtp_2} 27), +183 under
\code{dlygate4sd3_1}, the rest in gaps between cells. There is no new
rule, no metal or routing rule, and routing DRC is 0. \textbf{Accepted
under 45a by user decision 2026-10-26.} It is still not proven benign by
the foundry deck: Magic\textquotesingle s \code{nwell.4} is the
tap-connectivity check on an abstract view, and KLayout DRC, which would
settle it, is skipped on this host (8.5 GB, bead \bead{00ef}).

\textbf{\code{dft_scan_en} timing {[}M, all 9 corners{]}.} The TAP/pin
enable is launched by a TAP flop on the \textbf{falling} edge of
\code{tck} (\texttt{tck'}), so against the rising-edge capture inside
the \code{tck} image of the trees the worst hold through
\code{dft_scan_en} is \textbf{+43.961 ns} (nom\_tt; +39.4 to +46.0 ns
over the 9 corners) and setup \textbf{+42.157 ns} (+35.2 to +44.7). The
exposure raised in the bead notes (a tck-to-tck hold race on 18,803 SCE
pins) is not there; the repair step built a 28-buffer chain (about 10 ns
from the root to the farthest SCE), which matters for the scan-mode
clock period, not for the functional view. The same mechanism explains
the baseline\textquotesingle s +47 ns removal slack on the TAP reset
flops.

\textbf{Clock-mux insertion delay and skew {[}M{]}.} Cell delays on the
capture clock path: \code{u_pll_sub.u_ref_cm} mux + root buffer 0.470 ns
nom\_tt (0.474 baseline) and 0.840 ns max\_ss (0.848);
\code{u_cpu_pll_sub.u_ref_cm} 1.334 ns (1.460) and 2.511 ns (2.727).
\code{core_clk} latency/\allowbreak{}skew nom\_tt 2.350-2.897 / 0.547 ns
(baseline 2.326-2.822 / 0.495), max\_ss 4.183-5.163 / 0.980 (4.117-5.022
/ 0.905). \code{cpu_clk} latency/\allowbreak{}skew nom\_tt 2.067-3.424 /
1.357 ns (baseline 2.215-6.724 / 4.509), max\_ss 3.782-6.272 / 2.489
(4.056-11.897 / 7.841): the CPU domain tree is much better balanced, not
attributed.

\textbf{Scan-view hold on the CDC capture pins {[}M, nom\_tt, final
netlist + SPEF, \code{CVT_SDC_SCAN_VIEW=1}{]}.} 94 pins exempt in the
functional view (34 gray-pointer stage-0 D pins + 60 MBIST
tck-to-destination rows). Gray-pointer pins: \textbf{17 of 34 violate,
worst -3.083 ns} (j41m.5 measured 17/34 at -3.086..-3.240 on the earlier
netlist): unchanged by scan insertion. MBIST rows: \textbf{60 of 60
violate, -2.74 to -3.52 ns}, which is the worst scan-view hold of the
whole design. The masking of these flops in capture (user decision
2026-10-18) is \textbf{not implemented} and moves to bead \bead{j41m.15}
(user decision 2026-10-26): no ATPG exists yet and no hardware or
constraint change was made in this step. \textbf{Stage 2 is not
scan-hold-safe; no scan-hold-safe claim is made.}

\textbf{Max-cap {[}M,
\texttt{make\ sky130-\allowbreak{}soc-\allowbreak{}maxcap-\allowbreak{}check}{]}.}
With the waiver file as it stood after \prref{328}: 8 unwaived, 0 waived, 4
stale. User decisions of 2026-10-26 (\prref{330}), applied to
\src{pnr/sky130/soc/constraints/max_cap_waivers.json} and the plugin:

\begin{itemize}
\tightlist
\item
  Schema: a \code{driver} waiver may set
  \texttt{"stable":\allowbreak{}\ true} instead of \code{run_id} for a
  \textbf{hand-instantiated} driver whose name survives re-synthesis;
  the loader rejects \code{stable} on a tool-generated name
  (\code{_123_}, \code{netNNN}, \code{max_capNNN}, \code{wireNNN},
  \code{clkbuf_...}), the ceiling, corner list and justification still
  bind, and a stable waiver whose driver does not violate is reported
  STALE. The run-pinned path for synthesis names is unchanged. 22 new
  unit tests (\src{tb/tests/test_cvt_maxcap_waiver.py}, 81 pass).
\item
  \code{u-ref-cm-mux-clk-cpu-pll-sub} re-issued without \code{run_id}
  (stable driver \code{u_cpu_pll_sub.u_ref_cm.u_mux/X}, ceiling 0.45
  pF): scan run min/\allowbreak{}nom/\allowbreak{}max\_\allowbreak{}ss
  0.3659/\allowbreak{}0.3811/\allowbreak{}0.3912 pF vs limit 0.3397.
\item
  \code{sram-din0-22}: ceiling 0.35 -\textgreater{} 0.37 pF and
  \code{nom_ss} added (max\_ss 0.3617, nom\_ss 0.3395); new
  \code{sram-din0-23} (0.36 pF, max\_ss, measured 0.3478) and
  \code{sram-din0-26} (0.355 pF, max\_ss, measured 0.3426); the stale
  \code{sram-din0-29/30/31} deleted.
\item
  \code{_047147_/Y} (0.1388 pF vs 0.1383, +0.4 \%, synthesis-generated
  name, no macro load) stays \textbf{unwaived} and goes to bead
  \bead{e45j}.
\item
  Result on \code{RUN_2026-10-26_09-59-03} with \code{--strict-stale}: 7
  waived, \textbf{1 unwaived (\code{_047147_/Y})}, 0 stale. The
  previously shipped \code{RUN_2026-10-25_08-28-45} no longer passes the
  updated file (4 waived, 4 unwaived: \code{din0[29/30/31]} whose
  waivers were deleted, 2 stale \code{din0[23/26]}); the file now
  describes the scan configuration, the old run\textquotesingle s
  waivers are in git history.
\end{itemize}

\textbf{CDC budget check (j41m.5/\allowbreak{}y5mn) {[}M{]}.}
\code{check_cdc_timing_sky130.tcl} PASS at all 9 corners (22 groups,
every exception bound). max\_ss: \code{fifo_rd2wr} +9.891 and
\code{fifo_wr2rd} +9.875 ns of 12.5, MBIST groups +10.8 to +20.4 ns.
Negative control (every budget forced to 1 ps in a copy of the check
SDC): FAIL at nom\_tt and max\_ss, 22 of 22 groups violated.

\subsection{\texorpdfstring{\code{scan_en} is withheld until the
crypto scan-entry wipe is applied (bead \code{j41m.17}, user decision
2026-10-26)}{is withheld until the crypto scan-entry wipe is applied (bead , user decision 2026-10-26)}}\label{tm:ch:dft-architecture:1414-scan-en-is-withheld-until-the-crypto-scan-entry-wipe-is-applied-bead-j41m17-user-decision-2026-10-26}\origsec{14.14}

Closes the limit stated in 14.10 item 5 (crypto wipe section): the wipe
clear is functional D logic, the scan mux sits after it, so an edge with
\texttt{scan\_\allowbreak{}en\ =\allowbreak{}\ 1} loads \code{scan_in}
instead of the cleared value. Labels: \textbf{{[}R{]}} read from source,
\textbf{{[}M{]}} measured, \textbf{{[}A{]}} assumed.

\textbf{Design {[}R{]}.}

\begin{itemize}
\tightlist
\item
  \code{crypto_accel} has a new output \code{scan_wipe_done_o}
  (\src{rtl/periph/crypto_accel.sv}):
  \texttt{\textasciitilde{}((scan\_\allowbreak{}mode\_\allowbreak{}i\ \textasciicircum{}\ scan\_\allowbreak{}prev\_\allowbreak{}q)\ |\ (scan\_\allowbreak{}s1\_\allowbreak{}q\ \textasciicircum{}\ scan\_\allowbreak{}prev\_\allowbreak{}q)\ |\ wipe\_\allowbreak{}w)}
  with
  \texttt{wipe\_\allowbreak{}w\ =\allowbreak{}\ scan\_\allowbreak{}s2\_\allowbreak{}q\ \textasciicircum{}\ scan\_\allowbreak{}prev\_\allowbreak{}q}.
  It is 0 from the moment \code{scan_mode_i} differs from the settled
  value (the raw level is in the term, so it falls before any clock
  edge), stays 0 while the change walks
  \texttt{s1\ -\allowbreak{}>\ s2\ -\allowbreak{}>\ prev}, and rises
  just after edge 3, the same edge at which the \code{wipe_w}-qualified
  clears take effect. In functional mode it is a constant 1. It also
  drops for 3 \code{core_clk} edges after scan \textbf{exit} (the exit
  wipe), when \code{scan_en} is 0 anyway.
\item
  \code{soc_top} (\src{rtl/soc/soc_top.sv}, next to the
  \code{dft_scan_en} declaration):
  \texttt{dft\_\allowbreak{}scan\_\allowbreak{}en\ =\allowbreak{}\ dft\_\allowbreak{}scan\_\allowbreak{}en\_\allowbreak{}req\ \&\ crypto\_\allowbreak{}wipe\_\allowbreak{}done},
  where \code{dft_scan_en_req} is \code{dft_ctrl}\textquotesingle s
  \code{scan_en_o}. The \textbf{net named \code{dft_scan_en} is kept}
  for every consumer (the CPU macro\textquotesingle s \code{scan_en_i},
  the Stage 2 SE pins through \code{CVT_SCAN_EN_NET}, the STA queries,
  the existing tests); only its driver changed. \code{dft_ctrl},
  \code{dft_ctrl_ports} and \code{jtag_tap} are untouched. The gate
  covers pin mode and TAP-serial mode alike, because it sits after the
  \code{dft_ctrl} combiner.
\item
  Functional mode: \code{dft_scan_en} equals \code{dft_scan_en_req} once
  \code{scan_mode} has been 0 for 3 \code{core_clk} edges (the three
  edge-detect flops are unreset and track the level through reset).
  \code{scan_en} is 0 in functional mode, so the net stays 0.
\end{itemize}

\textbf{CDC / clocking rationale {[}R, A{]}.} No synchroniser was added,
deliberately.

\begin{enumerate}
\def\labelenumi{\arabic{enumi}.}
\tightlist
\item
  The gate matters only in scan mode, and in scan mode \code{core_clk},
  \code{cpu_core_clk} and every SE-pin flop run on the \textbf{one} test
  clock (\code{dft_clk_mux} on both PLL references, 14.2). The path
  \code{scan_prev_q} -\textgreater{} done -\textgreater{} SE is
  therefore a single-clock, single-cycle path launched at edge 3 and
  captured at the first shift edge, not a clock crossing. In functional
  mode the ANDed value is constant 0 and no scan flop reads it.
\item
  The only asynchronous ingredient is the raw \code{scan_mode_i} in the
  combinational term, the same quasi-static level that already drives
  the 256 key masks and that \src{tools/cdc/cdc_config.yml} lists under
  \code{reset_domains} (set before the first test-clock edge, not
  toggled while a clock runs). It is needed: with flops only,
  \code{done} would still be 1 at edge 1 and a tester that raised
  \code{scan_en} together with \code{scan_mode} could shift at edge 1.
  It can only move \code{done} 1 -\textgreater{} 0 (withholding), and
  \code{done} rises only through the synchronous \code{scan_prev_q}
  edge.
\item
  \code{cdc_sync} would add latency and a flop but fix nothing: its job
  is metastability from an asynchronous launch, and the launching flops
  here are on the test clock. A synchroniser per destination domain
  would also have to exist at 18,803 SE pins. Rejected.
\item
  Glitch behaviour: the three XOR terms share \code{scan_prev_q}, so
  when it toggles at edge 3 they all fall together and the NOR rises
  once, monotonically. Each of \code{scan_en_req}, \code{scan_mode},
  \code{scan_prev_q} changes alone in the protocol, and an AND with one
  changing input cannot glitch. The one ordering that is NOT protected:
  raising \code{scan_en} \textbf{before} \code{scan_mode} in pin mode
  (\texttt{scan\_\allowbreak{}en\ =\allowbreak{}\ 1} outside scan mode
  is outside the 14.1 contract) can produce a runt on the way down.
\item
  The netlist checks (\code{dft-scan-check}, the CDC gate) see no new
  register data input: \code{dft_scan_en} has no RTL sequential consumer
  before Stage 2 insertion (the macro\textquotesingle s \code{scan_en_i}
  is unused inside \code{rv32i_cpu_top}). {[}A until the checks are run,
  see the PR.{]}
\end{enumerate}

\textbf{Scan-entry protocol for the tester {[}R{]}.} Edge counts are
test-clock edges counted from the first edge after the mode change.

\begin{itemize}
\tightlist
\item
  \textbf{Pin mode (parallel).} Set
  \texttt{scan\_\allowbreak{}mode\_\allowbreak{}i\ =\allowbreak{}\ 1}
  with \texttt{scan\_\allowbreak{}en\_\allowbreak{}i\ =\allowbreak{}\ 0}
  and the test clock stopped. Apply \code{scan_clk_i} edges 1, 2, 3 (any
  \code{scan_en_i}; the gate forces the flops into capture). Edge 3
  clears the key and state. Keep
  \texttt{scan\_\allowbreak{}en\_\allowbreak{}i\ =\allowbreak{}\ 0}
  through \textbf{edge 4} (one spare edge so the SE tree has settled for
  the first shift; recommended, not enforced by hardware), then raise
  \code{scan_en_i}. Hardware guarantee: a shift can never occur at edges
  1 to 3, whatever the tester drives.
\item
  \textbf{TAP (serial).} After the \code{SCAN} instruction is updated,
  the test clock exists only in Capture-DR and Shift-DR. Do
  \textbf{three Capture-DR passes} (Select-DR, Capture-DR, Exit1-DR,
  Update-DR) before the first payload shift, then one more as the spare.
  Shift-DR edges taken earlier are accepted but are \textbf{capture}
  edges, not shifts: the gate holds \code{scan_en} low, so payload bits
  sent in those edges are lost. A tester that shifts at once is not
  harmed, it just has to send the data again.
\item
  \textbf{Not observable:} the wipe state is not readable over JTAG; the
  tester counts edges. A \code{scan_wipe_done} TDR bit is a possible
  follow-up, not part of this change.
\item
  \textbf{Exit} needs no rule: 3 functional \code{core_clk} edges after
  \code{scan_mode} falls wipe again, and the gate holds \code{scan_en}
  low meanwhile.
\end{itemize}

\textbf{Limits {[}R{]}.}

\begin{itemize}
\tightlist
\item
  \code{scan_out} shows each chain\textquotesingle s \textbf{last flop Q
  before any edge}. \textbf{Closed by gating (user decision 5a,
  2026-10-27)}: every SoC \code{scan_out_o[n]} is an \code{and2_2} of
  the chain\textquotesingle s last Q and the wipe-done net
  (\code{crypto_wipe_done}), built by \code{CVT.ScanInsert}
  (\code{CVT_SCAN_OUT_GATE_NET}) instead of the \code{buf_2}, so no
  chain end shows un-wiped crypto state before edge 3 whatever the chain
  order. Why gating and not "empty by chain order": chain 3 does end in
  a crypto flop (\code{u_crypto.g_sha.u_sha.block_i[182]}, i.e.
  \code{msg_q[182]}, in the fresh synthesis; \code{block_i[246]} in
  \code{RUN_2026-10-26_09-59-03}), so the gap was real and the end flop
  moves between insertions {[}M{]}. Chain ends of the fresh synthesis: 0
  \code{u_dma.linebuf[57][19]}, 1 \code{u_dma.linebuf[168][6]}, 2
  \code{u_bus.periph_axil_wdata[30]}, 3 crypto (above), 4
  \code{u_uart.baud_cnt_q[0]}, 5
  \code{u_apb_pll2_cdc.u_req_sync.sync_q[1]}; the two CPU-macro chains
  end in \code{u_core.u_csr.minstret_q[57]} and
  \code{u_mbist_ctrl.pause_cycles_q[21]} (no crypto in the macro, so
  they are not gated). The step fails, with no buffer fallback, if the
  net is absent or not singly driven;
  \texttt{check\_\allowbreak{}scan\_\allowbreak{}chains.\allowbreak{}py\ -\allowbreak{}-\allowbreak{}out-\allowbreak{}gate-\allowbreak{}net\ crypto\_\allowbreak{}wipe\_\allowbreak{}done}
  makes a missing gate an error (a buffer in its place is a tested
  mutant). \textbf{The net does not survive synthesis by itself} (it
  feeds only the \code{dft_scan_en} AND and is merged into anonymous
  logic: zero occurrences of the name in the netlist), so
  \code{soc_top.sv} carries \texttt{(*\ keep\ *)} on
  \code{crypto_wipe_done} (attribute only; with it the net survives with
  one driver) {[}M{]}.
\item
  A clock that never arrives (a stopped \code{scan_clk_i}, no
  Capture-DR) keeps \code{scan_en} low forever: the failure mode is
  "cannot shift", not "shifts early".
\item
  The key and \code{rk_q} are masked in scan mode and excluded from the
  chains regardless (14.2). The gate protects the \textbf{scanned}
  crypto flops (\code{s_q}, \code{msg_q}, SHA state,
  DOUT/\allowbreak{}DIGEST/\allowbreak{}IV).
\end{itemize}

\textbf{PD notes for the re-harden.} Measured in 14.14.1 below (the
\code{[A]} text that stood here is superseded).
\code{dft_scan_en_net_found} = 1 and the net \code{crypto_wipe_done}
survives with \texttt{(*\ keep\ *)}.

\textbf{Verification {[}M, 2026-10-26, local, Verilator{]}.}
\code{test_soc_scan_en_gate} 12/12 (new), \code{test_crypto} 53/53 (2
new), \code{soc_dft_scan} 13/13 and \code{soc_jtag} 13/13 after two
tests were changed to follow the protocol (they passed in the full suite
only because an earlier test left the wipe pipeline settled; each failed
alone). Mutants: gate removed, flops-only done, \code{wipe_w} term
dropped, done one edge early / late, gate tied 0 / tied 1 / inverted:
all killed, none equivalent; the 25-mutant \code{j41m.6} campaign still
25/25. \code{soc_all_ci} from clean: 911 passed, 0 failed, 54 skipped.
\code{dft-scan-check} Sky130 and ASAP7: 0 violations. CDC gate
(\texttt{make\ -\allowbreak{}C\ sim\ cdc\ CDC\_\allowbreak{}STRICT=\allowbreak{}1}):
PASS, counts identical to \code{f12c600}. Struct-select checker: 0
sites. Verilator zero-delay simulation cannot show analog glitches; the
glitch claim in the CDC rationale is by construction, and the tests
check that each signal changes once per phase. The "no residue" test
samples the scanned crypto flops at the first cycle \code{dft_scan_en}
is 1 (what a scan flop would shift out); the post-fill netlist was not
simulated.

\subsection{\texorpdfstring{Re-harden with the crypto
\code{scan_en} withhold and \code{scan_out_o} gating:
\code{RUN_2026-10-27_05-40-27} vs \code{RUN_2026-10-26_09-59-03} (bead
\code{j41m.17},
2026-10-27)}{Re-harden with the crypto  withhold and  gating:  vs  (bead , 2026-10-27)}}\label{tm:ch:dft-architecture:14141-re-harden-with-the-crypto-scan-en-withhold-and-scan-out-o-gating-run-2026-10-27-05-40-27-vs-run-2026-10-26-09-59-03-bead-j41m17-2026-10-27}\origsec{14.14.1}

Labels as above. Command:
\texttt{make\ -\allowbreak{}C\ pnr\ librelane-\allowbreak{}sky130-\allowbreak{}soc-\allowbreak{}noklayout}
(rccal default, flow \code{CVTClassicScan}, \code{CVT_SYNTH_GATE=warn},
KLayout DRC skipped as in the baseline), CPU macro = the committed scan
macro of 14.12 (not re-hardened; GDS/spice byte-identical to
\code{RUN_2026-10-26_05-39-04}). Wall time about 5.6 h (baseline 5.2 h).
\textbf{Result: flow exit with deferred errors Magic DRC, setup at
max\_ss (NPU, accepted, bead \bead{7dhp}) and max-cap (10 unwaived;
resolved by the waiver set below, user decision 2026-10-27). The
post-route CDC budget check failed at the three ss corners on the new
path and is fixed by a functional-view case analysis in the CDC-check
SDC (below).} Not a single-variable comparison: the RTL (wipe-done
output, AND into \code{dft_scan_en}, \texttt{(*\ keep\ *)}), the gated
\code{scan_out_o}, and the placement all changed.

{\small\setlength{\tabcolsep}{4pt}

{
\begin{longtable}[]{@{}
  >{\raggedright\arraybackslash}p{\dimexpr(\linewidth-4\tabcolsep)*447/1000\relax}
  >{\raggedright\arraybackslash}p{\dimexpr(\linewidth-4\tabcolsep)*267/1000\relax}
  >{\raggedright\arraybackslash}p{\dimexpr(\linewidth-4\tabcolsep)*267/1000\relax}@{}}
\toprule\noalign{}
\begin{minipage}[b]{\linewidth}\raggedright
\end{minipage} & \begin{minipage}[b]{\linewidth}\raggedright
this run
\end{minipage} & \begin{minipage}[b]{\linewidth}\raggedright
baseline
\end{minipage} \\
\midrule\noalign{}
\endhead
\bottomrule\noalign{}
\endlastfoot
LVS errors / routing DRC & 0 / 0 & 0 / 0 \\
Magic DRC (nwell.4 + met4.4a) & 10,959 (10,927 + 32) & 11,003 (10,971 +
32) \\
setup violators in the flow STA, max\_ss (other 8 corners) & \textbf{6},
WNS -0.153 ns, TNS -0.603 ns (0) & 0 (0) \\
hold violators, 9 corners & 0 & 0 \\
\code{core_clk} setup worst nom\_tt / nom\_ss / max\_ss (ns), replica &
+5.588 / +0.132 / \textbf{-0.153} & +5.671 / +0.156 / +0.002 \\
\code{cpu_clk} setup worst nom\_tt / max\_ss (ns), replica & +14.70 /
+5.19 & +15.92 / +7.76 \\
hold worst \code{core_clk} nom\_tt / max\_ss / max\_ff (ns) & +0.405 /
+0.231 / +0.268 & +0.415 / +0.223 / +0.258 \\
CPU macro input pins hold nom\_tt / max\_ss / max\_ff (ns) & +0.394 /
+1.272 / +0.169 & +0.568 / +1.457 / +0.274 \\
max-slew violators, max\_ss / nom\_tt (checker skipped in the flow) &
1,067 / 145 & 525 / 84 \\
max-cap violators in the flow STA, max\_ss & 6 & 5 \\
max-cap with the waiver file, \code{--strict-stale} & 10 unwaived, 3
stale & 1 unwaived, 0 stale \\
antenna pins / nets (diodes) & 382 / 290 (20,500) & 368 / 286
(17,642) \\
power / util / stdcell area & 119.10 mW / 46.64 \% / 2,070,340 um2 &
119.01 mW / 46.55 \% / 2,052,900 um2 \\
instances / timing-repair bufs / hold bufs & 332,973 / 28,726 / 213 &
328,829 / 28,045 / 111 \\
wirelength & 9.364 M & 9.214 M \\
\end{longtable}
}

}

The flow\textquotesingle s own setup figure is still the literal 0.000
latch artifact (14.11) at the other eight corners; the per-corner
replica STA (\src{pnr/scripts/sta_load_sky130_soc.tcl} with the
run\textquotesingle s SPEF, analysis SDC below) is the measurement.
Replica setup per corner, \code{core_clk} / \code{cpu_clk} (ns): nom\_tt
+5.588/\allowbreak{}+14.70, min\_tt +5.796/\allowbreak{}+15.18, max\_tt
+5.365/\allowbreak{}+14.24, nom\_ss +0.132/\allowbreak{}+5.87, min\_ss
+0.390/\allowbreak{}+6.60, max\_ss \textbf{-0.153}/+5.19, nom\_ff
+7.73/\allowbreak{}+17.89, min\_ff +7.91/\allowbreak{}+18.29, max\_ff
+7.55/\allowbreak{}+17.53.

\textbf{The 6 max\_ss setup violators {[}M{]}} are 6 endpoints
(\code{_169097_/D} -0.153, \code{_169104_/D} -0.132, \code{_169095_/D}
-0.110, \code{_169106_/D} -0.095, \code{_169111_/D} -0.079,
\code{_169116_/D} -0.033), all launched by the NPU weight-memory SRAM
macro outputs (\code{u_npu.g_on.u_wmem.u_sram_macro/dout0[12]},
\code{[8]}, a falling-edge half-cycle launch at 12.5 ns) into
\code{npu_weight_mem} pipeline flops (\code{sdfrtp_1}). They are
\textbf{not} the new crypto paths: the worst SCE endpoint at max\_ss is
+0.671 ns. User decision 2026-10-27: accepted, tracked in bead
\bead{7dhp}. The same class was seen in \code{RUN_2026-10-25_15-59-23}
(dout0{[}9{]}, -0.036 ns). The mid-flow resizer table (WNS -1.006 ns on
\code{_156429_/SCE}, seen in the log by the coordinator, not re-read
here) is a pre-route estimate {[}A{]}; the post-route SCE slack is
positive at all nine corners {[}M{]}.

\textbf{5(b) the two new paths, setup and hold, all nine corners {[}M,
replica STA, launch \code{core_clk} from the crypto edge-detect
flops{]}.} Analysis-only SDC (\src{tools/dft/make_scan_analysis_sdc.py}
writes it; queries in
\src{pnr/scripts/sta_scan_gate_paths_sky130_soc.tcl}; = the repo SDC
without the
\texttt{set\_\allowbreak{}false\_\allowbreak{}path\ -\allowbreak{}to\ scan\_\allowbreak{}out\_\allowbreak{}o}
line, plus
\texttt{set\_\allowbreak{}output\_\allowbreak{}delay\ 0\ -\allowbreak{}clock\ core\_\allowbreak{}clk\ scan\_\allowbreak{}out\_\allowbreak{}o};
the repo SDC is unchanged and still false-paths \code{scan_out_o}, so
the flow STA does not time that path).

{\small\setlength{\tabcolsep}{4pt}

{
\begin{longtable}[]{@{}
  >{\raggedright\arraybackslash}p{\dimexpr(\linewidth-6\tabcolsep)*186/1000\relax}
  >{\raggedright\arraybackslash}p{\dimexpr(\linewidth-6\tabcolsep)*258/1000\relax}
  >{\raggedright\arraybackslash}p{\dimexpr(\linewidth-6\tabcolsep)*350/1000\relax}
  >{\raggedright\arraybackslash}p{\dimexpr(\linewidth-6\tabcolsep)*186/1000\relax}@{}}
\toprule\noalign{}
\begin{minipage}[b]{\linewidth}\raggedright
corner
\end{minipage} & \begin{minipage}[b]{\linewidth}\raggedright
done -\textgreater{} \code{dft_scan_en} -\textgreater{} SE setup
\end{minipage} & \begin{minipage}[b]{\linewidth}\raggedright
done -\textgreater{} AND -\textgreater{} \code{scan_out_o} setup (worst
of 6)
\end{minipage} & \begin{minipage}[b]{\linewidth}\raggedright
done -\textgreater{} SE hold
\end{minipage} \\
\midrule\noalign{}
\endhead
\bottomrule\noalign{}
\endlastfoot
nom\_tt / min\_tt / max\_tt & +11.82 / +12.16 / +11.40 & +17.72 / +17.98
/ +17.47 & +1.665 / +1.863 / +1.531 \\
nom\_ss / min\_ss / max\_ss & +1.389 / +1.997 / \textbf{+0.671} & +11.93
/ +12.37 / +11.54 & +4.421 / +4.661 / +4.298 \\
nom\_ff / min\_ff / max\_ff & +16.01 / +16.25 / +15.69 & +19.97 / +20.15
/ +19.79 & +0.983 / +1.148 / +0.871 \\
\end{longtable}
}

}

Both pass as \code{core_clk}-to-\code{core_clk} paths. The max\_ss done
-\textgreater{} SE path has 23.2 ns of data delay in a 25 ns cycle: 3
hold-repair \code{dlygate4sd3_1} (4.2 ns) after the launch flop, then
the 28-buffer \code{max_cap} tree. The resizer cloned the enable AND
(\texttt{and2\_\allowbreak{}2\ clone13} drives 13 SCE pins from the same
inputs as
\texttt{and2\_\allowbreak{}4\ \_\allowbreak{}090536\_\allowbreak{}});
\code{check_scan_chains.py} now treats a gate with the same function and
input nets as the same root.

\textbf{CDC budget check: failed at the ss corners, fixed in the check
SDC {[}M{]}.} First result of
\texttt{make\ -\allowbreak{}C\ pnr\ check-\allowbreak{}cdc-\allowbreak{}timing-\allowbreak{}sky130-\allowbreak{}soc\ CDC\_\allowbreak{}CORNERS=\allowbreak{}all}
on this run: PASS at the six tt/ff corners, FAIL at nom\_ss, min\_ss and
max\_ss, one finding each: the domain-wide fallback
\texttt{core\_\allowbreak{}clk\ -\allowbreak{}>\ cpu\_\allowbreak{}clk}
(budget 0.9 x 25 ns = 22.5 ns). The path is the same done chain:
\code{_154808_} (crypto edge-detect flop, \code{core_clk})
-\textgreater{} 3 hold dlygates -\textgreater{} ... -\textgreater{}
\code{max_cap3478} -\textgreater{} \code{_156258_/SCE}, a scan flop of
chain 5 (the 981 \code{cpu_core_clk} flops): arrival 23.159 ns against
21.578 ns required, \textbf{-1.581 ns at max\_ss}. In the functional
clocking those 981 SE pins are a
\texttt{core\_\allowbreak{}clk\ -\allowbreak{}>\ cpu\_\allowbreak{}core\_\allowbreak{}clk}
crossing, but \code{dft_scan_en} is constant 0 in functional mode and in
scan mode both domains run on one test clock, so it is not a functional
crossing (the baseline had no such path: the enable came only from the
TAP and the pin).

User decision 2026-10-27: fix it in the CDC-check SDC with a
functional-view case analysis, not in RTL and not with an \code{*/SCE}
exemption. \code{sky130_soc_cdc_check.sdc} section 2b pins
\texttt{scan\_\allowbreak{}mode\_\allowbreak{}i\ =\allowbreak{}\ 0},
\texttt{scan\_\allowbreak{}en\_\allowbreak{}i\ =\allowbreak{}\ 0} and
the driver pin of the TAP serial-select flop
\code{u_dft_ctrl.tap_serial} = 0.
\textbf{\texttt{scan\_\allowbreak{}mode\_\allowbreak{}i\ =\allowbreak{}\ 0}
alone does not work {[}M{]}}: it leaves the same finding
(\code{dft_scan_en} is
\texttt{(tap\_\allowbreak{}serial\ ?\ tap\_\allowbreak{}scan\_\allowbreak{}en\ :\allowbreak{}\ pin\_\allowbreak{}scan\_\allowbreak{}en)\ \&\ crypto\_\allowbreak{}wipe\_\allowbreak{}done};
\code{scan_en_o} is not gated by scan\_mode and the TAP can raise
scan\_mode itself, so the request is not constant until the select and
the pin enable are pinned too). With the three, the request is a
constant 0, so the AND and its resizer clone are constant and the SE
tree drops out. Results on \code{RUN_2026-10-27_05-40-27} (no new PD
run; SDC-only change):

\begin{itemize}
\tightlist
\item
  \textbf{All nine corners PASS} (22 groups plus 2 domain-wide
  fallbacks, 36 rows per corner including the 4 skew rows).
\item
  \textbf{Before/\allowbreak{}after row comparison, 9 corners x 36 rows
  {[}M{]}}: the only status change is \code{domain_wide_core2cpu}
  VIOLATED -\textgreater{} MET at nom\_ss, min\_ss and max\_ss. Nothing
  else changes status; the one other numeric change is
  \code{apb_dbg_rst_status} slack up by 0.003 to 0.021 ns (the case
  analysis simplifies the reset mux in its cone), e.g. 17.155
  -\textgreater{} 17.165 ns at max\_ss. The domain-wide
  \texttt{core\_\allowbreak{}clk\ -\allowbreak{}>\ cpu\_\allowbreak{}clk}
  group is not emptied: its worst remaining path at max\_ss is a
  \code{core_clk} flop into a CPU-macro input, +8.961 ns.
\item
  \textbf{Negative control with the new SDC} (every budget forced to 1
  ps, nom\_tt and max\_ss): FAIL with \textbf{22 of 22 groups violated},
  the same 22 as before the change, so no MBIST, FIFO, bridge or PMU
  group is hidden by the case analysis.
\item
  \textbf{Scan view unaffected {[}M{]}}: with \code{CVT_SDC_SCAN_VIEW=1}
  the section is skipped (log line
  \texttt{CDC-\allowbreak{}FUNCTIONAL-\allowbreak{}VIEW:\allowbreak{}\ .\allowbreak{}.\allowbreak{}.\allowbreak{}\ NOT\ applied})
  and the original finding returns at max\_ss (-1.581 ns), i.e. the
  scan-view check still reports this path rather than suppressing it.
  \code{sky130_soc.sdc} is not modified.
\item
  \textbf{Recommendation for \code{sky130_soc.sdc} (not applied, a
  separate decision)}: do not add the \code{tap_serial} case analysis
  there yet. That SDC already pins \code{scan_mode_i}/\code{scan_en_i},
  and it declares \code{core_clk}/\code{cpu_clk} asynchronous, so the
  flow STA does not time the 981 cpu-domain SE pins but does time the
  core-domain ones: done -\textgreater{} SE is \textbf{+0.671 ns at
  max\_ss} and it is the flow\textquotesingle s only look at the SE
  tree. A constant \code{dft_scan_en} would remove all of them from the
  flow STA and from the post-route repair, which sized the tree and
  inserted its hold dlygates against exactly these paths. The proper
  home is the scan-mode SDC (\code{j41m.8} deliverable), which will time
  \code{dft_scan_en} on the test clock.
\end{itemize}

\textbf{CDC negative control, before the SDC change {[}M{]}.} Every
per-crossing and domain-wide budget forced to 1 ps in a copy of the
check SDC (\code{CDC_CHECK_MBIST_BUDGET_NS=0.001} for the MBIST rows):
FAIL with 22 findings at nom\_tt and at max\_ss, so all 22 groups are
bound and violated; the check is not vacuous. After the SDC change the
same control still gives 22 of 22 (above).

\textbf{Chain-gate check on the final netlist {[}M{]}.}
\texttt{check\_\allowbreak{}scan\_\allowbreak{}chains.\allowbreak{}py\ .\allowbreak{}.\allowbreak{}.\allowbreak{}\ -\allowbreak{}-\allowbreak{}out-\allowbreak{}gate-\allowbreak{}net\ crypto\_\allowbreak{}wipe\_\allowbreak{}done\ -\allowbreak{}-\allowbreak{}expect-\allowbreak{}chains\ 6\ -\allowbreak{}-\allowbreak{}expect-\allowbreak{}unscanned\ 922\ -\allowbreak{}-\allowbreak{}expect-\allowbreak{}length\ 3565,\allowbreak{}3565,\allowbreak{}3565,\allowbreak{}3565,\allowbreak{}3562,\allowbreak{}981}
on \code{49-openroad-fillinsertion/soc_top.nl.v}: exit 0, all six
\code{scan_out_o[n]} are an AND2 of the chain\textquotesingle s last
flop and \code{crypto_wipe_done}, 18,803 SCE pins trace to one root
(\code{dft_scan_en}, counting the resizer clone), the 922 unscanned
flops are explained (664 \code{tck_i}, 2 \code{scan_clk_i}, 256 crypto
keys).

\textbf{Max-cap, \code{--strict-stale} {[}M{]}: 10 unwaived, 3 stale;
investigation (user request 2026-10-27, nothing edited).}

{\footnotesize\setlength{\tabcolsep}{3pt}

{
\begin{longtable}[]{@{}
  >{\raggedright\arraybackslash}p{\dimexpr(\linewidth-4\tabcolsep)*327/1000\relax}
  >{\raggedright\arraybackslash}p{\dimexpr(\linewidth-4\tabcolsep)*327/1000\relax}
  >{\raggedright\arraybackslash}p{\dimexpr(\linewidth-4\tabcolsep)*327/1000\relax}@{}}
\toprule\noalign{}
\begin{minipage}[b]{\linewidth}\raggedright
net (driver)
\end{minipage} & \begin{minipage}[b]{\linewidth}\raggedright
cap vs limit (pF)
\end{minipage} & \begin{minipage}[b]{\linewidth}\raggedright
what is on it {[}M{]}
\end{minipage} \\
\midrule\noalign{}
\endhead
\bottomrule\noalign{}
\endlastfoot
\code{u_cpu_pll_sub.u_ref_cm.mux_clk} (\code{u_mux/X}, hand-instantiated
mux2\_4) & 0.4750 / 0.5077 / 0.5186
(min/\allowbreak{}nom/\allowbreak{}max\_\allowbreak{}ss) vs ceiling 0.45
(limit 0.3397) & mux2\_4 -\textgreater{} buf\_16 + 11 antenna diodes:
\textbf{identical loads in both runs}; coupling, see below \\
\code{_107570_/Y} (a21oi\_4) & 0.1580 (max\_ss), 0.1459 (nom\_ss) vs
0.1383 & 15 loads, 13 are antenna diodes \\
\code{hold19396/X}, \code{hold19401/X}, \code{hold19961/X}
(dlygate4sd3\_\allowbreak{}1) & 0.1113, 0.1046, 0.1020 vs 0.1007 & hold
dlygates on scan links: 1 SCD pin + 13, 6, 12 diodes \\
\code{_076179_/X} (or4\_4) & 0.3371 vs 0.3360 (+0.3 \%) & 7 functional
loads, no diodes (wire only) \\
\end{longtable}
}

}

\textbf{Why \code{mux_clk} grew 33 \% {[}M, SPEF, nom RC corner{]}.} Not
loads, not diodes: the net still drives one buf\_16 and 11 diodes.
Routed length grew only 6 \% (2,620 -\textgreater{} 2,775 um; met4 2,492
-\textgreater{} 2,441, met3 95 -\textgreater{} 303; the mux moved from
(3852, 2557) to (3646, 2505) while the root buffer stayed at (3954, 49),
2.5 mm away). The SPEF total went 0.356 -\textgreater{} 0.483 pF (max RC
0.366 -\textgreater{} 0.494): ground cap fell 0.290 -\textgreater{}
0.222 but \textbf{coupling rose 0.066 -\textgreater{} 0.261 pF, of which
0.158 pF is to \code{cpu_core_clk}}, the CPU domain\textquotesingle s
clock net, now routed alongside it for a long run
(baseline\textquotesingle s largest aggressor: 0.012 pF). So it is a
routing-adjacency accident between two clock nets, not an RTL or config
effect, and a cap-margin knob cannot reach it (hand-instantiated clock
driver, as 14.13 and \code{e45j} already found). A ceiling is exposed to
the next placement.

\textbf{Why max-slew doubled, 525 -\textgreater{} 1,067 violating pins
at max\_ss {[}M{]}.} 930 of the 1,067 are antenna-diode pins (baseline
419), i.e. +511 of the +542. They sit on few nets: 31 violating nets
(baseline 24), the worst driven by a \code{wire}-repair clkbuf\_16 with
48 to 55 diodes plus 4 \code{mux2_1.A1} loads. Nets carrying at least 30
diodes: \textbf{75 (baseline 45)}; by root signal:
\code{cpu_bridge_s_awaddr[]} 22 (10), \code{cpu_bridge_s_araddr[]} 8
(5), \code{cpu_bridge_s_awburst[]}/\code{awlen[]}/\code{arlen[]} 6 (9),
\code{mbist_rst_tck_n} 7 (7), \code{u_apb_pll_cdc.m_rst_sync_n} 2,
\code{core_rst_n} 2, plus 12 anonymous functional nets (4). The AXI
address/\allowbreak{}control bus leaving the CPU macro is the main
contributor. The cause is the diode count: antenna repair (step 40)
inserted \textbf{20,500 diodes (baseline 17,642, +16 \%)} against 1,411
violating nets before repair (1,358, +3.9 \%) and left 243 (219); by net
class the +2,858 is functional-logic nets +2,907, scan-link nets +180,
nets that also feed an SCD +4, single-flop-D nets -231, diode-only nets
-2. Other violating pins: \code{mux2.A1} 72, \code{sdfrtp.RESET_B} 16,
\code{o21ai.A1} 5, \code{or4.C} 4.

\textbf{Are the hold dlygates on scan links with 6 to 13 diodes the
common cause? {[}M{]} No.} They are the cause of those three max-cap
nets only (a dlygate4sd3\_\allowbreak{}1 has the lowest limit, 0.1007
pF). They are not new: the baseline has 98 scan-link nets driven by a
hold dlygate with diodes, up to 17 diodes (this run: 90, up to 13);
total scan-link diodes 637 -\textgreater{} 817. They are not the slew
driver (the slew nets are the AXI bus and reset nets above). The common
factor across the table is the antenna-diode pin capacitance on long
nets (\code{_107570_} and the three dlygates: 6 to 13 diodes each; the
14 \% excess of \code{_107570_} is diode pin cap), with two exceptions
that have a different cause: \code{mux_clk} (coupling) and
\code{_076179_} (plain wire, +0.3 \%).

\textbf{Applied (user decision 2026-10-27) {[}M{]}.}
\texttt{make\ -\allowbreak{}C\ pnr\ sky130-\allowbreak{}soc-\allowbreak{}maxcap-\allowbreak{}check\ SKY130\_\allowbreak{}SOC\_\allowbreak{}MAXCAP\_\allowbreak{}RUN=\allowbreak{}.\allowbreak{}.\allowbreak{}.\allowbreak{}/\allowbreak{}RUN\_\allowbreak{}2026-\allowbreak{}10-\allowbreak{}27\_\allowbreak{}05-\allowbreak{}40-\allowbreak{}27\ STRICT\_\allowbreak{}STALE=\allowbreak{}1}
now reports \textbf{10 waived, 0 unwaived, 0 stale: PASS}. The
wire-length arm
(\texttt{GRT\_\allowbreak{}DESIGN\_\allowbreak{}REPAIR\_\allowbreak{}MAX\_\allowbreak{}WIRE\_\allowbreak{}LENGTH\ =\allowbreak{}\ 1500},
proposal 1 below) is \textbf{not} run here; it is a follow-up under
\code{e45j} after the merge. The waiver file \code{max_cap_waivers.json}
is now:

\begin{itemize}
\tightlist
\item
  \code{u-ref-cm-mux-clk-cpu-pll-sub} (stable driver): ceiling 0.45
  -\textgreater{} \textbf{0.55 pF}, justification extended with the
  coupling finding (0.158 pF to \code{cpu_core_clk}).
\item
  \code{sram-din0-22/23/26-collar-mux}: \textbf{deleted} (stale).
\item
  \code{hold-dlygate-scd-link} (new, \textbf{load-class}):
  \code{cvt_maxcap_waiver.py} gains a \code{driver_class} key, "driver
  cell type starts with \code{sky130_fd_sc_hd__dlygate4sd3_} and its
  only non-antenna-diode load is one \code{sky130_fd_sc_hd__sdf*}
  \code{SCD} pin", ceiling \textbf{0.115 pF}, corners
  min/\allowbreak{}nom/\allowbreak{}max\_\allowbreak{}ss. The class is
  proven from the final netlist per violation (a second pass over the
  netlist collects every load of the violating nets); another cell type,
  any other functional load or a macro leaves the violation unwaived,
  and a class waiver that matches nothing is STALE. Covers
  \code{hold19396}, \code{hold19401}, \code{hold19961}.
\item
  \code{gen-a21oi-107570-run-2026-10-27} (run-pinned to
  \code{RUN_2026-10-27_05-40-27}, \code{_107570_/Y}): ceiling 0.160 pF
  (measured 0.1580 max\_ss, 0.1459 nom\_ss), corners nom\_ss and
  max\_ss.
\item
  \code{gen-or4-076179-run-2026-10-27} (run-pinned, \code{_076179_/X}):
  ceiling 0.338 pF (measured 0.3371), corner max\_ss. The two run-pinned
  entries go stale on the next re-synthesis by design.
  \src{tb/tests/test_cvt_maxcap_waiver.py} 100 pass: the class waiver is
  covered for accept, reject on a second non-diode load of another kind
  or a second SCD load or a non-SCD pin, reject on a macro load, reject
  on a different cell type, ceiling and corner binding, stale reporting
  (also when the class does not match), malformed entries, and the empty
  list staying strict; a mutant that makes the class test always true is
  killed by 5 tests.
\end{itemize}

\textbf{Proposals as made before the decision (proposal 1 deferred to
\code{e45j}).}

\begin{enumerate}
\def\labelenumi{\arabic{enumi}.}
\tightlist
\item
  \emph{Config, single variable:}
  \texttt{GRT\_\allowbreak{}DESIGN\_\allowbreak{}REPAIR\_\allowbreak{}MAX\_\allowbreak{}WIRE\_\allowbreak{}LENGTH\ =\allowbreak{}\ 1500}
  (existing \code{config_rccal_wl1500.json}) on this RTL vintage.
  Expected {[}A, from the \code{e45j} section 6 arm on
  \code{RUN_2026-10-08_23-12-19}: max\_ss slew 423 -\textgreater{} 102,
  antenna 205/273 -\textgreater{} 212/253,
  setup/\allowbreak{}hold/\allowbreak{}LVS/\allowbreak{}DRC unchanged,
  max-cap unchanged{]}: fewer long single-driver wires, so fewer
  diode-heavy nets and a lower slew count. It will not move
  \code{mux_clk}, and by \code{e45j} it did not change max-cap;
  confidence medium for slew, low for the cap list. A second candidate
  is \code{GRT_ANTENNA_MARGIN} 50 -\textgreater{} 35 (bead \bead{58q}
  data: 25 / 8 iterations gave 263 antenna pins vs 182 at 50 / 20, with
  max-slew 2,380 vs 3,199), a Pareto trade that gives up antenna count;
  try it after the wire-length arm.
\item
  \emph{Waivers, user\textquotesingle s call:} (a)
  \code{u-ref-cm-mux-clk-cpu-pll-sub} ceiling 0.45 -\textgreater{} 0.55
  pF (observed 0.5186; one neighbour added 0.13 pF this run, so a 2 \%
  margin would be fragile); (b) delete \code{sram-din0-22-collar-mux},
  \code{sram-din0-23-collar-mux}, \code{sram-din0-26-collar-mux}
  (stale); (c) the five tool-named nets cannot take \code{stable}
  waivers (the schema rejects generated names), so they would need
  run-pinned waivers that go stale on the next run. A stable alternative
  is a \textbf{load-class key} in \code{cvt_maxcap_waiver.py}: "driver
  is a hold dlygate whose only non-diode load is an \code{sdf*/SCD} pin,
  ceiling 0.115 pF" (covers \code{hold19396/19401/19961}, max observed
  0.1113), roughly 40 lines plus tests. \code{_107570_} (+14 \%) and
  \code{_076179_} (+0.3 \%) have no stable key.
\end{enumerate}

\section{\texorpdfstring{Stage 1b specification and implementation
record (bead \code{j41m.3}): JTAG
TAP}{Stage 1b specification and implementation record (bead ): JTAG TAP}}\label{tm:ch:dft-architecture:15-stage-1b-specification-and-implementation-record-bead-j41m3-jtag-tap}\origsec{15}

Scope: the TAP, its registers, the TAP-driven \code{dft_ctrl},
\code{soc_top} integration, checks. \textbf{No} scan insertion, ATPG,
MBIST or CPU-macro change. Decisions already taken: section 0 (decisions
5-7) and the user decisions of 2026-10-11 (N1 non-JEDEC placeholder
IDCODE, N2 compliant-minus-boundary-register, N4 TAP flops unscanned).
Labels as before. This section was written \textbf{before} the RTL as
its specification; the status table at its end is updated as each step
lands.

\subsection{Modules and ports}\label{tm:ch:dft-architecture:151-modules-and-ports}\origsec{15.1}

{\footnotesize\setlength{\tabcolsep}{3pt}

{
\begin{longtable}[]{@{}
  >{\raggedright\arraybackslash}p{\dimexpr(\linewidth-4\tabcolsep)*112/1000\relax}
  >{\raggedright\arraybackslash}p{\dimexpr(\linewidth-4\tabcolsep)*289/1000\relax}
  >{\raggedright\arraybackslash}p{\dimexpr(\linewidth-4\tabcolsep)*579/1000\relax}@{}}
\toprule\noalign{}
\begin{minipage}[b]{\linewidth}\raggedright
Module
\end{minipage} & \begin{minipage}[b]{\linewidth}\raggedright
File
\end{minipage} & \begin{minipage}[b]{\linewidth}\raggedright
Role
\end{minipage} \\
\midrule\noalign{}
\endhead
\bottomrule\noalign{}
\endlastfoot
\code{jtag_tap} & \src{rtl/soc/dft/jtag_tap.sv} & 16-state FSM, IR, TDR
bank, TDO, all in the TCK domain. No functional-clock logic. \\
\code{dft_ctrl} & \src{rtl/soc/dft/dft_ctrl.sv} & wraps \code{jtag_tap}
+ the 1a pin leg (\code{dft_ctrl_ports}, \textbf{unchanged}) + the TCK
scan-clock gate + the combiner. Same five outputs as
\code{dft_ctrl_ports} and the same four test inputs, plus the JTAG pins.
\code{soc_top} instantiates this instead of \code{dft_ctrl_ports}. \\
\end{longtable}
}

}

\code{soc_top} gains (all raw, no pad ring): \code{tck_i}, \code{tms_i},
\code{tdi_i}, \code{trst_ni} (in); \code{tdo_o}, \code{tdo_oe_o} (out).
Idle values used by every existing suite: \code{tck_i=0},
\code{tms_i=1}, \code{tdi_i=0}, \textbf{\code{trst_ni=0} (TAP held in
reset)}. \code{TDO} is never tri-stated in RTL (no \code{inout}); the
chip top builds the tristate from \code{tdo_oe_o}.

\subsection{TAP rules (IEEE 1149.1 where implemented)}\label{tm:ch:dft-architecture:152-tap-rules-ieee-11491-where-implemented}\origsec{15.2}

\begin{itemize}
\tightlist
\item
  16 states with the standard 4-bit encodings (TLR=4\textquotesingle hF,
  RTI=4\textquotesingle hC, SelDR=4\textquotesingle h7,
  CapDR=4\textquotesingle h6, ShDR=4\textquotesingle h2,
  Ex1DR=4\textquotesingle h1, PauDR=4\textquotesingle h3,
  Ex2DR=4\textquotesingle h0, UpDR=4\textquotesingle h5,
  SelIR=4\textquotesingle h4, CapIR=4\textquotesingle hE,
  ShIR=4\textquotesingle hA, Ex1IR=4\textquotesingle h9,
  PauIR=4\textquotesingle hB, Ex2IR=4\textquotesingle h8,
  UpIR=4\textquotesingle hD). State flops on TCK \textbf{rising} edge.
  \code{state_o[3:0]} is exported by \code{jtag_tap} for tests.
\item
  Reset: \code{trst_ni} low asynchronously puts every TAP flop in its
  reset state; release goes through one \code{cdc_reset_sync} (TCK
  domain, 2 stages, \code{scanmode_i=0}) so deassertion is synchronous
  to TCK. Five TCK rising edges with TMS=1 reach TLR from any state. The
  TAP is reset by nothing else (not \code{rst_n_i}, not any
  PLL/\allowbreak{}PMU/\allowbreak{}WDT reset). While the state flop is
  TLR, IR returns to IDCODE and every TDR update register and every
  output returns to its default at the next falling edge.
\item
  IR: 4 bits. Shift register on the rising edge (shifts in Shift-IR, LSB
  first, TDI enters the MSB), loaded with \textbf{\texttt{4'b0001}} in
  Capture-IR (the mandatory \code{...01}). The update register
  \code{ir_q} and the registered decodes change on the \textbf{falling}
  edge in Update-IR. IR reset value = IDCODE.
\item
  DR: the register selected by \code{ir_q} is loaded in Capture-DR,
  shifts in Shift-DR (rising edge), and its update stage (where it has
  one) changes on the falling edge in Update-DR. Selection is a one-hot
  AND-OR (no runtime-indexed read).
\item
  TDO: a flop launched on the \textbf{falling} edge; carries
  \code{ir_shift[0]} in Shift-IR, the selected DR\textquotesingle s LSB
  in Shift-DR, and \textbf{0 in every other state}. \code{tdo_oe_o} is a
  falling-edge flop that is 1 exactly in Shift-IR and Shift-DR.
\item
  Every \code{jtag_tap} output that reaches the SoC (\code{scan_mode_o},
  \code{serial_o}, \code{parallel_o}, \code{rst_assert_o},
  \code{scan_en_o}, \code{tdo_o}, \code{tdo_oe_o}) is a \textbf{direct
  falling-edge flop output computed from next-state values}: no
  combinational decode of IR/TDR bits downstream, so no output can
  glitch (they feed async resets, clock-mux selects and a clock-gate
  latch).
\end{itemize}

\subsection{Instruction map (private codes; FROZEN by user decision
2026-10-11, Q-OP: do not renumber)}\label{tm:ch:dft-architecture:153-instruction-map-private-codes-frozen-by-user-decision-2026-10-11-q-op-do-not-renumber}\origsec{15.3}

{\footnotesize\setlength{\tabcolsep}{3pt}

{
\begin{longtable}[]{@{}
  >{\raggedright\arraybackslash}p{\dimexpr(\linewidth-6\tabcolsep)*118/1000\relax}
  >{\raggedright\arraybackslash}p{\dimexpr(\linewidth-6\tabcolsep)*157/1000\relax}
  >{\raggedright\arraybackslash}p{\dimexpr(\linewidth-6\tabcolsep)*98/1000\relax}
  >{\raggedright\arraybackslash}p{\dimexpr(\linewidth-6\tabcolsep)*608/1000\relax}@{}}
\toprule\noalign{}
\begin{minipage}[b]{\linewidth}\raggedright
Opcode
\end{minipage} & \begin{minipage}[b]{\linewidth}\raggedright
Name
\end{minipage} & \begin{minipage}[b]{\linewidth}\raggedright
DR
\end{minipage} & \begin{minipage}[b]{\linewidth}\raggedright
Notes
\end{minipage} \\
\midrule\noalign{}
\endhead
\bottomrule\noalign{}
\endlastfoot
\texttt{4'h0} & EXTEST & BYPASS & \textbf{not implemented} (no boundary
register, N2); decodes to BYPASS \\
\texttt{4'h1} & IDCODE & 32 & reset default; DR =
\texttt{32'h15C0\_\allowbreak{}1001}: version 1, part \texttt{16'h5C01},
manufacturer \texttt{11'h000} (\textbf{invalid for JEDEC on purpose},
N1), LSB 1 \\
\texttt{4'h2} & SAMPLE/\allowbreak{}PRELOAD & BYPASS & \textbf{not
implemented} (N2) \\
\texttt{4'h3} & SCAN\_SEL & 3 & \code{[0]} PARALLEL, \code{[2:1]}
CHAIN\_GRP (0 all, 1 \code{clk_i} fabric, 2 \code{cpu_clk_i} fabric, 3
CPU macro). Update register, reset 0. Capture loads the current value
(readback). \textbf{No effect unless IR = SCAN.} CHAIN\_GRP has no
consumer until Stage 2 \\
\texttt{4'h4} & SCAN & 1 (stub) & enters scan mode (15.4). Until Stage 2
stitches chains the serial chain path is a 1-bit stub with BYPASS
timing; the single place to replace it is marked \code{STAGE2_CHAIN} in
\code{jtag_tap.sv} \\
\texttt{4'h5} & SCAN\_RST & 1 & \code{[0]} assert scan reset (level).
Update register, reset 0. Effective only while IR = SCAN \\
\texttt{4'h6} & MBIST\_\allowbreak{}CTRL & BYPASS & \textbf{reserved for
Stage 4}; decodes to BYPASS now \\
\texttt{4'h7} & MBIST\_\allowbreak{}STATUS & BYPASS & \textbf{reserved
for Stage 4}; decodes to BYPASS now \\
\texttt{4'hF} & BYPASS & 1 & Capture loads 0 \\
all others & & BYPASS & unknown opcodes -\textgreater{} BYPASS \\
\end{longtable}
}

}

\subsection{What the TAP drives, and how it combines with the 1a
pins}\label{tm:ch:dft-architecture:154-what-the-tap-drives-and-how-it-combines-with-the-1a-pins}\origsec{15.4}

Selecting SCAN (IR = \texttt{4'h4}) is the \textbf{only} thing that
enters TAP scan mode; TDR contents alone never disturb functional
operation (they are gated by IR = SCAN, and cleared by
TLR/\allowbreak{}TRST). With \code{tap_scan} = (IR==SCAN), \code{serial}
= tap\_scan \& \textasciitilde PARALLEL, and \code{pin_*} = the
unchanged 1a \code{dft_ctrl_ports} leg:

{\small\setlength{\tabcolsep}{4pt}

{
\begin{longtable}[]{@{}
  >{\raggedright\arraybackslash}p{\dimexpr(\linewidth-4\tabcolsep)*243/1000\relax}
  >{\raggedright\arraybackslash}p{\dimexpr(\linewidth-4\tabcolsep)*502/1000\relax}
  >{\raggedright\arraybackslash}p{\dimexpr(\linewidth-4\tabcolsep)*235/1000\relax}@{}}
\toprule\noalign{}
\begin{minipage}[b]{\linewidth}\raggedright
Output
\end{minipage} & \begin{minipage}[b]{\linewidth}\raggedright
Value
\end{minipage} & \begin{minipage}[b]{\linewidth}\raggedright
With the TAP idle
\end{minipage} \\
\midrule\noalign{}
\endhead
\bottomrule\noalign{}
\endlastfoot
\code{scan_mode_o}, \code{test_en_o} &
\texttt{pin\_\allowbreak{}scan\_\allowbreak{}mode\ |\ tap\_\allowbreak{}scan}
& = pin (bit-identical to 1a) \\
\code{scan_rst_no} &
\texttt{pin\_\allowbreak{}scan\_\allowbreak{}rst\_\allowbreak{}n\ \&\ \textasciitilde{}(tap\_\allowbreak{}scan\ \&\ SCAN\_\allowbreak{}RST[0])}
& = pin \\
\code{scan_en_o} &
\texttt{serial\ ?\ tap\_\allowbreak{}scan\_\allowbreak{}en\ :\allowbreak{}\ pin\_\allowbreak{}scan\_\allowbreak{}en}
& = pin \\
\code{test_clk_o} &
\texttt{serial\ ?\ tck\_\allowbreak{}scan\_\allowbreak{}gated\ :\allowbreak{}\ pin\_\allowbreak{}scan\_\allowbreak{}clk},
through the glitch-free switch of 15.5 & = pin \\
\end{longtable}
}

}

\code{tap_scan_en} is a falling-edge flop = (state==Shift-DR) \& serial.
\code{tck_scan_gated} is TCK gated by an \code{rv32i_clock_gate} whose
enable is (state in \{Capture-DR, Shift-DR\}) \& serial, so the scanned
flops see exactly the capture edge (the rising edge ending Capture-DR,
\code{scan_en}=0) and the shift edges (rising edges ending Shift-DR
cycles, \code{scan_en}=1), and \textbf{no} edge in
Run-Test/\allowbreak{}Idle, Pause or Update (no spurious capture).
PARALLEL=1 hands
\code{scan_clk_i}/\code{scan_en_i}/\code{scan_in_i}/\code{scan_out_o} to
the pins (the TAP still supplies \code{scan_mode}). The TAP
\textbf{never gates the pin path}: with the TAP idle, or in TLR, the 1a
pins work exactly as before. If a pin and the TAP are both in scan, the
TAP wins \code{test_clk}/\code{scan_en} while \code{serial} is set.

\subsection{\texorpdfstring{Clock/\allowbreak{}reset treatment and
the declared
crossings}{Clock/reset treatment and the declared crossings}}\label{tm:ch:dft-architecture:155-clockreset-treatment-and-the-declared-crossings}\origsec{15.5}

\begin{itemize}
\tightlist
\item
  TCK is an asynchronous third domain. Nothing in the TAP is clocked by
  a functional clock.
\item
  \textbf{Crossing 1, \code{trst_ni} -\textgreater{} TCK:} asynchronous
  assert, synchronous release through \code{cdc_reset_sync}
  (\code{u_trst_sync}, \code{jtag_tap.sv:113}; \code{u_sw_rst},
  \code{dft_ctrl.sv:137}). \textbf{What guards these two chains (bead
  \bead{j41m.11}, 2026-10-14).} \src{tools/cdc/cdc_config.yml} has
  \textbf{no} \code{required_synchronisers} entry for either chain, and
  one cannot be written with the existing gate. Measured with the CI
  flow (\texttt{make\ -\allowbreak{}C\ sim\ cdc}, pinned yosys): (1)
  cdc\_\allowbreak{}snitch classes a \code{cdc_reset_sync} chain
  \code{OK1} (the \code{:D} pins; stage 0 has D tied to 1, stage 1 sees
  only \code{tck_i}) and \code{OKX} (the \code{:R} async-clear pins, one
  source \code{trst_ni}), never \code{CDC}, and a
  \code{required_synchronisers} entry is satisfied only by \code{CDC}
  (or \code{BAD}+\code{magic} with \code{latch_launched}); (2)
  \code{cdc_snitch_proc.ys} runs \code{opt_merge}, which folds
  \code{u_tap.u_trst_sync} into \code{u_sw_rst} (same clock, same async
  clear, D = 1), so the report contains only
  \code{u_dft_ctrl.u_sw_rst.sync_q} and \code{u_trst_sync} is invisible
  to the gate by name. Bypassing \code{u_sw_rst} alone was caught only
  by accident (the raw \code{trst_ni} reaching \code{u_sync_a} combines
  with \code{scan_clk_i} because \code{trst_ni} is not in
  \code{reset_domains}), bypassing \code{u_trst_sync} alone passed. Do
  \textbf{not} add \code{trst_ni} to \code{reset_domains}: it would
  remove that accidental detection. The check is therefore structural,
  in the scan checker, which runs on a netlist that has had no
  \code{opt_merge}:
  \texttt{tools/\allowbreak{}dft/\allowbreak{}check\_\allowbreak{}scan\_\allowbreak{}clk\_\allowbreak{}rst.\allowbreak{}py\ -\allowbreak{}-\allowbreak{}expect-\allowbreak{}reset-\allowbreak{}sync\ SYNC=\allowbreak{}CONSUMER}
  (repeatable; \code{--reset-sync-min-flops}, default 2). It finds each
  chain by hierarchy
  (\code{u_dft_ctrl.u_tap.u_trst_sync=u_dft_ctrl.u_tap},
  \code{u_dft_ctrl.u_sw_rst=u_dft_ctrl.u_clk_sw}), requires at least two
  flops with async reset exactly \code{trst_ni} and clock exactly
  \code{tck_i}, requires the consumer to be reset from that chain, and
  fails if \code{trst_ni} reaches any other sequential cell without a
  synchroniser (the older TAP rule accepted a TAP flop reset straight
  from \code{trst_ni}, which is the hole this closes). Wired into
  \code{DFT_TAP_CHECK_ARGS} in \src{pnr/Makefile}, so
  \texttt{make\ -\allowbreak{}C\ pnr\ dft-\allowbreak{}scan-\allowbreak{}check}
  and \code{dft-scan-check-asap7} both run it. Measured: both PDK
  netlists 0 violations, 2 of 2 chains; scratch bypass of
  \code{u_sw_rst}, of \code{u_trst_sync}, and of both each fail on the
  Sky130 netlist naming the missing chain (the unmodified tree passes);
  25 pytest controls in \src{tb/tests/test_dft_scan_check.py}.
  \textbf{Limits, stated plainly:} (a) CI (bead \bead{j41m.12},
  2026-10-14): the \code{cdc} job of
  \src{.github/workflows/rtl-checks.yml} runs
  \texttt{make\ -\allowbreak{}C\ pnr\ dft-\allowbreak{}scan-\allowbreak{}check},
  \texttt{make\ -\allowbreak{}C\ pnr\ dft-\allowbreak{}scan-\allowbreak{}check-\allowbreak{}asap7}
  and
  \texttt{pytest\ tb/\allowbreak{}tests/\allowbreak{}test\_\allowbreak{}dft\_\allowbreak{}scan\_\allowbreak{}check.\allowbreak{}py}
  as blocking steps on the job\textquotesingle s pinned yosys and sv2v.
  The pytest step fails on any skipped test or fewer than 25 passes,
  because that file skips itself when it finds no yosys;
  \code{tests.yml} still has no yosys, so the file is skipped there and
  only the \code{cdc} job exercises it; (b) the check sees the RTL-level
  netlist, and a real synthesis merges the two identical chains into one
  (measured on \code{dft_ctrl} alone with \code{opt_merge}), so it
  proves the RTL structure, not that the post-synthesis netlist still
  has two chains. Earlier text (2026-10-13) said "declared a required
  synchroniser"; that was never true.
\item
  \textbf{The two \code{trst_ni} chains are duplicates on purpose (user
  decision 2026-10-14, bead \bead{j41m.12}).} \code{u_trst_sync}
  (\code{jtag_tap}) and \code{u_sw_rst} (\code{dft_ctrl}) are
  functionally identical: \code{STAGES} 2, clocked by \code{tck_i},
  async-cleared by \code{trst_ni}, scan pins tied off. Keeping both is
  intentional (per-module reset ownership; \code{jtag_tap} is left
  untouched). Synthesis \code{opt_merge} folds them into one chain
  (measured 73 -\textgreater{} 70 bit-flops on \code{dft_ctrl}), so a
  post-synthesis netlist is expected to show one chain, not two.
  \code{--expect-reset-sync} pins each chain by hierarchy on the
  pre-merge netlist, which is why it can still tell them apart.
\item
  \textbf{Crossing 2, TCK-launched static controls -\textgreater{}
  functional sinks:} \code{scan_mode}, \code{test_en}, \code{scan_rst},
  \code{scan_en} (falling-edge flops) feed only async-reset muxes,
  clock-mux selects, the clock-gate latch D and the crypto key mask.
  They are \textbf{deliberately not synchronised into any functional
  clock}: in scan mode the functional clock is stopped and replaced by
  the test clock, so a synchroniser in the destination domain would
  never advance, and the crypto mask (a data path into core-clock flops)
  must be effective before the first test-clock edge, not two edges
  later. This is the same \emph{quasi-static} contract as the 1a ports
  (set with the functional clocks stopped, not toggled while a clock
  runs). \textbf{How it is declared (corrected 2026-10-13; this text
  previously said an \code{accepted-risk} waiver):}
  \src{tools/cdc/cdc_config.yml} lists \code{tck_i} under
  \code{reset_domains}, with the justification as a comment; that
  removes the TCK source from the domain count of the registers behind
  these controls. It is \textbf{not} a \code{waivers:} entry (the
  file\textquotesingle s only two waivers are
  \texttt{kind:\allowbreak{}\ tool-\allowbreak{}limitation}, for the
  bridge reset and the APB response payload, neither for TCK) and not
  \code{accepted-risk}. The check that keeps it from being silent is
  \texttt{tools/\allowbreak{}dft/\allowbreak{}check\_\allowbreak{}scan\_\allowbreak{}clk\_\allowbreak{}rst.\allowbreak{}py\ -\allowbreak{}-\allowbreak{}tap-\allowbreak{}allowed-\allowbreak{}nets}
  (\texttt{make\ -\allowbreak{}C\ pnr\ dft-\allowbreak{}scan-\allowbreak{}check},
  allowed nets \code{tap_scan}, \code{tap_serial},
  \code{tap_rst_assert}, \code{test_clk_o}, \code{tdo_o},
  \code{tdo_oe_o}), which fails on any other TAP-launched net reaching
  non-TAP logic.

  \begin{itemize}
  \tightlist
  \item
    \textbf{Why these controls cannot be given a synchroniser or a FIFO
    (reasons, not omissions):} (1) the destination of \code{serial} is
    the \code{scan_clk_i} side of \code{dft_clk_switch} (\code{sel_a_i},
    \code{dft_ctrl.sv} \code{u_clk_sw}), and \code{scan_clk_i} idles LOW
    (tester contract 15.9.1), so a synchroniser or FIFO read side
    clocked there has no edges and would never advance (the deadlock
    15.9.2 describes for the latch input); (2) \code{scan_mode} selects
    its own destination clock: it is \code{sel_i} of the
    \code{dft_clk_mux} that picks the functional or the test clock in
    each \code{pll_subsystem} (\code{pll_subsystem.sv:131-134,192-195}),
    so a destination-clock synchroniser would be clocked by a clock that
    the signal itself is switching; (3) the zero-spare tester contract
    of 15.9.1 (the Update-IR to Capture-DR path gives exactly the TCK
    edges the clock-switch handshake needs) forbids adding latency to
    the TAP-to-clock-switch path. As built, the only synchronisers on
    TCK-related paths are \code{u_sync_a}/\code{u_sync_b} in
    \code{dft_clk_switch.sv:142-158} and the reset synchronisers
    \code{u_trst_sync} and \code{u_sw_rst}.
  \end{itemize}
\item
  \textbf{Test-clock mux (Q-CLK, resolved 2026-10-11 by user decision):
  a glitch-free enable-before-switch mux,
  \src{rtl/soc/dft/dft_clk_switch.sv}, instance \code{u_clk_sw} in
  \code{dft_ctrl}.} Sources: A = gated TCK (idles low), B = the
  \code{scan_clk_i} pin clock (assumed to idle LOW); select
  \code{serial} (a TCK falling-edge flop). Each source has an enable
  \textbf{latch} that is transparent while its own raw clock is low and
  holds while it is high (\code{en_a_q}, phase = raw \code{tck_i};
  \code{en_b_q}, phase = \code{scan_clk_i}), the two enables are
  cross-coupled through two \code{cdc_sync} (\code{u_sync_a} carries
  \texttt{\textasciitilde{}en\_\allowbreak{}b\_\allowbreak{}q} into
  \code{tck_i}, \code{u_sync_b} carries \code{en_a_q} into
  \code{scan_clk_i}), and
  \texttt{clk\_\allowbreak{}o\ =\allowbreak{}\ (clk\_\allowbreak{}a\ \&\ en\_\allowbreak{}a\_\allowbreak{}q)\ |\ (clk\_\allowbreak{}b\ \&\ en\_\allowbreak{}b\_\allowbreak{}q)};
  its reset (\code{u_sw_rst}, a second \code{cdc_reset_sync} on
  \code{trst_ni}, level-forced into the latches) selects B with A off,
  so the 1a pins-only operation (TAP in reset, TCK stopped) is
  unchanged. \textbf{Latches rather than the usual falling-edge enable
  flops, on purpose:} a latch lets the deselected source be switched OFF
  with no clock edge when it idles low (the normal tester state: one
  clock stopped), which flops cannot do. Measured by
  \code{test_dft_clk_switch} (19 tests, random-phase sweeps of both
  directions at three coprime period pairs, a
  runt/\allowbreak{}extra-edge/\allowbreak{}merge checker with a
  self-test) and run against a plain mux as the red control (17 of 19
  fail). \textbf{Which clock must be running for a switch:} (1) parallel
  -\textgreater{} serial (select A): \code{scan_clk_i} needs no edge if
  it idles low (stopped HIGH blocks it); \code{tck_i} must give 2 edges,
  which the TAP does by construction (Update-IR -\textgreater{}
  Select-DR -\textgreater{} Capture-DR is exactly enough: \textbf{zero
  spare}, do not shorten that path); (2) serial -\textgreater{} parallel
  (select B): the first pin pulse passes at once if \code{scan_clk_i}
  did not run while A was selected, otherwise after about 3
  \code{scan_clk_i} cycles; until then the output stays low with no
  extra edge; (3) \code{sel_a} must not reverse faster than one full
  handshake (2 cycles of the clock being enabled plus 2 of the one being
  retired; the TAP cannot reverse that fast, short dwell was seen to
  merge pulses in the unit bench); (4) an asynchronous
  TRST\_\allowbreak{}n/\allowbreak{}reset mid-pulse can leave a short
  LOW notch between two whole pulses and never shortens a pulse; scan
  state is destroyed by it anyway. Not modelled: metastability,
  STA/\allowbreak{}clock-gating checks of the latch and OR cells, a
  zero-width delta glitch inside one timestep (cocotb cannot see it).
\item
  Power-up: without a driven \code{trst_ni} the TAP flops power up
  unknown. \code{trst_ni} must be asserted at power-up (or TMS held 1
  for 5 TCK before relying on the TAP); at the chip top, tie it to a
  power-on reset. \textbf{Open question Q-TRST.}
\end{itemize}

\textbf{Deferred: MBIST control-word crossing (PROPOSAL, NOT
IMPLEMENTED; user decision 2026-10-13, bead \bead{j41m.10}).} The future
Stage 4 MBIST control word (section 6.2) is the one TCK-to-functional
crossing for which a FIFO has been proposed; it is not a scan control,
and its destinations would be the ungated functional clocks. The
intended mechanism is the depth-2 \code{cdc_gray_fifo} (merged in \prref{276}), one per destination domain, write side \code{tck_i}, read side
the \textbf{ungated} \code{core_clk} / \code{cpu_core_clk} (never a
gated clock). This \textbf{replaces} the per-bit \code{cdc_sync} plus
"wait \textgreater= 3 destination cycles" plan text of sections 6 and
9.3 as the intended direction, but the design is \textbf{not settled}:
15 open questions, each with a recommendation, are tracked on bead
\bead{j41m.10} and in the \prref{276} body (destination clock, one FIFO per
domain, write-strobe timing, overflow policy, DEPTH 2 vs 4, reset
architecture, TLR default delivery, SYNC\_\allowbreak{}STAGES, checker
and \code{cdc_config.yml} policy, SDC form, opcodes 6/7 leaving BYPASS,
scan testability of the unreset \code{mem_q}). They are not answered
here, and no RTL, checker, CDC-config or opcode change is made for it
before Stage 4. It does \textbf{not} apply to the existing scan controls
(crossing 2 above gives the reasons).

\subsection{Documented deviations from IEEE 1149.1 and from section
9}\label{tm:ch:dft-architecture:156-documented-deviations-from-ieee-11491-and-from-section-9}\origsec{15.6}

\begin{enumerate}
\def\labelenumi{\arabic{enumi}.}
\tightlist
\item
  No boundary register: EXTEST/\allowbreak{}SAMPLE/\allowbreak{}PRELOAD
  decode to BYPASS (N2). 1149.1 requires them; compliance is "minus
  boundary register". CLAMP/\allowbreak{}HIGHZ absent.
\item
  IDCODE manufacturer field \texttt{11'h000} is non-JEDEC (N1).
\item
  TAP flops are on no scan chain (N4); they are tested by
  IDCODE/\allowbreak{}TDR readback and the 1b suite.
\item
  SCAN\textquotesingle s serial chain path is a 1-bit stub until Stage 2
  (no chains exist to connect).
\item
  MBIST\_\allowbreak{}CTRL/\allowbreak{}MBIST\_\allowbreak{}STATUS
  opcodes are reserved and decode to BYPASS (MBIST is Stage 4).
\item
  Sections 8/9 said the direct scan pins are "ANDed with the TAP-held
  scan\_\allowbreak{}mode" so a floating pin cannot disturb functional
  mode. \textbf{Not implemented}: the 1a contract (pins work alone) and
  the standing decision that the TAP must not gate the parallel path
  take precedence; the pins are OR-ed. A floating \code{scan_mode_i} is
  therefore not protected by the TAP. Decision (2026-10-11, Q-PIN): keep
  OR; floating-pin protection is a pad pull-down at pad-ring time
  (section 9.5).
\item
  SCAN\_RST is a level bit; the tester pulses it by writing 1 then 0.
\item
  The crypto scan-entry scrub (14.3 item 1) is now done in
  \code{crypto_accel} hardware (\code{j41m.6}, 14.10); the TAP still
  offers no scrub of its own; \code{scan_en} is withheld in
  \code{soc_top} until the wipe completes (\code{j41m.17}, 14.14).
\end{enumerate}

\subsection{Tests and checks required of this stage}\label{tm:ch:dft-architecture:157-tests-and-checks-required-of-this-stage}\origsec{15.7}

Unit suite \code{test_jtag_tap} (module \code{jtag_tap} alone);
DFT-controller suite on \code{dft_ctrl}; SoC suite on
\code{tb_soc_pll}/\code{tb_soc_top}; mutation campaign on FSM
transitions, IR capture, opcode decode, TDR decode, TDO/OE timing, gate
enable, combiner; \code{check_scan_clk_rst.py} extended to cover
TAP-domain flops (clock = \code{tck_i} only, reset = \code{trst_ni}
through the sync only, unscanned by design) with a negative control; CDC
gate declarations; coverage ratchet; sv2v source lists; SDC constraints
text (TCK clock, asynchronous clock group, JTAG IO delays) in
\code{sky130_soc.sdc} only. \textbf{PD is not run in this stage.}

\subsection{Implementation record (2026-10-11, PR \#270)}\label{tm:ch:dft-architecture:158-implementation-record-2026-10-11-pr-270}\origsec{15.8}

Built as specified in 15.1 to 15.5: \src{rtl/soc/dft/jtag_tap.sv},
\src{rtl/soc/dft/dft_ctrl.sv} (instantiates the unchanged
\code{dft_ctrl_ports} as the pin leg), \code{soc_top} ports
\texttt{tck\_\allowbreak{}i\ tms\_\allowbreak{}i\ tdi\_\allowbreak{}i\ trst\_\allowbreak{}ni\ tdo\_\allowbreak{}o\ tdo\_\allowbreak{}oe\_\allowbreak{}o},
source lists (\src{pnr/Makefile} x2, \src{sim/Makefile},
\src{tb/cocotb/soc/Makefile} x3 DFT lists), port stubs
(\code{soc_top_ports.v}, both PDKs). Implementation choices inside the
specification, none of which change a pin or the instruction map:
\code{parallel_o} and \code{chain_grp_o} are gated by IR = SCAN exactly
like \code{serial_o}; Update-DR/\allowbreak{}Update-IR commit on the
falling edge after the rising edge that entered the Update state; TAP
flop count is 23 yosys cells (2 are the TRST synchroniser).

{\footnotesize\setlength{\tabcolsep}{3pt}

{
\begin{longtable}[]{@{}
  >{\raggedright\arraybackslash}p{\dimexpr(\linewidth-2\tabcolsep)*490/1000\relax}
  >{\raggedright\arraybackslash}p{\dimexpr(\linewidth-2\tabcolsep)*490/1000\relax}@{}}
\toprule\noalign{}
\begin{minipage}[b]{\linewidth}\raggedright
Check
\end{minipage} & \begin{minipage}[b]{\linewidth}\raggedright
Result {[}M{]}
\end{minipage} \\
\midrule\noalign{}
\endhead
\bottomrule\noalign{}
\endlastfoot
\code{test_jtag_tap} (unit, 16 tests x 3 TCK/\allowbreak{}system-clock
pairs 7/5, 13/3, 29/5 ns) & 48/48 pass \\
\code{test_dft_ctrl} & 7/7 pass \\
\code{test_soc_jtag} (tb\_\allowbreak{}soc\_\allowbreak{}pll) & 11/11
pass; first run 8/11, the 3 failures were test defects (absolute commit
thresholds that ignored the loop firmware\textquotesingle s
\textasciitilde0.28 IPC, a \code{key_q} deposit made while the
peripheral resets were still asserting, and a post-scan check that
forgot SCAN\_RST resets the crypto block), RTL unchanged \\
Hand mutants on TAP + \code{dft_ctrl} (IR capture value, Exit2-DR
transition, \code{scan_mode} from TDR bits, \code{scan_en} in
Capture-DR, ungated \code{chain_grp}/\code{rst_assert}, spurious gate
edge in Exit1-DR, pin-OR on \code{test_clk_o},
\texttt{scan\_\allowbreak{}mode\_\allowbreak{}o\ =\allowbreak{}\ pin\ \&\ tap},
...) & 19 applied, 19 killed; mutated copies lived outside the tree \\
\code{check_scan_clk_rst.py} + \src{tb/tests/test_dft_scan_check.py} &
14 pytest pass (8 new negative controls); both PDK netlists 0
violations: Sky130 1663 sequential cells / 28 scan-mode muxes, ASAP7
3749 / 29; 23 TAP-domain cells; scan mode net has the pin leg and a TAP
leg; undeclared-TAP-output control on the real netlist fails with 4
violations \\
CDC gate
(\texttt{make\ -\allowbreak{}C\ sim\ cdc\ CDC\_\allowbreak{}STRICT=\allowbreak{}1},
cdc\_\allowbreak{}snitch + cdc\_gate) & PASS, 0 unwaived, every waiver
and required synchroniser still matches; \code{tck_i} added to
\code{reset_domains} with the justification of 15.5 crossing 2, backed
by the \code{--tap-allowed-nets} check above (the gate can no longer see
a TCK-launched signal, the checker can) \\
\texttt{make\ -\allowbreak{}C\ sim\ lint}, Verible
(\src{.rules.verible_lint}), \code{check_struct_member_select.py},
\code{check_source_closure.py} (Sky130 and ASAP7 SoC) & clean / 0
findings / CLOSED, CLOSED; sv2v regenerates for both PDKs \\
Python QA as CI
(\texttt{ruff\ format\ -\allowbreak{}-\allowbreak{}check},
\texttt{ruff\ check}, \code{mypy}, \code{pylint}) & clean; pylint
9.55/10; \texttt{pytest\ tb/\allowbreak{}tests} pass (yosys-gated tests
run with the cdc-cache yosys) \\
SDC (\code{sky130_soc.sdc}: \texttt{create\_\allowbreak{}clock\ tck} 100
ns,
\texttt{set\_\allowbreak{}clock\_\allowbreak{}groups\ -\allowbreak{}asynchronous},
JTAG IO delays, \code{trst_ni} false path) & sources to completion in
OpenSTA against a port-only \code{soc_top} (tck clock and the six ports
resolve; the only warnings are the pre-existing missing-port ones from
the stub). \textbf{No STA sign-off and no PD run} \\
\end{longtable}
}

}

\textbf{Not verified:} any P\&R, CTS, timing or hold effect of the
TCK-gate / clock-select cells; STA with the TAP flop driving the
clock-mux selects (the functional-mode SDC cannot name that flop after
flattening, see the KNOWN GAP in \code{sky130_soc.sdc}); TCK frequency
limits; real-silicon power-up without a driven \code{trst_ni};
metastability (not modelable); the ASAP7 SDCs (not edited: Sky130 is the
target); \code{synth_gate/ci_smoke.sh} locally (not executable here; it
covers the CPU frontend, which this change does not touch; CI ran it
green); gate-level simulation of the TAP.

\textbf{Open questions still open (user decisions of 2026-10-11 closed
Q-PIN, Q-CLK and Q-OP and moved Q-SDC to a separate P4 bead):}

\begin{itemize}
\tightlist
\item
  Q-TRST: \code{trst_ni} must be asserted at power-up or 5 TCK with
  TMS=1 applied; tie it to a power-on reset at the chip top. Idle value
  used by the suites: asserted.
\item
  Q-OP (RESOLVED, frozen 2026-10-11, kept for the record): opcode values
  (15.3) are private codes chosen here; the doc left them open.
  SCAN\_RST was implemented as its own 1-bit instruction (the doc marks
  it optional),
  MBIST\_\allowbreak{}CTRL/\allowbreak{}MBIST\_\allowbreak{}STATUS
  reserved and decoded to BYPASS until Stage 4.
\item
  Q-STUB: SCAN\textquotesingle s serial data register is a 1-bit stub
  until Stage 2 (no chains exist), \code{CHAIN_GRP} has no consumer, and
  \code{scan_en}/\code{tap_scan_en} has no RTL consumer (its net is
  optimised away; add it to the \code{--tap-allowed-nets} list in
  \src{pnr/Makefile} when Stage 2 connects it).
\item
  Q-SDC (moved to its own P4 bead 2026-10-11): the ASAP7 SoC SDCs were
  not given TCK constraints; not edited here.
\item
  Crypto scan-entry scrub: closed in hardware by \code{j41m.6} (14.10).
  SCAN\_\allowbreak{}SEL/\allowbreak{}SCAN\_\allowbreak{}RST under SCAN
  also resets the crypto block through the scan reset, which is an
  additional, tester-driven scrub path and is not relied on.
\end{itemize}

\subsection{Q-CLK change record (2026-10-11): glitch-free
test-clock switch, supersedes the "plain mux" statements above}\label{tm:ch:dft-architecture:159-q-clk-change-record-2026-10-11-glitch-free-test-clock-switch-supersedes-the-plain-mux-statements-above}\origsec{15.9}

User decisions: Q-PIN keep OR (pad pull-downs at pad-ring time); Q-CLK
glitch-free enable-before-switch mux now; Q-OP frozen; Q-SDC separate
bead. Where 15.4/\allowbreak{}15.5/\allowbreak{}15.6/\allowbreak{}15.8
above or the earlier experience record say the test-clock select is a
plain mux, the "Test-clock mux" bullet of 15.5 is the current design.
Deviation from the request\textquotesingle s wording, flagged: the
per-source enables are latches (transparent while the own raw clock is
low), not falling-edge flops, because a flop-based enable cannot retire
a stopped deselected clock and would deadlock the common tester state;
the 2-FF cross-coupling, AND-OR output and reset-selects-pin-clock are
as requested.

{\footnotesize\setlength{\tabcolsep}{3pt}

{
\begin{longtable}[]{@{}
  >{\raggedright\arraybackslash}p{\dimexpr(\linewidth-2\tabcolsep)*490/1000\relax}
  >{\raggedright\arraybackslash}p{\dimexpr(\linewidth-2\tabcolsep)*490/1000\relax}@{}}
\toprule\noalign{}
\begin{minipage}[b]{\linewidth}\raggedright
Check
\end{minipage} & \begin{minipage}[b]{\linewidth}\raggedright
Result {[}M{]}
\end{minipage} \\
\midrule\noalign{}
\endhead
\bottomrule\noalign{}
\endlastfoot
\code{test_dft_clk_switch} (new) & 19/19; the same suite built with
\code{-GPLAIN_MUX=1}: 2/19 pass, 17 FAIL (runts, truncated pulses,
handover tests) \\
\code{test_jtag_tap} / \code{test_dft_ctrl} / \code{test_soc_jtag} /
\code{test_soc_dft_scan} (1a, unchanged) & 48/48, 16/16 (+9), 12/12
(+1), 10/10 \\
20 hand mutants on the mux (sync
dropped/\allowbreak{}single/\allowbreak{}wrong clock, latch phase, flop
instead of latch, no cross-coupling, flipped reset, ...) & 18 killed,
m16 equivalent (redundant reset term, forced by the
synchroniser\textquotesingle s own async reset), m20 a combinational
loop (caught by Verilator lint, not simulated) \\
\code{check_scan_clk_rst.py} on both PDK netlists
(\texttt{-\allowbreak{}-\allowbreak{}tap-\allowbreak{}alt-\allowbreak{}clock\ scan\_\allowbreak{}clk\_\allowbreak{}i},
new prefixes \code{u_clk_sw}, \code{u_sw_rst}) & 0 violations; 31
TAP-domain cells (was 23); Sky130 1671 sequential cells, ASAP7 3757;
removing the alt-clock option makes it fail on the two \code{u_sync_b}
flops; 15 pytest negative controls \\
CDC gate & PASS, 0 unwaived. The two handshake synchronisers show as
\texttt{BAD\ magic} in cdc\_\allowbreak{}snitch (the enable latch inputs
combine both clocks) and are documented in
\src{tools/cdc/cdc_config.yml}; see 15.9.2: root cause is a
cdc\_\allowbreak{}snitch convention, they are now declared with
\texttt{latch\_\allowbreak{}launched:\allowbreak{}\ true} under
\code{required_synchronisers} and proven present by the gate
(negative-controlled), clocking is also enforced by the checker above
and the unit suite \\
Verilator \code{-Wall}, Verible, \code{check_struct_member_select},
source closure (both SoC configs), sv2v (both PDKs), Python QA (ruff,
mypy, pylint 9.56) & clean / CLOSED / pass \\
\end{longtable}
}

}

Not verified (added): timing or STA of the
latch/\allowbreak{}AND/\allowbreak{}OR clock cells and their
clock-gating checks (Stage 2 scan-mode SDC), CTS behaviour of the extra
OR on the clock path, metastability. The Sky130 functional SDC carries a
note and still sources in OpenSTA.

\subsubsection{User decisions of 2026-10-12 and the tester
contract (stated plainly)}\label{tm:ch:dft-architecture:1591-user-decisions-of-2026-10-12-and-the-tester-contract-stated-plainly}\origsec{15.9.1}

The latch-based \code{dft_clk_switch} is \textbf{accepted}. The tester
contract is \textbf{accepted and part of the design}:

\begin{enumerate}
\def\labelenumi{\arabic{enumi}.}
\tightlist
\item
  \textbf{\code{scan_clk_i} must idle LOW} (a stopped-high
  \code{scan_clk_i} blocks every switch, both directions, until it
  falls), and must toggle for about 3 cycles after a serial
  -\textgreater{} parallel switch if it ran during scan (otherwise the
  first pin pulse passes at once).
\item
  \textbf{The TAP\textquotesingle s parallel -\textgreater{} serial
  entry has ZERO spare TCK edges.} The shortest legal TMS path Update-IR
  -\textgreater{} Select-DR -\textgreater{} Capture-DR gives exactly the
  2 TCK edges the handshake needs before the capture edge. Nobody may
  shorten that path or the handshake (for example by registering a
  signal in \code{dft_clk_switch} or the TAP). It is pinned by
  \texttt{test\_\allowbreak{}dft\_\allowbreak{}ctrl.\allowbreak{}test\_\allowbreak{}serial\_\allowbreak{}entry\_\allowbreak{}minimal\_\allowbreak{}path\_\allowbreak{}\{pin\_\allowbreak{}clock\_\allowbreak{}stopped,\allowbreak{}\ pin\_\allowbreak{}clock\_\allowbreak{}stopped\_\allowbreak{}after\_\allowbreak{}trst,\allowbreak{}\ pin\_\allowbreak{}clock\_\allowbreak{}running\}}
  (shortest TMS path; asserts that \code{en_a_q} rises between the 2nd
  TCK rise and the capture edge, and that the capture edge and every
  shift edge reach \code{test_clk_o}).
\item
  \code{sel_a} must not reverse faster than one full handshake (the TAP
  cannot).
\end{enumerate}

\subsubsection{\texorpdfstring{cdc\_\allowbreak{}snitch
\texttt{BAD\ magic} on the two handshake synchronisers: investigation
(2026-10-12)}{cdc\_snitch  on the two handshake synchronisers: investigation (2026-10-12)}}\label{tm:ch:dft-architecture:1592-cdc-snitch-bad-magic-on-the-two-handshake-synchronisers-investigation-2026-10-12}\origsec{15.9.2}

\textbf{Root cause (reproduced, not guessed).} \code{cdc_snitch.py}
(\code{check_bit}, \code{list_domains}) looks for flops by the substring
"DFF" in the cell type and classifies each flop input as \code{CDC} only
when the input cone has \textbf{exactly one} entry
(\texttt{len(dom\_\allowbreak{}in)\ =\allowbreak{}=\allowbreak{}\ 1})
and the register carries \code{magic_cdc}; any cone with two or more
entries is \code{BAD}, "doesn\textquotesingle t matter if they claim CDC
or not". \code{list_domains} recurses through every non-"DFF" cell and
appends one entry per \textbf{input pin} it reaches. A latch
(\texttt{\$\_\allowbreak{}DLATCH\_\allowbreak{}*}, ports E and D) is not
a "DFF" cell, so it is traversed like combinational logic and
contributes two entries (E = a module-input clock port, D = the enable
cone). Evidence, minimal reproducer (\code{snitch_repro.v}: four sources
in clock domain A into a \code{magic_cdc} 2-FF stage in domain B, run
through the repo\textquotesingle s own \code{cdc_snitch_proc.ys} +
\code{cdc_snitch.py}, yosys 0.64):

{
\begin{longtable}[]{@{}
  >{\raggedright\arraybackslash}p{\dimexpr(\linewidth-2\tabcolsep)*645/1000\relax}
  >{\raggedright\arraybackslash}p{\dimexpr(\linewidth-2\tabcolsep)*335/1000\relax}@{}}
\toprule\noalign{}
\begin{minipage}[b]{\linewidth}\raggedright
Source into the synchroniser
\end{minipage} & \begin{minipage}[b]{\linewidth}\raggedright
cdc\_\allowbreak{}snitch class
\end{minipage} \\
\midrule\noalign{}
\endhead
\bottomrule\noalign{}
\endlastfoot
A-domain flop Q (control) & \texttt{CDC\ magic} (1 x clk\_a) \\
\textbf{single-domain} latch (E = \textasciitilde clk\_a, D = A flop) &
\textbf{\texttt{BAD\ magic} (2 x clk\_a)} \\
AND of two A-domain flops (no latch at all) &
\textbf{\texttt{BAD\ magic} (2 x clk\_a)} \\
the same latch value registered by an A flop first & \texttt{CDC\ magic}
(1 x clk\_a) \\
\end{longtable}
}

So (i) the attribute survives (it is \code{magic} in all rows; it also
survives the flatten under \code{dft_ctrl}), (ii) the trigger is "more
than one pin in the launch cone", \textbf{a tool
limitation/\allowbreak{}convention} (a CDC source must be a single flop
output with no logic), not a lost attribute and not a config error, and
it fires for a perfectly clean single-clock latch. \textbf{Real
netlist.} \code{u_sync_b.sync_q[0]:D} (clk \code{scan_clk_i}):
\texttt{4\ x\ tck\_\allowbreak{}i} = E (raw TCK) plus \code{sel_a},
\texttt{\textasciitilde{}rst}, \code{b_off_in_a}, all TCK-domain: the
tool artifact only. \code{u_sync_a.sync_q[0]:D} (clk \code{tck_i}):
\texttt{2\ x\ scan\_\allowbreak{}clk\_\allowbreak{}i,\allowbreak{}\ 2\ x\ tck\_\allowbreak{}i}
= E (\code{scan_clk_i}) and \code{a_on_in_b} (B domain) plus
\code{sel_a}, reset (TCK domain). \textbf{Real design concern, found by
the investigation (not a tool artifact):} \code{en_b_q} is a latch whose
D contains \textbf{\code{sel_a_i} and the reset level, both launched in
the TCK domain, with no synchroniser into the \code{scan_clk_i} domain.}
That is the price of the stopped-clock property (a synchronised select
cannot retire a stopped B). It is glitch-free by the latch construction
(D can only matter while \code{scan_clk_i} is low), but it is not
metastability-safe: if \code{sel_a} rises within the latch aperture of a
\code{scan_clk_i} rising edge while \code{scan_clk_i} is TOGGLING,
\code{en_b_q} can go metastable while \code{clk_b} is high and the
output can take a runt. The accepted tester contract removes the case
(\code{scan_clk_i} idles low at every switch away from B, so the latch
is transparent and nothing is sampled); with a toggling
\code{scan_clk_i} the exposure is real but only at the one retire-B
instant, probability about (aperture / period). The simulation cannot
show it (cocotb has no metastability). \textbf{ACCEPTED RISK Q-META
(user decision 2026-10-12), contract-dependent:} select and reset reach
the \code{scan_clk_i}-side enable latch unsynchronised; it is safe ONLY
because \code{scan_clk_i} idles low at every switch (15.9.1); it cannot
be simulated. \textbf{Explicit check item for Stage 2 and for the ATPG
protocol definition:} (1) the tester/\allowbreak{}ATE protocol must keep
\code{scan_clk_i} low across every IR load that changes \code{serial}
(entering or leaving SCAN, and the SCAN\_SEL PARALLEL change followed by
SCAN), (2) STA of the scan-mode SDC must analyse the
\code{sel_a_i}/\allowbreak{}reset-to-\code{en_b_q} latch path
(setup/\allowbreak{}hold against \code{scan_clk_i}) and the clock-gating
checks on both enable latches, (3) re-evaluate whether to synchronise
select and reset into the \code{scan_clk_i} domain if the idles-low
contract is ever dropped. \textbf{Why the gate passed anyway, and what
was weak.} The gate counted these two registers as modelling artifacts:
after canonicalisation \code{tck_i} is in \code{reset_domains}, so
\code{u_sync_b} has no domains left and \code{u_sync_a} has one
(\code{PORT:scan_clk_i}); the gate only fails a \code{BAD} with two or
more distinct domains, and \code{required_synchronisers} only matches
\code{CDC}-classified registers. So the two synchronisers were
\textbf{not actually checked} by the gate (neither present nor correct);
15.9\textquotesingle s earlier statement that they could not be listed
was true only for the old gate code. \textbf{Applied (pure tooling, no
RTL or behaviour change, no check weakened).}
\src{tools/cdc/cdc_gate.py}: a \code{required_synchronisers} entry may
set \texttt{latch\_\allowbreak{}launched:\allowbreak{}\ true}; such an
entry is satisfied by a \code{CDC} register \textbf{or} a
\code{magic_cdc}-marked register that cdc\_\allowbreak{}snitch classed
\code{BAD}, and never by an unmarked \code{BAD} or an \code{OKX}.
\src{tools/cdc/cdc_config.yml} declares \code{dft-clk-switch-sync-a} and
\code{-sync-b} that way. Negative controls on the real report
(\code{soc_top_cdc.txt} of the current RTL): removing the
\code{u_sync_a} line, or stripping \code{magic} from the \code{u_sync_b}
line, makes the gate FAIL naming the entry; the unmodified report
PASSes; 4 new unit tests in \src{tb/tests/test_cdc_gate.py} (opted-in
BAD+magic passes, not-opted-in fails, unmarked BAD/OKX fail, stage
removed fails; 33/33). \textbf{Options not applied (they touch RTL or
weaken a check):}

\begin{itemize}
\tightlist
\item
  Register the latch outputs before they cross (a flop on the
  latch\textquotesingle s own clock in front of each synchroniser):
  cdc\_\allowbreak{}snitch would then classify CDC. Cost: the B-side
  flop only updates on \code{scan_clk_i} edges, so parallel
  -\textgreater{} serial with a stopped \code{scan_clk_i} would deadlock
  (the stopped-clock property is lost) and the zero-margin TAP entry
  gets worse. Not recommended.
\item
  Synchronise \code{sel_a}/reset into the B domain before the latch
  (removes the Q-META exposure, standard glitch-free mux): needs B edges
  to retire B, same deadlock with a stopped clock. Only worthwhile if
  the contract "\code{scan_clk_i} idles low at every switch" is ever
  dropped.
\item
  Upstream (BerkeleyLab/\allowbreak{}Bedrock) report: treat latch cells
  as sequential launch points, and count distinct domains rather than
  pins. A low-value report (the tool is "crude WIP" by its own header).
\item
  Alternate enforcement (what is applied plus
  \code{check_scan_clk_rst.py} TAP-prefix clocking and the unit suite).
\item
  Interaction with open \prref{274} (xtaa: \code{cdc_sync} gains
  \code{RESET_VAL}, edits \code{cdc_config.yml} and
  \code{CDC_SNITCH_POC.md} (\cref{tm:flow:cdc-snitch-poc})): both instances here keep the default reset
  value 0, which the switch relies on
  (\code{b_off_in_a}/\code{a_on_in_b} reset to 0 by design), so the RTL
  is unaffected as long as the default stays 0; the
  \code{cdc_config.yml} edits will conflict textually at the end of
  \code{required_synchronisers} and need a line-union.
\end{itemize}

\subsection{\texorpdfstring{Crypto scan wipe: hardware zeroize on
scan entry and exit (bead \code{j41m.6},
2026-10-17)}{Crypto scan wipe: hardware zeroize on scan entry and exit (bead , 2026-10-17)}}\label{tm:ch:dft-architecture:1410-crypto-scan-wipe-hardware-zeroize-on-scan-entry-and-exit-bead-j41m6-2026-10-17}\origsec{14.10}

Implements the user decision of 2026-10-15/\allowbreak{}17:
\code{crypto_accel} clears its key- and message-derived state on the
\code{scan_mode} edge, before any shift can happen, as a one-time wipe
that keeps the block scan-testable. Files:
\src{rtl/periph/crypto_accel.sv} (edge detect, wipe, bank clear),
\code{aes128_core.sv} and \code{sha256_core.sv} (new \code{wipe_i}). No
port was added to \code{crypto_accel} or \code{soc_top}.

\begin{enumerate}
\def\labelenumi{\arabic{enumi}.}
\tightlist
\item
  \textbf{Edge detect.}
  \texttt{wipe\_\allowbreak{}w\ =\allowbreak{}\ scan\_\allowbreak{}s2\_\allowbreak{}q\ \textasciicircum{}\ scan\_\allowbreak{}prev\_\allowbreak{}q}:
  a 2-FF synchroniser (\code{scan_s1_q}, \code{scan_s2_q}) plus one
  history flop, all in the crypto clock domain (\code{core_clk}).
  \textbf{A synchroniser is needed.} \code{scan_mode} is quasi-static
  but not synchronous to \code{core_clk}:
  \texttt{scan\_\allowbreak{}mode\_\allowbreak{}o\ =\allowbreak{}\ pin\ |\ tap\_\allowbreak{}scan}
  and \code{tap_scan} is a TCK-domain flop, and on exit the clock handed
  back is the functional PLL clock, which has no phase relationship to
  TCK. The raw level already feeds the 256 key/rk AND masks (a
  quasi-static data use); an \emph{edge} detector needs a clean
  synchronous level, and a metastable first stage must not reach 1,400
  flop D pins. The three flops are deliberately \textbf{unreset} so they
  track the input through reset: a reset (including \code{scan_rst_ni}
  pulsed inside scan mode) neither fabricates an entry edge on release
  nor hides the exit edge
  (\code{test_reset_neither_creates_nor_hides_a_scan_edge}). A reset
  shorter than 3 cycles can at worst produce a spurious wipe, the safe
  direction.
\item
  \textbf{Single-cycle synchronous clear, not an asynchronous clear.}
  \code{wipe_w} is one cycle wide and enters every wiped flop as
  \texttt{if\ (!rst\_\allowbreak{}n\ ||\ wipe\_\allowbreak{}w)}; the
  register-bank words (DOUT0-3, DIGEST0-7, IV0-3) are cleared through
  the bank\textquotesingle s hardware-write port, last in the
  \code{always_comb} so it overrides a same-cycle completion or
  \code{IV3} increment. Latency: \code{scan_mode} changes; edge 1
  samples it, edge 2 moves it to \code{scan_s2_q} (\code{wipe_w} high
  after this edge), \textbf{edge 3 clears}. On entry \code{core_clk} is
  the scan clock (\code{dft_clk_mux} selects it with the same
  \code{scan_mode}). It is not held: \code{wipe_w} is a one-cycle pulse
  per edge and a key written in scan mode persists
  (\code{test_scan_wipe_is_one_shot_so_scan_stays_testable}).
\item
  \textbf{Exit also wipes.} Recommended and implemented: leaving scan
  mode fires the same pulse, so anything loaded while in scan mode (a
  scan-loaded \code{s_q}, \code{msg_q}, IV, or a key written over APB in
  scan mode) never becomes a live key or live data. Tested by
  \code{test_scan_exit_wipes_scan_loaded_state}.
\item
  \textbf{\code{key_valid} drops} (\code{key_seen_q} is cleared), so
  after a wipe a start is the existing hang-free illegal path (done, no
  data). While \code{wipe_w} is high a START takes that path, an
  in-flight AES/SHA operation is aborted to done with no data
  (\texttt{C\_\allowbreak{}BUSY\ -\allowbreak{}>\ C\_\allowbreak{}DONE}),
  and KEY writes are dropped and latch \code{key_write_rejected}; DIN
  pushes in that cycle are cleared by the wipe. Software must not rely
  on a key or operation surviving a \code{scan_mode} change.
\item
  \textbf{What ATE must guarantee, and the limit. SUPERSEDED by 14.14
  (\code{j41m.17}): \code{soc_top} now ANDs \code{scan_en} with
  \code{crypto_accel}\textquotesingle s \code{scan_wipe_done_o}, so no
  shift is possible before edge 3 whatever the tester does; the text
  below is kept as history.} The clear is functional D logic. Scan
  insertion puts its scan mux \emph{after} that logic, so any edge with
  \texttt{scan\_\allowbreak{}en\ =\allowbreak{}\ 1} loads \code{scan_in}
  instead of the cleared value, and the last flop of a chain shows its Q
  at \code{scan_out} before the first edge. Therefore: after
  \code{scan_mode} rises, edges 1 to 3 must be capture-type
  (\texttt{scan\_\allowbreak{}en\ =\allowbreak{}\ 0}) and
  \code{scan_out} must not be sampled before edge 3 (a spare edge is
  recommended). On exit the functional clock restarts on its own and no
  rule applies. \textbf{This protects against software and against a
  cooperating tester; an adversarial tester that shifts at edge 1 reads
  the residue.} I looked for a design that needs no tester rule and
  found none that is both secure and clean under this
  repo\textquotesingle s checks: (a) an asynchronous clear needs a pulse
  generated from a flop (or a delay line), and
  \code{check_scan_clk_rst.py} requires every async reset cone to end at
  a primary input, a constant or a scan-mode mux; it would also lose the
  first shift bit and trip a scan DRC on a reset driven by a non-scan
  flop; (b) a level-held async clear is a reset held during scan, which
  the decision forbids. The robust fix is a handshake, not RTL inside
  the block: \code{dft_ctrl} withholds \code{scan_en_o} (and the first
  \code{scan_clk} edge) until \code{crypto_accel} reports the wipe done.
  That touches \code{dft_ctrl}/\code{soc_top} and the Stage 2 insertion
  (\code{j41m.8}), so it is left as a follow-up rather than done here;
  filed as a note on \code{j41m.8}.
\item
  \textbf{Interaction with the reset muxes.} The wipe does not use or
  alter \code{dft_rst_mux} / \code{scan_rst_ni}: in scan mode
  \texttt{core\_\allowbreak{}rst\_\allowbreak{}n\ =\allowbreak{}\ scan\_\allowbreak{}rst\_\allowbreak{}ni}
  and the existing synchronous reset still clears everything when the
  tester asserts it. \code{check_scan_clk_rst.py} is unchanged and
  passes (Sky130 sv2v netlist: 1,696 sequential cells, 294 async pins,
  28 scan-mode muxes, 0 violations; ASAP7: 3,782 / 294 / 29 / 0): no new
  async pin, no new clock.
\item
  \textbf{Flops wiped} (primary build, \code{SBOX_PARALLEL=16},
  \code{EN_AES=EN_SHA=1}): \code{key_q} 128, \code{key_seen_q} 4,
  \code{msg_q} 512, \code{ctr_op_q} 1; AES \code{s_q} 128, \code{rk_q}
  128, \code{rcon_q} 8, \code{round_q} 4, \code{state_q}; SHA \code{w_q}
  16x32, \code{a_q}..\code{h_q} 8x32, \code{hin_q} 256, \code{t_q} 7,
  \code{state_q}; wrapper \code{state_q}; bank words DOUT 128, DIGEST
  256, IV 128. \textbf{Not wiped:} \code{CTRL} (mode, IRQ enable,
  SHA\_CONT: not secret), sticky \code{done} and
  \code{key_write_rejected}, the three edge flops, and the APB
  bank\textquotesingle s other words (they hold no key or message data).
\item
  \textbf{Tests and mutation.} \code{test_crypto.py}: 8 scan-wipe tests
  (every flop class sampled in the first cycle after the clearing edge,
  because the idle pre-load of \code{s_q}/\code{rk_q} would refresh them
  a cycle later and hide a missing clear), plus
  \code{test_soc_dft_scan.py} wiring test through \code{soc_top}.
  Mutation campaign: \src{tools/verif/mutate_crypto_scan_wipe.py}
  (results in the PR description).
\end{enumerate}

\section{\texorpdfstring{Stage 4 specification (beads \code{j41m.14}
MBIST, \code{j41m.10} control-word
crossing)}{Stage 4 specification (beads  MBIST,  control-word crossing)}}\label{tm:ch:dft-architecture:16-stage-4-specification-beads-j41m14-mbist-j41m10-control-word-crossing}\origsec{16}

Specification written before the RTL; the status table at its end is
updated as each step lands. Labels as before: \textbf{{[}M{]}} measured,
\textbf{{[}R{]}} read from source, \textbf{{[}A{]}} assumed or
estimated. This section \textbf{supersedes} the Stage-4 plan text of
sections 6 and 10 and the "per-bit \code{cdc_sync}, wait \textgreater= 3
cycles" text of sections 6.2 and 9.3 (those are kept as history).
Decisions 4 and 5 (tester-only, TAP-driven) are unchanged.

\subsection{User decisions of 2026-10-17 and what this stage lands}\label{tm:ch:dft-architecture:160-user-decisions-of-2026-10-17-and-what-this-stage-lands}\origsec{16.0}

{\small\setlength{\tabcolsep}{4pt}

{
\begin{longtable}[]{@{}
  >{\raggedright\arraybackslash}p{\dimexpr(\linewidth-4\tabcolsep)*49/1000\relax}
  >{\raggedright\arraybackslash}p{\dimexpr(\linewidth-4\tabcolsep)*498/1000\relax}
  >{\raggedright\arraybackslash}p{\dimexpr(\linewidth-4\tabcolsep)*433/1000\relax}@{}}
\toprule\noalign{}
\begin{minipage}[b]{\linewidth}\raggedright
\#
\end{minipage} & \begin{minipage}[b]{\linewidth}\raggedright
Decision
\end{minipage} & \begin{minipage}[b]{\linewidth}\raggedright
Consequence
\end{minipage} \\
\midrule\noalign{}
\endhead
\bottomrule\noalign{}
\endlastfoot
D1 & Algorithm: March C- \textbf{plus a retention test} (write, pause,
read, programmable pause). Evaluate SMarchCHKB-family support. & Section
16.2 (algorithms), 16.3 (SMarch evaluation). \\
D2 & Scope: collars on \textbf{all} Sky130 SoC SRAM macros, including
the CPU-macro caches. & Section 16.1 (inventory), 16.7 (CPU macro). \\
D3 & Access: TAP-only start/\allowbreak{}status. Include the
\code{j41m.10} TCK -\textgreater{} core crossing (depth-2
\code{cdc_gray_fifo}). & Sections 16.4, 16.5. \\
\end{longtable}
}

}

\textbf{What lands in the PR that carries this section (phases, each
committed separately):} (1) this section; (2) RTL \code{mbist_ctrl},
\code{mem_collar}, the TAP data registers (opcodes 6/7), the crossing
\code{mbist_xing}, the integrating \code{soc_mbist}, with unit and
fault-injection suites; (3) integration in \code{soc_top} for the
\textbf{two SoC-level macros} (\code{sram_controller}, NPU weight
memory); (4) the \textbf{CPU-macro cache collars are designed here
(16.7) but NOT implemented in \code{rv32i_cpu_top}}: they change the
hard-macro boundary and need a CPU re-harden. Recommendation in 16.7:
land them with the \code{j41m.4} re-harden, not before.

\subsection{Memory inventory, reconciled with the
bead\textquotesingle s "x12" {[}R, re-checked 2026-10-17{]}}\label{tm:ch:dft-architecture:161-memory-inventory-reconciled-with-the-beads-x12-r-re-checked-2026-10-17}\origsec{16.1}

\code{grep} of \src{rtl/} for SRAM macro instantiations and the Sky130
SoC config (\src{pnr/sky130/soc/config.json},
\texttt{VERILOG\_\allowbreak{}DEFINES\ \_\allowbreak{}\_\allowbreak{}pnr\_\allowbreak{}\_\allowbreak{}\ SRAM\_\allowbreak{}SKY130},
\code{gpu_top_tieoff.sv}): \textbf{12 macro instances, no others}.

\begin{landscape}
\footnotesize\setlength{\tabcolsep}{3pt}

{
\begin{longtable}[]{@{}
  >{\raggedright\arraybackslash}p{\dimexpr(\linewidth-12\tabcolsep)*48/1000\relax}
  >{\raggedright\arraybackslash}p{\dimexpr(\linewidth-12\tabcolsep)*241/1000\relax}
  >{\raggedright\arraybackslash}p{\dimexpr(\linewidth-12\tabcolsep)*175/1000\relax}
  >{\raggedright\arraybackslash}p{\dimexpr(\linewidth-12\tabcolsep)*58/1000\relax}
  >{\raggedright\arraybackslash}p{\dimexpr(\linewidth-12\tabcolsep)*194/1000\relax}
  >{\raggedright\arraybackslash}p{\dimexpr(\linewidth-12\tabcolsep)*194/1000\relax}
  >{\raggedright\arraybackslash}p{\dimexpr(\linewidth-12\tabcolsep)*70/1000\relax}@{}}
\toprule\noalign{}
\begin{minipage}[b]{\linewidth}\raggedright
\#
\end{minipage} & \begin{minipage}[b]{\linewidth}\raggedright
Instance
\end{minipage} & \begin{minipage}[b]{\linewidth}\raggedright
Macro
\end{minipage} & \begin{minipage}[b]{\linewidth}\raggedright
Size
\end{minipage} & \begin{minipage}[b]{\linewidth}\raggedright
Ports used
\end{minipage} & \begin{minipage}[b]{\linewidth}\raggedright
Domain / clock
\end{minipage} & \begin{minipage}[b]{\linewidth}\raggedright
Lands in this PR
\end{minipage} \\
\midrule\noalign{}
\endhead
\bottomrule\noalign{}
\endlastfoot
1 & \code{u_sram.u_sram_macro} (\code{sram_controller.sv}) &
\code{sky130_sram_4kbyte_1rw1r_32x1024_8} & 1024 x 32 & port 0 write,
\textbf{port 1 read} (functional read path) & core, \code{core_clk} &
yes (soc\_top) \\
2 & \code{u_npu.g_on.u_wmem.u_sram_macro} (\code{npu_weight_mem.sv}) &
same & 1024 x 32 & port 0 RW, port 1 tied off & core, \code{core_clk} &
yes (soc\_top) \\
3-7 & icache \code{u_tag_sram}, 4 x \code{gen_data_sram[*].u_data_sram}
& \code{sky130_sram_1kbyte_1rw1r_32x256_8} & 256 x 32 each & port 0 RW,
port 1 tied off & CPU, \code{cpu_gated_clk}, \textbf{inside
\code{rv32i_cpu_top}} & designed (16.7) \\
8-12 & dcache \code{u_tag_sram}, 4 x \code{gen_data_sram[*].u_data_sram}
& same & 256 x 32 each & same & same & designed (16.7) \\
\end{longtable}
}

\end{landscape}

Reconciliation: "x12" = 2 SoC-level + 10 inside the CPU hard macro,
exactly section 3.1. Not memories for this purpose: the GPU
\code{shared_memory} (32 x \code{sram_1rw_128x32_asap7}) is tied off in
the Sky130 SoC and stays ASAP7-indicative; \code{dma_engine.linebuf},
FIFOs, regfile and cache valid/\allowbreak{}dirty arrays are flops and
go on scan chains (3.2). 147,456 bits in total. No spare rows or columns
exist, so the result is pass/fail with a fail record, not repair.

\subsection{Algorithms, encoding and architecture decision}\label{tm:ch:dft-architecture:162-algorithms-encoding-and-architecture-decision}\origsec{16.2}

\textbf{Architecture: a ROM-microcode sequencer, one per clock domain,
with selectable algorithm; not field-programmable, not hardwired per
algorithm.}

{\footnotesize\setlength{\tabcolsep}{3pt}

{
\begin{longtable}[]{@{}
  >{\raggedright\arraybackslash}p{\dimexpr(\linewidth-6\tabcolsep)*230/1000\relax}
  >{\raggedright\arraybackslash}p{\dimexpr(\linewidth-6\tabcolsep)*250/1000\relax}
  >{\raggedright\arraybackslash}p{\dimexpr(\linewidth-6\tabcolsep)*250/1000\relax}
  >{\raggedright\arraybackslash}p{\dimexpr(\linewidth-6\tabcolsep)*250/1000\relax}@{}}
\toprule\noalign{}
\begin{minipage}[b]{\linewidth}\raggedright
Option
\end{minipage} & \begin{minipage}[b]{\linewidth}\raggedright
Area on Sky130 {[}A, to be replaced by the measured number in 16.9{]}
\end{minipage} & \begin{minipage}[b]{\linewidth}\raggedright
Flexibility
\end{minipage} & \begin{minipage}[b]{\linewidth}\raggedright
Risk
\end{minipage} \\
\midrule\noalign{}
\endhead
\bottomrule\noalign{}
\endlastfoot
Hardwired FSM per algorithm & smallest for ONE algorithm; each extra
algorithm adds a state set & none after tape-out & each algorithm is its
own verification job \\
\textbf{ROM microcode, algorithm select (chosen)} & ROM is a constant
\code{case} function: synthesises to \textasciitilde100-300 gates,
\textbf{no flops}; sequencer \textasciitilde60-80 flops & six
algorithms, three backgrounds knobs, programmable pause, all selectable
through the TAP & one engine verified once; adding an algorithm is a ROM
row \\
Field-programmable microcode (loadable through the TAP) & + 22 elements
x \textasciitilde16 bits = \textasciitilde350 flops per domain
(\textasciitilde7.5 k um2 per domain at 21.3 um2 per \code{dfxtp_2},
{[}A{]}) plus a load register & maximum (silicon-debug algorithms) & a
corrupted or mis-loaded program silently becomes a different test; needs
its own readback check \\
\end{longtable}
}

}

Justification: memory is 88 \% of placed area and the test is
tester-only, so the controller must be small, but the cost of a loadable
program (\textasciitilde15 k um2 for two domains) is about the cost of
the whole chosen controller, and the algorithms the project needs are
known. If a silicon-debug algorithm is ever required the ROM becomes a
flop array behind the same sequencer (the sequencer already consumes a
microinstruction word; only its source changes).

\textbf{Sequencer.} One microinstruction = one March element: direction,
up to three read/\allowbreak{}write operations, a data-polarity per
operation, a \code{ck} flag (polarity XOR address bit 0, giving a
checkerboard), \code{pause_before}, and a \code{lane} kind (the
byte-lane element). The engine sweeps the element over the address
range, issues one operation per cycle, drains
\texttt{RD\_\allowbreak{}LAT\ +\allowbreak{}\ 1} cycles at each element
boundary (so background and expected data are constant while reads are
in flight), then fetches the next microinstruction.

\textbf{Data backgrounds} (word-oriented memory,
\texttt{log2(32)\ +\allowbreak{}\ 1\ =\allowbreak{}\ 6}):
\code{0x00000000}, \code{0x55555555}, \code{0x33333333},
\code{0x0F0F0F0F}, \code{0x00FF00FF}, \code{0x0000FFFF}. \code{w0}
writes the background \code{P}, \code{w1} writes
\texttt{\textasciitilde{}P}; with \code{ck} the polarity is XORed with
address bit 0. \texttt{bg\_\allowbreak{}sweep\ =\allowbreak{}\ 1}
repeats the selected algorithm for all six backgrounds.

\textbf{Algorithms} (3-bit \code{alg}; reserved codes finish at once
with \code{cfg_err}):

{\footnotesize\setlength{\tabcolsep}{3pt}

{
\begin{longtable}[]{@{}
  >{\raggedright\arraybackslash}p{\dimexpr(\linewidth-6\tabcolsep)*48/1000\relax}
  >{\raggedright\arraybackslash}p{\dimexpr(\linewidth-6\tabcolsep)*218/1000\relax}
  >{\raggedright\arraybackslash}p{\dimexpr(\linewidth-6\tabcolsep)*357/1000\relax}
  >{\raggedright\arraybackslash}p{\dimexpr(\linewidth-6\tabcolsep)*357/1000\relax}@{}}
\toprule\noalign{}
\begin{minipage}[b]{\linewidth}\raggedright
\code{alg}
\end{minipage} & \begin{minipage}[b]{\linewidth}\raggedright
Name
\end{minipage} & \begin{minipage}[b]{\linewidth}\raggedright
Elements
\end{minipage} & \begin{minipage}[b]{\linewidth}\raggedright
Detects (theory)
\end{minipage} \\
\midrule\noalign{}
\endhead
\bottomrule\noalign{}
\endlastfoot
0 & \code{MARCH_CM} & March C- (10N):
\texttt{up(w0)\ up(r0,\allowbreak{}w1)\ up(r1,\allowbreak{}w0)\ down(r0,\allowbreak{}w1)\ down(r1,\allowbreak{}w0)\ down(r0)}
& SAF, TF, AF, unlinked CFin/\allowbreak{}CFid/\allowbreak{}CFst \\
1 & \code{RETENTION} & \code{up(w0)}; PAUSE; \code{up(r0,w1)}; PAUSE;
\code{up(r1)} & data-retention faults (cell decays within the pause) \\
2 & \code{MARCH_CM_RET} (production default) & \code{MARCH_CM} then
\code{RETENTION} & both of the above \\
3 & \code{CKBD} & checkerboard
\texttt{up(wc)\ up(rc)\ up(w\textasciitilde{}c)\ up(r\textasciitilde{}c)},
then March C- on the checkerboard background & adjacent-cell shorts and
leakage, neighbourhood-sensitive coupling in logical space \\
4 & \code{LANE} & \code{up(w0)};
\texttt{up(w\ \textasciitilde{}P\ lane0,\allowbreak{}\ r,\allowbreak{}\ w\ \textasciitilde{}P\ lane1,\allowbreak{}\ r,\allowbreak{}\ w\ \textasciitilde{}P\ lane2,\allowbreak{}\ r,\allowbreak{}\ w\ \textasciitilde{}P\ lane3,\allowbreak{}\ r)};
\texttt{up(w\ P\ lane3,\allowbreak{}\ r,\allowbreak{}\ w\ P\ lane2,\allowbreak{}\ r,\allowbreak{}\ w\ P\ lane1,\allowbreak{}\ r,\allowbreak{}\ w\ P\ lane0,\allowbreak{}\ r)}
with \code{wmask} = one byte lane at a time, every read compared against
the cumulative expected word & write-mask (\code{wmask0}) lane faults,
which a full-word March cannot see. Both orders are needed: a lane whose
mask bit is stuck at 1 is visible only if some other lane is written
before it, so the forward pass exposes lanes 1-3 and the reverse pass
lanes 0-2 (\textbf{found by the fault-injection suite: the first,
forward-only version missed lane 0}) \\
5 & \code{FULL} & 2 + 3 + 4 in sequence & everything above \\
\end{longtable}
}

}

Retention pause:
\texttt{cycles\ =\allowbreak{}\ pause\_\allowbreak{}mant\ <<\ (2\ *\ pause\_\allowbreak{}exp)},
\code{pause_mant} 8 bits, \code{pause_exp} 3 bits, so 0 to 255
\textless\textless{} 14 = 4,177,920 cycles (104 ms at 40 MHz, 56 ms at
75 MHz). Longer pauses are made by the tester running the algorithm
repeatedly, or by \code{RETENTION} alone followed by re-reads (the
sticky fail record is per run).

\textbf{Microcode address map} (the value reported in \code{rec_elem}):
\code{MARCH_CM} 0-5, \code{RETENTION} 6-8, \code{CKBD} 9-18, \code{LANE}
19-21 (19 \code{up(w0)}, 20 forward lane sequence, 21 reverse lane
sequence). Algorithm ranges: 0 = 0-5, 1 = 6-8, 2 = 0-8, 3 = 9-18, 4 =
19-21, 5 = 0-21. Cycle count of one run = 1 + sum over its elements of
(1 + pause + nops x 2\^{}AW + 2), pinned by
\code{test_cycle_counts_pin_the_sequencing}.

\textbf{Port 1.} Memories with a functional read port 1 (\code{P1_MASK};
the SoC SRAM) read \textbf{both} ports in every read operation and
compare both against the same expected word, so the functional read path
(port 1) is tested in the same run at no extra cycles. Not covered:
simultaneous port-0 write / port-1 read of different addresses (the
section 6.1 wish list); stated as a limitation.

\subsection{SMarchCHKB-family evaluation (decision D1)}\label{tm:ch:dft-architecture:163-smarchchkb-family-evaluation-decision-d1}\origsec{16.3}

The user wrote \code{smarchchkd}. What is public and what is not: the
Siemens/\allowbreak{}Mentor Tessent algorithm family is named
\code{SMarch}, \code{SMarchCHKB}, \code{SMarchCHKBvcd},
\code{SMarchCHKBcil}, ...
(\href{https://www.scribd.com/document/355278009/Tessent-AppNote-Memory-Shared-BUS}{Tessent
AppNote MG576062 on Scribd} shows \code{SMarchCHKBcil} as the default
and discusses \code{SMarchCHKBvcd} and
bit-grouping/\allowbreak{}word-size limits); the generic industry name
is \code{SMarchCKBD}, "Serial March and Checkerboard", whose stated
purpose is "to detect memory failures using either fast row access or
fast column access"
(\href{https://www.design-reuse.com/articles/45915/memory-testing-self-repair-mechanism.html}{design-reuse
article}). \textbf{The element sequences of the Tessent algorithms are
proprietary and are not in any source this stage could reach}; the
article gives only the plain checkerboard
(\texttt{w\ c,\allowbreak{}\ r\ c,\allowbreak{}\ w\ \textasciitilde{}c,\allowbreak{}\ r\ \textasciitilde{}c},
up) and does not define SMarchCKBD\textquotesingle s sequence, coverage
or retention handling.

Decision: \textbf{implement an open SMarchCKBD-style algorithm
(\code{CKBD}, alg 3) and do not claim Tessent compatibility.} It
combines the two published ingredients (checkerboard and March) on one
engine: the checkerboard sequence above, then March C- with data
polarity tied to address parity (\code{ck}), swept over the six
backgrounds if \code{bg_sweep} is set. The closest named variant is
\code{SMarchCHKB}; \code{vcd} (probably "vertical column disturb" or
"verify column data", unverified) and \code{cil} (unverified) are NOT
implemented. The checkerboard is in \textbf{logical}
address/\allowbreak{}bit space: the physical row/\allowbreak{}column
mapping of the OpenRAM macros is scrambled and not in the
LEF/\allowbreak{}Liberty, so "adjacent" is a logical notion; the
6-background sweep plus address-parity polarity is the practical
substitute and is the stated limit of the claim. Coverage claims for
\code{CKBD} are what the fault-injection suite measures (coupling
between logically adjacent cells, address-parity-dependent faults), not
a Tessent figure. Adding a \code{vcd}/\code{cil}-like sequence later is
a ROM row (16.2), no sequencer change.

\subsection{Controller and collar interfaces}\label{tm:ch:dft-architecture:164-controller-and-collar-interfaces}\origsec{16.4}

\begin{lstlisting}[language={},caption={Diagram source (Mermaid)}]
flowchart LR
  TAP[jtag_tap, TCK] -->|MBIST_CTRL 29b, wr strobe| X[mbist_xing: one cdc_gray_fifo per domain]
  X -->|ctl, ungated core_clk| C1[mbist_ctrl core]
  X -->|ctl, ungated cpu_core_clk| C2[mbist_ctrl cpu, in CPU macro]
  C1 -->|broadcast ce/we/addr/din/wmask/rd1| K1[mem_collar sram_controller]
  C1 --> K2[mem_collar npu_weight_mem]
  K1 --- M1[4KB macro]
  K2 --- M2[4KB macro]
  M1 -->|dout0/dout1| C1
  M2 -->|dout0| C1
  C2 --> K3[10 collars I/D cache]
  C1 -->|status 69b, 2-FF + done delayed| TAP
  C2 -->|status 69b| TAP
\end{lstlisting}

\textbf{\code{mem_collar}} (pure input-side mux, on the functional side
of the \code{ifdef} arm, placed in the wrapper next to the macro
instance). \code{bist_en_i} selects:
\texttt{csb0\_\allowbreak{}m\ =\allowbreak{}\ bist\_\allowbreak{}en\ ?\ \textasciitilde{}b\_\allowbreak{}ce\ :\allowbreak{}\ csb0\_\allowbreak{}f},
\texttt{web0\_\allowbreak{}m\ =\allowbreak{}\ bist\_\allowbreak{}en\ ?\ \textasciitilde{}b\_\allowbreak{}we\ :\allowbreak{}\ web0\_\allowbreak{}f},
\texttt{wmask0\_\allowbreak{}m\ =\allowbreak{}\ bist\_\allowbreak{}en\ ?\ b\_\allowbreak{}wmask\ :\allowbreak{}\ wmask0\_\allowbreak{}f},
\code{addr0_m}, \code{din0_m} likewise; port 1 (only when
\code{HAS_P1}):
\texttt{csb1\_\allowbreak{}m\ =\allowbreak{}\ bist\_\allowbreak{}en\ ?\ \textasciitilde{}(b\_\allowbreak{}ce\ \&\ \textasciitilde{}b\_\allowbreak{}we\ \&\ b\_\allowbreak{}rd1)\ :\allowbreak{}\ csb1\_\allowbreak{}f},
\code{addr1_m}. Macro read data is \textbf{not} muxed:
\code{dout0}/\code{dout1} go to the functional logic unchanged and are
also exported to the controller (an extra fan-out, no extra mux on the
read path). \code{bist_en_i} is a flop output of the controller, so no
mux select can glitch. With
\texttt{bist\_\allowbreak{}en\ =\allowbreak{}\ 0} the collar is a
pass-through and the design is bit-identical to today (checked by the
existing \code{soc_all} with the collar present).

\textbf{\code{mbist_ctrl}} (parameters \code{N_MEM}, \code{AW},
\texttt{DW\ =\allowbreak{}\ 32},
\texttt{RD\_\allowbreak{}LAT\ =\allowbreak{}\ 2}, \code{P1_MASK}):
sequencer, data generator, one comparator per memory and port,
per-memory sticky \code{fail} and first-fail record. Broadcast outputs
are \textbf{registered}
(\texttt{b\_\allowbreak{}ce\ b\_\allowbreak{}we\ b\_\allowbreak{}wmask\ b\_\allowbreak{}addr\ b\_\allowbreak{}din\ b\_\allowbreak{}rd1}),
so the collar input pins are driven from flops. Read latency 2 is the
Sky130 macro contract {[}R,
\src{sim/sky130_sram_4kbyte_1rw1r_32x1024_8.sv}: inputs captured at the
posedge, operation at the following negedge, data sampled at the next
posedge; the same timing holds for the 1 KB model{]}, and is a parameter
checked by the unit suite against the real sim models.

\subsection{TAP registers, control word, crossing (j41m.10 folded
in)}\label{tm:ch:dft-architecture:165-tap-registers-control-word-crossing-j41m10-folded-in}\origsec{16.5}

\textbf{Opcodes (frozen, section 15.3, now implemented):} \texttt{4'h6}
\code{MBIST_CTRL}, \texttt{4'h7} \code{MBIST_STATUS}. Section 15.6
deviation 5 is retired; the unit suites that asserted BYPASS for 6/7 are
changed.

\code{MBIST_CTRL} (29 bits, LSB first, update register in
\code{jtag_tap}, Capture-DR reads back the current value, reset 0 =
idle):

{\small\setlength{\tabcolsep}{4pt}

{
\begin{longtable}[]{@{}
  >{\raggedright\arraybackslash}p{\dimexpr(\linewidth-4\tabcolsep)*82/1000\relax}
  >{\raggedright\arraybackslash}p{\dimexpr(\linewidth-4\tabcolsep)*175/1000\relax}
  >{\raggedright\arraybackslash}p{\dimexpr(\linewidth-4\tabcolsep)*723/1000\relax}@{}}
\toprule\noalign{}
\begin{minipage}[b]{\linewidth}\raggedright
Bits
\end{minipage} & \begin{minipage}[b]{\linewidth}\raggedright
Field
\end{minipage} & \begin{minipage}[b]{\linewidth}\raggedright
Meaning
\end{minipage} \\
\midrule\noalign{}
\endhead
\bottomrule\noalign{}
\endlastfoot
0 & \code{run} & level. 1 starts the selected algorithm when the
controller is idle; 0 aborts and clears all status \\
3:1 & \code{alg} & 16.2 \\
4 & \code{bg_sweep} & 1 = all six backgrounds \\
7:5 & \code{bg_sel} & background 0-5 (6, 7 = \code{cfg_err}) \\
15:8 & \code{pause_mant} & retention pause mantissa \\
18:16 & \code{pause_exp} & retention pause exponent \\
22:19 & \code{mask_sel} & memory whose first-failure XOR mask is
captured (latched at start) \\
26:23 & \code{rd_sel} & memory whose fail record is shown in the status
(live, may change after \code{done}) \\
27 & \code{dom_core} & enable the core-clock controller \\
28 & \code{dom_cpu} & enable the CPU-clock controller \\
\end{longtable}
}

}

\code{MBIST_STATUS} (read-only, captured at Capture-DR, 140 bits, LSB
first): bit 0 \code{tck_ovf} (sticky; control-word FIFO overflow), bit 1
\code{cpu_present}, bits 70:2 core block, bits 139:71 CPU block. One
69-bit block per domain:

{
\begin{longtable}[]{@{}
  >{\raggedright\arraybackslash}p{\dimexpr(\linewidth-2\tabcolsep)*99/1000\relax}
  >{\raggedright\arraybackslash}p{\dimexpr(\linewidth-2\tabcolsep)*881/1000\relax}@{}}
\toprule\noalign{}
\begin{minipage}[b]{\linewidth}\raggedright
Bits
\end{minipage} & \begin{minipage}[b]{\linewidth}\raggedright
Field
\end{minipage} \\
\midrule\noalign{}
\endhead
\bottomrule\noalign{}
\endlastfoot
0 & \code{done} (all selected tests complete; held until \code{run} =
0) \\
1 & \code{busy} \\
2 & \code{cfg_err} (reserved \code{alg} or \code{bg_sel}) \\
3 & \code{fail_any} \\
13:4 & \code{fail[9:0]} per-memory go/no-go (sticky) \\
23:14 & \code{rec_addr} first-fail word address of memory
\code{rd_sel} \\
28:24 & \code{rec_elem} element index (microcode address) \\
31:29 & \code{rec_slot} operation slot within the element \\
34:32 & \code{rec_bg} background index \\
36:35 & \code{rec_port} \{port 1 mismatch, port 0 mismatch\} \\
68:37 & \code{err_mask} XOR of read vs expected at the first failure of
memory \code{mask_sel} \\
\end{longtable}
}

Go/no-go per memory = \code{fail[i]}; "good" requires
\texttt{done\ =\allowbreak{}\ 1} and
\texttt{fail\_\allowbreak{}any\ =\allowbreak{}\ 0} and
\texttt{cfg\_\allowbreak{}err\ =\allowbreak{}\ 0}. A block that never
ran reads \texttt{done\ =\allowbreak{}\ 0}, so an absent or stuck
controller cannot read as a pass.

\textbf{Crossing (j41m.10).} Resolution of the 15 open questions of \prref{276} (the recommendation of each, taken because the user folded
\code{j41m.10} into this stage; any of them can still be overturned at
review):

{\footnotesize\setlength{\tabcolsep}{3pt}

{
\begin{longtable}[]{@{}
  >{\raggedright\arraybackslash}p{\dimexpr(\linewidth-2\tabcolsep)*86/1000\relax}
  >{\raggedright\arraybackslash}p{\dimexpr(\linewidth-2\tabcolsep)*894/1000\relax}@{}}
\toprule\noalign{}
\begin{minipage}[b]{\linewidth}\raggedright
Q
\end{minipage} & \begin{minipage}[b]{\linewidth}\raggedright
Resolution
\end{minipage} \\
\midrule\noalign{}
\endhead
\bottomrule\noalign{}
\endlastfoot
1, 2 & Applies to the MBIST word only; scan controls unchanged (section
15.5, crossing 2 stands) \\
3, 4 & One \code{cdc_gray_fifo} \textbf{per destination domain}; read
side the \textbf{ungated} \code{core_clk} / \code{cpu_core_clk}, never a
gated clock \\
5 & Write strobe is a falling-edge TAP flop set in Update-DR, so exactly
one rising TCK edge writes the word: \textbf{+1 TCK after Update-DR}.
Tester contract: after Update-DR of \code{MBIST_CTRL}, at least one more
TCK before relying on the write (any following Capture-DR satisfies
it) \\
6 & Overflow (strobe while \code{wr_ready} low) sets the sticky
\code{tck_ovf} in \code{MBIST_STATUS} bit 0; cleared by TRST or by an
\textbf{accepted} \texttt{run\ =\allowbreak{}\ 0} word (which includes
the Test-Logic-Reset default write). The TAP never stalls \\
7 & \textbf{DEPTH 2} with the overflow flag. Reason: the destination
only needs the latest word (\code{run} is a level, a held register), a
lost intermediate word is reported, and \code{cdc_gray_fifo} at DEPTH 2
is the variant \prref{276} validated; \code{async_axi_fifo} is not
involved \\
8 & Both FIFO resets are derived from \code{trst_ni} alone (write side:
the TAP\textquotesingle s \code{rst_tck_n}; read side: a
\code{cdc_reset_sync} of \code{trst_ni} into the destination clock, with
the scan-mode reset override). The controller is reset by that
\textbf{and} the functional domain reset, so \code{rst_n_i} can never
start MBIST, never touches the TAP, and cannot leave phantom FIFO
entries (functional reset does not reset the FIFO pointers) \\
9 & Entering Test-Logic-Reset (any path) writes the default word
(\texttt{run\ =\allowbreak{}\ 0}) once, so five TMS-high clocks abort a
running MBIST. The FIFO is empty or holds at most one stale word at that
time (DEPTH 2 is enough because the strobe is one per Update-DR and the
tester pace is slower than the destination) \\
10 & \code{SYNC_STAGES} 2 \\
11 & \code{check_scan_clk_rst.py} gets a narrow extension (flops under a
named prefix whose clock cone is exactly \code{tck_i} are TAP-domain),
\code{--tap-allowed-nets} and \code{--expect-tap-flops} updated;
\code{required_synchronisers} entries added for the two pointer
synchronisers; \code{tck_i} stays in \code{reset_domains} \\
12 &
\texttt{set\_\allowbreak{}max\_\allowbreak{}delay\ -\allowbreak{}datapath\_\allowbreak{}only}
is not implemented by OpenSTA
(\code{phase5_soc_multiclock.sdc:112-123}); the FIFO crossings follow
the form used there for the existing FIFOs. \textbf{Resolved by bead
\bead{y5mn} (2026-10-25):} budgets live in
\src{pnr/sky130/soc/constraints/sky130_soc_cdc_check.sdc} section 14
(gray pointers 12.5 ns = 0.5 x the faster period, mem\_q -\textgreater{}
holding register 22.5 ns, 69 status bits per domain 25 ns, the CPU macro
as launch cell with its 35 tied-0 status pins accounted for), checked by
\texttt{make\ -\allowbreak{}C\ pnr\ check-\allowbreak{}cdc-\allowbreak{}timing-\allowbreak{}sky130-\allowbreak{}soc},
negative control \code{CDC_CHECK_MBIST_BUDGET_NS}; \code{sky130_soc.sdc}
exempts the tck -\textgreater{} destination hold (pointer and payload)
as it does for the AXI FIFOs. \code{mbist_force} -\textgreater{}
\code{u_cpu_cg} stays an unsynchronised quasi-static control, not an SDC
row \\
13 & Opcodes 6/7 leave BYPASS in this stage \\
14 & Sections 6 and 9.3 carry a pointer to this section \\
15 & \code{mem_q} (unreset, written on \code{tck_i}) is in the TCK
domain beside the TAP flops: unscanned, covered by N4 (TDR readback:
write \code{MBIST_CTRL}, read it back) \\
\end{longtable}
}

}

\textbf{Status return crossing.}
\code{done}/\code{fail}/\allowbreak{}record/\allowbreak{}mask bits are
stable once \code{done} is high (the controller freezes them). Each
status bit goes through a 2-FF \code{cdc_sync} into \code{tck_i};
\code{done} and \code{busy} pass through \textbf{one extra TCK flop}, so
when the tester sees \texttt{done\ =\allowbreak{}\ 1} every data bit has
had a full TCK period to settle. Reading with
\texttt{done\ =\allowbreak{}\ 0} shows data that may be mid-update; that
is the documented meaning of "not done". A change of \code{rd_sel} after
\code{done} needs \textgreater= 3 TCK before the next Capture-DR
(documented tester rule). Likewise the control-word
FIFO\textquotesingle s write side learns that a slot was freed only
after the read pointer has crossed back into the TCK domain, which needs
TCK edges: \textbf{the tester keeps TCK running in
Run-Test/\allowbreak{}Idle} (found by
\code{test_overflow_is_reported_and_cleared}, which saw a second word
rejected while TCK was stopped).

\textbf{Reset and idle contract.}
\texttt{trst\_\allowbreak{}ni\ =\allowbreak{}\ 0} (the idle value of
every existing suite) holds the FIFO side and the controller in reset:
\texttt{bist\_\allowbreak{}en\ =\allowbreak{}\ 0}, all collars pass
through, no behaviour change. A functional reset leaves the controller
idle and clears the destination holding register, so the tester rewrites
\code{MBIST_CTRL}.

\subsection{\texorpdfstring{Clocking, scan\_\allowbreak{}mode,
PMU}{Clocking, scan\_mode, PMU}}\label{tm:ch:dft-architecture:166-clocking-scan-mode-pmu}\origsec{16.6}

\begin{itemize}
\tightlist
\item
  \textbf{Core controller and its two collars:} \code{core_clk}, the
  same clock as \code{sram_controller} and the NPU {[}R{]}. MBIST runs
  on the functional clock (at speed for the memories).
\item
  \textbf{CPU controller:} \code{cpu_gated_clk} (the CPU
  macro\textquotesingle s \code{clk_i}), fed by a control word whose
  FIFO read side is the \textbf{ungated} \code{cpu_core_clk}. The two
  are the same source and frequency (\code{cpu_gated_clk} is
  \code{cpu_core_clk} through \code{u_cpu_cg}); the existing design
  already treats signals from \code{cpu_core_clk} logic into
  \code{cpu_gated_clk} logic as one domain {[}R
  \code{soc_top.sv:540-660}{]}.
\item
  \textbf{MBIST clock force.} \code{u_cpu_cg} is also gated by the PMU.
  The CPU controller therefore needs the gate forced open while a
  CPU-domain test runs: the TAP-launched \code{mbist_clk_force}
  (quasi-static, a falling-edge TAP flop) is OR-ed into the
  gate\textquotesingle s \code{test_en}. Same quasi-static contract as
  \code{test_en} today (section 15.5, crossing 2). Precondition (tester
  contract): the CPU domain is powered and not isolated (PMU default
  after reset); MBIST does not drive the PMU. Core-domain memories are
  not on gated clocks.
\item
  \textbf{scan\_\allowbreak{}mode:} the controller\textquotesingle s and
  the destination FIFO\textquotesingle s asynchronous resets pass the
  same \code{dft_rst_mux} / \code{cdc_reset_sync} scan override as every
  other register, so in scan mode they are tester-controlled and
  \code{bist_en} cannot be 1 unless the chains put it there. The
  controller flops are ordinary scan flops. The TCK-side FIFO flops are
  TAP-domain (unscanned). In scan mode the SRAM macro is a black box
  exactly as before (section 6.3); MBIST adds no new scan-mode
  behaviour.
\item
  \textbf{Destructive test.} MBIST overwrites every tested memory. After
  a run the SoC needs a functional reset and memory reload (tester
  contract), the same as after scan.
\end{itemize}

\subsection{CPU-macro cache collars: design, proposal, re-harden
need}\label{tm:ch:dft-architecture:167-cpu-macro-cache-collars-design-proposal-re-harden-need}\origsec{16.7}

Ten memories (16.1 rows 3-12) sit inside \code{rv32i_cpu_top}; reaching
them needs a controller \textbf{inside} the macro (clock \code{clk_i} =
\code{cpu_gated_clk}) and these new ports. The existing 403 pins are
placed by \code{pin_order.cfg} with
\texttt{ERRORS\_\allowbreak{}ON\_\allowbreak{}UNMATCHED\_\allowbreak{}IO\ both}
{[}R \src{pnr/sky130/cpu/pin_order.README.md}{]}, so every new pin must
be added to it.

{
\begin{longtable}[]{@{}
  >{\raggedright\arraybackslash}p{\dimexpr(\linewidth-6\tabcolsep)*206/1000\relax}
  >{\raggedright\arraybackslash}p{\dimexpr(\linewidth-6\tabcolsep)*62/1000\relax}
  >{\raggedright\arraybackslash}p{\dimexpr(\linewidth-6\tabcolsep)*72/1000\relax}
  >{\raggedright\arraybackslash}p{\dimexpr(\linewidth-6\tabcolsep)*640/1000\relax}@{}}
\toprule\noalign{}
\begin{minipage}[b]{\linewidth}\raggedright
Port
\end{minipage} & \begin{minipage}[b]{\linewidth}\raggedright
Dir
\end{minipage} & \begin{minipage}[b]{\linewidth}\raggedright
Width
\end{minipage} & \begin{minipage}[b]{\linewidth}\raggedright
Source / sink
\end{minipage} \\
\midrule\noalign{}
\endhead
\bottomrule\noalign{}
\endlastfoot
\code{mbist_run_i} & in & 1 & CPU-domain holding register
(\texttt{run\ \&\ dom\_\allowbreak{}cpu}) \\
\code{mbist_cfg_i} & in & 26 &
\texttt{alg\ bg\_\allowbreak{}sweep\ bg\_\allowbreak{}sel\ pause\ mask\_\allowbreak{}sel\ rd\_\allowbreak{}sel} \\
\code{mbist_status_o} & out & 69 & the status block of 16.5 \\
\textbf{total} & & \textbf{96} & quasi-static, no timing path other than
the ordinary sync-to-clk \\
\end{longtable}
}

\code{mbist_clk_force} stays in \code{soc_top} (it drives
\code{u_cpu_cg}, outside the macro). Pin-count alternative if the pin
budget is contested: return the status in 16-bit slices selected by a
3-bit port (46 pins), at the cost of a multi-step tester readout and a
different CPU \code{MBIST_STATUS} layout; \textbf{not} chosen because it
would change the frozen TDR layout later.

Inside the macro: 10 \code{mem_collar} (tag and data of each cache,
collar between the cache logic and the macro instance in
\code{rv32i_icache.sv} / \code{rv32i_dcache.sv}), one \code{mbist_ctrl}
(\texttt{N\_\allowbreak{}MEM\ =\allowbreak{}\ 10},
\texttt{AW\ =\allowbreak{}\ 8},
\texttt{P1\_\allowbreak{}MASK\ =\allowbreak{}\ 0}), and the cache
modules gain \code{mbist_*} broadcast ports. The functional side of the
\code{SRAM_ASAP7} / FPGA / FreePDK45 arms stays untouched (the collar
sits only in the \code{SRAM_SKY130} arm and in a pass-through shim for
the others).

\textbf{Proposal: do NOT land the CPU-macro collars in the PR that
carries this section; land them with the \code{j41m.4} CPU-macro
scan-boundary re-harden.} Reasons: (1) both need the same boundary
change and the same \textasciitilde6-7 h CPU + SoC re-harden round
{[}section 5, A{]}; two rounds cost another 6-7 h; (2) both interact
with \code{pin_order.cfg} and the lock of lu09 (\#294), so one pin
change is safer than two; (3) the CPU timing margin is the thinnest in
the design (+0.174 ns, section 5) and a collar mux on ten SRAM
address/\allowbreak{}data pins plus a controller is exactly the kind of
change that may need the period relaxation of open item N7, which should
be measured once with scan muxes included; (4) everything the CPU side
needs is already built and unit-tested here (\code{mbist_ctrl} at
\texttt{N\_\allowbreak{}MEM\ =\allowbreak{}\ 10},
\texttt{AW\ =\allowbreak{}\ 8}, the CPU domain of
\code{mbist_xing}/\code{soc_mbist}, the TDR block), so the follow-up is
wiring and a re-harden, not new verification. \textbf{Re-harden is
required; none was run.} In this PR \code{soc_top} carries the
CPU-domain FIFO/\allowbreak{}holding-register path behind a parameter
(\texttt{MBIST\_\allowbreak{}CPU\ =\allowbreak{}\ 0}), the CPU status
block reads 0 and \texttt{cpu\_\allowbreak{}present\ =\allowbreak{}\ 0}.

\subsection{Verification required of this stage (non-negotiable per
the bead)}\label{tm:ch:dft-architecture:168-verification-required-of-this-stage-non-negotiable-per-the-bead}\origsec{16.8}

\begin{itemize}
\tightlist
\item
  \textbf{A BIST that always passes (or always fails) is the failure to
  rule out.} Fault-injecting SRAM models (same timing as the Sky130 sim
  models, equivalence-checked against them with faults off) support:
  stuck-at 0/1, transition up/down, inversion coupling, idempotent
  coupling, state coupling, address-decoder faults (no access, alias,
  multiple access), write-mask lane stuck, and retention decay with a
  configurable decay time.
\item
  An independent Python golden model of the algorithms (written from the
  algorithm definitions in 16.2, not from the ROM) predicts, for each
  (algorithm, fault), whether it is detected and the first-fail address
  / element / slot; the RTL record must match exactly. Faults an
  algorithm is NOT expected to detect (for example a retention fault
  under \code{MARCH_CM}) must pass. A no-fault run of every algorithm
  and background must pass with
  \texttt{fail\_\allowbreak{}any\ =\allowbreak{}\ 0,\allowbreak{}\ done\ =\allowbreak{}\ 1}.
\item
  Retention: a decay time T with pause \textless{} T passes and pause
  \textgreater{} T fails (the programmable pause is meaningful).
\item
  TAP-driven end-to-end test (\code{MBIST_CTRL} write, poll
  \code{MBIST_STATUS}, per-memory go/no-go and fail address through the
  TAP), including TLR abort, FIFO overflow flag, and the crossing at
  several TCK/core clock ratios.
\item
  Lint (Verilator -Wall), \code{dft-scan-check} for both PDKs with the
  pinned yosys, sv2v regeneration, the CDC gate, \code{soc_all_ci} plain
  and \code{FPGA=1}.
\end{itemize}

\subsection{Area {[}M, 2026-10-18{]}}\label{tm:ch:dft-architecture:169-area-m-2026-10-18}\origsec{16.9}

Method: \code{sv2v} + pinned yosys 0.62
\texttt{synth\ -\allowbreak{}flatten;\ dfflibmap;\ abc\ -\allowbreak{}liberty}
on \code{sky130_fd_sc_hd__tt_025C_1v80}, per block, no timing-driven
sizing (the real flow\textquotesingle s buffering and \code{dfxtp_2}
upsizing come on top). The estimate column is the section 6.3 / 16.9
estimate that this replaces.

{\small\setlength{\tabcolsep}{4pt}

{
\begin{longtable}[]{@{}
  >{\raggedright\arraybackslash}p{\dimexpr(\linewidth-6\tabcolsep)*524/1000\relax}
  >{\raggedright\arraybackslash}p{\dimexpr(\linewidth-6\tabcolsep)*144/1000\relax}
  >{\raggedright\arraybackslash}p{\dimexpr(\linewidth-6\tabcolsep)*253/1000\relax}
  >{\raggedright\arraybackslash}p{\dimexpr(\linewidth-6\tabcolsep)*59/1000\relax}@{}}
\toprule\noalign{}
\begin{minipage}[b]{\linewidth}\raggedright
Block
\end{minipage} & \begin{minipage}[b]{\linewidth}\raggedright
Estimate {[}A{]}
\end{minipage} & \begin{minipage}[b]{\linewidth}\raggedright
Measured {[}M{]}
\end{minipage} & \begin{minipage}[b]{\linewidth}\raggedright
Flops
\end{minipage} \\
\midrule\noalign{}
\endhead
\bottomrule\noalign{}
\endlastfoot
\code{mbist_ctrl}, core domain (N=2, AW=10, port 1 on memory 0) &
\textasciitilde6 k um2 & \textbf{13,314 um2} & 243 \\
\code{mbist_ctrl}, CPU domain (N=10, AW=8) & \textasciitilde15 k um2 &
\textbf{28,009 um2} & 398 \\
\code{mem_collar} AW=10 with port 1 / AW=8 without & \textasciitilde0.7
k um2 each & \textbf{653 / 513 um2} & 0 \\
\code{mbist_xing} (reset sync + depth-2 FIFO + holding register), per
domain & \textasciitilde1.5 k um2 & \textbf{3,562 um2} & 46 \\
\code{mbist_status_sync} (69 per-bit 2-FF + done/busy flop), per domain
& (not estimated) & \textbf{3,503 um2} & 140 \\
\code{jtag_tap} new vs \code{origin/main} (29 + 140
shift/\allowbreak{}update bits, decode, strobe) & \textasciitilde6 k um2
& \textbf{+8,217 um2} (2,620 -\textgreater{} 10,837) & +204 \\
\end{longtable}
}

}

Totals: \textbf{landed in this PR (core domain + 2 collars + TAP delta)
= 13,314 + 1,166 + 3,562 + 3,503 + 8,217 = 29,762 um2} (0.15 pp of the
20.36 M um2 core; 2.6 \% of the 1,150,630 um2 synthesised stdcell area
of section 8); \textbf{CPU-macro follow-up (controller + 10 collars +
xing + status) = 28,009 + 5,130 + 3,562 + 3,503 = 40,204 um2} (inside
the CPU macro: the controller and collars; in the SoC fabric: xing and
status). Everything together is about 70 k um2 (0.34 pp of the core),
above the 45-60 k um2 of section 6.3: the per-memory first-fail record
(23 flops per memory) and the 140 per-bit status synchronisers dominate.
Possible trims if area is contested {[}A, not tried{]}: keep one shared
first-fail record instead of one per memory (the per-memory go/no-go
bits stay), drop the redundant \code{pause_cycles_q} (22 flops).

\subsection{Implementation record (2026-10-18, PR \#306)}\label{tm:ch:dft-architecture:1610-implementation-record-2026-10-18-pr-306}\origsec{16.10}

\begin{itemize}
\tightlist
\item
  \textbf{RTL (landed):} \code{mbist_ctrl}, \code{mem_collar},
  \code{mbist_xing} + \code{mbist_ctl_hold}, \code{mbist_status_sync},
  \code{soc_mbist} (all \src{rtl/soc/dft/}); \code{jtag_tap} /
  \code{dft_ctrl} gain \code{MBIST_CTRL} / \code{MBIST_STATUS} and the
  write strobe; \code{soc_top} instantiates \code{soc_mbist} (core
  domain) and wires the collar bus into \code{sram_controller} and
  \code{npu_top} -\textgreater{} \code{npu_weight_mem}
  (SRAM\_\allowbreak{}SKY130 arm only; every other arm ties the new
  ports off). \code{soc_top} parameter \code{MBIST_CPU} (default 0)
  carries the CPU-domain crossing; the CPU macro is NOT changed.
\item
  \textbf{Found by the verification, fixed:} (1) the first LANE
  algorithm (forward lane order only) missed a write-mask
  lane-0-stuck-at-1 fault: a reverse lane pass (microcode address 21)
  was added. (2) The control-word FIFO\textquotesingle s write side
  learns of a freed slot only when TCK runs (read pointer must cross
  back): tester rule added to 16.5. (3) The Test-Logic-Reset default
  write occupies a FIFO slot until the destination clock reads it
  (visible only with a very slow destination clock).
\item
  \textbf{Verification:} \code{mbist_ctrl} (3 build points),
  \code{mbist_tap} (3 clock ratios), \code{mbist_soc_slice} (real
  wrappers), \code{mbist_fi_equiv} (FI model == real simulation model),
  the existing \code{jtag_tap} / \code{dft_ctrl} suites against the
  extended TAP model, \code{test_dft_scan_check.py} (+5 tests for the
  TCK-block rules). Results are in the PR description.
\item
  \textbf{Checker changes:} \src{tools/dft/check_scan_clk_rst.py}
  \code{--tck-prefix} / \code{--tck-sink} / \code{--expect-tck-flops}
  and the \code{SYNC=CONSUMER@any} form of \code{--expect-reset-sync} (a
  reset synchroniser in a functional clock domain); \src{pnr/Makefile}
  \code{DFT_TAP_CHECK_ARGS} carries the three declared TAP outputs and
  the two sink prefixes.
\item
  \textbf{Not done / open:} CPU-macro collars and re-harden (16.7); SDC
  for the new crossings (item 12 of the table in 16.5, STA owner);
  \code{required_synchronisers} entries in
  \src{tools/cdc/cdc_config.yml} for the FIFO pointer synchronisers (see
  the PR for the CDC gate result); scan testability of the TCK-side FIFO
  storage (unscanned by N4); metastability is not modelable here.
\end{itemize}

\subsection{\texorpdfstring{CPU-macro cache collars: RTL step
(bead \code{c54a}, 2026-10-18, PR
\#313)}{CPU-macro cache collars: RTL step (bead , 2026-10-18, PR \#313)}}\label{tm:ch:dft-architecture:1611-cpu-macro-cache-collars-rtl-step-bead-c54a-2026-10-18-pr-313}\origsec{16.11}

User decision 2026-10-18: combined re-harden. This step is RTL, tests
and pin planning only; the CPU macro re-harden and the one SoC re-harden
are PD steps that follow. It supersedes the "not implemented" statements
of 16.0 (4) and 16.7.

\textbf{Architecture {[}R, as implemented{]}}

\begin{itemize}
\tightlist
\item
  \textbf{Controller inside the macro} (not at \code{soc_top}): one
  \code{mbist_ctrl} (\texttt{N\_\allowbreak{}MEM\ =\allowbreak{}\ 10},
  \texttt{AW\ =\allowbreak{}\ 8},
  \texttt{P1\_\allowbreak{}MASK\ =\allowbreak{}\ 0}) in
  \code{rv32i_cpu_top}, ten \code{mem_collar}
  (\texttt{HAS\_\allowbreak{}P1\ =\allowbreak{}\ 0}) in
  \code{rv32i_icache} / \code{rv32i_dcache} between the cache logic and
  each macro, only in the \code{SRAM_SKY130} arm. Placing the controller
  at \code{soc_top} would have needed the ten macros\textquotesingle{}
  port-0 buses (about 10 x 44 pins) instead of 96 quasi-static pins.
  Memory index (status \code{fail} bit): 0 I\$ tag, 1-4 I\$ data, 5 D\$
  tag, 6-9 D\$ data.
\item
  \textbf{Reset-free controller} (\code{mbist_ctrl} parameter
  \texttt{ASYNC\_\allowbreak{}RST\ =\allowbreak{}\ 0}): the macro has
  only the synchronous \code{rst_n_i} and \code{check_scan_clk_rst.py}
  allows no async pin in it, so the controller has NO reset pin; its
  reset is folded into the run level
  (\texttt{run\ \&\ rst\_\allowbreak{}n\_\allowbreak{}i}) and
  \texttt{run\ =\allowbreak{}\ 0} flushes every flop within a few clocks
  (documented in \code{mbist_ctrl.sv}). A functional CPU reset therefore
  aborts a running test and can never start one. The SoC-fabric
  (core-domain) instance keeps
  \texttt{ASYNC\_\allowbreak{}RST\ =\allowbreak{}\ 1}, unchanged.
\item
  \textbf{Boundary: 96 pins}: \code{mbist_run_i} (1),
  \code{mbist_cfg_i[25:0]} (26), \code{mbist_status_o[68:0]} (69).
  \code{soc_mbist} drives the first two from the \code{cpu_core_clk}
  holding register and takes the third back into
  \code{mbist_status_sync}. Pin placement: \code{mbist_run_i} +
  \code{mbist_cfg_i} on the N edge, \code{mbist_status_o} on the S edge,
  in the middle of existing inter-pin gaps on the met2 track grid, so no
  existing pin moves (\code{pin_order.README.md} v5;
  \code{pin_template.def} PINS/NETS 407 -\textgreater{} 503, additions
  only).
\item
  \textbf{Clock force}: the Sky130 macro has no clock gate;
  \code{u_cpu_cg} in \code{soc_top} is forced open by the TAP-launched
  \code{mbist_force}. Found by
  \code{test_mbist_force_opens_the_closed_clock_gate}: the TAP drops the
  force at the Update-DR of the \texttt{run\ =\allowbreak{}\ 0} word, a
  few \code{cpu_core_clk} edges before that word reaches the macro, so
  with the PMU gate closed the controller would never see it and stay at
  \texttt{bist\_\allowbreak{}en\ =\allowbreak{}\ 1} behind a stopped
  clock. \code{cpu_cg_en} therefore also holds while the controller is
  busy (\texttt{mbist\_\allowbreak{}status\_\allowbreak{}o[0]\ |\ [1]},
  i.e. \texttt{st\_\allowbreak{}q\ !=\allowbreak{}\ IDLE}) and closes
  one gated edge after the controller has cleared. In the ASAP7 arm
  \code{test_en} is
  \texttt{scan\_\allowbreak{}mode\ |\ bist\_\allowbreak{}en} (the
  definition of 14.10); that arm has no collar, so the term is a no-op
  there today.
\item
  \textbf{Other arms}: default, FPGA and ASAP7 arms instantiate the same
  controller but the caches drive 0 as MBIST read data (no collar, no
  macro), so an MBIST run there reports FAIL on all ten bits and
  \texttt{done\ =\allowbreak{}\ 1}, never a vacuous pass, exactly like
  the core domain. \code{cpu_present} reads \code{MBIST_CPU}.
\item
  \textbf{\code{MBIST_CPU} default is 1} in \code{soc_top}. Consistency
  checked: Sky130 and ASAP7 SoC stubs carry the 96 ports,
  \src{pnr/Makefile} CPU source lists, \code{rtl/filelist/*.f}
  (\code{gen_filelist.py}), \code{phase5_soc_rtl_facts.json}
  (\texttt{gen-\allowbreak{}facts\ -\allowbreak{}-\allowbreak{}check}),
  UPF (\code{check_upf.py}: \code{u_cpu/mbist_status_o} classified
  \texttt{unclamped\ observability}). The SoC flow fails its macro pin
  match until the macro is re-hardened.
\end{itemize}

\textbf{Verification {[}M, 2026-10-18{]}}: Verilator -Wall adds no
warning on the default, Sky130, ASAP7 and FPGA arms of
\code{rv32i_cpu_top}; \code{dft-scan-check} (Sky130 SoC, ASAP7 SoC,
Sky130 CPU macro, ASAP7 CPU macro) 0 violations with the CPU macro still
at 0 async pins (the controller is reset-free), sequential words 944
-\textgreater{} 1033; \code{test_mbist_cpu_macro} 7/7 and 5/5 mutants
killed; \code{sim} rollup and the Sky130-arm icache / dcache /
cache\_\allowbreak{}integration suites green; the Sky130-arm RTL of the
macro is cycle-exact IDENTICAL to \code{origin/main} on the six gls
programs (collars in the loop, MBIST idle); \code{soc_all_ci} plain and
\code{FPGA=1} 831 passed / 0 failed / 41 skipped. Area (yosys synth +
abc, tt, SRAM macros black-boxed): 230,761 -\textgreater{} 256,402 um2
(+25,641, +11.1 \%; sequential +9,239), +3,402 cells, below the 33 k um2
estimate of 16.9 because the controller has no reset. Coarse grain: no
new \texttt{\$shiftx} / \texttt{\$shift}; +11 \texttt{\$pmux} and +1
\texttt{\$mem\_\allowbreak{}v2} (constant-label \code{case} functions
and the microcode ROM of \code{mbist_ctrl}, as in the core-domain
instance). NOT RUN: slang (not installed), timing, any PD step. See the
PR description for the counts. Suites: \code{mbist_cpu_macro} (the REAL
\code{rv32i_cpu_top} with \code{+define+SRAM_SKY130} behind the TAP,
fault-injecting 1 KB stand-ins), \code{mbist_fi_equiv_1k} (FI model
equals the real 1 KB simulation model), the extended \code{soc_jtag},
\code{dft-scan-check} x4.

\textbf{PD re-harden checklist (not run here)}

\begin{enumerate}
\def\labelenumi{\arabic{enumi}.}
\tightlist
\item
  CPU macro:
  \texttt{make\ -\allowbreak{}C\ pnr\ sky130-\allowbreak{}cpu-\allowbreak{}sv2v}
  then the Sky130 CPU flow with the committed \code{pin_order.cfg} +
  \code{pin_template.def} (strict template: it must report 503 pins);
  then
  \texttt{make\ -\allowbreak{}C\ pnr\ sky130-\allowbreak{}cpu-\allowbreak{}pin-\allowbreak{}lock-\allowbreak{}check\ NEW\_\allowbreak{}LEF=\allowbreak{}<run>/\allowbreak{}final/\allowbreak{}lef/\allowbreak{}rv32i\_\allowbreak{}cpu\_\allowbreak{}top.\allowbreak{}lef}
  must show no existing pin moved and the 96 new ones in place.
\item
  Commit the new \src{pnr/sky130/cpu/macro/} views (LEF, lib, GDS,
  SPICE, netlist) and regenerate the template only if a pin moved.
\item
  Timing: the collar adds one mux level on the address, data and enable
  pins of ten SRAM macros and \code{ss} reg-to-reg setup was already
  failing (r51p); measure, do not assume. Candidate relief if needed:
  register the \code{mem_collar} outputs inside \code{mbist_ctrl} is NOT
  possible (the functional side is the critical one); relax the period
  (item N7) or move the collar select earlier.
\item
  SoC: one re-harden with
  \texttt{MBIST\_\allowbreak{}CPU\ =\allowbreak{}\ 1};
  \code{phase5_soc_multiclock.sdc} / the Sky130 SoC SDC need the new
  CPU-domain MBIST crossings (item 12 of 16.5, same
  \texttt{set\_\allowbreak{}max\_\allowbreak{}delay\ -\allowbreak{}datapath\_\allowbreak{}only}
  gap as the core domain) and the \code{cpu_cg_en} feedback path (macro
  \code{mbist_status_o[1:0]} -\textgreater{} \code{u_cpu_cg} enable,
  same clock).
\item
  ASAP7 CPU macro views are not touched (shelved:
  ma7/\allowbreak{}lxv/\allowbreak{}2kn);
  \src{pnr/asap7/cpu/pin_order.cfg} is an explicit list that already
  lacked the scan pins and still lacks the MBIST pins: j41m.13 must add
  both when the ASAP7 macro is regenerated.
\end{enumerate}

\subsection{\texorpdfstring{CPU-macro collar-local compare (bead
\code{c54a},
2026-10-23)}{CPU-macro collar-local compare (bead , 2026-10-23)}}\label{tm:ch:dft-architecture:1612-cpu-macro-collar-local-compare-bead-c54a-2026-10-23}\origsec{16.12}

Two CPU-macro P\&R runs failed to route (global-route overflow 9,202 and
24,377 against 0 on the baseline; guides spread through 6 of the 10
cache macros, whose LEF blocks met1-met3). The diagnosis, inferred and
not measured per net, was that the macro-level \code{mbist_ctrl} is a
hub: it broadcast about 47 nets to the ten collars and took back 320
read-data wires (\code{dout0} of ten macros) across a 3.2 x 0.9 mm
array. User decision: compare in the collar.

\textbf{Change {[}R, as implemented{]}}

\begin{itemize}
\tightlist
\item
  \code{mem_collar} gets \code{LOCAL_CMP} (default 0, so the SoC
  \code{sram_controller} and NPU instances are unchanged). With 1, the
  collar rebuilds the expected word
  \texttt{bg\_\allowbreak{}pat(x\_\allowbreak{}bg)\ \textasciicircum{}\ expand4(x\_\allowbreak{}xm)}
  from an 8-net tag the controller broadcasts (\code{x_rd},
  \code{x_xm[3:0]}, \code{x_bg[2:0]}, all controller flop outputs) and
  compares its macro\textquotesingle s \code{dout0} (and \code{dout1}
  when \code{HAS_P1}) with it. The result is registered in the collar
  (no reset pin, same flush contract as the controller: \code{x_rd}
  falls within \code{RD_LAT} clocks of \texttt{run\ =\allowbreak{}\ 0})
  and returned as one strobe per read port (\code{cmp_fail0_o} /
  \code{cmp_fail1_o}).
\item
  \code{mbist_ctrl} gets \code{LOCAL_CMP}: it takes the strobes instead
  of \code{dout0_i} / \code{dout1_i}. The strobe is one clock later than
  the combinational compare, so the drain lasts
  \texttt{RD\_\allowbreak{}LAT\ +\allowbreak{}\ 1} cycles and the read
  tags have one more stage (address and slot), which keeps
  \texttt{rec\_\allowbreak{}addr\ /\allowbreak{}\ rec\_\allowbreak{}slot\ /\allowbreak{}\ rec\_\allowbreak{}elem\ /\allowbreak{}\ rec\_\allowbreak{}bg\ /\allowbreak{}\ rec\_\allowbreak{}port}
  and the sticky \code{fail} bits identical to the controller-side
  compare. \code{done} rises in the same cycle as the last \code{fail}
  bit.
\item
  \code{rv32i_cpu_top} instantiates the controller with
  \texttt{LOCAL\_\allowbreak{}CMP\ =\allowbreak{}\ 1} in the
  \code{SRAM_SKY130} arm only. The other arms have no collars and read
  constant 0: the controller compares that 0 itself
  (\texttt{LOCAL\_\allowbreak{}CMP\ =\allowbreak{}\ 0}, exactly as
  before), so an MBIST run there still FAILS and never passes vacuously.
\item
  Ports changed: \code{rv32i_core} and both caches lose
  \code{mbist_*_dout_o[159:0]} and gain \code{mbist_x_*_i} and
  \code{mbist_*_fail_o[4:0]}. The macro boundary (\code{rv32i_cpu_top}
  ports, 503 pins, \code{pin_template.def}) is unchanged.
\end{itemize}

\textbf{Timing choice (inferred, not measured)}: the SRAM \code{dout0}
nets are already half-cycle paths into the caches. The collar taps them
with one XOR2 per bit and registers only the 1-bit result, so
\code{dout0} carries one extra gate pin per bit, the same as the
controller-side XOR it replaces, but the wire to that pin is now local
instead of a run to the hub, and the half-cycle path ends at a flop
beside the macro. Registering \code{dout0} itself would add 32 flops per
collar (320) and a flop load on the same net, for no routing benefit.
The strobe then travels flop to flop in a full cycle. Needs an STA look
after the re-harden; none was run here.

\textbf{Diagnostic lost, stated exactly}: the XOR mask of the first
failing word (\code{status_o[68:37]}, the \code{err_mask} field of 16.5)
is no longer provided for the CPU macro: the failing data never reaches
the controller, so the 32 bits read 0 whatever \code{mask_sel} selects.
Everything else the status block carries is kept and is checked against
the golden model: per-memory sticky fail bits, \code{fail_any},
\code{done}, \code{busy}, \code{cfg_err}, and the first-fail record
(word address, element, slot, background, port). Consequence: a failing
macro, word, element and background are known but the failing bit
position inside the word is not, and the write-mask lane test still
detects a lane fault but cannot name the lane from the mask (the element
and slot still distinguish the forward and reverse lane passes). That is
acceptable for a go/no-go screen (this stage has no repair or redundancy
scheme that would consume a bit address), but not for bit-level failure
analysis; the core-domain controller (SoC macros) keeps
\texttt{LOCAL\_\allowbreak{}CMP\ =\allowbreak{}\ 0} and its mask. A
cheap way back, if wanted: a 32-flop first-fail capture per collar read
out serially (about 320 flops, +6 percent of the macro flop count, plus
a shift state in the controller). Not done. The 32
\code{mbist_status_o[68:37]} pins are now constant 0; they are NOT
removed (that changes the boundary and needs the user).

\textbf{Cost {[}M, yosys 0.62 generic gates on the sv2v netlist, base =
HEAD of this branch before the change{]}}: flops 5,285 -\textgreater{}
5,270 (-15: collars +10, extra tag stage +11, mask register and
\code{mask_sel} -36); the gate-level census band 5,220-5,370 of
\code{check_cpu_netlist_census.py} still holds. Nets between the
controller and the collars: 367 -\textgreater{} 65 (47 op-bus + 8 tag
broadcast, 10 returns).

\textbf{Update 2026-10-24: full-cycle capture (bead \bead{c54a})}: the
first collar compare ended in a rising-edge flop, a half-cycle path
because the Sky130 macro launches \code{dout} on the falling edge. In
the macro re-harden it failed setup at the ss corners only (icache -1.84
ns, dcache -0.94 ns). The strobe flop is now on the falling edge, so it
captures one full cycle after the launch; the controller launches the
expected-data tag from falling-edge copies on the same edge; a
rising-edge retime flop in the collar then drives the route to the
controller (a full-cycle flop-to-flop path). The strobe is two clocks
after the sample cycle: the drain is
\texttt{RD\_\allowbreak{}LAT\ +\allowbreak{}\ 2} cycles and the read
tags have two more stages. \code{dout} still carries one XOR2 input per
bit and no flop. Flops 5,270 -\textgreater{} 5,299 {[}M, yosys generic,
inside the census band{]}; 18 falling-edge flops, which need lock-up
handling at Stage 2 scan insertion. The timing argument is from the edge
relationships; it is not measured until the next macro re-harden. The
\texttt{LOCAL\_\allowbreak{}CMP\ =\allowbreak{}\ 0} logic is unchanged.
Suites on this RTL: \code{mbist_cpu_macro} 8/8, \code{mbist_ctrl} 5
configs, \code{mbist_tap} 24/24, \code{mbist_soc_slice} 3/3, both
\code{mbist_fi_equiv} 1/1.

\subsection{Parallel evaluation of Siemens Tessent MemoryBIST
(pointer)}\label{tm:ch:dft-architecture:1613-parallel-evaluation-of-siemens-tessent-memorybist-pointer}\origsec{16.13}

A Tessent 2025.4 MemoryBIST flow for the same two SoC-level SRAM macros
(TCD memory description, generated and inserted controller, VCS
simulation of Tessent\textquotesingle s own patterns, a fault campaign
over the 105-fault catalogue of \code{mbist_golden.py}, and area against
16.9 measured the same way) lives in \src{dft/tessent/mbist/}. It does
not change or replace anything in this section. Results and a
recommendation (keep the in-house MBIST; Tessent is a possible licensed
cross-check): \src{docs/design/TESSENT_MBIST_EVALUATION.md} (\cref{tm:dec:tessent-mbist-evaluation}).

\section{Tessent scan + ATPG evaluation (pointer)}\label{tm:ch:dft-architecture:17-tessent-scan-atpg-evaluation-pointer}\origsec{17}

A Siemens Tessent 2025.4 scan-insertion and ATPG evaluation of the
Sky130 SoC fabric (6 chains of \textasciitilde3,044, stuck-at 94.17 \%
test coverage / 1,574 patterns, transition 81.28 \% / 2,567 patterns, 20
patterns simulated clean on VCS with a negative control) is in
\src{docs/design/TESSENT_ATPG_FLOW.md} (\cref{tm:flow:tessent-atpg}), code in \src{dft/tessent/atpg/}.
It follows sections 8, 14, 15 and decision N4 unchanged, is not
back-annotated into PD, and does not modify the RTL.

\section{\texorpdfstring{Stage 3 ATPG tool spike record (bead
\code{j41m.15},
2026-10-28)}{Stage 3 ATPG tool spike record (bead , 2026-10-28)}}\label{tm:ch:dft-architecture:18-stage-3-atpg-tool-spike-record-bead-j41m15-2026-10-28}\origsec{18}

Labels: \textbf{{[}M{]}} measured (file or command named),
\textbf{{[}R{]}} read from source, \textbf{{[}A{]}} assumed or inferred.
Code: \src{tools/dft/atpg/} (README there); heavy outputs under
\code{/nobackup/claude_sim_build/j41m15/}, none in git.
\textbf{Evaluation only: no RTL and no PD config was touched.} Box (user
decision 2026-10-28, N12): about 5 working days, \textbf{12 GB RAM and
24 h wall per run}, stop at the first rung that breaks a limit. Each
Quaigh run is under \src{tools/dft/atpg/run_capped.sh}
(\texttt{systemd-\allowbreak{}run\ -\allowbreak{}-\allowbreak{}user\ -\allowbreak{}-\allowbreak{}scope\ -\allowbreak{}p\ MemoryMax=\allowbreak{}12G\ -\allowbreak{}p\ MemorySwapMax=\allowbreak{}0},
\texttt{timeout\ 24h},
\texttt{/\allowbreak{}usr/\allowbreak{}bin/\allowbreak{}time\ -\allowbreak{}v}),
which writes \code{metrics.json} (rc, wall, max RSS, cgroup peak,
\code{oom_kill}).

\textbf{Result in one paragraph.} Quaigh 0.0.6 completes the three
module rungs (\code{npu}, \code{crypto}, \code{dma}) and \textbf{breaks
the 12 GB box on the fabric rung}: the real Stage 2 post-fill netlist
(\code{RUN_2026-10-28_05-14-30}, 412,472 collapsed faults) is OOM-killed
at 12 GiB after 679 s with 29,376 patterns stored (78 \% progress). No
fabric-level deterministic coverage number exists. A random-pattern-only
grade of the same cut gives 69.59 \% (a lower bound, not an ATPG
result). Per decision 8 and the strict list of 2.4, the 99.0 \% / 98.0
\% goal is therefore \textbf{not demonstrated at SoC scale}; a restated
target is proposed in 18.8 and \textbf{not adopted} (superseded: 99 / 98
adopted provisionally 2026-10-30, 18.8). \textbf{Update (bead
\bead{j41m.18}, 18.9): with a sparse-row patch, Quaigh completes the
fabric rung in the same box (5.12 GiB, 50 min): 99.16 \% fault / 99.95
\% test coverage on the collapsed, excluded universe.}

\subsection{\texorpdfstring{What Quaigh 0.0.6 does (Step 0) {[}R:
\code{quaigh-0.0.6/src/atpg.rs}, \code{sim/fault.rs}, \code{equiv.rs},
\code{cmd.rs}{]}}{What Quaigh 0.0.6 does (Step 0) {[}R: , , , {]}}}\label{tm:ch:dft-architecture:181-what-quaigh-006-does-step-0-r-quaigh-006srcatpgrs-simfaultrs-equivrs-cmdrs}\origsec{18.1}

\begin{itemize}
\tightlist
\item
  \textbf{"unobservable" is a proof.}
  \code{find_pattern_detecting_fault} builds the good-vs-faulty miter
  and calls \code{prove} (Kissat); \code{Unsat} returns \code{None} and
  the fault is counted \code{unobservable}; \code{Interrupted} panics,
  so a timeout can never be mistaken for a proof. Only these faults are
  subtracted in test coverage. The count is printed only in the tqdm
  progress line (\code{unobservable=N}), not in a summary line;
  \code{atpg_coverage_report.py} reads the last occurrence.
\item
  \textbf{Universe.} The collapsed list \code{Fault::all_unique}
  (equivalence collapsing of single-fanout and
  xor/\allowbreak{}buf/\allowbreak{}and-input faults) unless
  \code{--with-redundant-faults}. Not comparable one-to-one with
  Tessent\textquotesingle s representative faults.
\item
  \textbf{Unresolved faults exist.}
  \texttt{detected\ +\allowbreak{}\ unobservable\ <\ total} is possible:
  the Stage 0 \code{timer} known answer is 4336 detected, 351
  unobservable, 5 neither (4692). They stay undetected in both numbers.
\item
  \textbf{Sequential BLIF is accepted} (\code{.latch}); \code{atpg}
  exposes the flops itself (\code{expose_dff}: Q appended to the inputs,
  D to the outputs, in \code{.latch} order).
\item
  \textbf{\code{--num-random} / \code{--num-cycles} do not tune the
  combinational flow.} They switch to "random patterns only": no SAT, no
  grading, written to the file. The plan\textquotesingle s "retry the
  fabric rung with a larger \code{--num-random}" therefore does not
  exist as an option. Also \code{--num-random} sizes patterns from the
  primary inputs only and does not expose flops (a patterns-vs-network
  width assertion fails on a \code{.latch} BLIF), so the random grade in
  18.3 uses the explicit full-scan view.
\item
  \textbf{Memory is O(patterns x faults).} \code{TestPatternGenerator}
  keeps
  \texttt{pattern\_\allowbreak{}detections:\allowbreak{}\ Vec<Vec<bool>>}
  (one byte per pattern-fault pair) and the compression step adds
  \code{fault_to_patterns} / \code{pattern_to_faults} index vectors.
  Each hard fault adds 64 patterns (\code{add_random_patterns_from}).
  {[}A: the per-phase peak was not measured; the fabric kill point
  matches the byte matrix, see 18.3.{]}
\item
  \textbf{\texttt{quaigh\ equiv} on a sequential BLIF is a 1-cycle check
  from the reset state, not a full-scan proof}
  (\code{check_equivalence_bounded}). An early version of
  \code{regrade.py} used it for the "provably untestable" reading of
  undetected samples; that was wrong and was corrected (the
  combinational view is used, and the proof is repeated with
  Yosys\textquotesingle{} SAT solver, 18.5).
\end{itemize}

\subsection{Method {[}R/M{]}}\label{tm:ch:dft-architecture:182-method-rm}\origsec{18.2}

{\footnotesize\setlength{\tabcolsep}{3pt}

{
\begin{longtable}[]{@{}
  >{\raggedright\arraybackslash}p{\dimexpr(\linewidth-2\tabcolsep)*305/1000\relax}
  >{\raggedright\arraybackslash}p{\dimexpr(\linewidth-2\tabcolsep)*675/1000\relax}@{}}
\toprule\noalign{}
\begin{minipage}[b]{\linewidth}\raggedright
Step
\end{minipage} & \begin{minipage}[b]{\linewidth}\raggedright
What
\end{minipage} \\
\midrule\noalign{}
\endhead
\bottomrule\noalign{}
\endlastfoot
Module rungs & \code{synth_module.sh}: \code{sv2v} with the Sky130 SoC
define set and file list
(\texttt{make\ -\allowbreak{}C\ pnr\ sky130-\allowbreak{}soc-\allowbreak{}sv2v-\allowbreak{}inputs}),
then Yosys 0.46 \texttt{synth\ -\allowbreak{}flatten},
\code{dfflibmap}/\texttt{abc\ -\allowbreak{}liberty} with the Sky130 tt
Liberty. Not the LibreLane recipe, so cell counts differ from the PD
netlist. The Sky130 4 KB SRAM stays a black box. \\
Fabric rung & the post-fill netlist
\code{.../RUN_2026-10-28_05-14-30/49-openroad-fillinsertion/soc_top.nl.v}
(6 fabric chains, 18,803 scan flops, 922 unscanned).
\code{netlist_prep.py} deletes
fill/\allowbreak{}decap/\allowbreak{}tap/\allowbreak{}diode cells
(1,265,632) and applies the exclusions of 18.4. \\
Cut (\code{mk_cut.py}) & \code{read_liberty} without \code{-lib} (every
Sky130 cell becomes logic or a flop), a wrapper that ties the scan-mode
capture constants, \code{flatten}, \code{async2sync} + \code{dffunmap}
(no async flop is left), latches -\textgreater{} flops, \textbf{no
\code{opt_dff}} (it deletes flops whose D folds to a constant; a scan
flop stays loadable),
\texttt{abc\ -\allowbreak{}g\ AND,\allowbreak{}NAND,\allowbreak{}OR,\allowbreak{}NOR,\allowbreak{}XOR,\allowbreak{}XNOR,\allowbreak{}ANDNOT,\allowbreak{}ORNOT,\allowbreak{}MUX},
\code{write_blif} with \code{.latch} for every flop, then a dead-logic
drop. \texttt{opt\ -\allowbreak{}full} was found to delete every flop
and is not used. \\
Capture constants (fabric) &
\texttt{scan\_\allowbreak{}mode\_\allowbreak{}i=\allowbreak{}1,\allowbreak{}\ scan\_\allowbreak{}en\_\allowbreak{}i=\allowbreak{}0,\allowbreak{}\ trst\_\allowbreak{}ni=\allowbreak{}0,\allowbreak{}\ tck\_\allowbreak{}i=\allowbreak{}0,\allowbreak{}\ scan\_\allowbreak{}rst\_\allowbreak{}ni=\allowbreak{}1};
\code{scan_clk_i} is the only clock. Module rungs: \code{rst_n=1} (and
\code{scan_mode_i=1} for crypto). \\
Black boxes & CPU macro \code{rv32i_cpu_top} and the SRAM macros:
outputs become free pseudo-inputs, inputs become observed outputs
(\texttt{expose\ -\allowbreak{}evert}). Quaigh has no X state, so
coverage of the logic around them is an \textbf{upper bound}. A
black-box clock pin has no data driver and is dropped from the outputs
(1 on the fabric: \code{u_cpu.clk_i}). \\
Clock logic & clock-network logic that only reaches flop clock pins is
dropped (Quaigh ignores \code{.latch} clocks, and left in, its faults
would inflate "proven untestable"); the one clock-gate cell without a
Liberty function (\code{sdlclkp_4}) is dropped and counted. \\
\end{longtable}
}

}

Quaigh counts after its own \code{expose_dff} {[}M, logs{]}: fabric
19,409 inputs, 19,238 outputs, 262,016 gates, 1,346,520 faults, 412,472
collapsed; flops in the cut: 18,804.

\subsection{\texorpdfstring{Ladder {[}M: each \code{metrics.json}
and \code{q/report.json} under
\code{/nobackup/claude_sim_build/j41m15/<rung>/}{]}}{Ladder {[}M: each  and  under {]}}}\label{tm:ch:dft-architecture:183-ladder-m-each-metricsjson-and-qreportjson-under-nobackupclaude-sim-buildj41m15rung}\origsec{18.3}

\begin{landscape}
\footnotesize\setlength{\tabcolsep}{3pt}

{
\begin{longtable}[]{@{}
  >{\raggedright\arraybackslash}p{\dimexpr(\linewidth-22\tabcolsep)*230/1000\relax}
  >{\raggedright\arraybackslash}p{\dimexpr(\linewidth-22\tabcolsep)*52/1000\relax}
  >{\raggedright\arraybackslash}p{\dimexpr(\linewidth-22\tabcolsep)*67/1000\relax}
  >{\raggedright\arraybackslash}p{\dimexpr(\linewidth-22\tabcolsep)*48/1000\relax}
  >{\raggedright\arraybackslash}p{\dimexpr(\linewidth-22\tabcolsep)*93/1000\relax}
  >{\raggedright\arraybackslash}p{\dimexpr(\linewidth-22\tabcolsep)*48/1000\relax}
  >{\raggedright\arraybackslash}p{\dimexpr(\linewidth-22\tabcolsep)*59/1000\relax}
  >{\raggedright\arraybackslash}p{\dimexpr(\linewidth-22\tabcolsep)*56/1000\relax}
  >{\raggedright\arraybackslash}p{\dimexpr(\linewidth-22\tabcolsep)*56/1000\relax}
  >{\raggedright\arraybackslash}p{\dimexpr(\linewidth-22\tabcolsep)*70/1000\relax}
  >{\raggedright\arraybackslash}p{\dimexpr(\linewidth-22\tabcolsep)*96/1000\relax}
  >{\raggedright\arraybackslash}p{\dimexpr(\linewidth-22\tabcolsep)*104/1000\relax}@{}}
\toprule\noalign{}
\begin{minipage}[b]{\linewidth}\raggedright
Rung
\end{minipage} & \begin{minipage}[b]{\linewidth}\raggedright
Flops in cut
\end{minipage} & \begin{minipage}[b]{\linewidth}\raggedright
Collapsed faults
\end{minipage} & \begin{minipage}[b]{\linewidth}\raggedright
Detected
\end{minipage} & \begin{minipage}[b]{\linewidth}\raggedright
Proven untestable (SAT)
\end{minipage} & \begin{minipage}[b]{\linewidth}\raggedright
Unresolved
\end{minipage} & \begin{minipage}[b]{\linewidth}\raggedright
Fault coverage
\end{minipage} & \begin{minipage}[b]{\linewidth}\raggedright
Test coverage
\end{minipage} & \begin{minipage}[b]{\linewidth}\raggedright
Patterns kept
\end{minipage} & \begin{minipage}[b]{\linewidth}\raggedright
Wall
\end{minipage} & \begin{minipage}[b]{\linewidth}\raggedright
Peak RAM (cgroup)
\end{minipage} & \begin{minipage}[b]{\linewidth}\raggedright
Completed
\end{minipage} \\
\midrule\noalign{}
\endhead
\bottomrule\noalign{}
\endlastfoot
\code{npu} (\code{npu_top}, SRAM black-boxed) & 1,169 & 81,126 & 80,905
& 220 & 1 & 99.728 \% & 99.9988 \% & 1,199 & 56 s & 1.21 GiB & yes \\
\code{crypto} (\code{crypto_accel}, 256 key flops excluded) & 2,755 &
93,178 & 91,235 & 1,943 & 0 & 97.915 \% & 100.000 \% & 1,578 & 149 s &
1.68 GiB & yes \\
\code{dma} (\code{dma_engine}) & 9,088 & 137,172 & 136,969 & 203 & 0 &
99.852 \% & 100.000 \% & 13,432 & 845 s & \textbf{11.42 GiB} (RSS 11.89
GB) & yes, 0.6 GiB under the cap \\
\code{fabric} (Stage 2 post-fill netlist, CPU macro + SRAMs black-boxed)
& 18,804 & 412,472 & n/a & 428 proven at the kill & n/a & n/a & n/a &
n/a & 679 s to the kill & \textbf{12.0 GiB, \code{oom_kill=1}} &
\textbf{NO: broke the box} \\
\end{longtable}
}

\end{landscape}

Fabric detail {[}M \code{fabric/q/run.log}, \code{metrics.json}{]}:
after the random phase 274,941 of 412,472 faults (66.66 \%) were
detected; the deterministic phase had reached 78 \% of its progress bar
with 29,376 patterns stored when the kernel killed the run. 29,376 x
412,472 B = 12.1 GB, which is the cap: the byte matrix of 18.1 alone
fills it, before the compression step that follows. {[}A: completion
would need more than 12 GiB; the pattern count at completion is unknown.
The 24 h limit was not the binding one.{]} The cap, not the wall time,
is the failure.

Fabric random-pattern-only grade {[}M \code{fabric/rand/grade_4096/}{]}:
4,096 random patterns (\texttt{quaigh\ atpg\ -\allowbreak{}r}, then
\code{atpg-report} on the full-scan view, 377,138 faults in that view
because \code{expose_dff} adds buffers): \textbf{262,454 detected =
69.59 \%}, 169 s, 1.9 GB. This is a floor for a pattern set (no pattern
targets a hard fault, nothing is proven), shown only to bound what the
random phase alone buys; the universe differs from the 412,472 above, so
the two percentages are not interchangeable.

\code{dma} passing by 0.6 GiB means the module-rung method is also at
its edge: a block of \textasciitilde9 k flops is about the largest
Quaigh can finish under this cap.

\subsection{\texorpdfstring{Exclusion manifest {[}M
\code{fabric/exclusions.json},
\code{hold_mask_1028.json}{]}}{Exclusion manifest {[}M , {]}}}\label{tm:ch:dft-architecture:184-exclusion-manifest-m-fabricexclusionsjson-hold-mask-1028json}\origsec{18.4}

State with every coverage number: the fault universe \textbf{excludes}
the logic and flops below.

{\footnotesize\setlength{\tabcolsep}{3pt}

{
\begin{longtable}[]{@{}
  >{\raggedright\arraybackslash}p{\dimexpr(\linewidth-4\tabcolsep)*327/1000\relax}
  >{\raggedright\arraybackslash}p{\dimexpr(\linewidth-4\tabcolsep)*327/1000\relax}
  >{\raggedright\arraybackslash}p{\dimexpr(\linewidth-4\tabcolsep)*327/1000\relax}@{}}
\toprule\noalign{}
\begin{minipage}[b]{\linewidth}\raggedright
Exclusion
\end{minipage} & \begin{minipage}[b]{\linewidth}\raggedright
Count
\end{minipage} & \begin{minipage}[b]{\linewidth}\raggedright
Treatment in the cut
\end{minipage} \\
\midrule\noalign{}
\endhead
\bottomrule\noalign{}
\endlastfoot
\code{tck_i}-clocked flops (TAP, MBIST TDRs, incl.
negedge-behind-inverter) & \textbf{664} & not scanned, tested through
JTAG (user decision 2026-10-26). Reset-type flops (554 of the 666 TAP +
clock-switch flops) are removed and Q is tied to the reset value
(\code{trst_ni=0}). \\
\code{scan_clk_i}-clocked flops (clock switch) & \textbf{2} & same \\
no-reset flops among the 666 (power-up unknown) & 112 & kept as flops
whose D is their own Q (never observed); Yosys drops them (constant
clock \code{tck_i=0}) so they enter the cut as constants, not as free
inputs. {[}M: the cut holds 18,804 flops = 18,803 scan flops + 1; A: the
constant value{]}. This differs from the plan\textquotesingle s "flagged
free inputs" wording. \\
crypto key flops (\code{u_crypto.key_q} 128, \code{u_aes.rk_q} 128;
decision 3) & \textbf{256} & removed, Q tied 0; scan-mode masks make
their cone constant. Their exclusive D fan-in is unobserved and falls
out of the universe. \\
scan-view hold-masked capture flops & \textbf{94} & Q stays a
pseudo-input (still shift-controllable), D tied 0 so the capture value
is not observed \\
black boxes & CPU macro (\code{u_cpu}), SRAM (\code{u_sram},
\code{u_npu...u_sram_macro}) & free outputs / observed inputs, upper
bound (no X) \\
physical cells (fill, decap, tap, antenna diode) & 1,265,632 & deleted
(no logic) \\
\end{longtable}
}

}

Unscanned-flop accounting {[}M{]}: 19,725 flops = 18,803 scan cells +
922 unscanned; the 922 = 664 + 2 + 256 + 0 unexplained
(\texttt{other\_\allowbreak{}unscanned\_\allowbreak{}flops:\allowbreak{}\ 0}).

\textbf{Hold mask} (\code{derive_hold_mask.py}). The 94 flops are
selected by \textbf{stable synchroniser Q-net names}, never by instance
id: 34
\texttt{u\_\allowbreak{}cpu\_\allowbreak{}axi\_\allowbreak{}cdc.\allowbreak{}u\_\allowbreak{}\{aw,\allowbreak{}w,\allowbreak{}ar,\allowbreak{}b,\allowbreak{}r\}\_\allowbreak{}fifo.\allowbreak{}u\_\allowbreak{}\{wr\_\allowbreak{}ptr\_\allowbreak{}to\_\allowbreak{}rd,\allowbreak{}rd\_\allowbreak{}ptr\_\allowbreak{}to\_\allowbreak{}wr\}.\allowbreak{}sync\_\allowbreak{}q[0][i]}
(cdc\_\allowbreak{}gray\_\allowbreak{}axi), 4
\texttt{u\_\allowbreak{}mbist.\allowbreak{}\{g\_\allowbreak{}cpu.\allowbreak{}u\_\allowbreak{}xing\_\allowbreak{}cpu|u\_\allowbreak{}xing\_\allowbreak{}core\}.\allowbreak{}u\_\allowbreak{}fifo.\allowbreak{}u\_\allowbreak{}wr\_\allowbreak{}ptr\_\allowbreak{}to\_\allowbreak{}rd.\allowbreak{}sync\_\allowbreak{}q[0][i]}
(mbist\_\allowbreak{}ptr\_\allowbreak{}sync), 4 \code{...ctl_w[i]}
(mbist\_\allowbreak{}ctl\_\allowbreak{}hold), 26
\code{u_mbist.core_cfg_w[i]}
(mbist\_\allowbreak{}core\_\allowbreak{}cfg), 26 \code{mb_cpu_cfg[i]}
(mbist\_\allowbreak{}cpu\_\allowbreak{}cfg). Validated \textbf{94/94}
against
\code{/nobackup/claude_sim_build/j41m8_sta/new_nom_tt_exempt.d/exempt_pins.txt}
on \code{RUN_2026-10-26_09-59-03} (the STA\textquotesingle s own list: 0
missing, 0 extra) and re-derived on \code{RUN_2026-10-28_05-14-30} (94,
same split). Negative control: dropping the \code{mbist_cpu_cfg} rule
gives 26 missing and exit 1. Pytest:
\src{tb/tests/test_dft_atpg_hold_mask.py}.

\subsection{Sanity checks}\label{tm:ch:dft-architecture:185-sanity-checks}\origsec{18.5}

{\footnotesize\setlength{\tabcolsep}{3pt}

{
\begin{longtable}[]{@{}
  >{\raggedright\arraybackslash}p{\dimexpr(\linewidth-2\tabcolsep)*490/1000\relax}
  >{\raggedright\arraybackslash}p{\dimexpr(\linewidth-2\tabcolsep)*490/1000\relax}@{}}
\toprule\noalign{}
\begin{minipage}[b]{\linewidth}\raggedright
Check
\end{minipage} & \begin{minipage}[b]{\linewidth}\raggedright
Result
\end{minipage} \\
\midrule\noalign{}
\endhead
\bottomrule\noalign{}
\endlastfoot
Stage 0 \code{timer} known answer {[}M \code{timer_s0}{]} &
\texttt{quaigh\ atpg} on the committed Stage 0 BLIF: \textbf{4336/4692 =
92.41 \%}, \code{unobservable=351}, 358 patterns, as in 2.2. Reproduced.
The same \code{timer} through the new Sky130-cell pipeline gives a
different universe (2,840 faults, 97.50 \% fault coverage, 99.86 \% test
coverage): different netlist (mapped, pruned), not a contradiction. \\
Independent re-grade, method & \code{regrade_vl.py}: the sampled faults
are injected as select-controlled muxes into a copy of the cut,
\textbf{Yosys} \code{read_blif}/\code{write_verilog} defines the logic,
\textbf{Verilator} simulates every Quaigh pattern with each fault
selected, a fault is detected when an output differs. For undetected
faults \textbf{Yosys\textquotesingle{} SAT}
(\texttt{miter\ -\allowbreak{}equiv} +
\texttt{sat\ -\allowbreak{}prove}, a different solver from
Quaigh\textquotesingle s Kissat) decides untestable vs coverage hole.
Negative controls: an all-zero pattern must detect clearly fewer of the
same faults; an injected wrong verdict must fail the comparison
(pytest). \code{regrade.py} additionally cross-checks with its own
Python simulator and \texttt{quaigh\ simulate}. \\
\code{npu} (250 faults, 241 modelled) & 199 detected + 42 undetected;
all 42 undetected proven untestable by Yosys SAT; \textbf{0
disagreements} with the Python simulator; all-zero control 22 vs 199
detected \\
\code{crypto} (250) & 200 detected + 50 undetected; all 50 proven
untestable by Yosys SAT; \textbf{0 disagreements}; control 40 vs 200 \\
\code{dma} (250, 249 modelled) & 200 detected + 49 undetected;
\textbf{45 proven untestable by Yosys SAT, 4 testable but missed by the
13,432 patterns} (all four are sa0 on single-fanout AND-gate outputs
feeding a NOR/\allowbreak{}ANDNOT); \textbf{0 disagreements} with the
Python simulator on detection; control 37 vs 200. See the full-universe
grade below. \\
\code{dma} full uncollapsed universe {[}M \code{dma/full/}{]} &
\texttt{quaigh\ atpg-\allowbreak{}report\ -\allowbreak{}-\allowbreak{}with-\allowbreak{}redundant-\allowbreak{}faults}
on the same 13,432 patterns: \textbf{451,725 / 452,588 = 99.81 \%} (863
undetected), 165 s, 6.4 GB. The collapsed list gave 136,969 / 137,172 =
99.85 \% with 203 proven. \\
\code{timer} (80, 78 modelled) & 59 + 19, 19/19 proven untestable, 0
disagreements \\
Quaigh \code{simulate} cross-check & \code{regrade.py} (npu) 200
detected + 50 undetected agree with \texttt{quaigh\ simulate} on
fault-injected copies (16 min). Its "provably untestable" column
predates the full-scan correction of 18.1 and is superseded by the Yosys
SAT column above. \\
Reading & A sample of 250 faults per rung is a spot check of the
pattern-to-output simulation and of the untestable proofs, not a recount
of Quaigh\textquotesingle s fault list (Quaigh reports totals only). The
unresolved fault and the \code{timer} 5 are not covered. The fabric has
no pattern set to re-grade. \\
\textbf{Finding: "test coverage 100 \%" is a statement about
Quaigh\textquotesingle s collapsed list} & On \code{dma} 4 of 49 sampled
undetected net faults are testable (Yosys SAT) yet missed. Whether those
four are members of Quaigh\textquotesingle s collapsed list was not
established {[}A: the collapse rules in \code{fault.rs} treat input pins
after signal inversion, and the mapping from a BLIF net fault to a
Quaigh pin fault was not reconstructed{]}. Magnitude {[}A{]}: 863
uncollapsed faults are undetected (0.19 \% of 452,588); at 4/49 of them
testable that is roughly 70 faults, so test coverage on the uncollapsed
list is about 99.98 \%, not 100.00 \%, with large sampling error.
Reported numbers in 18.3 are Quaigh\textquotesingle s own and keep the
label "collapsed list". \code{npu} and \code{crypto} showed 0 holes in
42 and 50 sampled undetected faults. \\
\end{longtable}
}

}

\subsection{\texorpdfstring{Delta against Tessent (94.17 \% test /
93.25 \% fault, \code{TESSENT_ATPG_FLOW.md} (\cref{tm:flow:tessent-atpg}) 4.2, 4.3)
{[}A{]}}{Delta against Tessent (94.17 \% test / 93.25 \% fault,  4.2, 4.3) {[}A{]}}}\label{tm:ch:dft-architecture:186-delta-against-tessent-9417-test-9325-fault-tessent-atpg-flowmd-42-43-a}\origsec{18.6}

No like-for-like number exists (the fabric rung failed), so the
comparison is by class, with nothing reclassified to close a gap:

{\footnotesize\setlength{\tabcolsep}{3pt}

{
\begin{longtable}[]{@{}
  >{\raggedright\arraybackslash}p{\dimexpr(\linewidth-6\tabcolsep)*221/1000\relax}
  >{\raggedright\arraybackslash}p{\dimexpr(\linewidth-6\tabcolsep)*259/1000\relax}
  >{\raggedright\arraybackslash}p{\dimexpr(\linewidth-6\tabcolsep)*259/1000\relax}
  >{\raggedright\arraybackslash}p{\dimexpr(\linewidth-6\tabcolsep)*242/1000\relax}@{}}
\toprule\noalign{}
\begin{minipage}[b]{\linewidth}\raggedright
Factor
\end{minipage} & \begin{minipage}[b]{\linewidth}\raggedright
Tessent run
\end{minipage} & \begin{minipage}[b]{\linewidth}\raggedright
This spike
\end{minipage} & \begin{minipage}[b]{\linewidth}\raggedright
Direction of the effect
\end{minipage} \\
\midrule\noalign{}
\endhead
\bottomrule\noalign{}
\endlastfoot
Netlist & older pre-Stage-2 yosys-synthesised netlist (18,263 scan
cells, \code{edfxtp} flops) & the real Stage 2 post-fill netlist & not
comparable cell for cell \\
Fault model & Sky130 cell pins, 661,690 representative faults & generic
gates after \texttt{abc\ -\allowbreak{}g}, 412,472 collapsed faults &
different universe; neither is the other\textquotesingle s subset \\
X handling & X-aware; black boxes and unscanned flops are X sources
(AU.BB 1.05 \%, 6,915 faults) & no X: black-box outputs are free inputs
& spike is \textbf{optimistic} around the CPU macro, SRAMs and NPU \\
TAP/MBIST flops (666) and cones & in the denominator, undetected (AU.TC
1.27 \%, part of AU.PC) & excluded from the universe (user decision
2026-10-26) & spike \textbf{higher} by that share, which Tessent counts
as loss \\
Crypto key mask (AU.PC 2,598 representative faults) & undetected, in the
denominator & key flops and exclusive cone removed. {[}A: the 1,943
SAT-proven untestable faults of the \code{crypto} rung (2.09 \% of
93,178) are the masked legs and constant cones; their cause was not
itemised per gate.{]} They are subtracted in test coverage and shown in
fault coverage & test coverage hides it, fault coverage shows it \\
Redundant / unclassified & RE 0.98 \%, AU 1.11 \%, PU/SEQ/UO 0.19 \% &
"proven untestable" (SAT) + "unresolved" & equivalent buckets \\
Scan path faults & {[}A: whether Tessent\textquotesingle s 661,690
include SI/SE/SCD pin faults was not checked{]} & scan muxes folded
under \code{scan_en_i=0}; SI/SE faults not in the universe & spike
\textbf{higher} if Tessent counts them \\
Completion & 94.17 \% in 4 min 55 s, 1.18 GB & no completed fabric
result; 12 GiB exhausted at 679 s & Tessent scales;
Quaigh\textquotesingle s data structure does not \\
\end{longtable}
}

}

What follows {[}A{]}: nothing in the module rungs contradicts a high
ceiling (97.9 to 99.9 \% fault coverage on blocks, 100 \% test coverage
after SAT-proven untestable faults), and the Tessent shortfall is
attributed to black boxes (not demonstrated here, since this spike
cannot see X), TAP/MBIST tied cells (excluded here), and the crypto
mask. The spike does \textbf{not} show that an open tool can reach 99.0
\% test coverage on the SoC: that statement has no measurement behind
it.

\subsection{Limits and what was not done {[}A unless noted{]}}\label{tm:ch:dft-architecture:187-limits-and-what-was-not-done-a-unless-noted}\origsec{18.7}

\begin{itemize}
\tightlist
\item
  Generic-gate faults, not Sky130 cell-pin faults; the final number must
  be re-simulated on cell pins (follow-up bead).
\item
  Stuck-at only. Delay faults (N10) are untouched; Quaigh has no
  transition model.
\item
  Patterns are not mapped to chain order, no SVF/TAP replay, no
  gate-level simulation of the patterns on the scan netlist (follow-up
  bead). Quaigh\textquotesingle s scan-path faults (shift, flush, SI/SE)
  are not graded.
\item
  The SRAM macros are not shadowed (6.3): the largest recoverable class
  in the Tessent data.
\item
  The CPU macro (5,299 scan flops) is a black box; its own scan netlist
  as an extra rung was not run.
\item
  Fault (AUCOHL) was not built: the user allowed it only as a fallback
  when Quaigh fails a rung; it did fail the fabric rung, but
  Fault\textquotesingle s own engine is random patterns plus an
  \code{iverilog} fault simulation (2.2) and Atalanta/\allowbreak{}PODEM
  are non-commercial, so a build was judged not worth the time-box
  without the owner\textquotesingle s decision (listed in 18.8).
\item
  \code{quaigh} was not patched. The O(patterns x faults) byte matrix is
  a data-structure limit: bit-packing it (8x) or splitting the fault
  list is a change to the tool, not a parameter. \textbf{Superseded by
  18.9 (bead \bead{j41m.18}): a sparse-row patch fixed it and the fabric
  rung completes.}
\item
  Whole-SoC fault coverage (the 98 \% of 2.1) needs the fault count
  inside the excluded cones. A/B estimate {[}M count, A reading{]}: the
  same fabric cut built with
  \texttt{netlist\_\allowbreak{}prep.\allowbreak{}py\ -\allowbreak{}-\allowbreak{}keep-\allowbreak{}unscanned}
  (the 922 unscanned flops kept as ordinary full-scan flops,
  \code{fabric_noexcl/}) has \textbf{416,278} collapsed faults against
  \textbf{412,472} with the exclusions applied, i.e. the exclusions
  remove about \textbf{3,806 faults (0.91 \% of the larger list)}. If
  every one of them counted as undetected, the whole-SoC fault coverage
  could not exceed 412,472 / 416,278 = \textbf{99.09 \%} even with
  perfect detection elsewhere; the 98 \% goal leaves about 1.1 points
  for everything else. The two cuts are not a subset pair (Quaigh gate
  counts 262,016 vs 261,111), so this is an estimate of the order of
  magnitude, in line with the 0.7 \% + 0.2-0.4 \% budget of 2.1.
\end{itemize}

\subsection{Proposed restated target (NOT adopted: the owner
decides)}\label{tm:ch:dft-architecture:188-proposed-restated-target-not-adopted-the-owner-decides}\origsec{18.8}

\textbf{ADOPTED 2026-10-30 (owner, bead \bead{j41m.23}); supersedes
items 1-3 below, which stay as history.} Item 1 no longer holds: the
fabric now has a number (18.9). Item 2, the per-block fallback, is
dropped. Item 3 was resolved by route (b). Item 4 changes: \textbf{99.0
\% test / 98.0 \% fault is adopted as the provisional Stage 3 target},
applied to the fabric list, with the CPU macro as a separate rung at the
same target. The sign-off conditions are in 2.1.

\begin{enumerate}
\def\labelenumi{\arabic{enumi}.}
\tightlist
\item
  \textbf{Do not state one SoC stuck-at number.} The fabric rung failed
  in the 12 GB box; no deterministic fabric coverage was measured.
\item
  \textbf{Per-block targets (the fallback of 2.4)}, measured with this
  flow: stuck-at test coverage \textgreater= 99.0 \% and fault coverage
  \textgreater= 97.0 \% per block on blocks of \textless=
  \textasciitilde9 k flops (npu 100.00 / 99.73, crypto 100.00 / 97.92,
  dma 100.00 / 99.85 are the evidence), each with the exclusion
  statement of 18.4 and the SAT-proven untestable count itemised.
\item
  \textbf{For the SoC number choose one of}: (a) a commercial or patched
  engine run (Tessent has a measured 94.17 \% on an older netlist and
  finishes in minutes), (b) Quaigh patched to a bit-packed detection
  matrix, then re-run the fabric rung, (c) partition the fabric by
  hierarchy into cuts that fit the box and report per block. Options (b)
  and (c) keep the open-tool route; (a) is the only one with a measured
  SoC-scale result.
\item
  Keep 99.0 / 98.0 as a goal (decision 8), not a commitment.
\end{enumerate}

\textbf{Owner decisions 2026-10-29 (bead \bead{j41m.18}):} route
\textbf{(b)} chosen and executed (result in 18.9); \textbf{Fault
(AUCOHL) is dropped} as a fallback (random patterns plus an
\code{iverilog} fault simulation add nothing over a completing Quaigh);
the restated target above is still \textbf{not adopted}: the owner
decides it with the 18.9 number in hand (follow-up bead).
\textbf{Decided 2026-10-30 (bead \bead{j41m.23}): adopted provisionally,
see the note at the top of this section and 2.1.}

\subsection{\texorpdfstring{Patched Quaigh: fabric result (bead
\code{j41m.18},
2026-10-29)}{Patched Quaigh: fabric result (bead , 2026-10-29)}}\label{tm:ch:dft-architecture:189-patched-quaigh-fabric-result-bead-j41m18-2026-10-29}\origsec{18.9}

\textbf{Patch {[}R/M{]}.}
\src{tools/dft/atpg/quaigh/0001-sparse-pattern-detections.patch} (Quaigh
\code{0.0.6+cvt1}, built by
\texttt{tools/\allowbreak{}dft/\allowbreak{}build\_\allowbreak{}quaigh.\allowbreak{}sh\ patched}
into \code{/nobackup/claude_sim_build/j41m18/qroot}; the stock binary is
kept). \code{pattern_detections} becomes, per pattern, the ascending
list of detected fault indices (\code{u32}) instead of one byte per
(pattern, fault). Reason, read from \code{atpg.rs}: only the first 256
random patterns are graded against every fault; each later 64-pattern
batch is graded against the still-undetected faults only, so its dense
rows were almost all zeros. \code{compress_patterns} builds
\code{fault_to_patterns} from the sparse rows in the same ascending
order and keeps the \code{max_by} tie-break, so pattern selection is
unchanged. \code{patterns} (one byte per input) is left dense. No other
file changed.

\textbf{Equivalence gate {[}M: \texttt{cmp\ -\allowbreak{}s} of the
pattern files; each \code{metrics.json} under
\texttt{/\allowbreak{}nobackup/\allowbreak{}claude\_\allowbreak{}sim\_\allowbreak{}build/\allowbreak{}\{j41m15,\allowbreak{}j41m18\}/\allowbreak{}}{]}.}
Same cuts (symlinked from j41m.15), same seed; only the binary differs.

{
\begin{longtable}[]{@{}
  >{\raggedright\arraybackslash}p{\dimexpr(\linewidth-6\tabcolsep)*275/1000\relax}
  >{\raggedright\arraybackslash}p{\dimexpr(\linewidth-6\tabcolsep)*198/1000\relax}
  >{\raggedright\arraybackslash}p{\dimexpr(\linewidth-6\tabcolsep)*310/1000\relax}
  >{\raggedright\arraybackslash}p{\dimexpr(\linewidth-6\tabcolsep)*198/1000\relax}@{}}
\toprule\noalign{}
\begin{minipage}[b]{\linewidth}\raggedright
Rung
\end{minipage} & \begin{minipage}[b]{\linewidth}\raggedright
Pattern file vs stock
\end{minipage} & \begin{minipage}[b]{\linewidth}\raggedright
Peak RAM stock -\textgreater{} patched (cgroup)
\end{minipage} & \begin{minipage}[b]{\linewidth}\raggedright
Wall stock -\textgreater{} patched
\end{minipage} \\
\midrule\noalign{}
\endhead
\bottomrule\noalign{}
\endlastfoot
\code{timer} (Stage 0 known answer) & identical & 42 MiB -\textgreater{}
9 MiB & 1.5 s -\textgreater{} 1.4 s \\
\code{npu} & identical & 1.21 GiB -\textgreater{} 0.08 GiB & 56 s
-\textgreater{} 42 s \\
\code{crypto} & identical & 1.68 GiB -\textgreater{} 0.12 GiB & 149 s
-\textgreater{} 138 s \\
\code{dma} & identical & 11.42 GiB -\textgreater{} 0.94 GiB & 845 s
-\textgreater{} 269 s \\
\end{longtable}
}

\textbf{Fabric rung {[}M: \code{j41m18/fabric/q/metrics.json},
\code{report.json}, \code{run.log}{]}.} Same cut as 18.3
(\code{RUN_2026-10-28_05-14-30}, exclusions of 18.4, CPU macro and SRAMs
black-boxed), under \code{run_capped.sh} (12 GiB, 24 h):
\textbf{completed}, 2,977 s, cgroup peak \textbf{5.12 GiB},
\code{oom_kill=0}.

\begin{landscape}
\small\setlength{\tabcolsep}{4pt}

{
\begin{longtable}[]{@{}
  >{\raggedright\arraybackslash}p{\dimexpr(\linewidth-12\tabcolsep)*143/1000\relax}
  >{\raggedright\arraybackslash}p{\dimexpr(\linewidth-12\tabcolsep)*80/1000\relax}
  >{\raggedright\arraybackslash}p{\dimexpr(\linewidth-12\tabcolsep)*199/1000\relax}
  >{\raggedright\arraybackslash}p{\dimexpr(\linewidth-12\tabcolsep)*96/1000\relax}
  >{\raggedright\arraybackslash}p{\dimexpr(\linewidth-12\tabcolsep)*128/1000\relax}
  >{\raggedright\arraybackslash}p{\dimexpr(\linewidth-12\tabcolsep)*120/1000\relax}
  >{\raggedright\arraybackslash}p{\dimexpr(\linewidth-12\tabcolsep)*215/1000\relax}@{}}
\toprule\noalign{}
\begin{minipage}[b]{\linewidth}\raggedright
Collapsed faults
\end{minipage} & \begin{minipage}[b]{\linewidth}\raggedright
Detected
\end{minipage} & \begin{minipage}[b]{\linewidth}\raggedright
Proven untestable (SAT)
\end{minipage} & \begin{minipage}[b]{\linewidth}\raggedright
Unresolved
\end{minipage} & \begin{minipage}[b]{\linewidth}\raggedright
Fault coverage
\end{minipage} & \begin{minipage}[b]{\linewidth}\raggedright
Test coverage
\end{minipage} & \begin{minipage}[b]{\linewidth}\raggedright
Patterns generated / kept
\end{minipage} \\
\midrule\noalign{}
\endhead
\bottomrule\noalign{}
\endlastfoot
412,472 & 409,022 & 3,230 & 220 & \textbf{99.16 \%} & \textbf{99.95 \%}
& 218,944 / 23,627 \\
\end{longtable}
}

\end{landscape}

The random phase is the one 18.3 recorded (256 patterns, 274,941
detected). 218,944 generated patterns x 412,472 B would have been about
90 GB in the stock byte matrix, so the stock tool could not have
finished this rung with any cap this host offers.

\textbf{Independent re-grade {[}M:
\code{j41m18/fabric/regrade_own.json}, \code{regrade_vl.json}; method of
18.5{]}.} 250 sampled faults (200 detected, 50 undetected per Quaigh):
Verilator confirms all 200 detected, \textbf{0 disagreements}; Yosys SAT
proves \textbf{all 50 undetected untestable, 0 coverage holes}; all-zero
control detects 27 vs 200. \code{regrade.py} (own simulator) agrees;
\texttt{quaigh\ simulate} was skipped (too slow at this size). Wall 5
min + 51 min, peak 4.1 GB. The 220 unresolved faults are not in the
sample.

\textbf{Reading {[}A{]}.} Every caveat of 18.6/18.7 still applies:
generic-gate collapsed faults (not Sky130 cell pins), stuck-at only, no
X (the coverage around the CPU macro, SRAMs and NPU macro is an upper
bound), 922 unscanned flops and their exclusive cones outside the
universe (whole-SoC fault coverage ceiling about 99.09 \% from 18.7),
patterns not mapped to chains or replayed on the scan netlist. The dma
finding of 18.5 (testable faults missed outside Quaigh\textquotesingle s
collapsed list) was not repeated here: the uncollapsed
\texttt{atpg-\allowbreak{}report\ -\allowbreak{}-\allowbreak{}with-\allowbreak{}redundant-\allowbreak{}faults}
grade of the fabric patterns was not run. Against decision 8: fabric
test coverage 99.95 \% \textgreater= 99.0 \% on this
flow\textquotesingle s universe; fault coverage 99.16 \% on the excluded
universe, whose whole-SoC equivalent, counting every excluded fault as
undetected against the \code{fabric_noexcl} list of 18.7, is about
409,022 / 416,278 = \textbf{98.26 \%} {[}A: the two cuts are not a
subset pair, so this is an estimate{]}. Neither is a sign-off number
until the owner restates the target (18.8) and the cell-pin
re-simulation is done.

\textbf{Uncollapsed grade (bead \bead{j41m.23}, 2026-10-30) {[}M:
\texttt{/\allowbreak{}nobackup/\allowbreak{}claude\_\allowbreak{}sim\_\allowbreak{}build/\allowbreak{}j41m23/\allowbreak{}\{dma\_\allowbreak{}full,\allowbreak{}fabric\_\allowbreak{}full\}/\allowbreak{}metrics.\allowbreak{}json},
\code{run.log}{]}.}
\texttt{quaigh\ atpg-\allowbreak{}report\ -\allowbreak{}-\allowbreak{}with-\allowbreak{}redundant-\allowbreak{}faults}
(patched binary, under \code{run_capped.sh}) on the 23,627 kept fabric
patterns. Calibration first: the same command on the j41m.15 \code{dma}
patterns gives \textbf{451,725 / 452,588}, identical to the stock binary
(18.5), with a cgroup peak of 0.61 GiB against stock\textquotesingle s
6.1 GiB. Report mode grades every pattern against every fault, so the
sparse rows hold every detection, but detection density is low enough
that they still beat the byte matrix.

\begin{landscape}
\small\setlength{\tabcolsep}{4pt}

{
\begin{longtable}[]{@{}
  >{\raggedright\arraybackslash}p{\dimexpr(\linewidth-12\tabcolsep)*236/1000\relax}
  >{\raggedright\arraybackslash}p{\dimexpr(\linewidth-12\tabcolsep)*96/1000\relax}
  >{\raggedright\arraybackslash}p{\dimexpr(\linewidth-12\tabcolsep)*96/1000\relax}
  >{\raggedright\arraybackslash}p{\dimexpr(\linewidth-12\tabcolsep)*105/1000\relax}
  >{\raggedright\arraybackslash}p{\dimexpr(\linewidth-12\tabcolsep)*140/1000\relax}
  >{\raggedright\arraybackslash}p{\dimexpr(\linewidth-12\tabcolsep)*140/1000\relax}
  >{\raggedright\arraybackslash}p{\dimexpr(\linewidth-12\tabcolsep)*166/1000\relax}@{}}
\toprule\noalign{}
\begin{minipage}[b]{\linewidth}\raggedright
Universe
\end{minipage} & \begin{minipage}[b]{\linewidth}\raggedright
Faults
\end{minipage} & \begin{minipage}[b]{\linewidth}\raggedright
Detected
\end{minipage} & \begin{minipage}[b]{\linewidth}\raggedright
Undetected
\end{minipage} & \begin{minipage}[b]{\linewidth}\raggedright
Fault coverage
\end{minipage} & \begin{minipage}[b]{\linewidth}\raggedright
Wall
\end{minipage} & \begin{minipage}[b]{\linewidth}\raggedright
Peak RAM (cgroup)
\end{minipage} \\
\midrule\noalign{}
\endhead
\bottomrule\noalign{}
\endlastfoot
fabric, collapsed (above) & 412,472 & 409,022 & 3,450 & 99.16 \% & 2,977
s (ATPG) & 5.12 GiB \\
fabric, \textbf{uncollapsed} & 1,346,520 & 1,332,673 & 13,847 &
\textbf{98.97 \%} & 919 s & 1.96 GiB \\
\end{longtable}
}

\end{landscape}

Test coverage on the uncollapsed list cannot be computed: Quaigh proves
untestable faults only on the collapsed list, so the 13,847 mix
redundant and missed faults (the dma reading of 18.5). The uncollapsed
figure is 0.19 points below the collapsed one (dma: 0.04). {[}A{]}
Scaling it by the exclusion ratio of 18.7 (412,472 / 416,278) gives
about \textbf{98.07 \%} whole-fabric fault coverage on the uncollapsed
universe: above the adopted 98.0 \% (2.1), with a margin of under 0.1
point, which a cell-pin re-grade (\code{j41m.20}) can erase. Both
universes meet the 98.0 \% fault coverage target only provisionally.
\end{relocated}


\chapter{Tessent flows}\label{tm:ch:tessent-flows}

This chapter describes the Tessent flows for test insertion and pattern generation.

\begin{relocated}{docs/design/TESSENT_ATPG_FLOW.md}
% Converted from docs/design/TESSENT_ATPG_FLOW.md by tools/docs/md2tex.py; edit this file, not the original.
\section{Tessent scan insertion + ATPG evaluation flow (Sky130 SoC
fabric)}\label{tm:flow:tessent-atpg}

DFT Stage 2 evaluation with Siemens Tessent Shell 2025.4 (epic
\code{j41m}; beads \code{j41m.8} scan insertion and \code{j41m.5}
scan-mode SDC are the real Stage 2 deliverables, this flow is the
evaluation of them). Code: \src{dft/tessent/atpg/}. \textbf{Evaluation
only: nothing here is back-annotated into PD, no RTL is modified, and
the architecture decisions of \code{DFT_ARCHITECTURE.md} (\cref{tm:ch:dft-architecture}) (test ports,
\code{scan_mode}/\code{scan_en}, reset and clock muxes,
\code{dft_clk_switch}, N4 TAP flops on no chain, CPU macro as a black
box, crypto key mask and wipe) are followed, not re-decided.}

Labels: \textbf{{[}M{]}} measured (file named), \textbf{{[}R{]}} read
from source, \textbf{{[}A{]}} assumed/\allowbreak{}inferred.

\subsection{How to run}\label{tm:flow:tessent-atpg:1-how-to-run}\origsec{1}

\begin{lstlisting}[language=bash]
cd dft/tessent/atpg
make all YOSYS=<yosys> SV2V=<sv2v>        # synth, scan, atpg (stuck-at)
make atpg-trans                            # transition-delay (LOC), 26 min
BUILD=$PWD/build sh scripts/vcs_sim.sh build/patterns/soc_stuck_serial.v $PWD/build/vcs_stuck   # VCS pattern check
make clean
\end{lstlisting}

Targets: \code{synth} (sv2v netlist -\textgreater{}
sky130\_\allowbreak{}fd\_\allowbreak{}sc\_\allowbreak{}hd gate netlist,
yosys), \code{cells} (Liberty -\textgreater{} Tessent cell library),
\code{scan} (insertion + scan DRC + DRC triage), \code{atpg} (stuck-at),
\code{atpg-trans}, \code{all}, \code{clean}. Everything is generated
into \code{dft/tessent/atpg/build/} (gitignored); small summary reports
are committed in \src{dft/tessent/atpg/reports/}.

Host requirements: \code{tessent} on PATH with \code{MGLS_LICENSE_FILE},
\code{yosys} \textgreater= 0.62, \code{sv2v}, VCS for the pattern check.
\textbf{The Sky130 PDK is not installed on the evaluation host}, so
\code{scripts/fetch_lib.sh} downloads
\code{sky130_fd_sc_hd__tt_025C_1v80.lib} from the OpenROAD-flow-scripts
sky130hd platform (sha256 pinned in the script). This is the same
Liberty the OpenROAD/\allowbreak{}LibreLane flows read, but it was NOT
taken from this repo\textquotesingle s PD runs. Tessent is started with
\texttt{<\ /\allowbreak{}dev/\allowbreak{}null} and \code{-replace} (see
section 8).

\subsection{Flow and decisions}\label{tm:flow:tessent-atpg:2-flow-and-decisions}\origsec{2}

{\footnotesize\setlength{\tabcolsep}{3pt}

{
\begin{longtable}[]{@{}
  >{\raggedright\arraybackslash}p{\dimexpr(\linewidth-4\tabcolsep)*99/1000\relax}
  >{\raggedright\arraybackslash}p{\dimexpr(\linewidth-4\tabcolsep)*440/1000\relax}
  >{\raggedright\arraybackslash}p{\dimexpr(\linewidth-4\tabcolsep)*440/1000\relax}@{}}
\toprule\noalign{}
\begin{minipage}[b]{\linewidth}\raggedright
Step
\end{minipage} & \begin{minipage}[b]{\linewidth}\raggedright
What
\end{minipage} & \begin{minipage}[b]{\linewidth}\raggedright
Decision / reason
\end{minipage} \\
\midrule\noalign{}
\endhead
\bottomrule\noalign{}
\endlastfoot
sv2v &
\texttt{make\ -\allowbreak{}C\ pnr\ sky130-\allowbreak{}soc-\allowbreak{}sv2v}
regenerated from current \code{main} (60 files,
\texttt{\_\allowbreak{}\_\allowbreak{}pnr\_\allowbreak{}\_\allowbreak{}\ SRAM\_\allowbreak{}SKY130\ USE\_\allowbreak{}SKY130\_\allowbreak{}CLKCELLS})
& A committed synthesised netlist exists only in untracked PD run
directories; none matches current RTL (TAP, MBIST, clock cells are
newer). Regenerated, not reused. \\
synth & yosys 0.64, \texttt{synth\ -\allowbreak{}flatten},
\code{dfflibmap}/\code{abc} against the tt Liberty, latch techmap,
\code{hilomap} (\code{scripts/synth.ys.in}) & CPU macro and the two SRAM
macros stay black boxes. \textbf{Not the PD netlist}: yosys picks
\code{edfxtp_1} (enable flop) for 10,174 flops where the LibreLane
recipe produced \code{dfxtp} + mux. Area/flop counts therefore differ
from the PD runs. 17m44s, 1.1 GB {[}M \code{reports/runtime.txt}{]}. \\
fix-up & \code{scripts/fixup_gate.py}: delete the six
\code{scan_out_o[5:0]} placeholder ties, rename the \texttt{\$paramod}
CPU macro & 14.4 contract step (a), done on the generated netlist (RTL
untouched). \\
cell library & \code{lib2v.py} (Liberty -\textgreater{} zero-delay
Verilog models) -\textgreater{} LibComp -\textgreater{} \code{.atpglib}
-\textgreater{} \code{write_cell_library} + \code{read_liberty}
-\textgreater{} \code{.tcelllib} with a curated
\code{dft_cell_selection} & See section 3. \\
scan &
\texttt{set\_\allowbreak{}context\ dft\ -\allowbreak{}scan\ -\allowbreak{}no\_\allowbreak{}rtl},
\texttt{add\_\allowbreak{}scan\_\allowbreak{}mode\ soc\_\allowbreak{}scan\ -\allowbreak{}chain\_\allowbreak{}count\ 6\ -\allowbreak{}single\_\allowbreak{}clock\_\allowbreak{}domain\_\allowbreak{}chains\ on\ -\allowbreak{}lockup\_\allowbreak{}cell\_\allowbreak{}type\ latch}
bound to the existing \code{scan_in_i[5:0]}/\code{scan_out_o[5:0]},
\code{insert_test_logic} & 6 fabric chains per
DFT\_\allowbreak{}ARCHITECTURE section 8 (5 \code{clk_i} + 1
\code{cpu_clk_i}); chains 6,7 belong to the CPU macro and are not
stitched here. In test mode every flop is on \code{scan_clk_i} (clock
muxes), so the tool sees one domain and balances by length, not by
domain. \\
constraints & \code{scan_mode_i}=1;
\texttt{clk\_\allowbreak{}i\ cpu\_\allowbreak{}clk\_\allowbreak{}i\ tck\_\allowbreak{}i\ tdi\_\allowbreak{}i\ trst\_\allowbreak{}ni}=0;
\texttt{tms\_\allowbreak{}i\ rst\_\allowbreak{}n\_\allowbreak{}i\ cpu\_\allowbreak{}rst\_\allowbreak{}n\_\allowbreak{}i}=1;
clocks \code{scan_clk_i} (off 0) and \code{scan_rst_ni} (off 1); scan
enable \code{scan_en_i} & TAP held in reset (\code{trst_ni}=0) so the
pin leg of \code{dft_ctrl} is the only scan control. \textbf{Deviation}:
the SE net is the port \code{scan_en_i}, not the internal
\code{dft_scan_en} (\code{soc_top.sv:360-382}); they are equal with the
TAP idle. TAP-serial scan is therefore not exercised. \\
ATPG & \texttt{set\_\allowbreak{}context\ patterns\ -\allowbreak{}scan},
\code{import_scan_mode}, stuck-at then transition
(\texttt{set\_\allowbreak{}pattern\_\allowbreak{}type\ -\allowbreak{}sequential\ 2},
LOC) & Single thread
(\texttt{add\_\allowbreak{}processors\ -\allowbreak{}number} not
accepted, section 8). \\
\end{longtable}
}

}

\subsection{Cell library (Liberty import) {[}M{]}}\label{tm:flow:tessent-atpg:3-cell-library-liberty-import-m}\origsec{3}

\code{read_liberty} alone does \textbf{not} give Tessent a usable
library: it only attaches attributes (scan pins, \code{test_cell},
clock-gating) to models that must come from \code{read_verilog} or
\code{read_cell_library}, and \code{read_cell_library} wants
\code{.atpglib}/\code{.tcelllib}, not Verilog. Pipeline
(\texttt{make\ cells}):

\begin{enumerate}
\def\labelenumi{\arabic{enumi}.}
\tightlist
\item
  \code{scripts/lib2v.py}: Liberty
  \code{function}/\code{ff}/\code{latch}/\code{statetable}
  -\textgreater{} plain Verilog-2001 models. 415 of 428 cells modelled.
  \textbf{Not modelled (13):}
  \texttt{ebufn\_\allowbreak{}\{1,\allowbreak{}2,\allowbreak{}4,\allowbreak{}8\}},
  \texttt{einvn\_\allowbreak{}\{0,\allowbreak{}1,\allowbreak{}2,\allowbreak{}4,\allowbreak{}8\}},
  \texttt{einvp\_\allowbreak{}\{1,\allowbreak{}2,\allowbreak{}4,\allowbreak{}8\}}
  (three-state; the SoC has none, \code{reports/lib2v_report.txt}). Bare
  buffers are emitted as \code{buf} gates, because
  \texttt{assign\ X\ =\allowbreak{}\ A} becomes a \code{_wire} primitive
  that Tessent will not accept as a buffer.
\item
  \code{libcomp} -\textgreater{} \code{sky130_hd.atpglib}: 415 models,
  \textbf{0 errors}, \code{lcVerify} "ALL PASSED for all patterns" (415
  models) {[}M \code{build/libcomp/libcomp.log}{]}. Note this checks the
  atpglib against \emph{my generated Verilog}, not against vendor models
  (none available on this host).
\item
  \code{read_liberty} warnings (kept, not fixed):
  \code{dlclkp_*}/\code{sdlclkp_*} lack \code{clock_gate_test_pin}, so
  Tessent will not use them as test-logic clock gaters (the SoC clock
  gate is hand-instantiated and already has \code{SCE}); the 13 tristate
  cells have no match; \code{probe*}, \code{lpflow_*} buffers fail the
  buffer check and lose the attribute (unused).
\item
  \code{scripts/mk_tcelllib.py} replaces the auto-created
  \code{dft_cell_selection} (which picks \code{lpflow_*} low-power cells
  and \code{_4} drive strengths) with a curated \code{sky130hd}
  selection (\code{buf_2}, \code{inv_2}, \code{mux2_2}, \code{dlxtn_1}
  lockup latch, \code{sdf*} scan cells, ...). Scan equivalents for every
  flop in the netlist exist in the Liberty
  (\code{dfxtp_1}-\textgreater{}\code{sdfxtp_1},
  \code{dfrtp_1}-\textgreater{}\code{sdfrtp_1},
  \code{dfstp_*}-\textgreater{}\code{sdfstp_*},
  \code{edfxtp_1}-\textgreater{}\code{sedfxtp_1}).
\end{enumerate}

\subsection{Results {[}M{]}}\label{tm:flow:tessent-atpg:4-results-m}\origsec{4}

\subsubsection{Scan insertion and DRC}\label{tm:flow:tessent-atpg:41-scan-insertion-and-drc}\origsec{4.1}

6 chains of 3044, 3044, 3044, 3044, 3044, 3043 cells (\textbf{18,263
scan cells}, max imbalance 1 cell; the doc\textquotesingle s section 8
proposal was \textasciitilde3,511 for 5 fabric chains) via existing
\code{scan_in_i[5:0]}/\code{scan_out_o[5:0]}, 0 top-level ports added,
12 instances added (the tool reports 6 retiming elements; the other 6
are not itemised), flush test passes in the ATPG-time DRC
(\texttt{Chain\ =\allowbreak{}\ soc\_\allowbreak{}scan\_\allowbreak{}chain*\ successfully\ traced}).
ScanDEF written (\code{build/soc_top.scanorder},
\code{write_scan_order}). 1m50s, 1.17 GB. The test procedure file is
inside the STIL (\code{soc_stuck.stil} Procedures block).

{\footnotesize\setlength{\tabcolsep}{3pt}

{
\begin{longtable}[]{@{}
  >{\raggedright\arraybackslash}p{\dimexpr(\linewidth-6\tabcolsep)*331/1000\relax}
  >{\raggedright\arraybackslash}p{\dimexpr(\linewidth-6\tabcolsep)*70/1000\relax}
  >{\raggedright\arraybackslash}p{\dimexpr(\linewidth-6\tabcolsep)*141/1000\relax}
  >{\raggedright\arraybackslash}p{\dimexpr(\linewidth-6\tabcolsep)*437/1000\relax}@{}}
\toprule\noalign{}
\begin{minipage}[b]{\linewidth}\raggedright
Rule
\end{minipage} & \begin{minipage}[b]{\linewidth}\raggedright
Count
\end{minipage} & \begin{minipage}[b]{\linewidth}\raggedright
Handling
\end{minipage} & \begin{minipage}[b]{\linewidth}\raggedright
Triage (by Q-net RTL hierarchy, \code{reports/drc_triage.txt})
\end{minipage} \\
\midrule\noalign{}
\endhead
\bottomrule\noalign{}
\endlastfoot
S1 unstable non-scan cells when clocks off & 417 & demoted to warning &
243 \code{u_dft_ctrl.u_tap}, 122 \code{u_mbist.u_stsync_core}, 29
\code{mbist_ctl}, 8 \code{u_mbist.u_xing_core}, 4 \code{u_clk_sw}, rest
single TAP outputs. \textbf{All TAP/MBIST TCK-domain flops: decision N4
(TAP flops on no chain).} \\
S2 clock capture ability & 471 & demoted to warning & same set (+64
\code{u_mbist.u_xing_core}): these flops are clocked by \code{tck_i},
held 0. Not a fix: the tool otherwise refuses to leave setup mode. \\
D5 non-scan memory element & 483 & warning & 415 TIE-0 (TAP flops held
by \code{trst_ni}), 6 TIE-1, 56 INIT-X, 6 TLA (transparent latches).
Same blocks. \\
C4 clock may be affected by captured data & 731 & note & all on
\code{scan_rst_ni}: the reset cone contains flops (\code{pll_locked_o},
\code{cpu_pll_locked_o}, CDC reset-sync flops). Mux selects are constant
in test mode, so likely structural pessimism; \textbf{not proven},
section 7. \\
S7 / E5 & 13 / 308 & note & S7: scannable cells tied to constants
(reset-synchroniser stages, \code{u_crypto.scan_s1_q}). E5: X
propagation from black boxes. \\
Any rule with \code{error} handling after demotion & 0 & & \\
\end{longtable}
}

}

No S1/S2 violator is outside the declared TAP/\allowbreak{}MBIST-TCK
set, so no exclusion was made to clear a report. Nothing needed an RTL
change to proceed.

\subsubsection{Coverage and patterns}\label{tm:flow:tessent-atpg:42-coverage-and-patterns}\origsec{4.2}

\begin{landscape}
\small\setlength{\tabcolsep}{4pt}

{
\begin{longtable}[]{@{}
  >{\raggedright\arraybackslash}p{\dimexpr(\linewidth-12\tabcolsep)*187/1000\relax}
  >{\raggedright\arraybackslash}p{\dimexpr(\linewidth-12\tabcolsep)*60/1000\relax}
  >{\raggedright\arraybackslash}p{\dimexpr(\linewidth-12\tabcolsep)*100/1000\relax}
  >{\raggedright\arraybackslash}p{\dimexpr(\linewidth-12\tabcolsep)*107/1000\relax}
  >{\raggedright\arraybackslash}p{\dimexpr(\linewidth-12\tabcolsep)*133/1000\relax}
  >{\raggedright\arraybackslash}p{\dimexpr(\linewidth-12\tabcolsep)*213/1000\relax}
  >{\raggedright\arraybackslash}p{\dimexpr(\linewidth-12\tabcolsep)*180/1000\relax}@{}}
\toprule\noalign{}
\begin{minipage}[b]{\linewidth}\raggedright
Fault model
\end{minipage} & \begin{minipage}[b]{\linewidth}\raggedright
Faults
\end{minipage} & \begin{minipage}[b]{\linewidth}\raggedright
Test coverage
\end{minipage} & \begin{minipage}[b]{\linewidth}\raggedright
Fault coverage
\end{minipage} & \begin{minipage}[b]{\linewidth}\raggedright
ATPG effectiveness
\end{minipage} & \begin{minipage}[b]{\linewidth}\raggedright
Patterns
\end{minipage} & \begin{minipage}[b]{\linewidth}\raggedright
CPU / wall
\end{minipage} \\
\midrule\noalign{}
\endhead
\bottomrule\noalign{}
\endlastfoot
Stuck-at & 661,690 & \textbf{94.17 \%} & 93.25 \% & 99.99 \% &
\textbf{1,574} & 292 s / 4m55s, 1.18 GB \\
Transition (LOC, 2 pulses) & 608,006 & \textbf{81.28 \%} & 80.42 \% &
99.31 \% & \textbf{2,567} (2,564 clock-sequential) & 1,587 s / 26m29s,
1.22 GB \\
\end{longtable}
}

\end{landscape}

Targets in the DFT skill are 99 \% SAF / 95 \% transition;
\textbf{neither is met}, and the project\textquotesingle s own goal
(\code{DFT_ARCHITECTURE} 2.1) is 99 \% test / 98 \% fault. The shortfall
is explained, not reclassified (section 4.3). Reports:
\code{reports/stuck_statistics.rpt},
\code{reports/trans_statistics.rpt}. Shift time at 6 x 3,044: 1,574 x
3,045 = 4.8 M cycles (0.19 s at 25 MHz); transition 7.8 M cycles {[}A,
arithmetic{]}.

\subsubsection{Where the stuck-at coverage goes (5.83 \% not
detected)}\label{tm:flow:tessent-atpg:43-where-the-stuck-at-coverage-goes-583-not-detected}\origsec{4.3}

{\footnotesize\setlength{\tabcolsep}{3pt}

{
\begin{longtable}[]{@{}
  >{\raggedright\arraybackslash}p{\dimexpr(\linewidth-4\tabcolsep)*327/1000\relax}
  >{\raggedright\arraybackslash}p{\dimexpr(\linewidth-4\tabcolsep)*327/1000\relax}
  >{\raggedright\arraybackslash}p{\dimexpr(\linewidth-4\tabcolsep)*327/1000\relax}@{}}
\toprule\noalign{}
\begin{minipage}[b]{\linewidth}\raggedright
Class (faults, \% of 661,690)
\end{minipage} & \begin{minipage}[b]{\linewidth}\raggedright
Cause
\end{minipage} & \begin{minipage}[b]{\linewidth}\raggedright
Attribution (\code{reports/stuck_fault_triage.txt}, representative
faults)
\end{minipage} \\
\midrule\noalign{}
\endhead
\bottomrule\noalign{}
\endlastfoot
AU.BB black boxes, 6,915 (1.05 \%) & CPU macro and 2 SRAM macros are X
sources/\allowbreak{}sinks; no shadow logic exists
(DFT\_\allowbreak{}ARCHITECTURE 6.3 shadow flops are not built) &
\code{u_cpu_axi_cdc} 3,097 (the CPU AXI boundary), \code{u_sram} 460,
\code{u_npu} 395, CPU-facing glue \\
AU.PC pin constraints, 14,707 (2.22 \%) & logic dead under the test-mode
constants & \code{u_crypto} 2,598 (the \code{key_q}/\code{rk_q} AND
masks, \texttt{\textasciitilde{}scan\_\allowbreak{}mode}: by design),
\code{u_mbist} 419, \code{u_dft_ctrl} 304, functional mux legs of the
clock/\allowbreak{}reset muxes \\
AU.TC tied cells, 8,407 (1.27 \%) & TAP/\allowbreak{}MBIST-TCK flops
held in reset make their cones constant & \code{u_dft_ctrl} 3,101,
\code{u_mbist} 852 (N4) \\
AU unclassified, 7,318 (1.11 \%) & not broken down by the tool &
\code{u_crypto} 1,155, \code{u_npu} 438, \code{u_mbist} 362,
I2C/\allowbreak{}TRNG/\allowbreak{}SPI/\allowbreak{}PWM/\allowbreak{}GPIO
100-270 each (CDC synchroniser cones, constant-reset flops) \\
RE redundant 6,474 (0.98 \%), PU 950, SEQ 280, UO 88 (ATPG abort) & &
\code{u_mbist} 168 SEQ \\
\end{longtable}
}

}

The 477 TAP/MBIST flops off the chains are N4 by decision (about 2.6 \%
of the sequential cells); they account for most of AU.TC and part of
AU.PC {[}A, by attribution, not by an A/B run{]}. Crypto key-mask loss
(2,598 representative faults, plus the key-fanout cones) is the
quantified cost of decision 3.

\textbf{Transition}: 81.28 \% with 89,242 unclassified AU (14.7 \%) is
not interpretable as a sign-off number. There are no timing exceptions,
no OCC, no SDC, the CDC crossings are not masked, and launch-on-capture
uses the single ATE clock \code{scan_clk_i}. It shows the flow works and
the order of magnitude; section 9 lists what a real transition run
needs.

\subsection{Pattern verification (VCS) {[}M{]}}\label{tm:flow:tessent-atpg:5-pattern-verification-vcs-m}\origsec{5}

\code{scripts/vcs_sim.sh}: VCS Y-2026.03, the scan-inserted netlist
\code{soc_top_scan.v}, the generated zero-delay cell models, X-output
stand-ins for the black boxes (\code{mk_sim_stubs.py}), and the Tessent
serial testbench (chain test + \textbf{20 sampled patterns of 1,574},
\texttt{-\allowbreak{}sample\ 20}).

{\footnotesize\setlength{\tabcolsep}{3pt}

{
\begin{longtable}[]{@{}
  >{\raggedright\arraybackslash}p{\dimexpr(\linewidth-2\tabcolsep)*490/1000\relax}
  >{\raggedright\arraybackslash}p{\dimexpr(\linewidth-2\tabcolsep)*490/1000\relax}@{}}
\toprule\noalign{}
\begin{minipage}[b]{\linewidth}\raggedright
Run
\end{minipage} & \begin{minipage}[b]{\linewidth}\raggedright
Result
\end{minipage} \\
\midrule\noalign{}
\endhead
\bottomrule\noalign{}
\endlastfoot
good netlist & \texttt{Begin\ chain\ test} / \texttt{End\ chain\ test} /
\textbf{\texttt{No\ error\ between\ simulated\ and\ expected\ patterns}}
(2.80 ms simulated, 2m34s, 153 MB) \\
negative control: one mid-file scan link \code{.SCD()} tied to
\texttt{1'b0} &
\textbf{\texttt{Error\ between\ simulated\ and\ expected\ patterns}}
(detected) \\
\end{longtable}
}

}

NOT RUN: the full 1,574-pattern set, the transition patterns,
parallel-mode patterns, and any simulation against vendor Sky130 cell
models (the cell models are derived from the same Liberty the ATPG used,
so a Liberty misreading would be common-mode). The negative control
proves the check can fail; it does not prove coverage.

\subsection{\texorpdfstring{Runtime and memory {[}M
\code{reports/runtime.txt}{]}}{Runtime and memory {[}M {]}}}\label{tm:flow:tessent-atpg:6-runtime-and-memory-m-reportsruntimetxt}\origsec{6}

yosys 17m44s / 1.10 GB; cell library \textless{} 3 min; scan insertion +
DRC 1m50s / 1.17 GB; stuck-at ATPG 4m55s / 1.18 GB; transition ATPG
26m29s / 1.22 GB; VCS compile \textasciitilde1 min, run 2m34s / 153 MB.
Host: 8 cores shared with a SpyGlass job; Tessent ran single-threaded.

\subsection{Testability findings (proposals; no RTL changed)}\label{tm:flow:tessent-atpg:7-testability-findings-proposals-no-rtl-changed}\origsec{7}

{\footnotesize\setlength{\tabcolsep}{3pt}

{
\begin{longtable}[]{@{}
  >{\raggedright\arraybackslash}p{\dimexpr(\linewidth-6\tabcolsep)*48/1000\relax}
  >{\raggedright\arraybackslash}p{\dimexpr(\linewidth-6\tabcolsep)*311/1000\relax}
  >{\raggedright\arraybackslash}p{\dimexpr(\linewidth-6\tabcolsep)*311/1000\relax}
  >{\raggedright\arraybackslash}p{\dimexpr(\linewidth-6\tabcolsep)*311/1000\relax}@{}}
\toprule\noalign{}
\begin{minipage}[b]{\linewidth}\raggedright
\#
\end{minipage} & \begin{minipage}[b]{\linewidth}\raggedright
Finding
\end{minipage} & \begin{minipage}[b]{\linewidth}\raggedright
Where
\end{minipage} & \begin{minipage}[b]{\linewidth}\raggedright
Proposal
\end{minipage} \\
\midrule\noalign{}
\endhead
\bottomrule\noalign{}
\endlastfoot
1 & \code{scan_out_o[5:0]} placeholder tie must be deleted before
insertion & \src{rtl/soc/soc_top.sv:421} & Stage 2 insertion step (this
flow deletes it in the netlist) \\
2 & SE is the internal TAP-or-pin net, not the raw port &
\src{rtl/soc/soc_top.sv:360-382} & bind SE to \code{dft_scan_en} in real
Stage 2; TAP-serial ATPG mode is untested here \\
3 & 477-483 TAP/\allowbreak{}MBIST-TCK flops on no chain (N4) &
\src{rtl/soc/dft/jtag_tap.sv}, \code{mbist_xing.sv},
\code{mbist_status_sync.sv} & accepted; they plausibly cost about 1.5 pp
of stuck-at coverage together with their cones {[}A: AU.TC 1.27 \% +
part of AU.PC, not isolated by a separate run{]}. A separate TAP scan
mode would recover it (N4 option) \\
4 & 731 C4 notes: \code{scan_rst_ni} reset cones contain flops &
\src{rtl/soc/pll/pll_subsystem.sv:121-127},
\code{pll_clkgen_stub.sv:52-57} (\code{locked_q}),
\code{cdc_reset_sync.sv} & check with the Tessent C4 report that the
muxes with constant select block the flop legs; if not, reset-pin faults
need \code{scan_rst_ni} pulses the tool will not generate \\
5 & Crypto mask makes the key cone untestable by design &
\src{rtl/periph/crypto_accel.sv:491},
\src{rtl/periph/aes128_core.sv:461} & accepted (decision 3); quantified
here: 2,598 AU.PC representative faults \\
6 & CPU AXI boundary, SRAM pins: no shadow/\allowbreak{}bypass logic &
\src{rtl/soc/soc_top.sv} (\code{u_cpu}, \code{u_sram}, \code{u_npu}) &
the section 6.3 shadow flops are the biggest recoverable block (AU.BB
1.05 \% plus their downstream cones) \\
7 & Latches without a scan equivalent: \code{dlxtn_1} in the clock gate
and \code{dft_clk_switch} enables, 6 transparent latches (TLA) &
\src{rtl/mem/rv32i_clock_gate.sv:75},
\src{rtl/soc/dft/dft_clk_switch.sv} & none needed
(clock-gate/\allowbreak{}test-clock internals); recorded \\
\end{longtable}
}

}

\subsection{Notes on tooling (for the next person)}\label{tm:flow:tessent-atpg:8-notes-on-tooling-for-the-next-person}\origsec{8}

\begin{itemize}
\tightlist
\item
  \texttt{tessent\ -\allowbreak{}shell\ -\allowbreak{}dofile\ x.\allowbreak{}tcl}
  \textbf{hangs in interactive mode after any aborted dofile}; always
  pass \texttt{<\ /\allowbreak{}dev/\allowbreak{}null}. \code{-logfile}
  is not overwritten without \code{-replace} (a stale log looks like a
  result).
\item
  Start-up is 15 s on an idle host and was up to 9 minutes on a loaded
  one.
\item
  The \code{-scan} legacy commands (\code{add_scan_groups},
  \code{add_scan_chains}, \code{run}) exist but \code{insert_test_logic}
  in \texttt{set\_\allowbreak{}context\ dft\ -\allowbreak{}scan} did
  nothing with them ("No test logic was inserted"). The working flow is
  the modern one:
  \texttt{-\allowbreak{}no\_\allowbreak{}rtl\ -\allowbreak{}design\_\allowbreak{}id\ gate},
  \code{.tcelllib} with a \code{dft_cell_selection},
  \code{add_scan_mode}, \code{analyze_scan_chains},
  \code{insert_test_logic},
  \texttt{write\_\allowbreak{}design\ -\allowbreak{}tsdb}, then
  \texttt{set\_\allowbreak{}context\ patterns\ -\allowbreak{}scan} +
  \code{read_design} + \code{import_scan_mode}
  (\code{/home/EDA/mentor/tessent2025/tessent_2025.4/share/UsageExamples/}
  \code{tessent_example_flat_flow_v2025.4.tgz} is the reference).
  \texttt{write\_\allowbreak{}design\ -\allowbreak{}tsdb} fails if the
  TSDB exists: the Makefile removes it.
\item
  \texttt{add\_\allowbreak{}processors\ -\allowbreak{}number\ N} is
  rejected in 2025.4 (one thread was used);
  \texttt{set\_\allowbreak{}multiprocessing\_\allowbreak{}options\ -\allowbreak{}multithreading\ on}
  is accepted.
\item
  The auto-generated \code{dft_cell_selection} is unusable for Sky130
  (low-power cells); a buffer, inverter, mux and lockup latch are
  mandatory or insertion aborts.
\item
  \texttt{set\_\allowbreak{}drc\_\allowbreak{}handling\ S1\ warning} /
  \texttt{S2\ warning} is required to leave setup mode with unscanned
  TCK-domain flops.
\item
  \code{pgrep}-style waiting and \code{sleep} loops are blocked in agent
  sessions; use run-in-background and completion notifications.
\end{itemize}

\subsection{Gap list toward real Stage 2}\label{tm:flow:tessent-atpg:9-gap-list-toward-real-stage-2}\origsec{9}

\begin{enumerate}
\def\labelenumi{\arabic{enumi}.}
\tightlist
\item
  \textbf{PD back-annotation}: the scan netlist is not placed. Needs the
  ScanDEF into ODB (the OpenROAD DFT stage in \src{tools/dft/} used a
  different stitcher), scan re-ordering, a buffer tree for the 18 k-sink
  SE net (Tessent connects \code{scan_en_i} straight to every scan cell
  {[}A: not measured{]}), hold repair on 18 k Q-\textgreater SI links
  (\code{GRT_RESIZER_HOLD_SLACK_MARGIN} 0.3 ns), and the
  \code{sdf*}/\code{sedfxtp} cells allowed by the flow\textquotesingle s
  \code{DONT_USE_CELLS} (unverified).
\item
  \textbf{Netlist realism}: re-run on the LibreLane synthesis netlist
  (it uses \code{dfxtp}+mux, not \code{edfxtp}) and on the PDK-installed
  Liberty and vendor Verilog models.
\item
  \textbf{Scan-mode SDC} (\code{j41m.5}): shift (\code{scan_clk_i},
  \code{scan_en_i}=1) and capture (\code{scan_en_i}=0, LOC pulses)
  constraints; the functional SDC pins \code{scan_mode_i}=0.
\item
  \textbf{CPU macro wrapper} (\code{j41m.4}): the 4,880-flop macro is a
  black box here. Its scan chains (indices 6,7), SRAM collars and shadow
  logic must exist and be modelled (a Tessent core
  description/\code{add_core_instance} with the macro\textquotesingle s
  chain info, or the macro netlist flattened in) before an SoC coverage
  number means anything.
\item
  \textbf{SRAM shadow/\allowbreak{}bypass}
  (DFT\_\allowbreak{}ARCHITECTURE 6.3) and \textbf{TAP-serial} mode (SE
  = \code{dft_scan_en}).
\item
  \textbf{At-speed}: OCC/clock controller or ATE bursts, timing-aware
  transition ATPG with SDC/STA exceptions, CDC masking,
  \code{en_latch}/\code{dft_clk_switch} handling, the Q-META tester
  contract (\code{scan_clk_i} low across IR loads).
\item
  \textbf{Crypto wipe protocol} (\code{j41m.6} limit, 14.10 item 5):
  three capture-type edges after \code{scan_mode} rises; Tessent test
  setup does not know this. The first patterns here shift immediately.
\item
  \textbf{Compression} (TestKompress/\allowbreak{}EDT) was not attempted
  (not cheap on this netlist; flat 6 x 3,044 is already modest).
  \code{mttestkompress_c} is licensed.
\item
  Convert the patterns for an ATE/JTAG path (SVF/STIL via TAP) and run
  full-set and parallel-mode simulation.
\item
  Resolve the 731 C4 notes (finding 4) and the Tessent "unclassified AU"
  (14.7 \% of transition faults).
\end{enumerate}

\src{tools/dft/check_scan_clk_rst.py} /
\texttt{make\ -\allowbreak{}C\ pnr\ dft-\allowbreak{}scan-\allowbreak{}check}
were not affected: no RTL, netlist generator or Makefile under
\src{pnr/} was changed.
\end{relocated}

\begin{relocated}{docs/design/TESSENT_MBIST_EVALUATION.md}
% Converted from docs/design/TESSENT_MBIST_EVALUATION.md by tools/docs/md2tex.py; edit this file, not the original.
\section{Siemens Tessent MemoryBIST vs the in-house MBIST (parallel
evaluation)}\label{tm:dec:tessent-mbist-evaluation}

Branch \code{feat/tessent-mbist}, flow in \src{dft/tessent/mbist/}. Epic
\code{j41m} (DFT), follows \src{docs/design/DFT_ARCHITECTURE.md} (\cref{tm:ch:dft-architecture}) section
16 (the in-house MBIST). \textbf{Nothing under \src{rtl/}, \src{pnr/} or
\src{tb/} was changed and the in-house MBIST is untouched}: the Tessent
hardware is generated into a separate netlist
(\code{dft/tessent/mbist/build/}, gitignored). The decision between the
two is the user\textquotesingle s; section 9 is a recommendation only.

Labels as in the architecture document: \textbf{{[}M{]}} measured
(command or file named), \textbf{{[}R{]}} read from a source file,
\textbf{{[}A{]}} assumed or inferred (reasoning given).

Tools: Tessent Shell 2025.4 (licence checkout verified with a smoke
run), VCS Y-2026.03, yosys 0.64+345 (the repo pins 0.62; see 6.4), sv2v
0.0.13.1 behind the RHEL8 glibc wrapper.

\subsection{What was built}\label{tm:dec:tessent-mbist-evaluation:1-what-was-built}\origsec{1}

{\footnotesize\setlength{\tabcolsep}{3pt}

{
\begin{longtable}[]{@{}
  >{\raggedright\arraybackslash}p{\dimexpr(\linewidth-2\tabcolsep)*490/1000\relax}
  >{\raggedright\arraybackslash}p{\dimexpr(\linewidth-2\tabcolsep)*490/1000\relax}@{}}
\toprule\noalign{}
\begin{minipage}[b]{\linewidth}\raggedright
Item
\end{minipage} & \begin{minipage}[b]{\linewidth}\raggedright
File
\end{minipage} \\
\midrule\noalign{}
\endhead
\bottomrule\noalign{}
\endlastfoot
Memory description (TCD) for \code{sky130_sram_4kbyte_1rw1r_32x1024_8} &
\src{dft/tessent/mbist/lib/sky130_sram_4kbyte_1rw1r_32x1024_8.tcd_memory}
(+ \code{.v_interface}, the interface-only Verilog view) \\
Evaluation top: the unmodified \code{sram_controller} and
\code{npu_weight_mem}, one clock, in-house \code{mb_*} inputs tied
inactive & \src{dft/tessent/mbist/rtl/mbist_eval_top.sv} \\
Dofiles: generate (DRC + DftSpecification), insert (hardware, ICL,
patterns), sim (VCS) & \src{dft/tessent/mbist/dofiles/} \\
Makefile: \code{mbist-generate}, \code{mbist-insert}, \code{mbist-sim},
\code{mbist-area}, \code{mbist-evidence}, \code{all}, \code{clean} &
\src{dft/tessent/mbist/Makefile} \\
Fault-injecting memory wrapper (reuses
\src{tb/cocotb/soc/sram_fi_core.sv}, plusargs set the fault) &
\src{dft/tessent/mbist/tb/sky130_sram_fi_tess.sv} \\
Fault campaign (same fault universe as the in-house golden model) and
area scripts & \code{scripts/fi_campaign.py}, \code{scripts/area.sh},
\code{scripts/area_breakdown.py} \\
Small committed results & \src{dft/tessent/mbist/evidence/} \\
\end{longtable}
}

}

Flow, as run {[}M{]}
(\texttt{make\ clean\ \&\&\ make\ mbist-\allowbreak{}generate\ mbist-\allowbreak{}insert\ mbist-\allowbreak{}sim}
from scratch reproduced the numbers below):

\begin{enumerate}
\def\labelenumi{\arabic{enumi}.}
\tightlist
\item
  \code{sv2v} of the same SoC sources the Sky130 flow uses, with
  \texttt{-\allowbreak{}-\allowbreak{}define=\allowbreak{}SRAM\_\allowbreak{}SKY130\ -\allowbreak{}-\allowbreak{}define=\allowbreak{}\_\allowbreak{}\_\allowbreak{}pnr\_\allowbreak{}\_\allowbreak{}}
  (and \textbf{not} \code{USE_SKY130_CLKCELLS}: those hand-instantiated
  clock cells are a black box to Tessent; see 7).
\item
  Tessent \texttt{set\_\allowbreak{}context\ dft\ -\allowbreak{}rtl},
  \code{tcd_memory} library, \code{read_verilog},
  \texttt{add\_\allowbreak{}clocks\ clk\ -\allowbreak{}period\ 25ns},
  \code{check_design_rules} (clean), \code{create_dft_specification},
  \code{process_dft_specification}, \code{extract_icl},
  \code{create_patterns_specification}.
\item
  \texttt{run\_\allowbreak{}testbench\_\allowbreak{}simulations\ -\allowbreak{}simulator\ vcs}
  on Tessent\textquotesingle s own generated testbenches, against the
  project\textquotesingle s behavioural model
  \src{sim/sky130_sram_4kbyte_1rw1r_32x1024_8.sv}, then against the
  fault-injecting model.
\item
  yosys
  \texttt{synth;\ dfflibmap;\ abc\ -\allowbreak{}liberty\ sky130\_\allowbreak{}fd\_\allowbreak{}sc\_\allowbreak{}hd\_\allowbreak{}\_\allowbreak{}tt\_\allowbreak{}025C\_\allowbreak{}1v80}
  on the Tessent-inserted design and on the in-house blocks, the method
  of section 16.9 of the architecture document.
\end{enumerate}

\subsection{Memory library: what is in the TCD and what was assumed}\label{tm:dec:tessent-mbist-evaluation:2-memory-library-what-is-in-the-tcd-and-what-was-assumed}\origsec{2}

Authored from the macro\textquotesingle s Liberty
(\code{pnr/sky130/soc/macro/*_TT_1p8V_25C.lib}), the OpenRAM model
beside it, and the project\textquotesingle s simulation model.
\textbf{Assumptions} (also in the file header):

{\footnotesize\setlength{\tabcolsep}{3pt}

{
\begin{longtable}[]{@{}
  >{\raggedright\arraybackslash}p{\dimexpr(\linewidth-6\tabcolsep)*48/1000\relax}
  >{\raggedright\arraybackslash}p{\dimexpr(\linewidth-6\tabcolsep)*311/1000\relax}
  >{\raggedright\arraybackslash}p{\dimexpr(\linewidth-6\tabcolsep)*311/1000\relax}
  >{\raggedright\arraybackslash}p{\dimexpr(\linewidth-6\tabcolsep)*311/1000\relax}@{}}
\toprule\noalign{}
\begin{minipage}[b]{\linewidth}\raggedright
\#
\end{minipage} & \begin{minipage}[b]{\linewidth}\raggedright
Assumption
\end{minipage} & \begin{minipage}[b]{\linewidth}\raggedright
Basis
\end{minipage} & \begin{minipage}[b]{\linewidth}\raggedright
Risk
\end{minipage} \\
\midrule\noalign{}
\endhead
\bottomrule\noalign{}
\endlastfoot
A1 & Both ports are synchronous, read latency one clock, no output
register (\texttt{DataOutStage\ None}, \texttt{OperationSet\ sync}). The
macro captures inputs at the rising edge and reads at the following
falling edge, so data is valid half a clock later and sampled at the
next rising edge. & {[}R{]}
\src{sim/sky130_sram_4kbyte_1rw1r_32x1024_8.sv} header; Liberty
\texttt{timing\_\allowbreak{}type\ :\allowbreak{}\ falling\_\allowbreak{}edge}
on \code{dout0/dout1} & The half-cycle internal timing is invisible at
the pins in RTL simulation. Tessent\textquotesingle s at-speed timing
assumptions for the real macro (SDC/\allowbreak{}MinHold) are
\textbf{not validated}. \\
A2 & \code{csb0/csb1/web0} active low, \code{wmask0} active high, one
bit per 8-bit lane (\code{GroupWriteEnable}, \texttt{BitGrouping\ 8}) &
{[}R{]} Liberty, model & none \\
A3 & Port 1 is read only (\texttt{LogicalPorts\ 1R1RW}:
\textbf{\code{[nR][nW][nRW]} order}, \texttt{1RW\ 1R} and \code{1RW1R}
are rejected) & {[}R{]} Liberty & none \\
A4 & No transparent/\allowbreak{}bypass mode
(\texttt{TransparentMode\ None}); the macro is a black box in scan mode,
as in the in-house flow & \code{DFT_ARCHITECTURE.md} (\cref{tm:ch:dft-architecture}) 6.3 & A Tessent
ATPG flow around the macro would want a bypass; not evaluated. \\
A5 & The physical row/\allowbreak{}column map is \textbf{not known}
(OpenRAM scrambles it; LEF/\allowbreak{}Liberty do not say). The address
map is a logical split
\texttt{Row[7:\allowbreak{}0]\ =\allowbreak{}\ Address[9:\allowbreak{}2]},
\texttt{Column[1:\allowbreak{}0]\ =\allowbreak{}\ Address[1:\allowbreak{}0]},
chosen because Tessent \textbf{silently forces \code{BitGrouping} to the
word size (32)} on a memory with no column address bits, which would
drop the byte-lane write-mask test. & Tessent note at TCD read & All
neighbourhood algorithms (checkerboard, the SMarchCHKB family) work in
logical space only, the same limit as the in-house \code{CKBD}
(architecture 16.3). A physical map would need the OpenRAM scrambler
documentation. \\
A6 & \texttt{MinHold\ 1},
\texttt{MilliWattsPerMegaHertz\ 0.\allowbreak{}3} (the latter is only a
power-budget input; Tessent reports
\texttt{power\_\allowbreak{}mw\ 24.\allowbreak{}00} for the run) &
placeholder & Power number is not measured; do not use it. \\
\end{longtable}
}

}

\subsection{Comparison, Tessent vs in-house}\label{tm:dec:tessent-mbist-evaluation:3-comparison-tessent-vs-in-house}\origsec{3}

Both test the two SoC-level 4 KB macros of \code{soc_top} in parallel
from one controller on \code{core_clk} (25 ns, 40 MHz).

{\footnotesize\setlength{\tabcolsep}{3pt}

{
\begin{longtable}[]{@{}
  >{\raggedright\arraybackslash}p{\dimexpr(\linewidth-4\tabcolsep)*327/1000\relax}
  >{\raggedright\arraybackslash}p{\dimexpr(\linewidth-4\tabcolsep)*327/1000\relax}
  >{\raggedright\arraybackslash}p{\dimexpr(\linewidth-4\tabcolsep)*327/1000\relax}@{}}
\toprule\noalign{}
\begin{minipage}[b]{\linewidth}\raggedright
Aspect
\end{minipage} & \begin{minipage}[b]{\linewidth}\raggedright
In-house (\code{mbist_ctrl} + \code{mem_collar}, TAP opcodes 6/7)
\end{minipage} & \begin{minipage}[b]{\linewidth}\raggedright
Tessent MemoryBIST 2025.4 (generated)
\end{minipage} \\
\midrule\noalign{}
\endhead
\bottomrule\noalign{}
\endlastfoot
Algorithms available in the evaluated build & March C-
(\code{MARCH_CM}), retention (\code{RETENTION}), both
(\code{MARCH_CM_RET}), checkerboard + March C- (\code{CKBD}), byte-lane
(\code{LANE}), \code{FULL}; six data backgrounds, optional sweep;
selected at run time through \code{MBIST_CTRL} {[}R{]} & Library of 22
algorithms in the install
(\code{lib/tools/etv/membistipg_algos/*.algo}): LVMarchCMinus,
LVMarchX/\allowbreak{}Y/\allowbreak{}LA,
LVCheckerboard1X1/\allowbreak{}4X4, LVGal*, LVWalkingPat,
LVRowBar/\allowbreak{}ColumnBar,
LVAddress/\allowbreak{}DataInterconnect, LVBitSurroundDisturb, LVMasest,
ReadOnly, SMarch, SMarchCHKB,
SMarchCHKBci/\allowbreak{}cil/\allowbreak{}vcd. The run evaluated 6
(LVMarchCMinus, SMarch, SMarchCHKB, SMarchCHKBci, SMarchCHKBcil,
LVCheckerboard1X1), all selectable at run time (controller type
\emph{Programmable}, "HardProgrammable" in the generation log)
{[}M{]} \\
March C- & yes & yes (\code{LVMARCHCMINUS}); detection on the 105-fault
universe identical to the in-house \code{MARCH_CM} (97/105 per memory,
5.1) \\
Retention & \code{RETENTION}/\code{MARCH_CM_RET}, programmable pause 0
to 4.18 M cycles & \textbf{No retention algorithm in the library}; a
\code{ParallelRetentionTest} pattern (\code{parallel_retention_time},
run mode \code{hw_default}) pauses between write and read. Detection of
a 5000-clock decay with an 8000-clock pause demonstrated (5.3).
User-defined retention would be a soft algorithm with the delay counter
{[}A: not built{]} \\
Checkerboard & logical-space \code{CKBD} & \code{LVCHECKERBOARD1X1};
\code{SMarchCHKB*} family (names match the user\textquotesingle s
\code{smarchchkd}) \\
User-defined algorithm & ROM row in \code{mbist_ctrl} (RTL change) &
Soft algorithm instructions (\code{soft_instruction_count},
\code{AlgorithmResourceOptions}) or a custom \code{.algo} file {[}R,
syntax{]}; \textbf{not exercised} \\
Fault classes (measured, 5.2) & SAF, TF, AF
(none/\allowbreak{}alias/\allowbreak{}multi),
CFin/\allowbreak{}CFid/\allowbreak{}CFst inter-word, intra-word coupling
with \code{FULL}, write-mask lane (only \code{LANE}/\code{FULL}), port-1
stuck, retention & SAF, TF, AF, CFin/\allowbreak{}CFid/\allowbreak{}CFst
inter-word 100 \% with every March-class algorithm, intra-word coupling
75 to 100 \% depending on algorithm, port-1 stuck, retention;
\textbf{write-mask lane stuck-at-1 not detected by any library
algorithm} (8/16, 5.2) \\
Test time, March C-, both macros in parallel (cycles, 25 ns) &
\textbf{10,259} (0.26 ms) & \textbf{24,768} (0.62 ms): the controller
makes one pass per logical port ("Test Port" table in the generation
log: port A, then port B); the 1R port of the NPU weight memory is tied
off inside \code{npu_weight_mem} and the pass is wasted there {[}M{]} \\
Test time, other algorithms (cycles) & \code{RETENTION} 4,106 + 2
pauses; \code{MARCH_CM_RET} 14,364 + 2 pauses; \code{CKBD} 14,367;
\code{LANE} 17,418; \code{FULL} 46,147 + 2 pauses; x6 with background
sweep {[}R formula pinned by
\code{test_cycle_counts_pin_the_sequencing}{]} & SMarchCHKBcil
\textbf{90,164} (default); SMarchCHKBci about 73.8 k; SMarchCHKB about
57.5 k; SMarch about 45.2 k; LVCheckerboard1X1 about 16.6 k (VCS time
minus about 52 us IJTAG setup); retention pattern about 15.8 k + 8 k
pause {[}M, see 5.1{]} \\
Area, sky130\_\allowbreak{}fd\_\allowbreak{}sc\_\allowbreak{}hd, same
flow (um2) & \textbf{29,733} landed (controller 13,061; two collars
1,191; TCK-\textgreater core crossing 3,761; status return 3,503; TAP
delta 8,217) & \textbf{25,850} (full, 6 algorithms) / \textbf{23,750}
(March C- only) / \textbf{27,050} (full with the Tessent-inserted TAP);
details 6 \\
Access / integration & Custom private TAP opcodes 6/7 (\code{MBIST_CTRL}
29 b, \code{MBIST_STATUS} 140 b), control word crosses TCK to
\code{core_clk} through a depth-2 \code{cdc_gray_fifo}; tester software
is plain register writes and polling & IJTAG (IEEE 1687): eight raw
ports
\texttt{ijtag\_\allowbreak{}\{tck,\allowbreak{}reset,\allowbreak{}ce,\allowbreak{}se,\allowbreak{}ue,\allowbreak{}sel,\allowbreak{}si,\allowbreak{}so\}}
on a SIB network, or a Tessent-inserted 1149.1 TAP
(\texttt{-\allowbreak{}host\_\allowbreak{}scan\_\allowbreak{}interface\_\allowbreak{}type\ tap}:
\texttt{tck\ tms\ tdi\ tdo\ trst\ tdo\_\allowbreak{}en}, +1.2 k um2).
Patterns, ICL and PDL come out of the tool. Not connectable to the
existing \code{jtag_tap} without a new IR opcode and an ICL for it
(7.2) \\
Repair / BIRA & none (pass/\allowbreak{}fail record), by design: the
OpenRAM macros have no spare rows or columns & MemoryBisr, BIRA and
fuse-map options exist (\code{RepairOptions}, \code{MemoryBisr}),
\textbf{not applicable to these macros} (no redundancy) and not
exercised \\
Diagnosis & one first-failure record per run: address, element, slot,
background, port, 32-bit XOR mask; per-memory sticky go/no-go &
per-memory go (\code{compare_memory_go}) and per-comparator go-id,
stop-on-error with 12-bit counter, failure-count limit,
bitmap/\allowbreak{}ATE diagnosis flows
(\code{CharacterizationDebugOptions}); the bitmap tooling was
\textbf{not run} \\
Effort and maintenance & RTL (\code{mbist_ctrl} 598 lines, collar,
crossing, status, TAP changes), golden model and fault suites already
written and running in CI; open source toolchain only; CPU-macro cache
collars still to do (16.7) & TCD about 100 lines + three dofiles; the
11.5 k-line controller is regenerated, not maintained; \textbf{needs a
Tessent MemoryBIST licence} (features \code{mtmemorybist_c} etc. present
on this host, absent in CI); one TCD per macro (the 1 KB cache macro
would need another); algorithm library maintained by the vendor \\
\end{longtable}
}

}

\subsection{Rules, violations and warnings triaged}\label{tm:dec:tessent-mbist-evaluation:4-rules-violations-and-warnings-triaged}\origsec{4}

{\footnotesize\setlength{\tabcolsep}{3pt}

{
\begin{longtable}[]{@{}
  >{\raggedright\arraybackslash}p{\dimexpr(\linewidth-6\tabcolsep)*48/1000\relax}
  >{\raggedright\arraybackslash}p{\dimexpr(\linewidth-6\tabcolsep)*311/1000\relax}
  >{\raggedright\arraybackslash}p{\dimexpr(\linewidth-6\tabcolsep)*311/1000\relax}
  >{\raggedright\arraybackslash}p{\dimexpr(\linewidth-6\tabcolsep)*311/1000\relax}@{}}
\toprule\noalign{}
\begin{minipage}[b]{\linewidth}\raggedright
\#
\end{minipage} & \begin{minipage}[b]{\linewidth}\raggedright
Tool message
\end{minipage} & \begin{minipage}[b]{\linewidth}\raggedright
Cause
\end{minipage} & \begin{minipage}[b]{\linewidth}\raggedright
Resolution
\end{minipage} \\
\midrule\noalign{}
\endhead
\bottomrule\noalign{}
\endlastfoot
1 &
\texttt{Unrecognized\ syntax\ ('1RW1R').\allowbreak{}\ Correct\ syntax:\allowbreak{}\ [<int>R][<int>W][<int>RW]}
& Port-type order in \code{LogicalPorts} &
\texttt{LogicalPorts\ :\allowbreak{}\ 1R1RW;} \\
2 & Note: "memory has no column address bits ... BitGrouping changed to
32" & TCD declared all address bits as row bits & Logical split
\code{Row[7:0]/Column[1:0]} (assumption A5) so \texttt{BitGrouping\ 8}
is kept \\
3 & \code{DFT_C1-1} memory clock not properly sourced, tracing stops at
\code{u_pll_sub/u_stub/out_clk_o} (then at
\code{u_pll_sub/u_ref_cm/.../muxinst/Z}) on \code{soc_top} & See section
7. First the clock path contains the black-box
\code{sky130_fd_sc_hd__mux2_4} (\code{USE_SKY130_CLKCELLS}); with that
define removed it stops at the behavioural \code{dft_clk_mux} &
Evaluation top \code{mbist_eval_top} with one \code{clk} port \\
4 & Duplicate \code{sky130_sram_4kbyte_1rw1r_32x1024_8} in simulation
(sv2v carries the blackbox stub into the design file) & Stub module and
simulation model have the same name & Stub removed from the sv2v input
list; Tessent reads \code{lib/*.v_interface} with
\texttt{-\allowbreak{}interface\_\allowbreak{}only\ -\allowbreak{}exclude\_\allowbreak{}from\_\allowbreak{}file\_\allowbreak{}dictionary} \\
5 & \texttt{Warning\ [PPHW01]} x5 in the controller generation log:
\code{NumberOfDataRegisterBits} not specified (32 bits assumed);
\texttt{DataPolarityEnable:\allowbreak{}ON\ is\ ignored\ since\ the\ physical\ data\ map\ does\ not\ exist}
for LVMARCHCMINUS and LVCHECKERBOARD1X1 (x2); group write enable port
"will be enabled" by operation set SYNC (x2, one per memory) &
Consequence of A5 (no physical map) and of the 1RW+1R TCD & Accepted.
The data-polarity warning means the two LV algorithms run without the
physical-neighbour data inversion: logical-space only \\
6 &
\texttt{Warning:\allowbreak{}\ Rule\ FN1\ violation\ occurs\ 3770\ times},
\texttt{Rule\ FP13\ .\allowbreak{}.\allowbreak{}.\allowbreak{}\ 71\ times}
at \code{extract_icl} after insertion & Netlist-sanity rules
re-evaluated on the generated RTL; \code{check_design_rules} on the
input design reports none & \textbf{Not triaged to root cause};
non-blocking (ICL extraction, patterns and VCS simulation all pass).
Needs a look at the rule text before any production use \\
7 & Default patterns use
\texttt{reduced\_\allowbreak{}address\_\allowbreak{}count\ :\allowbreak{}\ on}
(simulated in 65 us) & Tessent\textquotesingle s generated simulation
patterns shorten the address range & Own patterns specification with the
option off for every algorithm (\code{dofiles/mbist_sim.do}) so the
fault campaign runs the full 1024 addresses \\
8 &
\texttt{Error:\allowbreak{}\ Unrecognized\ argument\ '-\allowbreak{}pattern\_\allowbreak{}id'}
for \code{process_patterns_specification} & Argument exists only on
\code{run_testbench_simulations} & Fixed in
\code{dofiles/mbist_sim.do} \\
9 & yosys: \texttt{\$finish\ executed} while reading the sv2v output &
\texttt{\$fatal} elaboration checks of \code{sram_controller} at its
default \code{MEM_WORDS=4096} &
\texttt{read\_\allowbreak{}verilog\ -\allowbreak{}defer} (as the
repo\textquotesingle s synthesis gate does) \\
\end{longtable}
}

}

\subsection{Simulation and fault-injection results {[}M{]}}\label{tm:dec:tessent-mbist-evaluation:5-simulation-and-fault-injection-results-m}\origsec{5}

\subsubsection{Fault-free simulation on VCS}\label{tm:dec:tessent-mbist-evaluation:51-fault-free-simulation-on-vcs}\origsec{5.1}

Tessent-generated testbenches, the project\textquotesingle s behavioural
SRAM model, full address count, per-memory go compare:

{
\begin{longtable}[]{@{}
  >{\raggedright\arraybackslash}p{\dimexpr(\linewidth-4\tabcolsep)*330/1000\relax}
  >{\raggedright\arraybackslash}p{\dimexpr(\linewidth-4\tabcolsep)*198/1000\relax}
  >{\raggedright\arraybackslash}p{\dimexpr(\linewidth-4\tabcolsep)*452/1000\relax}@{}}
\toprule\noalign{}
\begin{minipage}[b]{\linewidth}\raggedright
Pattern (algorithm)
\end{minipage} & \begin{minipage}[b]{\linewidth}\raggedright
Result
\end{minipage} & \begin{minipage}[b]{\linewidth}\raggedright
Simulated time (incl. about 52 us IJTAG setup)
\end{minipage} \\
\midrule\noalign{}
\endhead
\bottomrule\noalign{}
\endlastfoot
ICLNetwork (IJTAG network verify) & pass, 0 miscompares & 10 us \\
LVMARCHCMINUS & pass, 0 miscompares & 671.7 us \\
SMARCH & pass & 1182.0 us \\
SMARCHCHKB & pass & 1488.8 us \\
SMARCHCHKBCI & pass & 1897.0 us \\
SMARCHCHKBCIL & pass & 2302.2 us \\
LVCHECKERBOARD1X1 & pass & 467.1 us \\
Retention (200 us pause) & pass & 647.1 us \\
\end{longtable}
}

\texttt{unscheduled\ 0\ queued\ 0\ running\ 0\ pass\ 8\ fail\ 0}. The
same eight patterns also pass against the fault-injecting model with all
faults off (\texttt{pass\ 8\ fail\ 0}); the in-house suite proves that
model equal to the real simulation model
(\code{test_fi_model_equals_real_model}). The default signoff patterns
of \texttt{make\ mbist-\allowbreak{}insert} (reduced address count) also
pass, 3 of 3.

\subsubsection{Fault injection}\label{tm:dec:tessent-mbist-evaluation:52-fault-injection}\origsec{5.2}

The in-house fault catalogue
(\code{tb/cocotb/soc/mbist_golden.py::catalogue(10)}: 105 faults at
AW=10) was injected into \textbf{each} memory (\code{sr} =
\code{sram_controller} macro, \code{wm} = NPU weight memory) through the
fault-injecting model, plus two read-port-1 stuck-at faults on \code{sr}
(the weight memory ties port 1 off). 1,272 simulations, all from one
compiled \code{simv} per pattern; a fault is \textbf{detected} if the
testbench reports more than zero miscompares (the fault-free run reports
zero, checked first). Counts below are pooled over both memories (212 =
2 x 105 + 2).

\begin{landscape}
\small\setlength{\tabcolsep}{4pt}

{
\begin{longtable}[]{@{}
  >{\raggedright\arraybackslash}p{\dimexpr(\linewidth-12\tabcolsep)*181/1000\relax}
  >{\raggedright\arraybackslash}p{\dimexpr(\linewidth-12\tabcolsep)*143/1000\relax}
  >{\raggedright\arraybackslash}p{\dimexpr(\linewidth-12\tabcolsep)*86/1000\relax}
  >{\raggedright\arraybackslash}p{\dimexpr(\linewidth-12\tabcolsep)*114/1000\relax}
  >{\raggedright\arraybackslash}p{\dimexpr(\linewidth-12\tabcolsep)*133/1000\relax}
  >{\raggedright\arraybackslash}p{\dimexpr(\linewidth-12\tabcolsep)*143/1000\relax}
  >{\raggedright\arraybackslash}p{\dimexpr(\linewidth-12\tabcolsep)*181/1000\relax}@{}}
\toprule\noalign{}
\begin{minipage}[b]{\linewidth}\raggedright
Fault class
\end{minipage} & \begin{minipage}[b]{\linewidth}\raggedright
LVMARCHCMINUS
\end{minipage} & \begin{minipage}[b]{\linewidth}\raggedright
SMARCH
\end{minipage} & \begin{minipage}[b]{\linewidth}\raggedright
SMARCHCHKB
\end{minipage} & \begin{minipage}[b]{\linewidth}\raggedright
SMARCHCHKBCI
\end{minipage} & \begin{minipage}[b]{\linewidth}\raggedright
SMARCHCHKBCIL
\end{minipage} & \begin{minipage}[b]{\linewidth}\raggedright
LVCHECKERBOARD1X1
\end{minipage} \\
\midrule\noalign{}
\endhead
\bottomrule\noalign{}
\endlastfoot
AFalias & 6/6 & 6/6 & 6/6 & 6/6 & 6/6 & 4/6 \\
AFmulti & 6/6 & 6/6 & 6/6 & 6/6 & 6/6 & 4/6 \\
AFnone & 6/6 & 6/6 & 6/6 & 6/6 & 6/6 & 6/6 \\
CFid & 40/40 & 40/40 & 40/40 & 40/40 & 40/40 & 8/40 \\
CFid (intra-word) & 4/8 & 4/8 & 6/8 & 6/8 & 6/8 & 0/8 \\
CFin & 20/20 & 20/20 & 20/20 & 20/20 & 20/20 & 8/20 \\
CFin (intra-word) & 4/4 & 4/4 & 4/4 & 4/4 & 4/4 & 0/4 \\
CFst & 40/40 & 40/40 & 40/40 & 40/40 & 40/40 & 20/40 \\
CFst (intra-word) & 4/8 & 4/8 & 8/8 & 8/8 & 8/8 & 4/8 \\
P1STUCK & 2/2 & 2/2 & 2/2 & 2/2 & 2/2 & 2/2 \\
SAF & 28/28 & 28/28 & 28/28 & 28/28 & 28/28 & 28/28 \\
TF & 28/28 & 28/28 & 28/28 & 28/28 & 28/28 & 28/28 \\
WMASK & 8/16 & 8/16 & 8/16 & 8/16 & 8/16 & 8/16 \\
\textbf{all} & \textbf{196/212} & \textbf{196/212} & \textbf{202/212} &
\textbf{202/212} & \textbf{202/212} & \textbf{120/212} \\
\end{longtable}
}

\end{landscape}

{
\begin{longtable}[]{@{}lllll@{}}
\toprule\noalign{}
Fault class & MARCH\_CM & CKBD & LANE & FULL \\
\midrule\noalign{}
\endhead
\bottomrule\noalign{}
\endlastfoot
AFalias & 3/3 & 3/3 & 3/3 & 3/3 \\
AFmulti & 3/3 & 3/3 & 3/3 & 3/3 \\
AFnone & 3/3 & 3/3 & 3/3 & 3/3 \\
CFid & 20/20 & 20/20 & 5/20 & 20/20 \\
CFid (intra-word) & 2/4 & 2/4 & 2/4 & 4/4 \\
CFin & 10/10 & 10/10 & 5/10 & 10/10 \\
CFin (intra-word) & 2/2 & 2/2 & 2/2 & 2/2 \\
CFst & 20/20 & 20/20 & 15/20 & 20/20 \\
CFst (intra-word) & 2/4 & 2/4 & 2/4 & 4/4 \\
SAF & 14/14 & 14/14 & 14/14 & 14/14 \\
TF & 14/14 & 14/14 & 14/14 & 14/14 \\
WMASK & 4/8 & 4/8 & 8/8 & 8/8 \\
\textbf{all} & \textbf{97/105} & \textbf{97/105} & \textbf{76/105} &
\textbf{105/105} \\
\end{longtable}
}

The in-house columns are the \textbf{golden-model prediction} for the
same faults (the repo\textquotesingle s fault-injection suites assert
that the RTL reproduces the golden model exactly,
\code{test_mbist_ctrl.py}), not a new RTL simulation.

Readings:

\begin{itemize}
\tightlist
\item
  Tessent \code{LVMARCHCMINUS} and the in-house \code{MARCH_CM} detect
  the same faults: 97 of 105 per memory. Both miss the same classes:
  half the intra-word coupling cases and the write-mask faults that a
  full-word March cannot see.
\item
  The SMarchCHKB family adds intra-word coupling (up to 100 of 105 per
  memory) and \code{LVCHECKERBOARD1X1} alone is weak
  (120/\allowbreak{}212), as expected of a checkerboard without a March.
\item
  \textbf{Write-mask lane faults:} every Tessent library algorithm
  leaves the stuck-at-1 mask-lane cases undetected (8/16). The in-house
  \code{LANE}/\code{FULL} algorithms detect all 8 per memory. This is
  the one fault class where the in-house algorithm set is stronger, and
  it is a property of the library algorithms plus the
  \texttt{group\_\allowbreak{}write\_\allowbreak{}enable\_\allowbreak{}control\ :\allowbreak{}\ always\_\allowbreak{}on}
  default {[}A: a system-controlled mask
  (\texttt{group\_\allowbreak{}write\_\allowbreak{}enable\_\allowbreak{}control\ :\allowbreak{}\ system})
  with a soft algorithm might close it; not tried{]}.
\item
  In-house \code{FULL} (all of the above with a background sweep)
  detects 105/105. No Tessent algorithm run reached 105/105.
\end{itemize}

\subsubsection{Retention}\label{tm:dec:tessent-mbist-evaluation:53-retention}\origsec{5.3}

{\footnotesize\setlength{\tabcolsep}{3pt}

{
\begin{longtable}[]{@{}
  >{\raggedright\arraybackslash}p{\dimexpr(\linewidth-4\tabcolsep)*437/1000\relax}
  >{\raggedright\arraybackslash}p{\dimexpr(\linewidth-4\tabcolsep)*360/1000\relax}
  >{\raggedright\arraybackslash}p{\dimexpr(\linewidth-4\tabcolsep)*183/1000\relax}@{}}
\toprule\noalign{}
\begin{minipage}[b]{\linewidth}\raggedright
Case
\end{minipage} & \begin{minipage}[b]{\linewidth}\raggedright
\code{RETENTION} pattern (200 us = 8,000-clock pause)
\end{minipage} & \begin{minipage}[b]{\linewidth}\raggedright
LVMARCHCMINUS (no pause)
\end{minipage} \\
\midrule\noalign{}
\endhead
\bottomrule\noalign{}
\endlastfoot
Cell decays after 5,000 clocks (longer than the longest write-to-read
gap of a pause-free March, about 2,100 clocks, shorter than the pause) &
\textbf{detected} on both memories & not detected \\
Cell decays after 200,000 clocks (longer than the pause) & not detected
(pause too short, as it must be) & not detected \\
\end{longtable}
}

}

So the pause is meaningful in both directions. The in-house test has the
same property (its programmable pause is also pinned by the
repo\textquotesingle s fault-injection suite).

\subsection{Area}\label{tm:dec:tessent-mbist-evaluation:6-area}\origsec{6}

\subsubsection{Method {[}M{]}}\label{tm:dec:tessent-mbist-evaluation:61-method-m}\origsec{6.1}

\code{scripts/area.sh}: yosys
\texttt{synth\ -\allowbreak{}top;\ dfflibmap\ -\allowbreak{}liberty\ L;\ abc\ -\allowbreak{}liberty\ L;\ opt\_\allowbreak{}clean\ -\allowbreak{}purge;\ stat\ -\allowbreak{}liberty\ L}
with
\texttt{L\ =\allowbreak{}\ sky130\_\allowbreak{}fd\_\allowbreak{}sc\_\allowbreak{}hd\_\allowbreak{}\_\allowbreak{}tt\_\allowbreak{}025C\_\allowbreak{}1v80.\allowbreak{}lib},
no timing-driven sizing, hierarchy kept
(\texttt{read\_\allowbreak{}verilog\ -\allowbreak{}defer}). In-house
blocks are synthesised alone with the parameters of \code{soc_mbist}
(\texttt{N\_\allowbreak{}MEM=\allowbreak{}2,\allowbreak{}\ AW=\allowbreak{}10,\allowbreak{}\ P1\_\allowbreak{}MASK=\allowbreak{}2'b01}).
The Tessent number is the \textbf{difference} between the whole
evaluation top after insertion and before (\code{mbist_eval_top}, memory
macro as black box), so it includes the memory interfaces, comparators,
controller, BAP and SIBs.

\subsubsection{Results (um2)}\label{tm:dec:tessent-mbist-evaluation:62-results-um2}\origsec{6.2}

{\small\setlength{\tabcolsep}{4pt}

{
\begin{longtable}[]{@{}
  >{\raggedright\arraybackslash}p{\dimexpr(\linewidth-8\tabcolsep)*323/1000\relax}
  >{\raggedright\arraybackslash}p{\dimexpr(\linewidth-8\tabcolsep)*266/1000\relax}
  >{\raggedright\arraybackslash}p{\dimexpr(\linewidth-8\tabcolsep)*182/1000\relax}
  >{\raggedright\arraybackslash}p{\dimexpr(\linewidth-8\tabcolsep)*73/1000\relax}
  >{\raggedright\arraybackslash}p{\dimexpr(\linewidth-8\tabcolsep)*136/1000\relax}@{}}
\toprule\noalign{}
\begin{minipage}[b]{\linewidth}\raggedright
Block
\end{minipage} & \begin{minipage}[b]{\linewidth}\raggedright
In-house
\end{minipage} & \begin{minipage}[b]{\linewidth}\raggedright
Tessent full
\end{minipage} & \begin{minipage}[b]{\linewidth}\raggedright
Tessent min
\end{minipage} & \begin{minipage}[b]{\linewidth}\raggedright
Tessent + TAP
\end{minipage} \\
\midrule\noalign{}
\endhead
\bottomrule\noalign{}
\endlastfoot
Baseline (\code{sram_controller} + \code{npu_weight_mem} + in-house
collars) & 7,073 & 7,073 & 7,073 & 7,073 \\
Same design after insertion & n/a & 32,923 & 30,823 & 34,139 \\
\textbf{Inserted MBIST delta} & see below & \textbf{25,850} &
\textbf{23,750} & \textbf{27,066} \\
Controller (sequencer, address/\allowbreak{}data generators, FSM) &
\code{mbist_ctrl} 13,061 (243 flops) & about 14,380 & about 12,260 &
about 14,380 \\
Memory interfaces + comparators + status & 2 collars 1,191 (656 + 536) &
about 8,090 & about 8,090 & about 8,090 \\
Access (TAP delta / BAP + SIBs / Tessent TAP) & TAP delta 8,217
{[}recorded, 16.9, not re-measured{]} & about 2,180 (BAP 1,549 + SIB
628) & about 2,180 & about 3,380 (+1,200 TAP) \\
CDC (TCK to core, status return) & crossing 3,761 + status 3,503 & not
separately present {[}A{]} & & \\
\textbf{Total} & \textbf{29,733} & \textbf{25,850} & \textbf{23,750} &
\textbf{27,066} \\
\end{longtable}
}

}

Notes: (a) the Tessent breakdown is the sum of per-module areas and
accounts for about 96 \% of the delta, the rest is not attributed to a
module {[}A: multiply-instantiated small cells{]}; (b) ten
\texttt{\$\_\allowbreak{}DLATCH\_\allowbreak{}N\_\allowbreak{}} in the
Tessent clock-gate cells are not mapped by this flow (about 150 um2
missing); (c) "min" is
\texttt{AdvancedOptions/\allowbreak{}algorithm\ LVMARCHCMINUS} with no
extra algorithms, the controller is still reported \emph{Programmable}
and Tessent still lists a second library algorithm in its log, so "min"
is not a hard-wired single-algorithm controller {[}M:
\code{evidence/generation/min_c1.generation_log}{]}; (d) the in-house
\code{mbist_ctrl} re-measured here is 13,061 against 13,314 recorded in
16.9 (yosys 0.64 against 0.62, \code{mem_collar} 656/536 against
653/513); (e) all numbers are before the real flow\textquotesingle s
buffering and upsizing.

The two are comparable in size. Tessent is 13 \% smaller than the
in-house total with six algorithms selectable and per-memory go/no-go,
20 \% smaller in the March-C- only configuration, and 9 \% smaller with
its own 1149.1 TAP. The in-house figure also pays for a TCK-to-core
crossing and a 140-bit status mirror that the Tessent IJTAG interface
does not have in this form {[}A: how Tessent\textquotesingle s BAP
crosses from \code{ijtag_tck} to the BIST clock was not studied; its
generated \code{posedge_synchronizer_reset} cells show it has
synchronisers{]}.

\subsection{Integration findings}\label{tm:dec:tessent-mbist-evaluation:7-integration-findings}\origsec{7}

\subsubsection{\texorpdfstring{\code{soc_top} as the Tessent
target}{as the Tessent target}}\label{tm:dec:tessent-mbist-evaluation:71-soc-top-as-the-tessent-target}\origsec{7.1}

Insertion into \code{soc_top} itself failed at \code{check_design_rules}
with \code{DFT_C1-1}: the memory clock \code{core_clk} reaches the
macros through \code{dft_clk_mux}, whose select is \code{dft_scan_mode},
an OR of the primary input \code{scan_mode_i} and TAP flops of
\code{dft_ctrl}. Tessent can only constrain \textbf{primary inputs}
(\code{add_input_constraints} rejects internal nets:
\texttt{'dft\_\allowbreak{}scan\_\allowbreak{}mode'\ is\ not\ a\ primary\ input}),
so it cannot prove the mux transparent. Tried:
\texttt{add\_\allowbreak{}input\_\allowbreak{}constraints\ scan\_\allowbreak{}mode\_\allowbreak{}i\ -\allowbreak{}c0},
additionally declaring \code{scan_clk_i} as a clock, a pseudo primary
input at \code{u_pll_sub/ref_clk_o}: none removed the error {[}M, run
interactively; logs not kept{]}. With \code{USE_SKY130_CLKCELLS} the
trace first stops at the black-box library mux. That the cause is the
non-constant select is inferred from the message and the three failed
workarounds {[}A{]}. Practical consequence: running Tessent on the full
SoC needs either an RTL/\allowbreak{}constraint change (a static
test-mode input for that mux) or an internal-pin clock declaration,
which would move the clock root for simulation. This evaluation
therefore uses \code{mbist_eval_top}, which holds the two real SRAM
users with one clock.

\subsubsection{Access to the existing TAP}\label{tm:dec:tessent-mbist-evaluation:72-access-to-the-existing-tap}\origsec{7.2}

The in-house \code{jtag_tap} has no IJTAG instruction. Hosting
Tessent\textquotesingle s SIB network behind it needs (a) an additional
IR opcode (the existing numbering is frozen by decision Q-OP, so this is
additive, not a renumbering) selecting a 1687 host data register, (b) an
ICL for \code{jtag_tap} so Tessent can retarget patterns, and (c) a
decision on the second \code{dft_ctrl}/TCK clock crossing. None of this
was built. The two options evaluated instead: raw IJTAG ports (8 pins, a
sibling of the TAP) and Tessent\textquotesingle s own TAP (6 pins, +1.2
k um2), which would duplicate the project\textquotesingle s TAP.

\subsection{Not done, and limits}\label{tm:dec:tessent-mbist-evaluation:8-not-done-and-limits}\origsec{8}

\begin{itemize}
\tightlist
\item
  Gate-level simulation, SDC for the generated hardware, timing of the
  controller and at-speed behaviour of the real macro: not run.
\item
  Tessent on \code{soc_top}, on the 1 KB CPU-macro cache SRAMs, and
  combination with the in-house collar left in place: not run.
\item
  The \code{smarchchkd} request: the library contains \code{SMarchCHKB},
  \code{SMarchCHKBci}, \code{SMarchCHKBcil} and \code{SMarchCHKBvcd};
  the element sequences are licensed content in the install and are
  neither copied nor reproduced here. \code{SMarchCHKBvcd} was not
  generated.
\item
  Physical row/\allowbreak{}column mapping unknown (A5), so no claim
  about physical-neighbour coverage in either flow.
\item
  Fault models are the functional ones of \code{sram_fi_core.sv}; no
  analogue/\allowbreak{}resistive defects.
\item
  Parallel jobs: the fault campaign ran two \code{simv} processes at a
  time; Tessent ran with
  \texttt{set\_\allowbreak{}multiprocessing\_\allowbreak{}options\ -\allowbreak{}processors\ 2}.
\end{itemize}

\subsection{Recommendation (the decision is the
user\textquotesingle s)}\label{tm:dec:tessent-mbist-evaluation:9-recommendation-the-decision-is-the-users}\origsec{9}

\begin{enumerate}
\def\labelenumi{\arabic{enumi}.}
\tightlist
\item
  \textbf{Keep the in-house MBIST for the Sky130 SoC.} It already meets
  the section 16 requirements, is verified by golden models and fault
  suites that run in open-source CI, detects every fault of the shared
  catalogue (105/105 with \code{FULL}) including write-mask lane faults
  that no Tessent library algorithm detected, tests both ports in one
  pass (about 2.4 times fewer cycles for March C-), and needs no
  licence.
\item
  \textbf{Do not replace it with Tessent on this evidence.}
  Tessent\textquotesingle s area is only 13 to 20 \% lower, it has no
  retention algorithm in the library, its access network does not plug
  into the existing TAP without RTL changes, and it cannot be run on
  \code{soc_top} as is (7.1).
\item
  \textbf{Where Tessent would earn its place:} (a) if the project ever
  moves to a memory with redundancy and needs BIRA/\allowbreak{}BISR,
  (b) if an ATE flow needs standard ICL/PDL retargeting, STIL/WGL
  patterns and the bitmap/\allowbreak{}diagnosis tooling, (c) if a
  vendor-qualified algorithm set (SMarchCHKB family) is a customer
  requirement. For (c) the cheapest route is adding the Tessent flow in
  \src{dft/tessent/mbist/} as a \textbf{cross-check} in a licensed
  environment: the fault campaign script and the TCD are ready, and it
  already shows the algorithm sets agree on the common faults.
\item
  \textbf{Cheap follow-ups either way:} try
  \texttt{group\_\allowbreak{}write\_\allowbreak{}enable\_\allowbreak{}control\ :\allowbreak{}\ system}
  to see whether Tessent can close the write-mask gap; build the
  physical address map if OpenRAM documents the scrambler; and note that
  a Tessent controller needs a per-macro TCD, so the CPU-macro 1 KB
  cache SRAM (not covered here) would need a second one.
\end{enumerate}

\subsection{Reproducing}\label{tm:dec:tessent-mbist-evaluation:10-reproducing}\origsec{10}

\begin{lstlisting}[language={}]
cd dft/tessent/mbist
# host overrides (not committed): host/env.mk  (SV2V := <sv2v wrapper>)   host/env.sh  (YOSYS, SV2V, SKY130_LIB)
source ../../../tools/eda/sim_env.sh            # VCS licence/env
make mbist-generate mbist-insert mbist-sim mbist-area mbist-evidence
\end{lstlisting}

Licence: \code{MGLS_LICENSE_FILE=27017@eda.ic}. Runtime on this host:
generate + insert about 2 min, simulation about 2 min, fault campaign
about 10 min, area about 3 min.
\end{relocated}

\chapter{FPGA prototyping}\label{tm:ch:fpga-prototyping}

This chapter describes the FPGA prototype. It gives the board choice and the measured results.

\begin{relocated}{docs/FPGA_BOARD_SELECTION.md}
% Converted from docs/FPGA_BOARD_SELECTION.md by tools/docs/md2tex.py; edit this file, not the original.
\section{\texorpdfstring{FPGA board selection (bead \code{hlrm}, GH
\#245)}{FPGA board selection (bead , GH \#245)}}\label{tm:dec:fpga-board-selection}

Status: \textbf{recommendation, not a purchase decision.} Written
2026-10-14. It builds on the yosys resource counts in
\code{FPGA_STAGE0_EVALUATION.md} (\cref{tm:dec:fpga-stage0-evaluation})
§2. fmax exists as an ECP5 lower bound and as a real Vivado Artix-7
result for the CPU-only build (§5, "Measured: ECP5 lower bound" and
"Measured: Artix-7 XC7A200T (Vivado 2026.1)"); the "fits" for the full
SoC with GPU stays an estimate unless that subsection says otherwise.

\subsection{How much FPGA the design needs}\label{tm:dec:fpga-board-selection:1-how-much-fpga-the-design-needs}\origsec{1}

Xilinx 7-series LUT6 counts from yosys 0.62 \code{synth_xilinx} (Stage
0, measured), plus assumed overheads:

{\small\setlength{\tabcolsep}{4pt}

{
\begin{longtable}[]{@{}
  >{\raggedright\arraybackslash}p{\dimexpr(\linewidth-8\tabcolsep)*122/1000\relax}
  >{\raggedright\arraybackslash}p{\dimexpr(\linewidth-8\tabcolsep)*148/1000\relax}
  >{\raggedright\arraybackslash}p{\dimexpr(\linewidth-8\tabcolsep)*271/1000\relax}
  >{\raggedright\arraybackslash}p{\dimexpr(\linewidth-8\tabcolsep)*232/1000\relax}
  >{\raggedright\arraybackslash}p{\dimexpr(\linewidth-8\tabcolsep)*206/1000\relax}@{}}
\toprule\noalign{}
\begin{minipage}[b]{\linewidth}\raggedright
Build
\end{minipage} & \begin{minipage}[b]{\linewidth}\raggedright
Today (measured)
\end{minipage} & \begin{minipage}[b]{\linewidth}\raggedright
+ bring-up / debug logic (\textasciitilde5 K, assumed)
\end{minipage} & \begin{minipage}[b]{\linewidth}\raggedright
+ 25 \% future RTL growth (assumed)
\end{minipage} & \begin{minipage}[b]{\linewidth}\raggedright
Device size at \textless= 70 \% LUT use
\end{minipage} \\
\midrule\noalign{}
\endhead
\bottomrule\noalign{}
\endlastfoot
CPU only (no GPU) & 38.6 K & \textasciitilde44 K & \textasciitilde55 K &
\textasciitilde78 K \\
Full SoC with GPU & 62.4 K + \textasciitilde5 K LUT-RAM &
\textasciitilde72 K & \textasciitilde90 K & \textasciitilde120-130 K \\
\end{longtable}
}

}

\begin{itemize}
\tightlist
\item
  \textbf{Bring-up logic} means a debug bridge (JTAG-to-AXI or UART),
  one or two ILA cores, clocking and an IO wrapper. The \textasciitilde5
  K figure is an assumption, not measured.
\item
  \textbf{25 \% growth} is headroom for future RTL (a larger NPU, more
  IP, a bigger SRAM), not a planned feature list.
\item
  \textbf{70 \% LUT} is a common ceiling for routability and timing
  closure.
\item
  \textbf{Margin in our favour:} Vivado typically reports 10-30 \% fewer
  LUTs than yosys (inference, not measured here).
\item
  \textbf{GPU shared memory:} 7-series mapped the 16 KB GPU shared
  memory to LUT-RAM (\textasciitilde5 K LUT sites). Forcing block RAM
  (\texttt{ram\_\allowbreak{}style\ =\allowbreak{}\ "block"}) would
  remove that.
\item
  \textbf{Not the limit:} block RAM (\textasciitilde9 RAMB36-eq) and DSP
  (19-44) are small next to every candidate below.
\item
  \textbf{Main SRAM is 4 KB today}
  (\texttt{SRAM\_\allowbreak{}MEM\_\allowbreak{}WORDS\ =\allowbreak{}\ 1024}).
  Firmware validation may want much more; every candidate has enough
  block RAM for tens to hundreds of KB.
\end{itemize}

\subsection{Candidates, with Singapore price and stock}\label{tm:dec:fpga-board-selection:2-candidates-with-singapore-price-and-stock}\origsec{2}

Prices are distributor snapshots of unknown date, \textbf{before
Singapore\textquotesingle s 9 \% GST}. Confirm before ordering.

{\footnotesize\setlength{\tabcolsep}{3pt}

{
\begin{longtable}[]{@{}
  >{\raggedright\arraybackslash}p{\dimexpr(\linewidth-8\tabcolsep)*286/1000\relax}
  >{\raggedright\arraybackslash}p{\dimexpr(\linewidth-8\tabcolsep)*86/1000\relax}
  >{\raggedright\arraybackslash}p{\dimexpr(\linewidth-8\tabcolsep)*126/1000\relax}
  >{\raggedright\arraybackslash}p{\dimexpr(\linewidth-8\tabcolsep)*126/1000\relax}
  >{\raggedright\arraybackslash}p{\dimexpr(\linewidth-8\tabcolsep)*355/1000\relax}@{}}
\toprule\noalign{}
\begin{minipage}[b]{\linewidth}\raggedright
Board (FPGA)
\end{minipage} & \begin{minipage}[b]{\linewidth}\raggedright
LUTs
\end{minipage} & \begin{minipage}[b]{\linewidth}\raggedright
Full SoC with growth
\end{minipage} & \begin{minipage}[b]{\linewidth}\raggedright
CPU only with growth
\end{minipage} & \begin{minipage}[b]{\linewidth}\raggedright
Singapore price and stock
\end{minipage} \\
\midrule\noalign{}
\endhead
\bottomrule\noalign{}
\endlastfoot
\textbf{Digilent Nexys Video} (XC7A200T, -1) & 134.6 K &
\textbf{\textasciitilde62 \%, fits} & \textasciitilde37 \% &
\textbf{S\$873.49, 83 in stock} at DigiKey SG (17-week lead time when
out); RS SG had 3 \\
Digilent Genesys 2 (XC7K325T, -2) & 203.8 K & \textasciitilde41 \%, fits
easily & \textasciitilde25 \% &
**S\(1,529.99, 46 in stock** at DigiKey SG (18-week lead time); Mouser SG S\)1,527.61,
53 in stock \\
AMD Kria KV260 (Zynq UltraScale+ XCK26) & \textasciitilde117 K &
\textasciitilde72 \%, at the limit & \textasciitilde43 \% & No Singapore
listing found. DigiKey US US\(283.86, 420 in stock (about S\)380 at
\textasciitilde1.34 S\(/US\), before shipping and GST) \\
Digilent Arty A7-100T (XC7A100T) & 63.4 K & does not fit &
\textasciitilde87 \%, too tight & No Singapore listing found. DigiKey US
US\$299 \\
Radiona ULX3S-85F (Lattice ECP5, open toolchain) & 83.6 K (LUT4) & does
not fit & \textasciitilde60 \%, fits & No Singapore distributor. Crowd
Supply US\(235 + US\)18 worldwide shipping \\
\end{longtable}
}

}

XCK26 resources (Farnell listing): 256 K system logic cells, 144 block
RAM, 64 UltraRAM, 1.2 K DSP.

\subsection{Recommendation}\label{tm:dec:fpga-board-selection:3-recommendation}\origsec{3}

\textbf{Digilent Nexys Video (XC7A200T).} It is the cheapest board found
in stock in Singapore that holds the whole SoC, GPU included, with room
for growth.

\begin{itemize}
\tightlist
\item
  \textbf{vs Genesys 2:} more headroom and a faster speed grade, but
  \textasciitilde75 \% more expensive for margin the design does not
  need yet. Pick it only if fmax turns out to be the problem (the
  measured Artix-7 fmax, \textasciitilde74-107 MHz in Vivado, does not
  trigger that; section 5). Re-checked against the Kintex-7 measurement
  (section 5, "Measured: Kintex-7"): about 1.7-1.9x faster than Artix-7,
  which does not change this.
\item
  \textbf{vs KV260:} cheapest by far; its ARM processor could load
  firmware and act as a debug host, and its UltraRAM could hold a much
  larger SoC SRAM. But the full SoC sits at the 70 \% ceiling once it
  grows, PL pin access for the peripherals is limited, and no Singapore
  stock was found.
\item
  \textbf{Arty A7-100T and ULX3S-85F:} CPU-only boards. The Arty has no
  room to grow; the ULX3S is the only fully open toolchain option.
\item
  \textbf{Tools:} Nexys Video\textquotesingle s XC7A200T is covered by
  the free Vivado ML Standard edition (verify at download time).
\end{itemize}

\subsection{Second clock source (add-on)}\label{tm:dec:fpga-board-selection:4-second-clock-source-add-on}\origsec{4}

Testing the CDC crossings properly needs two \textbf{independent}
clocks; two outputs of one MMCM are frequency-locked and do not exercise
an asynchronous crossing (\code{FPGA_STAGE0_EVALUATION.md} (\cref{tm:dec:fpga-stage0-evaluation}) §7, decision
9.5). None of the candidates was confirmed to have a second user
oscillator, so plan on an add-on.

No Digilent-branded programmable-oscillator Pmod was found. Options:

{\footnotesize\setlength{\tabcolsep}{3pt}

{
\begin{longtable}[]{@{}
  >{\raggedright\arraybackslash}p{\dimexpr(\linewidth-6\tabcolsep)*214/1000\relax}
  >{\raggedright\arraybackslash}p{\dimexpr(\linewidth-6\tabcolsep)*255/1000\relax}
  >{\raggedright\arraybackslash}p{\dimexpr(\linewidth-6\tabcolsep)*255/1000\relax}
  >{\raggedright\arraybackslash}p{\dimexpr(\linewidth-6\tabcolsep)*255/1000\relax}@{}}
\toprule\noalign{}
\begin{minipage}[b]{\linewidth}\raggedright
Option
\end{minipage} & \begin{minipage}[b]{\linewidth}\raggedright
What
\end{minipage} & \begin{minipage}[b]{\linewidth}\raggedright
Price (snapshot)
\end{minipage} & \begin{minipage}[b]{\linewidth}\raggedright
Notes
\end{minipage} \\
\midrule\noalign{}
\endhead
\bottomrule\noalign{}
\endlastfoot
\textbf{A. Adafruit Si5351A breakout (\#2045)} (recommended) &
I2C-programmable clock generator, 3 outputs, 8 kHz to 150+ MHz, own 25
MHz crystal, 3.3 V regulator, level-shifted I2C & about £7.70-7.90 (The
Pi Hut), €12.95 (Bits \& Parts), A\$14.60 (Core Electronics); no
Singapore listing found & Not Pmod form factor: wire it to a Pmod header
(jumpers or a Pmod breakout). Its own crystal makes its outputs
independent of the board oscillator. It can be programmed by the
SoC\textquotesingle s own I2C controller or a small FPGA-side init
block. Headers need soldering on some clones. \\
B. Fixed canned oscillator on a Pmod breakout & A DIP/SMD crystal
oscillator on a Pmod prototyping/\allowbreak{}breakout board & a few
dollars & Cheapest; one fixed frequency per part; hand assembly. \\
C. Maxim DS1086LPMB1 & Pmod-format programmable oscillator board &
unknown & Listed in a 2012 Maxim collateral PDF; likely obsolete,
availability unconfirmed. \\
\end{longtable}
}

}

\textbf{Clock-capable pins matter.} A clock entering the FPGA on a
non-clock-capable pin cannot reach an MMCM/BUFG on the dedicated route;
Vivado then needs
\texttt{CLOCK\_\allowbreak{}DEDICATED\_\allowbreak{}ROUTE\ FALSE}, which
adds skew and jitter (usually acceptable for a CDC soak clock, not for a
performance clock). For Nexys Video, a Digilent forum thread
(user-reported, \textbf{unverified}) says Pmod JA is single-ended, not
clock-capable and limited to \textless10 MHz, while JB/JC go to
differential pairs. Before buying, check the Nexys Video master XDC /
schematic for which Pmod pins land on MRCC/SRCC sites. The
board\textquotesingle s FMC LPC connector carries dedicated
clock-capable \code{CLK*_M2C} pins and is the clean alternative if no
Pmod pin qualifies.

\subsection{Clock-speed sanity check: tool availability (checked
2026-10-14)}\label{tm:dec:fpga-board-selection:5-clock-speed-sanity-check-tool-availability-checked-2026-10-14}\origsec{5}

{\footnotesize\setlength{\tabcolsep}{3pt}

{
\begin{longtable}[]{@{}
  >{\raggedright\arraybackslash}p{\dimexpr(\linewidth-4\tabcolsep)*170/1000\relax}
  >{\raggedright\arraybackslash}p{\dimexpr(\linewidth-4\tabcolsep)*405/1000\relax}
  >{\raggedright\arraybackslash}p{\dimexpr(\linewidth-4\tabcolsep)*405/1000\relax}@{}}
\toprule\noalign{}
\begin{minipage}[b]{\linewidth}\raggedright
Tool
\end{minipage} & \begin{minipage}[b]{\linewidth}\raggedright
Available?
\end{minipage} & \begin{minipage}[b]{\linewidth}\raggedright
Notes
\end{minipage} \\
\midrule\noalign{}
\endhead
\bottomrule\noalign{}
\endlastfoot
Vivado & \textbf{Yes since 2026-10-15} (v2026.1,
\code{/nobackup/FPGA_SW/2026.1/Vivado}, free BASIC licence; installed
after this table was first written, the original note was: not on the
host, not in \code{/opt}, \code{/tools} or the nix store, and not in
nixpkgs (proprietary). A free Vivado ML Standard install needed an AMD
account download (\textasciitilde100 GB installed; \code{/nobackup} has
\textasciitilde600 GB free). & \\
nextpnr 0.11.1 (nixpkgs) & Yes, via
\texttt{nix\ build\ nixpkgs\#nextpnr} & Architectures: \code{ecp5},
\code{ice40}, \code{nexus}, \code{generic}, himbaechel \code{gowin} /
\code{ng-ultra} / \code{gatemate}. \textbf{No Xilinx 7-series}
(nextpnr-xilinx / prjxray are not in nixpkgs). \\
trellis (ECP5 database) & Yes, nixpkgs \code{trellis} & Needed by
\code{nextpnr-ecp5}. \\
openFPGALoader 1.1.1 & Yes, nixpkgs & Board programming only. \\
yosys & 0.69 in nixpkgs; Stage 0 used 0.62 from the nix store & Neither
is in the project flake (\code{FPGA_STAGE0_EVALUATION.md} (\cref{tm:dec:fpga-stage0-evaluation}) §1). \\
\end{longtable}
}

}

Consequences:

\begin{itemize}
\tightlist
\item
  A \textbf{real post-route fmax for the recommended Artix-7} needs
  Vivado.
\item
  With open tools available today, the only post-route timing is
  \textbf{nextpnr-ecp5 on an LFE5U-85F}. ECP5 is a slower family than
  Artix-7, so that figure is a lower bound for Artix-7, and only the
  CPU-only build fits.
\end{itemize}

\subsubsection{Measured: ECP5 lower bound (2026-10-14, CPU-only build,
post place-and-route)}\label{tm:dec:fpga-board-selection:measured-ecp5-lower-bound-2026-10-14-cpu-only-build-post-place-and-route}

\textbf{Read this as a LOWER BOUND for the recommended Artix-7
(XC7A200T), not an estimate of it.} ECP5 is a slower family than
Artix-7. It was run because it is the only post-route timing the open
tools can give today (decision on bead \bead{hlrm}: "ECP5 lower bound
now, Vivado later"). No Xilinx timing exists yet.

\textbf{Result (nextpnr-ecp5, LFE5U-85F-6BG381C, CPU-only SoC,
\texttt{-\allowbreak{}-\allowbreak{}freq\ 60}, 3 seeds, post-route "Max
frequency"):}

{
\begin{longtable}[]{@{}
  >{\raggedright\arraybackslash}p{\dimexpr(\linewidth-8\tabcolsep)*573/1000\relax}
  >{\raggedright\arraybackslash}p{\dimexpr(\linewidth-8\tabcolsep)*74/1000\relax}
  >{\raggedright\arraybackslash}p{\dimexpr(\linewidth-8\tabcolsep)*74/1000\relax}
  >{\raggedright\arraybackslash}p{\dimexpr(\linewidth-8\tabcolsep)*74/1000\relax}
  >{\raggedright\arraybackslash}p{\dimexpr(\linewidth-8\tabcolsep)*185/1000\relax}@{}}
\toprule\noalign{}
\begin{minipage}[b]{\linewidth}\raggedright
Clock domain (what it clocks)
\end{minipage} & \begin{minipage}[b]{\linewidth}\raggedright
Seed 1
\end{minipage} & \begin{minipage}[b]{\linewidth}\raggedright
Seed 2
\end{minipage} & \begin{minipage}[b]{\linewidth}\raggedright
Seed 3
\end{minipage} & \begin{minipage}[b]{\linewidth}\raggedright
Min .. max
\end{minipage} \\
\midrule\noalign{}
\endhead
\bottomrule\noalign{}
\endlastfoot
\code{cpu_gated_clk} (CPU core, caches, CPU-side CDC; ASIC target 1282
MHz) & 47.95 & 51.66 & 48.33 & \textbf{47.95 .. 51.66 MHz} \\
\code{core_clk} (bus, peripherals, CRYPTO, NPU, DMA, SRAM; ASIC target
571 MHz) & 47.35 & 52.88 & 55.22 & \textbf{47.35 .. 55.22 MHz} \\
\code{cpu_core_clk} (ungated CPU-domain sync/\allowbreak{}reset logic,
174 FFs) & 132.03 & 148.06 & 122.80 & 122.80 .. 148.06 \\
\code{tck_i} (JTAG/\allowbreak{}DFT) & 101.32 & 122.22 & 109.05 & 101.32
.. 122.22 \\
\code{clk_i} (harness shell, 290 FFs) & 229.25 & 223.56 & 231.86 &
(harness, not SoC) \\
\end{longtable}
}

Both SoC domains \textbf{failed} the 60 MHz target in all three seeds,
so the two numbers that matter are the achieved Fmax of the two big
domains: about \textbf{48-52 MHz (CPU)} and \textbf{47-55 MHz (fabric)}.
Seed spread is about +-5 MHz. Only one target (60 MHz) was run for the
three seeds; the placer is timing-driven toward the target, so a
different target would move the numbers (not measured). Against the ASIC
clocks this is a prototype scale factor of roughly 25x (CPU, 1282 MHz)
and 10-12x (fabric, 571 MHz). It sits in the upper half of the Stage 0
depth-proxy bracket (ECP5 no-GPU, 38-61 MHz;
\code{FPGA_STAGE0_EVALUATION.md} (\cref{tm:dec:fpga-stage0-evaluation}) 2.5), which that proxy had labelled
"order of magnitude only".

\textbf{Critical paths (post-route, copied from the nextpnr reports;
module from the sv2v source line of the start flop):}

{\footnotesize\setlength{\tabcolsep}{3pt}

{
\begin{longtable}[]{@{}
  >{\raggedright\arraybackslash}p{\dimexpr(\linewidth-6\tabcolsep)*91/1000\relax}
  >{\raggedright\arraybackslash}p{\dimexpr(\linewidth-6\tabcolsep)*296/1000\relax}
  >{\raggedright\arraybackslash}p{\dimexpr(\linewidth-6\tabcolsep)*296/1000\relax}
  >{\raggedright\arraybackslash}p{\dimexpr(\linewidth-6\tabcolsep)*296/1000\relax}@{}}
\toprule\noalign{}
\begin{minipage}[b]{\linewidth}\raggedright
Domain
\end{minipage} & \begin{minipage}[b]{\linewidth}\raggedright
Path (seed 1)
\end{minipage} & \begin{minipage}[b]{\linewidth}\raggedright
Shape
\end{minipage} & \begin{minipage}[b]{\linewidth}\raggedright
Seeds 2 and 3
\end{minipage} \\
\midrule\noalign{}
\endhead
\bottomrule\noalign{}
\endlastfoot
\code{cpu_gated_clk} & \code{rv32i_core} \code{mem_trap_entry_r} flop
(hazard unit) -\textgreater{} \code{stall_ex1_ex2} -\textgreater{}
\code{flush_ex1a_ex1b} -\textgreater{} \code{ex1b_comb} -\textgreater{}
pipeline-register input mux & 11 LUT levels, 3.12 ns logic + 17.73 ns
routing (85 \% routing) & Same start flop; endpoint is the
\code{u_icache.valid_array} bit reset (12 levels, 3.78 ns logic + 15.58
/ 16.91 ns routing) \\
\code{core_clk} & \code{axil_to_apb} \code{apb_paddr} flop
-\textgreater{} APB decode -\textgreater{} UART register bank write
-\textgreater{} sync-reset pin of a UART \code{apb4_register_bank} flop
& 11 LUT levels, 3.57 ns logic + 17.56 ns routing (83 \% routing) & Seed
2: CRYPTO \code{aes128_core} S-box (\code{sub_bytes}) cone, 15 levels,
4.18 + 14.73 ns. Seed 3: I2C \code{apb4_register_bank} clock-divider
-\textgreater{} \code{i2c_bit_engine.tick_q}, 22 levels (carry chain),
7.82 + 10.29 ns \\
\end{longtable}
}

}

Reading: the CPU limiter is the \textbf{trap / stall / flush control
cone} (the high-fanout control class the ASIC work also fought), not the
ALU, forwarding mux, cache tag compare or NPU MAC. The fabric limiter is
\textbf{not one path}: three seeds gave three different modules (UART
bank, AES, I2C), all within \textasciitilde8 MHz of each other, i.e. a
flat wall of many near-critical paths. Both are \textasciitilde85 \%
wire delay on a design this size, so seed noise is large and a vendor
tool\textquotesingle s placer may do better or worse.

\textbf{Utilisation (identical in every seed, \code{85k} device):}
TRELLIS\_\allowbreak{}COMB 39,004 / 83,640 (\textbf{46 \%}; 34,290 logic
LUT4 + 3,582 carry + 272 RAM + 136 RAMW), TRELLIS\_\allowbreak{}FF
14,639 (\textbf{17 \%}), DP16KD 17 / 208 (8 \%), MULT18X18D 19 / 156 (12
\%), DCCA (global clock buffers) 5 / 56, TRELLIS\_\allowbreak{}IO 6 /
365. Routing completed in every run (no unrouted nets);
nextpnr\textquotesingle s router1 reports no congestion map, so "no
congestion problem" is inferred from completion, not measured. The
harness adds +118 LUT4 and +290 FF to the Stage 0 count (34,172 /
14,349).

\textbf{Method:} \src{tools/fpga/stage0_pnr_ecp5.sh} (sv2v with
\texttt{-\allowbreak{}-\allowbreak{}define=\allowbreak{}FPGA\ -\allowbreak{}-\allowbreak{}define=\allowbreak{}\_\allowbreak{}\_\allowbreak{}pnr\_\allowbreak{}\_\allowbreak{}}
-\textgreater{} yosys 0.62
\texttt{synth\_\allowbreak{}ecp5\ -\allowbreak{}flatten\ -\allowbreak{}run\ :\allowbreak{}check},
the Stage 0 settings, now with iopad insertion -\textgreater{}
\texttt{nextpnr-\allowbreak{}ecp5\ -\allowbreak{}-\allowbreak{}85k\ -\allowbreak{}-\allowbreak{}package\ CABGA381\ -\allowbreak{}-\allowbreak{}speed\ 6\ -\allowbreak{}-\allowbreak{}freq\ 60\ -\allowbreak{}-\allowbreak{}seed\ N\ -\allowbreak{}-\allowbreak{}timing-\allowbreak{}allow-\allowbreak{}fail\ -\allowbreak{}-\allowbreak{}ignore-\allowbreak{}loops}).
Summaries via \src{tools/fpga/stage0_pnr_report.py}. Outputs in
\code{/nobackup/fpga_ecp5_fmax/} (not in the repo). Versions, all from
nix: nextpnr 0.11.1
\code{/nix/store/z9dl8x2hkcnwa71b9fwjsqxdhgx7q4xc-nextpnr-0.11.1},
prjtrellis \code{trellis-unstable-2025-01-30}
\code{/nix/store/y7jvwg7amb7nk951clb1sl4k055wrnfc-trellis-unstable-2025-01-30}
(database \code{lgnv5k4789pidnpsw9pwsf9vv84k3066}), yosys 0.62 and sv2v
0.0.13.1 as in Stage 0. Nothing installed outside nix; \src{flake.nix}
untouched. Cost: yosys 3:29 / 1.7 GB; nextpnr 8:25 (seed 1, quiet host),
42:14 (seed 2) and 19:07 (seed 3) with other jobs loading the host (load
average \textasciitilde13), peak RSS 1.33-1.35 GB each.

\textbf{What the harness cut} (\src{tools/fpga/ecp5_timing_harness.sv};
the core RTL is unchanged): \code{soc_top} has 292 port bits (103 in,
189 out) against 205 user IO on CABGA381, so it cannot be pinned. The 99
non-clock inputs come from one serial shift register on a single pin,
the 189 outputs are registered and XOR-folded into one pin (so no output
cone is optimised away); \code{clk_i}, \code{cpu_clk_i},
\code{scan_clk_i}, \code{tck_i} stay real pins. \textbf{Not timed:}
every pad-level path (input to register, register to output), IO
standards, banks, pin placement (nextpnr placed the six IO freely, no
LPF). The shells add \textasciitilde300 FFs (2 \%).

\textbf{How the gated clock and the clock tree were treated (caveats
that matter):}

\begin{itemize}
\tightlist
\item
  Nothing here is a real clock tree. \code{core_clk} is \code{clk_i}
  through the \code{dft_clk_mux} \textbf{LUT}, \code{cpu_core_clk} is
  \code{cpu_clk_i} through a LUT, and \code{cpu_gated_clk} is the
  behavioural latch + AND \code{rv32i_clock_gate} (Stage 0 option D,
  unchanged). nextpnr promoted all five clock nets to global clocks (5
  DCCA) and timed each as its own domain; it derives the LUT delay (0.24
  ns) into the clock-to-clock relation. It does \textbf{not} time the
  latch enable path.
\item
  yosys maps the four latches (\code{u_cpu_cg}, \code{u_tck_cg} clock
  gates; \code{dft_clk_switch} \code{en_a_q} / \code{en_b_q}) as LUT
  feedback, i.e. four combinational loops. \textbf{nextpnr refuses to
  run without \code{--ignore-loops}} ("Timing analysis failed due to
  combinational loops"); the loops are therefore cut, not analysed. A
  LUT-feedback latch as a clock gate is unproven in hardware; the real
  prototype needs one of the §6 gating options of
  \code{FPGA_STAGE0_EVALUATION.md} (\cref{tm:dec:fpga-stage0-evaluation}).
\item
  \texttt{-\allowbreak{}-\allowbreak{}freq\ 60} applied to
  \textbf{every} domain; the ASIC\textquotesingle s different per-domain
  targets are not expressed.
\item
  \textbf{Cross-domain paths are timed as synchronous} paths at the 60
  MHz relation. The \#94 SDC
  (\texttt{set\_\allowbreak{}clock\_\allowbreak{}groups\ -\allowbreak{}asynchronous},
  \texttt{set\_\allowbreak{}max\_\allowbreak{}delay\ -\allowbreak{}datapath\_\allowbreak{}only})
  is \textbf{not translated}; the CDC FIFOs and bridges are timed like
  ordinary logic. The worst cross-domain figures nextpnr printed (44.99
  / 47.49 / 51.89 MHz in seeds 1 / 2 / 3: \code{tck_i} falling edge
  -\textgreater{} \code{core_clk}, a JTAG/DFT path, and in seed 3 a
  harness shift-register flop -\textgreater{} \code{core_clk} reset
  fanout) are in the same band as the domain Fmax and are not separate
  limiters. The relations between the two main SoC domains
  (\code{core_clk} \textless-\textgreater{} \code{cpu_gated_clk})
  printed only max-delay values (seed 1: 2.77 ns and 1.84 ns), no
  failing frequency.
\item
  No hold analysis exists in nextpnr; no hold number is produced.
\end{itemize}

\textbf{With the GPU:} one attempt (seed 1, same flow and harness)
\textbf{did not place}. Packing reported 84,953 LUT4 (101 \%) and 85,909
TRELLIS\_\allowbreak{}COMB of 83,640 (\textbf{102 \%}; 67,369 logic LUT4
+ 4,888 carry + 8,464 RAM LUT4 = 80 \% of the RAM-capable sites), 19,202
FF (22 \%), 49 DP16KD, and the placer stopped with "ERROR: Unable to
find legal placement for all cells, design is probably at utilisation
limit" (about 2 min of nextpnr). So the with-GPU build does not fit an
LFE5U-85F, consistent with the Stage 0 sizing, and it has \textbf{no}
timing number. No Artix-7 conclusion follows from this (that device is
larger).

\textbf{Not verified:} any Xilinx/\allowbreak{}Vivado timing (the
Artix-7 figure is unmeasured; ECP5 -\textgreater{} Artix-7 speed-up is
not quantified here); a different \code{--freq} target or
\code{router2}; post-synthesis simulation of the netlist; a real clock
gate; pad timing; power; hold.

\subsubsection{Measured: Artix-7 XC7A200T (Vivado 2026.1)}\label{tm:dec:fpga-board-selection:measured-artix-7-xc7a200t-vivado-20261}

Real vendor result for the recommended part: \textbf{Vivado v2026.1
(lin64, Build 6511674), \code{xc7a200tsbg484-1}, free BASIC licence},
post-route, non-project batch flow. Reproduce with
\src{tools/fpga/vivado_fmax.sh} (+ \code{vivado_fmax.tcl},
\code{vivado_fmax_constraints.tcl}, \code{vivado_fmax_report.py});
\code{VIVADO_SETTINGS} must point at the install\textquotesingle s
\code{settings64.sh}. Outputs live in \code{/nobackup/fpga_vivado_fmax/}
(not in the repo). The Genesys 2 comparison is the next subsection (bead
\bead{hlrm.1}).

\textbf{Result, CPU-only build (no GPU; the Sky130 tie-off replaces
\src{rtl/gpu}), clock targets \code{clk_i} 10 ns (100 MHz) and
\code{cpu_clk_i} 7 ns (143 MHz), one run:}

\begin{landscape}
\footnotesize\setlength{\tabcolsep}{3pt}

{
\begin{longtable}[]{@{}
  >{\raggedright\arraybackslash}p{\dimexpr(\linewidth-12\tabcolsep)*260/1000\relax}
  >{\raggedright\arraybackslash}p{\dimexpr(\linewidth-12\tabcolsep)*48/1000\relax}
  >{\raggedright\arraybackslash}p{\dimexpr(\linewidth-12\tabcolsep)*122/1000\relax}
  >{\raggedright\arraybackslash}p{\dimexpr(\linewidth-12\tabcolsep)*113/1000\relax}
  >{\raggedright\arraybackslash}p{\dimexpr(\linewidth-12\tabcolsep)*50/1000\relax}
  >{\raggedright\arraybackslash}p{\dimexpr(\linewidth-12\tabcolsep)*126/1000\relax}
  >{\raggedright\arraybackslash}p{\dimexpr(\linewidth-12\tabcolsep)*260/1000\relax}@{}}
\toprule\noalign{}
\begin{minipage}[b]{\linewidth}\raggedright
Domain (ASIC target)
\end{minipage} & \begin{minipage}[b]{\linewidth}\raggedright
Target
\end{minipage} & \begin{minipage}[b]{\linewidth}\raggedright
WNS (setup)
\end{minipage} & \begin{minipage}[b]{\linewidth}\raggedright
TNS / failing endpoints
\end{minipage} & \begin{minipage}[b]{\linewidth}\raggedright
WHS (hold)
\end{minipage} & \begin{minipage}[b]{\linewidth}\raggedright
Estimated fmax = 1/(T - WNS)
\end{minipage} & \begin{minipage}[b]{\linewidth}\raggedright
Limiting path
\end{minipage} \\
\midrule\noalign{}
\endhead
\bottomrule\noalign{}
\endlastfoot
\code{clk_i} = \code{core_clk} fabric: bus, peripherals, CRYPTO, NPU,
DMA, SRAM (571 MHz) & 10.0 ns & \textbf{-2.085 ns} & -1152 ns / 1824 of
23,568 & +0.058 ns & \textbf{\textasciitilde83 MHz} & NPU MAC array: a
PE weight register -\textgreater{} \code{acc_q} of column 2; 19 levels
(12 CARRY4, 4 LUT2, 2 LUT4, 1 LUT3), 6.67 ns logic + 5.40 ns route \\
\code{cpu_clk_i} = gated CPU core (1282 MHz) & 7.0 ns & \textbf{-2.332
ns} & -3867 ns / 4839 of 13,583 & +0.117 ns & \textbf{\textasciitilde107
MHz} & \code{rv32i_core} EX2 \code{ex_mem.pc} -\textgreater{} EX1a/EX1b
\code{alu_result} (ALU carry chain); second family: the same flop
-\textgreater{} \code{u_csr} \code{mcycle}/\code{minstret} clock
enables; 14 levels (6-7 CARRY4, 5 LUT6), 3.77 ns logic + 5.51 ns
route \\
\code{tck_i} (JTAG TAP) & 10.0 ns & +0.000 ns (latch borrowing) & 0 &
+0.160 ns & not limited (\textgreater= 100 MHz) & TAP \code{serial_o}
-\textgreater{} clock-gate / clock-switch latch D \\
\end{longtable}
}

\end{landscape}

Read as: \textbf{Artix-7 -1, CPU-only SoC: about 83 MHz fabric and about
107 MHz CPU after a full place-and-route}, against about 48-55 MHz on
the ECP5 lower bound (so the lower bound was about 1.7x pessimistic for
the fabric and about 2x for the CPU) and the Stage 0 depth-proxy bracket
of 58-99 MHz (\code{FPGA_STAGE0_EVALUATION.md} (\cref{tm:dec:fpga-stage0-evaluation}) 2.5; the measured fabric
number sits inside it, the CPU one just above it). Prototype scale
factor vs the ASIC clocks: about 7x (fabric, 571 -\textgreater{}
\textasciitilde83 MHz) and about 12x (CPU, 1282 -\textgreater{}
\textasciitilde107 MHz). \textbf{Both domains missed their target, so
the estimate is what the router achieved while still trying; a
\texttt{WNS\ >=\allowbreak{}\ 0} domain would only have given a lower
bound} (in an earlier 10 ns/10 ns run the CPU domain met timing at
+0.033 ns, i.e. it stopped optimising at 100 MHz; the 7 ns target is
what exposed \textasciitilde107 MHz). Only one target per domain was run
in the reported configuration, so the sensitivity of the estimate to the
target is not measured. Single run, no seed sweep.

\textbf{Cross-domain crossings (reported separately, bounded by the
translated CDC budgets, see below):} \code{clk_i} -\textgreater{}
\code{cpu_clk_i} WNS -1.444 ns (153 of 1258 endpoints fail) and
\code{cpu_clk_i} -\textgreater{} \code{clk_i} WNS -1.052 ns (15 of 506),
both against a
\texttt{set\_\allowbreak{}max\_\allowbreak{}delay\ -\allowbreak{}datapath\_\allowbreak{}only\ 0.\allowbreak{}9\ x\ destination\ period}
bound (6.3 ns and 9.0 ns). This is the ASIC check-SDC budget applied
literally; it is tighter than a plain asynchronous cut, and with a 7 ns
CPU clock the budget for crossings into the CPU domain is only 6.3 ns.
Not separate limiters of the two domain numbers; they are listed so
nothing is hidden.

\textbf{Utilisation (post-route, \code{xc7a200tsbg484-1}):} Slice LUTs
18,981 of 133,800 available (\textbf{14.2 \%}; 18,696 logic + 285
LUT-RAM), Slice registers 13,206 of 267,600 (4.9 \%, 4 of them latches),
F7 / F8 muxes 640 / 130, block RAM tiles 12 of 365 (3.3 \%: 1 RAMB36 +
22 RAMB18), DSP48E1 2 of 740 (0.3 \%), bonded IOB 6 of 285, BUFGCTRL 6
of 32. (The post-route report lists 133,800 LUTs available against
134,600 in the part query and the synthesis report; the reason was not
investigated.) \textbf{Versus yosys Stage 0} (CPU-only, 38,554 LUT1-6,
14,079 FF, 9 RAMB36-equivalents, 19 DSP48E1): Vivado used \textbf{about
0.49x the LUTs} (the "10-30 \% fewer" assumption in section 1 was too
conservative: it is \textasciitilde51 \% fewer here), about the same
flops, 12 vs \textasciitilde9 block-RAM tiles and 2 vs 19 DSPs (the 4x4
INT8 NPU multipliers went to LUTs and CARRY4 chains, not DSP48E1). All
blocks are present in the hierarchy report; the harness was checked and
swept nothing away ("Trimming check" below). This does not move the
section 1 sizing for the full SoC (that depends on the GPU, see below),
but it removes the "38.6 K LUTs, 70 \% rule" worry for the CPU-only
build on this part: it uses 14 \%.

\textbf{Method (what was and was not done):}

\begin{itemize}
\tightlist
\item
  Vivado\textquotesingle s own SystemVerilog frontend read the RTL
  directly (\texttt{read\_\allowbreak{}verilog\ -\allowbreak{}sv}, one
  call per file,
  \texttt{synth\_\allowbreak{}design\ -\allowbreak{}verilog\_\allowbreak{}define\ FPGA\ -\allowbreak{}verilog\_\allowbreak{}define\ \_\allowbreak{}\_\allowbreak{}pnr\_\allowbreak{}\_\allowbreak{}},
  default \texttt{flatten\_\allowbreak{}hierarchy\ rebuilt}); the sv2v
  output was \textbf{not} needed. File list = the Stage 0 list. One
  deviation, no committed RTL change: Vivado 2026.1 rejects the
  \texttt{input\ logic\ x} ports of the Sky130 GPU tie-off
  \src{pnr/sky130/soc/gpu_top_tieoff.sv} under its own
  \texttt{`default\_\allowbreak{}nettype\ none} ({[}Synth 8-6735{]},
  Verilator / sv2v / yosys accept it), so the flow reads a copy of that
  one file with the directive line removed. The tie-off only exists in
  the no-GPU build.
\item
  IO: the existing \src{tools/fpga/ecp5_timing_harness.sv} is
  vendor-neutral and was used \textbf{unchanged} (4 clock pins,
  \code{din}, \code{dout}; 99 inputs from one shift register, 189
  outputs folded into one pin). In-context synthesis, \code{opt_design},
  \code{place_design}, \code{phys_opt_design}, \code{route_design},
  default directives, 4 threads. Pad-level paths, IO standards and pin
  placement are \textbf{not} timed (Vivado placed the 6 pins itself, no
  XDC pin constraints). The harness shells add \textasciitilde300 FFs.
\item
  Clocks: \code{create_clock} on \code{clk_i}, \code{cpu_clk_i},
  \code{scan_clk_i}, \code{tck_i}; no generated clocks (Vivado
  propagates the port clock through the DFT clock-mux LUT and the
  clock-gate AND).
\item
  Translated from
  \texttt{phase5\_\allowbreak{}soc\_\allowbreak{}multiclock\{,\allowbreak{}\_\allowbreak{}check\}.\allowbreak{}sdc}
  (the run prints one
  \texttt{VFMAX-\allowbreak{}CDC\ TRANSLATED\ /\allowbreak{}\ CUT\ /\allowbreak{}\ MISS}
  line per exception; \code{report_exceptions} is saved): domain-wide
  \texttt{set\_\allowbreak{}max\_\allowbreak{}delay\ -\allowbreak{}datapath\_\allowbreak{}only\ 0.\allowbreak{}9\ x\ destination\ period}
  both ways (Vivado form of the check SDC;
  \texttt{set\_\allowbreak{}clock\_\allowbreak{}groups\ -\allowbreak{}asynchronous}
  would override and cut these); the 0.5 x destination-period
  gray-pointer budgets on all 10 pointer crossings of the five
  \code{u_cpu_axi_cdc} channel FIFOs (directions as in the check SDC:
  aw/w/ar written in the CPU domain, b/r the other way); the
  \code{u_apb_pll2_cdc} / \code{u_apb_dbg_cdc} \code{req_q} /
  \code{ack_q} level crossings and \code{u_cpu_clk_dis_sync} (all bound,
  1-2 flops each); \code{set_false_path} to the 18 async set/clear pins
  of the \code{cdc_reset_sync} flops; \code{scan_clk_i} physically
  exclusive with the two functional clocks.
\item
  \textbf{Functional mode} (needed, and found the hard way): the SDC
  resolves the DFT clock muxes with \code{set_case_analysis} on the
  \code{scan_mode_i} / \code{scan_en_i} ports. In the harness those are
  flops and the mux is a LUT, so in a first run \code{tck_i} and
  \code{scan_clk_i} propagated into \code{core_clk} and timed every
  functional flop three times (a bogus \code{tck_i} WNS of -2.85 ns came
  from a TAP falling-edge flop into a CRYPTO register). Fix:
  \texttt{set\_\allowbreak{}clock\_\allowbreak{}sense\ -\allowbreak{}stop\_\allowbreak{}propagation}
  for \code{tck_i} and \code{scan_clk_i} at the one LUT that merges them
  into \code{core_clk}; \code{tck_i} stays async to the functional
  clocks, so TCK -\textgreater{} functional crossings are \textbf{cut}.
\item
  Harness-shell cuts, so no fake path is timed: \code{in_sr} flops that
  model SDC-false-pathed ports (\code{rst_n_i}, \code{cpu_rst_n_i},
  \code{apb_*}, \code{uart_rx_i}, \code{spi_miso_i}, the DFT inputs) and
  the \code{in_sr} -\textgreater{} CPU-domain and CPU-domain
  -\textgreater{} output-shell crossings are \code{set_false_path}, as
  the SDC false-paths the port equivalents. \code{gpio_in_i} /
  \code{i2c_*_i} shell flops keep their \code{clk_i} launch.
\item
  \textbf{Not translated:} port \code{set_input_delay} /
  \code{set_output_delay}, the SRAM \code{-hold} false paths, the GPU /
  CPU-APB \code{set_multicycle_path} (they exist for ASIC hard-macro
  boundaries; here CPU and GPU are plain RTL), clock uncertainty /
  latency / transition (the ASIC clock-tree model; Vivado computes its
  own), \code{set_case_analysis} itself (replaced by the clock-sense
  stop above), and any hold exception.
\end{itemize}

\textbf{Clock gate and DFT latches (kept behavioural, as Stage 0 option
D):}

\begin{itemize}
\tightlist
\item
  Vivado inferred 4 latches (\code{LDCE}, 3 with inverted \code{G}): the
  \code{u_cpu_cg} and \code{u_tck_cg} clock-gate enables and the two
  \code{dft_clk_switch} \code{en_a_q} / \code{en_b_q} (the third GPU
  gate does not exist in the CPU-only build). No combinational loop
  exists (\code{check_timing} loops = 0, latch loops = 0), unlike the
  four loops yosys/\allowbreak{}nextpnr had to cut with
  \code{--ignore-loops}.
\item
  Every LUT-driven clock net was \textbf{promoted to a global buffer}:
  \code{core_clk} (the DFT clock-mux LUT) -\textgreater{} one BUFG with
  8,071 clock pins, \code{cpu_gated_clk} (latch + AND, as a LUT)
  -\textgreater{} one BUFG with 4,808 pins, \code{cpu_core_clk}
  -\textgreater{} one BUFG with 58 pins, plus the three port clocks; 6
  BUFGCTRL of 32 in total. So the gated clock is \textbf{not} a
  fabric-routed net: it is a LUT feeding a BUFG, the LUT sits in the
  clock path (extra insertion delay, included in
  Vivado\textquotesingle s skew numbers; measured worst-path clock skew
  on the CPU critical path is -0.05 ns), and the gate is \textbf{not} a
  BUFGCE/\allowbreak{}glitch-free clock enable in hardware.
\item
  \code{report_methodology}: 5 x \textbf{TIMING-14 "LUT on the clock
  tree" (Critical Warning)} (the DFT muxes and the gate), 17 x LUTAR-1
  (LUT drives an async reset), 13 x SYNTH-6, TIMING-9 / TIMING-10 once
  each (unknown CDC logic / missing synchroniser property), LATCH-1
  advisory. They are warnings about the structure, not timing failures.
\item
  The enable latch \textbf{is timed}: \code{u_cpu_cg/en_latch_reg/D} is
  a \code{cpu_clk_i} fall-edge latch endpoint with 0.000 ns slack
  \textbf{by time borrowing} (2.058 ns borrowed of a 3.5 ns half period,
  source \code{u_cpu_clk_dis_sync} through one LUT4). It did not limit
  either domain.
\item
  The \emph{effect of gating on the reported numbers} was \textbf{not
  isolated} (no run with the gate removed, which would need an RTL arm;
  that is the open clock-gating decision). The evidence for "small" is
  that the CPU critical path starts and ends on gated flops with skew of
  -0.05 ns and positive hold slack (+0.117 ns); the evidence against
  being fully clean is TIMING-14.
\end{itemize}

\textbf{Cost (CPU-only):} synthesis 4:07 wall / 2.49 GB peak RSS; place
and route 10:06 wall / 2.73 GB; both well inside the host limits.

\textbf{With the GPU (full SoC), same flow, same targets (\code{clk_i}
10 ns, \code{cpu_clk_i} 7 ns), one run: it fits, comfortably.}

{\small\setlength{\tabcolsep}{4pt}

{
\begin{longtable}[]{@{}
  >{\raggedright\arraybackslash}p{\dimexpr(\linewidth-4\tabcolsep)*425/1000\relax}
  >{\raggedright\arraybackslash}p{\dimexpr(\linewidth-4\tabcolsep)*158/1000\relax}
  >{\raggedright\arraybackslash}p{\dimexpr(\linewidth-4\tabcolsep)*398/1000\relax}@{}}
\toprule\noalign{}
\begin{minipage}[b]{\linewidth}\raggedright
\end{minipage} & \begin{minipage}[b]{\linewidth}\raggedright
CPU-only
\end{minipage} & \begin{minipage}[b]{\linewidth}\raggedright
Full SoC with GPU
\end{minipage} \\
\midrule\noalign{}
\endhead
\bottomrule\noalign{}
\endlastfoot
\code{clk_i} fabric WNS / estimated fmax & -2.085 ns / \textasciitilde83
MHz & \textbf{-3.452 ns / \textasciitilde74 MHz} (4527 of 56,511
endpoints fail) \\
\code{cpu_clk_i} CPU WNS / estimated fmax & -2.332 ns /
\textasciitilde107 MHz & \textbf{-3.183 ns / \textasciitilde98 MHz}
(6267 of 13,588 fail) \\
Hold WHS (worst) & +0.058 ns & \textbf{-0.112 ns, 8 endpoints} (see
below) \\
\code{clk_i} -\textgreater{} \code{cpu_clk_i} / \code{cpu_clk_i}
-\textgreater{} \code{clk_i} crossings WNS & -1.444 / -1.052 ns & -1.454
/ -1.229 ns \\
Slice LUTs (post-route) & 18,981 (14.2 \%) & \textbf{40,586 of 133,800
(30.3 \%)}: 37,485 logic + 3,101 LUT-RAM \\
Slice registers & 13,206 (4.9 \%) & 17,988 (6.7 \%), 5 latches \\
Block RAM tiles & 12 (3.3 \%) & 28 of 365 (7.7 \%): 1 RAMB36 + 54
RAMB18 \\
DSP48E1 & 2 (0.3 \%) & 26 of 740 (3.5 \%) \\
BUFGCTRL & 6 & 7 of 32 \\
Wall time / peak RSS (synthesis + P\&R in one session for the GPU build)
& 4:07 + 10:06 / 2.7 GB & 21:23 / 3.28 GB \\
\end{longtable}
}

}

\begin{itemize}
\tightlist
\item
  \textbf{GPU fit verdict on XC7A200T: fits with a large margin.} 30 \%
  of the LUTs, 8 \% of the block RAM, 4 \% of the DSPs. That is half of
  the "\textasciitilde62 \%" of section 2, which came from yosys LUT
  counts plus 25 \% growth and bring-up margin; today\textquotesingle s
  measured SoC is 30 \%, so even after +25 \% growth and +5 K bring-up
  logic (both assumed, not measured) the estimate is about 42 \%.
\item
  \textbf{GPU shared memory went to block RAM in Vivado}:
  \code{shared_memory} is 32 RAMB18 (plus 9.6 K LUTs of mux and control
  around them); yosys had chosen LUT-RAM (\textasciitilde5 K LUT sites)
  for the same RTL. The GPU vector register file went to LUT-RAM (2,816
  LUT-RAM sites in \code{gpu_top}). So the Stage 1 "GPU shared-memory
  RAM style" decision is not forced on this part: Vivado already picks
  block RAM.
\item
  \textbf{Timing with the GPU is worse, and the CPU number moved even
  though the CPU RTL did not.} The fabric limiter is a GPU
  warp-scheduler \code{warp_done} flop -\textgreater{} DMA
  \code{state_q} clock enable (16 levels, 2.79 ns logic + 9.51 ns route
  = 77 \% wire, clock skew -0.91 ns); the CPU limiter in this run is EX2
  \code{ex_mem.mem_size} -\textgreater{} I-cache \code{refill_buf} clock
  enable (12 levels, 2.0 ns logic + 7.8 ns route). The CPU domain gained
  a worse WNS (-2.33 -\textgreater{} -3.18 ns, \textasciitilde107
  -\textgreater{} \textasciitilde98 MHz) with identical CPU logic, so
  **the placement-to-placement spread of a single run is of the order of
  1 ns (\textasciitilde10 \%)**; the honest reading is "fabric
  \textasciitilde74-83 MHz, CPU \textasciitilde98-107 MHz" rather than
  any single figure, and one run per build is all that was done (no seed
  or target sweep).
\item
  \textbf{Hold:} worst -0.112 ns on 8 endpoints (THS -0.49 ns). The
  worst one is a harness shell flop \code{in_sr[82]}
  (\code{gpio_in_i[31]}) -\textgreater{} the GPIO input synchroniser,
  with 1.52 ns of clock skew: the shell flop sits behind a plain
  \code{clk_i} BUFG while the synchroniser sits behind the
  \code{core_clk} LUT + BUFG (extra insertion delay). The
  SDC\textquotesingle s
  \texttt{set\_\allowbreak{}false\_\allowbreak{}path\ -\allowbreak{}hold\ -\allowbreak{}from\ [all\_\allowbreak{}inputs]}
  was \textbf{not} translated, so this is a harness / pad-hold artifact;
  the other seven endpoints were not individually inspected. It does
  show a real structural point: a LUT in the clock path adds about 1.5
  ns of insertion delay against a direct BUFG.
\item
  \textbf{Clock gating / latches:} 5 \code{LDCE} (the CPU, GPU and TCK
  clock-gate enables and the two clock-switch latches); every LUT-driven
  clock net again became a BUFG (the GPU gated clock is its own BUFG
  with 5,462 pins); 8 x TIMING-14 (Critical Warning) instead of 5; no
  combinational loops.
\end{itemize}

\textbf{Trimming check (is Vivado\textquotesingle s 2x-smaller LUT count
and 2-vs-19 DSP real, or did logic disappear?) - verdict: no material
trimming; the numbers are valid for the design as written.} Evidence,
all measured unless marked inference:

\begin{enumerate}
\def\labelenumi{\arabic{enumi}.}
\tightlist
\item
  \textbf{Harness did not trim anything.} \code{soc_top} synthesised
  alone, out of context
  (\texttt{vivado\_\allowbreak{}fmax.\allowbreak{}sh\ .\allowbreak{}.\allowbreak{}.\allowbreak{}\ ooc}:
  no harness, every port kept, nothing tied off) gives CPU-only
  \textbf{18,687 LUT / 12,957 FF / 12 BRAM tiles / 2 DSP} and with GPU
  \textbf{40,017 LUT / 17,763 FF / 28 BRAM tiles / 26 DSP}. The
  in-context runs (18,700 / 13,239 and 40,038 / 18,046) differ by
  exactly the harness shells (+282 / +283 FF, +13 / +21 LUT). The 6-pin
  harness neither created constants nor orphaned outputs.
\item
  \textbf{Removal messages are small and behaviour-preserving.}
  Synthesis log, CPU-only: 37 flops "unused, removed" (31 are the high
  bits of the NPU requantiser product register, absorbed into the
  DSP48E1 P register, per the \texttt{DSP\ Report} lines), 7 unused CSR
  registers (\code{rv32i_csr_file}: two trimmed from 32 to 4 bits), 92
  flops with a constant-0 input propagated (59 in \code{u_crypto}), and
  \textgreater=100 duplicate registers merged (Synth 8-3886 hit its
  100-message print limit; the printed ones are pipeline duplicates such
  as \code{id_ex_reg.rs2_addr[4:0]} == \code{instruction[24:20]}).
  \code{opt_design} then removed \textbf{192 cells in total} and
  propagated no constants. Nothing block-sized is removed.
\item
  \textbf{Every block is present with its state.} Flip-flops per block,
  Vivado (post-synthesis hierarchy) vs yosys (the same sv2v netlist,
  non-flattened \code{synth_xilinx}, one instance per module assumed):
  CPU 4,717 vs 6,284 (-25 \%, the merged duplicates), CRYPTO 2,799 vs
  3,007, NPU 1,076 vs 1,011, DMA 801 vs 545, I2C 374 vs 620, GPIO 288 vs
  290, PWM 260 vs 288, SPI 205 vs 323 (other peripherals not tabulated);
  SoC total 12,957 vs 14,079 (-8 \%). A trimmed block would have lost
  its flops; none did. (Per-block \textbf{LUT} counts cannot be compared
  the same way: with \texttt{flatten\_\allowbreak{}hierarchy\ rebuilt}
  Vivado attributes LUTs by pre-optimisation names, so for example the
  APB read-mux LUTs of CRYPTO appear under \code{u_axil_apb_bridge}.
  Only the totals are meaningful.)
\item
  \textbf{The DSP gap is a mapping choice.} The only multipliers Vivado
  put in DSP48E1 are the NPU requantiser product (\code{prod_d_w} = 32 x
  17 -\textgreater{} 2 cascaded DSPs:
  \texttt{DSP\ Report\ .\allowbreak{}.\allowbreak{}.\allowbreak{}\ (PCIN>>17)+\allowbreak{}A*B}).
  The 16 INT8 x INT8 PE multipliers in \code{npu_mac_array}
  (\texttt{16'(w\_\allowbreak{}q)\ *\ 16'(\$signed(a))}, i.e. 8-bit
  operands behind a sign extension) went to LUT fabric plus a 12-CARRY4
  chain, which is the fabric-fmax limiting path above and the reason the
  NPU costs 1,977 LUTs. yosys put the same 16 + 3 (requantiser)
  multipliers, plus one in \code{i2c_controller} (not traced), into 19
  DSP48E1. Vivado clearly does infer DSPs for wide multiplies: the GPU
  adds 24 DSP48E1 in Vivado (26 - 2) against 25 in yosys (44 - 19).
\item
  \textbf{Where the 2x LUT gap comes from is only partly shown.} Shown:
  S-boxes of AES went to 8 RAMB18 ROMs in Vivado while yosys
  \code{proc_rom} made them logic; duplicates merged (-1.1 K FF and
  their LUTs); wide-mux packing (640 F7

  \begin{itemize}
  \tightlist
  \item
    130 F8 vs 3,271 + 872 in yosys). Inference only, not isolated per
    cause: Vivado\textquotesingle s LUT6 packing and cross-boundary
    optimisation versus ABC\textquotesingle s \code{synth_xilinx} result
    (yosys\textquotesingle s own non-flattened run is \textbf{bigger}
    than its flattened one, 15.9 K + 5.1 K INV LUTs for the CPU alone,
    so 38.5 K is itself a flatten-dependent figure). The GPU increment
    is similar in both tools (+21.3 K LUT in Vivado, +23.8 K in yosys),
    so the gap is in the common base, not the GPU. Consequence: the
    utilisation and fmax numbers above stand. The LUT ratio should be
    read as "Vivado needs roughly half to two thirds of the LUTs yosys
    reports for this RTL", not as a bound for other RTL.
  \end{itemize}
\end{enumerate}

\textbf{What the measurement changes in this document:}

\begin{itemize}
\tightlist
\item
  Section 1\textquotesingle s "Vivado typically 10-30 \% fewer LUTs than
  yosys (inference)" is replaced by a measurement:
  \textbf{\textasciitilde51 \% fewer for CPU-only (18,981 vs 38,554) and
  \textasciitilde35 \% fewer for the full SoC (40,586 LUT + LUT-RAM vs
  62,353 logic LUT + \textasciitilde5 K LUT-RAM sites)}. The yosys-based
  sizing in sections 1 and 2 is therefore conservative by about 2x for
  LUTs, and the "fits" column above is pessimistic.
\item
  Section 2\textquotesingle s "Arty A7-100T does not fit the full SoC"
  is \textbf{no longer true on LUTs alone}: 40.6 K of 63.4 K is
  \textasciitilde64 \% today, and \textasciitilde80 \% with the 25 \%
  growth, so it fits today, with no room for bring-up logic and growth
  (plus the other limits not measured here: its pins; its 135 block RAM
  tiles and 240 DSPs would cover the 28 and 26 used here, from the
  device data sheet, not checked in Vivado). The recommendation below
  does not change, because the user wants growth headroom and debug
  logic, and the Nexys Video still has \textasciitilde2x the LUTs for a
  price that was already decided to be acceptable; it is noted as an
  alternative only if cost dominates.
\item
  Timing: on both Artix-7 builds the prototype clocks are
  \textbf{\textasciitilde74-107 MHz}, a 7-17x scale factor against the
  ASIC targets. That is fast enough for firmware validation and the CDC
  soak, and well above the ECP5 lower bound; it does \textbf{not} need
  the Genesys 2\textquotesingle s faster speed grade, so the "pick it
  only if fmax turns out to be the problem" condition of section 3 is
  not triggered by this measurement.
\end{itemize}

\textbf{NOT verified in the Artix-7 numbers:} any other speed grade or
part (Kintex-7 / Genesys 2 is measured in the next subsection); a sweep
of clock targets or seeds (one run per build, \textasciitilde1 ns spread
observed); post-route simulation; power; pad-level timing and hold
(harness shells only); \code{phys_opt_design} after routing and any
non-default directive (a longer-effort flow may do better); the effect
of removing the behavioural clock gate; the other 7 hold endpoints (not
inspected); any run on hardware.

\subsubsection{\texorpdfstring{Measured: Kintex-7 XC7K325T-2 (Vivado
2026.1), the Genesys 2 part (bead
\code{hlrm.1})}{Measured: Kintex-7 XC7K325T-2 (Vivado 2026.1), the Genesys 2 part (bead )}}\label{tm:dec:fpga-board-selection:measured-kintex-7-xc7k325t-2-vivado-20261-the-genesys-2-part-bead-hlrm1}

Same flow as the Artix-7 run above, only the part changed:
\textbf{\code{xc7k325tffg900-2}}, Vivado v2026.1, same RTL
(origin/\allowbreak{}main \code{7875da3}; no RTL change since the
Artix-7 run, only bead notes), same harness, same filelist, same
constraints script, same flow (\code{opt_design}, \code{place_design},
\code{phys_opt_design}, \code{route_design}, default directives),
\texttt{maxThreads\ 4}. \code{tools/fpga/vivado_fmax.*} needed
\textbf{no change}: the part was already the \code{VFMAX_PART} variable
and the Artix-7 default is untouched. Invocation:
\texttt{VFMAX\_\allowbreak{}PART=\allowbreak{}xc7k325tffg900-\allowbreak{}2\ VFMAX\_\allowbreak{}PERIOD\_\allowbreak{}CPU=\allowbreak{}7\ vivado\_\allowbreak{}fmax.\allowbreak{}sh\ .\allowbreak{}\ <out>\ \{nogpu|gpu\}\ .\allowbreak{}.\allowbreak{}.\allowbreak{}\ 10}
(CPU-only: \code{synth} then \code{impl}; with GPU: \code{both}, as for
Artix-7). Outputs in \code{/nobackup/fpga_vivado_fmax_k7/} (not in the
repo). Every figure below is {[}M{]} (read from the Vivado reports of
this session) unless marked {[}A{]} (assumed / inferred) or {[}R{]}
(read from another document).

\textbf{Step 1: like-for-like targets (\code{clk_i} 10 ns,
\code{cpu_clk_i} 7 ns), one run per build. Result: the Kintex-7 met
every constraint, so these are lower bounds, not achieved fmax} (the
placer stops optimising once the target is met; the Artix-7 run missed
all four, so its figures were what the router achieved while still
trying).

{\footnotesize\setlength{\tabcolsep}{3pt}

{
\begin{longtable}[]{@{}
  >{\raggedright\arraybackslash}p{\dimexpr(\linewidth-6\tabcolsep)*155/1000\relax}
  >{\raggedright\arraybackslash}p{\dimexpr(\linewidth-6\tabcolsep)*191/1000\relax}
  >{\raggedright\arraybackslash}p{\dimexpr(\linewidth-6\tabcolsep)*370/1000\relax}
  >{\raggedright\arraybackslash}p{\dimexpr(\linewidth-6\tabcolsep)*263/1000\relax}@{}}
\toprule\noalign{}
\begin{minipage}[b]{\linewidth}\raggedright
Build / domain
\end{minipage} & \begin{minipage}[b]{\linewidth}\raggedright
Artix-7 -1 WNS {[}R{]} (est. fmax)
\end{minipage} & \begin{minipage}[b]{\linewidth}\raggedright
Kintex-7 -2 WNS (lower bound 1/\allowbreak{}(T-WNS))
\end{minipage} & \begin{minipage}[b]{\linewidth}\raggedright
WNS swing
\end{minipage} \\
\midrule\noalign{}
\endhead
\bottomrule\noalign{}
\endlastfoot
CPU-only, \code{clk_i} fabric & -2.085 ns (\textasciitilde83 MHz) &
\textbf{+1.411 ns} (\textgreater= 116 MHz) & +3.50 ns \\
CPU-only, \code{cpu_clk_i} & -2.332 ns (\textasciitilde107 MHz) &
\textbf{+0.088 ns} (\textgreater= 145 MHz; the reported WNS is the
clock-gate enable latch, borrowing; the worst register-to-register path
has +0.748 ns, \textgreater= 160 MHz) & +2.42 ns (\textgreater= 3.08 ns
on the register path) \\
With GPU, \code{clk_i} fabric & -3.452 ns (\textasciitilde74 MHz) &
\textbf{+0.367 ns} (\textgreater= 104 MHz) & +3.82 ns \\
With GPU, \code{cpu_clk_i} & -3.183 ns (\textasciitilde98 MHz) &
\textbf{+0.102 ns} (\textgreater= 145 MHz, same latch endpoint) & +3.29
ns \\
\end{longtable}
}

}

\textbf{Step 2 (an addition to the brief, needed because step 1
saturated): tighter targets, \code{clk_i} 7 ns (143 MHz) and
\code{cpu_clk_i} 5 ns (200 MHz)}, implementation only on the same
synthesis checkpoints (so only the constraint changes), one run per
build. This is what measures an achieved fmax, the same way the Artix-7
7 ns CPU target did. It is a different target pair from the Artix-7 run,
so it is \textbf{not} a strict like-for-like; step 1 is the strict one.

{\footnotesize\setlength{\tabcolsep}{3pt}

{
\begin{longtable}[]{@{}
  >{\raggedright\arraybackslash}p{\dimexpr(\linewidth-8\tabcolsep)*144/1000\relax}
  >{\raggedright\arraybackslash}p{\dimexpr(\linewidth-8\tabcolsep)*49/1000\relax}
  >{\raggedright\arraybackslash}p{\dimexpr(\linewidth-8\tabcolsep)*166/1000\relax}
  >{\raggedright\arraybackslash}p{\dimexpr(\linewidth-8\tabcolsep)*277/1000\relax}
  >{\raggedright\arraybackslash}p{\dimexpr(\linewidth-8\tabcolsep)*344/1000\relax}@{}}
\toprule\noalign{}
\begin{minipage}[b]{\linewidth}\raggedright
Build / domain
\end{minipage} & \begin{minipage}[b]{\linewidth}\raggedright
Target
\end{minipage} & \begin{minipage}[b]{\linewidth}\raggedright
WNS / failing endpoints
\end{minipage} & \begin{minipage}[b]{\linewidth}\raggedright
Estimated fmax = 1/(T-WNS)
\end{minipage} & \begin{minipage}[b]{\linewidth}\raggedright
Limiting path (start -\textgreater{} end; logic)
\end{minipage} \\
\midrule\noalign{}
\endhead
\bottomrule\noalign{}
\endlastfoot
CPU-only, \code{clk_i} fabric & 7.0 ns & \textbf{+0.245 ns}, 0 of 23,444
(met) & \textbf{\textgreater= 148 MHz} (lower bound; true limit not
reached) & NPU MAC: \code{ain_q} -\textgreater{} column-3 \code{acc_q},
20 levels (13 CARRY4, 5 LUT2), 3.21 ns logic + 3.57 ns route \\
CPU-only, \code{cpu_clk_i} & 5.0 ns & \textbf{-0.566 ns}, 1807 of 13,589
& \textbf{\textasciitilde180 MHz} & \code{u_cpu_pmu_rst_sync} flop
-\textgreater{} synchronous-reset pin of \code{id_ex_reg}
(\code{rs2_data}, \code{csr_op}...), 4 levels
(LUT3/\allowbreak{}5/\allowbreak{}6/\allowbreak{}6), 0.40 ns logic +
3.71 ns route, clock skew -1.12 ns: a \textbf{reset-distribution} path,
not the ALU path that limited Artix-7 \\
With GPU, \code{clk_i} fabric & 7.0 ns & \textbf{-0.955 ns}, 269 of
56,540 & \textbf{\textasciitilde126 MHz} & GPU warp scheduler
\code{rr_ptr} -\textgreater{} \code{u_gif_adapter} \code{ar_addr_q}, 16
levels (9 LUT6, 1 MUXF7), 1.07 ns logic + 5.65 ns route, skew -1.10
ns \\
With GPU, \code{cpu_clk_i} & 5.0 ns & \textbf{-0.257 ns}, 243 of 13,607
& \textbf{\textasciitilde190 MHz} & \code{ex1c_ex1b_reg} \code{pc}
-\textgreater{} \code{u_csr} \code{mhpmcounter4} clock enable, 16 levels
(8 CARRY4, 4 LUT6), 1.34 ns logic + 3.58 ns route \\
\end{longtable}
}

}

(Step-1 critical paths, for the record: CPU-only \code{clk_i} worst is
\code{u_pll_sub} stub \code{locked_q} -\textgreater{} SHA \code{w_q}
reset pin, 2 levels, 7.4 ns of routing, +1.411 ns slack; with GPU it is
a GPU scheduler flop (\code{rr_ptr}) -\textgreater{} DMA \code{state_q}
clock-enable path, the same family that limited Artix-7 (17 levels, 1.48
ns logic + 6.98 ns route, skew -0.93 ns); \code{cpu_clk_i} is the
clock-gate enable latch in both, then \code{mem_wb_reg.pc}
-\textgreater{} \code{u_csr} \code{mcause} clock enable (15 levels, 7
CARRY4) at +0.748 ns CPU-only.)

\textbf{Prototype scale factor} (ASIC 571 MHz fabric / 1282 MHz CPU),
using the step-2 numbers: fabric \textasciitilde126 MHz with GPU (4.5x)
and

\begin{quote}
= 148 MHz CPU-only (\textless= 3.9x); CPU \textasciitilde190 MHz with
GPU (6.7x) and \textasciitilde180 MHz CPU-only (7.1x). Artix-7 was
7-17x.
\end{quote}

\textbf{Kintex-7 vs Artix-7, does the difference exceed the
\textasciitilde1 ns run-to-run spread?} Yes, on every comparison, with
the caveat that the spread was observed on Artix-7 only and is not
re-measured here (one run per build and target on Kintex-7 as well).

\begin{itemize}
\tightlist
\item
  Strict like-for-like (step 1, same 10/7 ns targets): the WNS swing is
  \textbf{+2.4 to +3.8 ns} (\textgreater= 3.1 ns on the CPU register
  path), i.e. 2.4x to 3.8x the 1 ns spread, and it flips all four
  domains from failing to meeting timing.
\item
  Achieved periods (step 2 against the Artix-7 estimates): with GPU,
  fabric 8.0 ns vs 13.5 ns (-5.5 ns, 1.7x) and CPU 5.3 ns vs 10.2 ns
  (-4.9 ns, 1.9x); CPU-only, fabric \textless= 6.8 ns vs 12.1 ns
  (\textgreater= 5.3 ns, \textgreater= 1.8x) and CPU 5.6 ns vs 9.3 ns
  (-3.8 ns, 1.7x). Because the Kintex-7 targets were tighter and tighter
  targets push the placer harder, a part of this gap may be the target,
  not the device {[}A{]}; the strict step-1 swing does not have that
  issue.
\item
  Reading: \textbf{Kintex-7 -2 is roughly 1.7x to 1.9x faster than
  Artix-7 -1 on this design} (fabric \textasciitilde126 MHz+ vs
  \textasciitilde74-83 MHz, CPU \textasciitilde180-190 MHz vs
  \textasciitilde98-107 MHz), and its fabric limiter moves around (NPU
  MAC, GPU scheduler -\textgreater{} bus) just as on Artix-7. Single
  runs, so read the ratio as 1.7-1.9x +/- the placement spread, not as a
  precise factor.
\end{itemize}

\textbf{Utilisation (post-route, \code{xc7k325tffg900-2}, step-1 runs;
step-2 within +/-2 \% of these):}

{\small\setlength{\tabcolsep}{4pt}

{
\begin{longtable}[]{@{}
  >{\raggedright\arraybackslash}p{\dimexpr(\linewidth-4\tabcolsep)*202/1000\relax}
  >{\raggedright\arraybackslash}p{\dimexpr(\linewidth-4\tabcolsep)*412/1000\relax}
  >{\raggedright\arraybackslash}p{\dimexpr(\linewidth-4\tabcolsep)*367/1000\relax}@{}}
\toprule\noalign{}
\begin{minipage}[b]{\linewidth}\raggedright
\end{minipage} & \begin{minipage}[b]{\linewidth}\raggedright
CPU-only
\end{minipage} & \begin{minipage}[b]{\linewidth}\raggedright
With GPU
\end{minipage} \\
\midrule\noalign{}
\endhead
\bottomrule\noalign{}
\endlastfoot
Slice LUTs & 18,375 of 203,800 (\textbf{9.0 \%}; 18,090 logic + 285
LUT-RAM) & 39,518 (\textbf{19.4 \%}; 36,417 logic + 3,101 LUT-RAM) \\
Slice registers & 13,143 (3.2 \%), 4 latches & 17,950 (4.4 \%), 5
latches \\
Block RAM tiles & 12 of 445 (2.7 \%) & 28 (6.3 \%): 1 RAMB36 + 54
RAMB18 \\
DSP48E1 & 2 of 840 (0.2 \%) & 26 (3.1 \%) \\
BUFGCTRL & 6 of 32 & 7 of 32 \\
Artix-7 for reference {[}R{]} & 18,981 LUT (14.2 \%), 13.2 K FF, 12
BRAM, 2 DSP & 40,586 LUT (30.3 \%), 18.0 K FF, 28 BRAM, 26 DSP \\
\end{longtable}
}

}

The mapping is the same as on Artix-7 (same block RAM, DSP and register
counts; LUT count 3 \% lower CPU-only and 3 \% lower with GPU, a packing
difference between the two device families {[}A{]}). The Kintex-7 simply
has 1.5x the LUTs.

\textbf{Hold:} CPU-only is clean in both steps (WHS +0.061 / +0.066 ns
on \code{clk_i}, +0.088 / +0.094 ns on \code{cpu_clk_i}). \textbf{With
GPU it is not, and this one is not a harness-shell artifact:} step 1
\code{clk_i} WHS \textbf{-0.312 ns on 118 endpoints} (THS -13.9 ns);
step 2 \textbf{-0.206 ns on 203 endpoints} (THS -19.9 ns). The worst
endpoints are internal bus-to-GPU registers (\code{u_axil_ic} /
\code{u_periph_bridge} / \code{u_sram} \code{buf_data_q} -\textgreater{}
\code{u_gpu} \code{ar_addr_q}, \code{aw_addr_q}, \code{if_id_q}), with
\textbf{+0.92 to +1.10 ns of clock skew} (destination later than
source). Cause {[}A, not isolated{]}: the GPU clock is the behavioural
LUT clock gate promoted to its own BUFG, so its insertion delay is later
than the plain \code{clk_i} BUFG (8 x TIMING-14 "LUT on the clock tree",
the same structural finding as Artix-7, where the skew was about 1.5
ns). On Artix-7 only 8 endpoints (-0.112 ns) failed hold and the worst
was a harness shell flop, so the Kintex-7 hold picture is worse and
different. It does not affect the fmax numbers, but it is a real finding
for the clock-gate decision (bead \bead{hlrm.2}): on hardware this would
need a hold fix (balanced clock, a vendor clock-enable primitive instead
of a LUT gate, or removing the gate on the prototype), and the run was
not repeated with a post-route \code{phys_opt_design}, which may or may
not close it.

\textbf{Cost (all within the host limits; MemAvailable was 12.6-13.9 GB
before every job, no OOM, no waiting needed):}

{
\begin{longtable}[]{@{}lll@{}}
\toprule\noalign{}
Job & Wall time & Peak RSS \\
\midrule\noalign{}
\endhead
\bottomrule\noalign{}
\endlastfoot
CPU-only synthesis & 3:56 & 2.49 GB \\
CPU-only P\&R, 10/7 ns & 5:59 & 3.14 GB \\
With GPU, synthesis + P\&R in one session, 10/7 ns & 15:08 & 4.07 GB \\
CPU-only P\&R, 7/5 ns & 8:31 & 2.83 GB \\
With GPU P\&R, 7/5 ns & 12:13 & 3.09 GB \\
\end{longtable}
}

(Artix-7 for reference {[}R{]}: 4:07 + 10:06 CPU-only, 21:23 with GPU.
Kintex-7 was faster for the same flow, probably because it met timing at
the first target so the router and \code{phys_opt} had less to do
{[}A{]}.)

\textbf{What this does to the recommendation:} \textbf{nothing; section
3 stands (Nexys Video).} Section 3 says to pick the Genesys 2 "only if
fmax turns out to be the problem". Kintex-7 is measurably faster (about
1.7-1.9x), but Artix-7 already gives \textasciitilde74-107 MHz, a 7-17x
scale factor against the ASIC clocks, which the Artix-7 conclusion above
already judged fast enough for firmware validation and the CDC soak
{[}A, judgement carried over, not re-evaluated{]}, so fmax is not the
problem. The Genesys 2 costs about 75 \% more (section 2 {[}R{]}) for
headroom that was not needed (19 \% of its LUTs with the GPU, against 30
\% on the XC7A200T; the XC7A200T has ample margin). The one new input is
the GPU-build hold behaviour above, which is a clock-gate design issue,
not a part selection issue.

\textbf{NOT verified in the Kintex-7 numbers:} a seed sweep or more than
one run per build and target (the Artix-7 \textasciitilde1 ns spread is
assumed to apply, not re-measured); the true fmax of the two domains
that still met timing at the tightest target tried (CPU-only fabric

\begin{quote}
= 148 MHz); the effect of \code{phys_opt_design} after routing and
non-default directives; fixing or characterising the GPU-build hold
endpoints; pad-level timing, power, any hardware run; the remaining
cause analysis for the CPU-only \code{cpu_clk_i} reset-path limiter (it
is a 1,807-endpoint family, the top three endpoints inspected); the
Genesys 2 board itself (pinout, clocks, price) is untouched.
\end{quote}

\begin{enumerate}
\def\labelenumi{\arabic{enumi}.}
\tightlist
\item
  Board (recommendation: Nexys Video).
\item
  Clock-speed sanity check: done for the CPU-only build. ECP5 lower
  bound \textasciitilde48-52 MHz CPU, \textasciitilde47-55 MHz fabric;
  Artix-7 XC7A200T-1 in Vivado 2026.1 \textasciitilde107 MHz CPU,
  \textasciitilde83 MHz fabric (above). Kintex-7 XC7K325T-2 (Genesys 2)
  \textasciitilde180-190 MHz CPU, \textasciitilde126 MHz+ fabric
  (subsection above).
\item
  Second clock source (recommendation: Si5351A breakout), after
  confirming clock-capable pins on the chosen board.
\item
  The Stage 1 decisions already recorded on bead \bead{hlrm}: GPU
  shared-memory RAM style, clock-gating strategy, strip DFT on the
  prototype, IO wrapper and \code{boot_rom} image owner.
\end{enumerate}

\subsection{Sources}\label{tm:dec:fpga-board-selection:sources}

\begin{itemize}
\tightlist
\item
  \href{https://www.digikey.sg/en/products/detail/digilent-inc/410-316/5456481}{DigiKey
  SG -- Nexys Video 410-316}
\item
  \href{https://sg.rs-online.com/web/p/fpga-development-tools/1346477}{RS
  Singapore -- Nexys Video}
\item
  \href{https://www.digikey.sg/en/products/detail/digilent-inc/410-300/5730862}{DigiKey
  SG -- Genesys 2 410-300}
\item
  \href{https://www.mouser.sg/c/semiconductors/engineering-development-tools/programmable-logic-ic-development-tools?pg=8}{Mouser
  SG -- programmable logic dev tools}
\item
  \href{https://sg.element14.com/digilent/410-300/dev-board-fpga/dp/2528808}{element14
  SG -- Genesys 2}
\item
  \href{https://boards.fpgadeveloper.com/boards/Genesys-2}{FPGA
  Developer -- Genesys 2}
\item
  \href{https://www.digikey.com/en/products/detail/amd/SK-KV260-G/13985269}{DigiKey
  US -- SK-KV260-G}
\item
  \href{https://fr.farnell.com/amd/sk-kv260-g-ed/kit-de-dem-64-bits-arm-cortex/dp/4719664}{Farnell
  -- SK-KV260-G-ED (XCK26 resources)}
\item
  \href{https://www.digikey.com/en/products/detail/digilent-inc/410-319-1/9445912}{DigiKey
  US -- Arty A7-100T}
\item
  \href{https://crowdsupply.com/radiona/ulx3s}{Crowd Supply -- ULX3S}
\item
  \href{https://learn.adafruit.com/adafruit-si5351-clock-generator-breakout/overview}{Adafruit
  -- Si5351A breakout overview}
\item
  \href{https://thepihut.com/products/adafruit-si5351a-clock-generator-breakout-board-8khz-to-160mhz}{The
  Pi Hut -- Adafruit Si5351A}
\item
  \href{https://www.bitsandparts.nl/en/Clock-generator-Si5351A-Breakout-8KHz-to-160MHz-Adafruit-2045-p100227}{Bits
  \& Parts -- Adafruit 2045}
\item
  \href{https://core-electronics.com.au/adafruit-si5351a-clock-generator-breakout-board-8khz-to-160mhz.html}{Core
  Electronics -- Adafruit 2045}
\item
  \href{https://wcm-cce.cldnet.analog.com/media/en/news-marketing-collateral/solutions-bulletins-brochures/maxim-analog-essentials-collection.pdf}{Maxim
  analog essentials (DS1086LPMB1)}
\item
  \href{https://forum.digilent.com/topic/17773-nexys-video-pmod-connectors/}{Digilent
  forum -- Nexys Video Pmod connectors}
\item
  \href{https://reference.digilentinc.com/programmable-logic/nexys-video/reference-manual}{Nexys
  Video reference manual}
\end{itemize}
\end{relocated}

\begin{relocated}{docs/FPGA_STAGE0_EVALUATION.md}
% Converted from docs/FPGA_STAGE0_EVALUATION.md by tools/docs/md2tex.py; edit this file, not the original.
\section{\texorpdfstring{FPGA Stage 0 evaluation (bead \code{hlrm}, GH
\#245, PR \#279)}{FPGA Stage 0 evaluation (bead , GH \#245, PR \#279)}}\label{tm:dec:fpga-stage0-evaluation}

Feasibility only: open tools, \textbf{no board, no vendor toolchain, no
place-and-route, no timing, no hardware.} Everything below is either
\emph{measured} (tool output quoted, source named) or \emph{inferred}
(marked). Nothing here is a decision; the open human decisions are
collected in section 9.

\subsection{Method}\label{tm:dec:fpga-stage0-evaluation:1-method}\origsec{1}

\begin{enumerate}
\def\labelenumi{\arabic{enumi}.}
\tightlist
\item
  \textbf{One define, \code{FPGA}}, selects block-RAM-inferable memories
  at every macro / flat-array site (section 3). With \code{FPGA}
  undefined the design is unchanged (section 4).
\item
  \textbf{File list} = the cocotb \code{PLL_SOURCES} list parsed out of
  \src{tb/cocotb/soc/Makefile} (cannot drift from what the regression
  elaborates) + \src{rtl/mem/rv32i_clock_gate.sv} (that list never names
  it: Verilator finds it through \code{-I}, sv2v has no \code{-I} and
  would silently blackbox it) + \code{rtl/fpga/*.sv}.
  \src{tools/fpga/stage0_filelist.sh}.
\item
  \textbf{No-GPU} = \code{rtl/gpu/*} replaced by
  \src{pnr/sky130/soc/gpu_top_tieoff.sv}, the \emph{same} synthesisable
  tie-off the Sky130 SoC flow uses (\src{pnr/Makefile}
  \code{SKY130_SOC_SV_FILES}). No new exclusion mechanism. Note the
  Sky130 flow also blackboxes the CPU as a hard macro; \textbf{here the
  CPU stays real}, so "no-GPU" is CPU + peripherals + CRYPTO + NPU + DMA
  + SRAM.
\item
  \textbf{Synthesis}: sv2v
  (\texttt{-\allowbreak{}-\allowbreak{}define=\allowbreak{}FPGA\ -\allowbreak{}-\allowbreak{}define=\allowbreak{}\_\allowbreak{}\_\allowbreak{}pnr\_\allowbreak{}\_\allowbreak{}})
  -\textgreater{} yosys \code{read_verilog},
  \texttt{hierarchy\ -\allowbreak{}check\ -\allowbreak{}top\ soc\_\allowbreak{}top},
  \texttt{synth\_\allowbreak{}xilinx\ -\allowbreak{}family\ xc7\ -\allowbreak{}flatten\ -\allowbreak{}noiopad}
  or
  \texttt{synth\_\allowbreak{}ecp5\ -\allowbreak{}flatten\ -\allowbreak{}noiopad},
  \code{stat}, \code{write_json}. \texttt{ulimit\ -\allowbreak{}v} 10
  GB, one job at a time, outputs under
  \code{/nobackup/fpga_stage0/synth} (not in the repo).
  \src{tools/fpga/stage0_synth.sh}. \code{-noiopad}: no
  IOBUF/\allowbreak{}IBUF insertion, so port buffers are not in the
  counts.
\item
  \textbf{Resource + depth numbers} come from the yosys JSON via
  \src{tools/fpga/stage0_report.py}.
\item
  \code{__pnr__} is the define every synthesis flow in this repo already
  uses; it zero-initialises \code{boot_rom} (see section 7).
\end{enumerate}

\subsubsection{Tool versions and where they came from}\label{tm:dec:fpga-stage0-evaluation:tool-versions-and-where-they-came-from}

{\footnotesize\setlength{\tabcolsep}{3pt}

{
\begin{longtable}[]{@{}
  >{\raggedright\arraybackslash}p{\dimexpr(\linewidth-6\tabcolsep)*295/1000\relax}
  >{\raggedright\arraybackslash}p{\dimexpr(\linewidth-6\tabcolsep)*113/1000\relax}
  >{\raggedright\arraybackslash}p{\dimexpr(\linewidth-6\tabcolsep)*351/1000\relax}
  >{\raggedright\arraybackslash}p{\dimexpr(\linewidth-6\tabcolsep)*221/1000\relax}@{}}
\toprule\noalign{}
\begin{minipage}[b]{\linewidth}\raggedright
Tool
\end{minipage} & \begin{minipage}[b]{\linewidth}\raggedright
Version
\end{minipage} & \begin{minipage}[b]{\linewidth}\raggedright
Source
\end{minipage} & \begin{minipage}[b]{\linewidth}\raggedright
In the project flake?
\end{minipage} \\
\midrule\noalign{}
\endhead
\bottomrule\noalign{}
\endlastfoot
sv2v & 0.0.13.1 &
\code{/nix/store/bknj130bjxz018c73yawkjmbzjhppqbc-sv2v-0.0.13.1}
(hard-coded in \src{pnr/Makefile} \texttt{SV2V\ ?=\allowbreak{}}) &
\textbf{No.} Not in any \src{flake.nix} devShell. \\
yosys & 0.62 (git 7326bb7) &
\code{/nix/store/f1q0w7rd0a4ny4hqvfxlhs4cmariidcy-yosys-0.62} (found in
the nix store; presumably pulled in by the LibreLane environment:
inference) & \textbf{No.} \\
Verilator & 5.048 & project flake \code{default} devShell & Yes \\
nextpnr (any), prjtrellis, prjxray, openFPGALoader & - & \textbf{absent}
(searched \code{/nix/store}) & No \\
\end{longtable}
}

}

Nothing was installed. The shared external LibreLane install was not
edited. \textbf{A flake change is needed before Stage 1}: today yosys
and sv2v are reachable only by absolute store path. A small \code{fpga}
devShell (yosys, sv2v, nextpnr-ecp5 + prjtrellis, openFPGALoader) would
make this reproducible; it is \emph{not} done in this PR (decision 9.6).

\subsection{Results}\label{tm:dec:fpga-stage0-evaluation:2-results}\origsec{2}

\subsubsection{\texorpdfstring{Resource table (measured: yosys 0.62
\code{stat} / JSON, \code{-flatten}, no
P\&R)}{Resource table (measured: yosys 0.62  / JSON, , no P\&R)}}\label{tm:dec:fpga-stage0-evaluation:21-resource-table-measured-yosys-062-stat-json-flatten-no-pr}\origsec{2.1}

{\small\setlength{\tabcolsep}{4pt}

{
\begin{longtable}[]{@{}
  >{\raggedright\arraybackslash}p{\dimexpr(\linewidth-8\tabcolsep)*130/1000\relax}
  >{\raggedright\arraybackslash}p{\dimexpr(\linewidth-8\tabcolsep)*215/1000\relax}
  >{\raggedright\arraybackslash}p{\dimexpr(\linewidth-8\tabcolsep)*340/1000\relax}
  >{\raggedright\arraybackslash}p{\dimexpr(\linewidth-8\tabcolsep)*164/1000\relax}
  >{\raggedright\arraybackslash}p{\dimexpr(\linewidth-8\tabcolsep)*130/1000\relax}@{}}
\toprule\noalign{}
\begin{minipage}[b]{\linewidth}\raggedright
\end{minipage} & \begin{minipage}[b]{\linewidth}\raggedright
Xilinx 7-series, no-GPU
\end{minipage} & \begin{minipage}[b]{\linewidth}\raggedright
Xilinx 7-series, with GPU
\end{minipage} & \begin{minipage}[b]{\linewidth}\raggedright
ECP5, no-GPU
\end{minipage} & \begin{minipage}[b]{\linewidth}\raggedright
ECP5, with GPU
\end{minipage} \\
\midrule\noalign{}
\endhead
\bottomrule\noalign{}
\endlastfoot
LUT (logic) & 38,554 (LUT1-6) & 62,353 & 34,172 (LUT4) & 67,427
(LUT4) \\
+ LUT-RAM & 46 RAM32M (\textasciitilde184 LUT sites) & 46 RAM32M + 704
RAM64M + 1024 RAM128X1S (\textasciitilde5,048 LUT sites) & 68
TRELLIS\_\allowbreak{}DPR16X4 & 2,116 TRELLIS\_\allowbreak{}DPR16X4 \\
FF & 14,079 & 19,665 & 14,349 & 18,911 \\
Latch cells & 4 (LDCE) & 5 (LDCE) & see 2.4 & see 2.4 \\
Block RAM & 15 RAMB18 + 1 RAMB36 (= 9 RAMB36-eq) & same & 17 DP16KD (18
Kb) & 49 DP16KD \\
DSP & 19 DSP48E1 & 44 DSP48E1 & 19 MULT18X18D & 45 MULT18X18D \\
Carry & 852 CARRY4 & 1,194 CARRY4 & 1,731 CCU2C (2 bit each) & 2,372
CCU2C \\
Wide-mux & 3,271 MUXF7 / 872 MUXF8 & 4,472 / 941 & 6,220 PFUMX / 1,407
L6MUX21 & 9,505 / 512 \\
yosys wall / peak RSS & 5:38 / 2.0 GB & 8:36 / 2.1 GB & 3:05 / 1.6 GB &
9:42 / 2.1 GB \\
\end{longtable}
}

}

Caveats on the table: yosys/ABC LUT counts are typically 10-30 \%
\emph{worse} than Vivado/\allowbreak{}Quartus (\textbf{inference}, not
measured here); LUT-RAM sites are my conversion from primitive counts
(RAM32M = 4 sites, RAM64M = 4, RAM128X1S = 2); ECP5 DPR16X4 costs
roughly 2-3 slices each (\textbf{inference}, from memory of the ECP5
docs). Latch cells and unmapped logic are in the totals as yosys
reported them.

\subsubsection{\texorpdfstring{Which memories inferred block RAM
(measured:
\texttt{mapping\ memory\ .\allowbreak{}.\allowbreak{}.\allowbreak{}}
lines in the
logs)}{Which memories inferred block RAM (measured:  lines in the logs)}}\label{tm:dec:fpga-stage0-evaluation:22-which-memories-inferred-block-ram-measured-mapping-memory-lines-in-the-logs}\origsec{2.2}

{\small\setlength{\tabcolsep}{4pt}

{
\begin{longtable}[]{@{}
  >{\raggedright\arraybackslash}p{\dimexpr(\linewidth-6\tabcolsep)*362/1000\relax}
  >{\raggedright\arraybackslash}p{\dimexpr(\linewidth-6\tabcolsep)*140/1000\relax}
  >{\raggedright\arraybackslash}p{\dimexpr(\linewidth-6\tabcolsep)*309/1000\relax}
  >{\raggedright\arraybackslash}p{\dimexpr(\linewidth-6\tabcolsep)*169/1000\relax}@{}}
\toprule\noalign{}
\begin{minipage}[b]{\linewidth}\raggedright
Memory
\end{minipage} & \begin{minipage}[b]{\linewidth}\raggedright
Bits
\end{minipage} & \begin{minipage}[b]{\linewidth}\raggedright
Xilinx 7
\end{minipage} & \begin{minipage}[b]{\linewidth}\raggedright
ECP5
\end{minipage} \\
\midrule\noalign{}
\endhead
\bottomrule\noalign{}
\endlastfoot
I-cache tag + 4 data banks (5 x 256x32) & 5 x 8 Kb & 5 x BRAM SDP
(RAMB18) & 5 x PDPW16KD \\
D-cache tag + 4 data banks (5 x 256x32) & 5 x 8 Kb & 5 x BRAM SDP & 5 x
PDPW16KD \\
\code{sram_controller} main SRAM (1024x32, byte-write) & 32 Kb & 1 x
BRAM TDP (RAMB36) & 1 x DP16KD pair (2 x 18 Kb) \\
NPU weight memory (4 x 256x32) & 4 x 8 Kb & 4 x BRAM SDP & 4 x
PDPW16KD \\
DMA \code{linebuf} (not an FPGA arm; inferred from the existing RTL) &
small & BRAM SDP (RAMB18) & PDPW16KD \\
DMA \code{q_mem}, 5 x CDC FIFO \code{mem_q} & small & LUT-RAM & DPR16X4
(LUT-RAM) \\
GPU shared memory: 32 banks x 128x32 (16 KB total) & 32 x 4 Kb &
\textbf{LUT-RAM, 32 x RAM128X1S-based (about 5 K LUT sites)} &
\textbf{BRAM (32 x PDPW16KD)} \\
GPU vector register file \code{rf} (8 lanes) & 8 x ... & LUT-RAM &
LUT-RAM \\
CPU register file & 31 x 32 x 3 read ports & flops/LUT logic (not mapped
as \texttt{mapping\ memory}) & same \\
\end{longtable}
}

}

\textbf{Finding:} the vendor mappers disagree on the same RTL for the
GPU shared memory (7-series chose LUT-RAM, ECP5 chose BRAM), and
7-series LUT-RAM there costs \textasciitilde5 K LUT sites. Forcing
\texttt{(*\ ram\_\allowbreak{}style\ =\allowbreak{}\ "block"\ *)} in
\code{fpga_sram_1rw} for the 128-deep case would move 7-series to 32
RAMB18. I did \textbf{not} do this (it is a
sizing/\allowbreak{}portability choice, section 9.3). \code{proc_rom}
turned 31 (no-GPU) / 33 (GPU) constant-\code{case} memories (AES/SHA
S-boxes and constants) into logic; none became block RAM.

\subsubsection{Blackboxes / undefined modules}\label{tm:dec:fpga-stage0-evaluation:23-blackboxes-undefined-modules}\origsec{2.3}

\texttt{hierarchy\ -\allowbreak{}check\ -\allowbreak{}top\ soc\_\allowbreak{}top}
passed in all four jobs and the log greps for blackbox / undefined /
missing are empty (measured). The earlier hazard is real though: before
\code{rv32i_clock_gate.sv} was added to the list explicitly, Verilator
lint reported \code{MODMISSING} for it (3 instances) and sv2v would have
silently blackboxed it. With IOBUF stripped (\code{-noiopad}) and no
vendor primitives instantiated in RTL, nothing else is left undefined.

\subsubsection{\texorpdfstring{Latches (measured:
\texttt{Latch\ inferred} lines)}{Latches (measured:  lines)}}\label{tm:dec:fpga-stage0-evaluation:24-latches-measured-latch-inferred-lines}\origsec{2.4}

Six \code{proc_dlatch} inferences, identical in all four jobs:
\code{rv32i_clock_gate.en_latch} (the ICG latch),
\code{dft_clk_switch.en_a_q} / \code{en_b_q} (intentional: glitch-free
clock switch), and three \code{sv2v_autoblock_N.w} temporaries in
\code{rv32i_dcache} / \code{rv32i_icache}. The final cell count is 4
LDCE (no-GPU) / 5 LDCE (GPU) on 7-series, which matches the three
\code{rv32i_clock_gate} instances (\code{u_cpu_cg}, \code{u_gpu_cg},
\code{u_tck_cg}) + the two switch latches with the GPU gate removed by
the tie-off; the three \code{autoblock.w} latches do not appear to
survive as cells (\textbf{UNVERIFIED}: I did not trace them past
\code{proc}). They exist in the shared RTL, not in the \code{FPGA} arms.

\subsubsection{fmax: NOT MEASURED (ECP5 lower bound measured later)}\label{tm:dec:fpga-stage0-evaluation:25-fmax-not-measured-ecp5-lower-bound-measured-later}\origsec{2.5}

\begin{quote}
A post-route ECP5 lower bound now exists: CPU-only, LFE5U-85F-6, nextpnr
0.11.1, about 48-52 MHz (CPU domain) and 47-55 MHz (fabric domain) over
3 seeds; see
\code{FPGA_BOARD_SELECTION.md} (\cref{tm:dec:fpga-board-selection}) §5.
A real Artix-7 XC7A200T Vivado 2026.1 post-route result
(\textasciitilde107 MHz CPU, \textasciitilde83 MHz fabric, CPU-only) is
in the same section, "Measured: Artix-7 XC7A200T".
\end{quote}

There is \textbf{no nextpnr}, no Vivado, no Quartus. yosys has no
placement or routing. What exists is a \emph{depth proxy}: the longest
combinational path in LUT-level equivalents (LUT = 1, MUXF = 0.5, CARRY4
= 0.25; cut at
FF/\allowbreak{}latch/\allowbreak{}BRAM/\allowbreak{}DSP/\allowbreak{}ports),
from the mapped JSON:

{
\begin{longtable}[]{@{}
  >{\raggedright\arraybackslash}p{\dimexpr(\linewidth-6\tabcolsep)*173/1000\relax}
  >{\raggedright\arraybackslash}p{\dimexpr(\linewidth-6\tabcolsep)*154/1000\relax}
  >{\raggedright\arraybackslash}p{\dimexpr(\linewidth-6\tabcolsep)*375/1000\relax}
  >{\raggedright\arraybackslash}p{\dimexpr(\linewidth-6\tabcolsep)*279/1000\relax}@{}}
\toprule\noalign{}
\begin{minipage}[b]{\linewidth}\raggedright
\end{minipage} & \begin{minipage}[b]{\linewidth}\raggedright
Depth (LUT-eq)
\end{minipage} & \begin{minipage}[b]{\linewidth}\raggedright
Bracket (depth x ns/level + overhead)
\end{minipage} & \begin{minipage}[b]{\linewidth}\raggedright
Assumption
\end{minipage} \\
\midrule\noalign{}
\endhead
\bottomrule\noalign{}
\endlastfoot
7-series, no-GPU & 20.5 & 58 .. 99 MHz & 0.45-0.80 ns/level + 0.9 ns \\
7-series, GPU & 20.5 & 58 .. 99 MHz & same \\
ECP5, no-GPU & 19.0 & 38 .. 61 MHz & 0.80-1.30 ns/level + 1.3 ns \\
ECP5, GPU & 23.5 & 31 .. 50 MHz & same \\
\end{longtable}
}

Limits: the ns/level figures are my assumptions, not device data; it
ignores fan-out and routing congestion (this design has large flat
fan-outs, e.g. the ASIC work\textquotesingle s 836-fanout counter-enable
cone); it treats all clocks together (two clock ports, \code{clk_i} and
\code{cpu_clk_i}); the ECP5 runs cut 4 / 6 combinational loops
(clock-gate / clock-mux latch paths) so the path may be understated
there. \textbf{Treat "tens of MHz" as the order of magnitude and nothing
more.} An ASIC target of 1282 MHz (CPU) / 571 MHz (fabric) implies a
prototype scale factor of roughly 10-40x; all later measurements must be
labelled with it.

\subsection{The FPGA arms (what was written, by whom, what each
infers)}\label{tm:dec:fpga-stage0-evaluation:3-the-fpga-arms-what-was-written-by-whom-what-each-infers}\origsec{3}

Written directly by the fpga-orchestrator (not sub-delegated to the
rtl-design orchestrator); lint and checker below.

{\footnotesize\setlength{\tabcolsep}{3pt}

{
\begin{longtable}[]{@{}
  >{\raggedright\arraybackslash}p{\dimexpr(\linewidth-4\tabcolsep)*248/1000\relax}
  >{\raggedright\arraybackslash}p{\dimexpr(\linewidth-4\tabcolsep)*366/1000\relax}
  >{\raggedright\arraybackslash}p{\dimexpr(\linewidth-4\tabcolsep)*366/1000\relax}@{}}
\toprule\noalign{}
\begin{minipage}[b]{\linewidth}\raggedright
Site
\end{minipage} & \begin{minipage}[b]{\linewidth}\raggedright
Arm
\end{minipage} & \begin{minipage}[b]{\linewidth}\raggedright
Infers
\end{minipage} \\
\midrule\noalign{}
\endhead
\bottomrule\noalign{}
\endlastfoot
\src{rtl/fpga/fpga_sram_1rw.sv} (new) & stand-in for every 32-bit 1RW
macro; ports = the macro/sim contract; dout holds on write and when
\code{csb0} is high (identical to \src{sim/sram_1rw_256x32_verilator.v},
the model the regression uses) & single-port block RAM, sync read, no
reset, one address \\
\src{rtl/fpga/fpga_sdp_ram_be.sv} (new) & byte-enabled simple-dual-port,
registered read, read-first on a same-address collision & block RAM
SDP/TDP with byte write \\
\code{rv32i_icache.sv}, \code{rv32i_dcache.sv} & \texttt{`elsif\ FPGA}
arm for tag and data
(\texttt{fpga\_\allowbreak{}sram\_\allowbreak{}1rw\ \#(.\allowbreak{}AW(8))})
& BRAM \\
\code{npu_weight_mem.sv} & \texttt{`elsif\ FPGA} in the banked arm: 4 x
\code{fpga_sram_1rw}; the shared 4:1 mux and 2-cycle read-latency
normalisation are untouched & BRAM \\
\code{gpu/shared_memory.sv} & \texttt{`ifdef\ FPGA} around the bank
instance (\code{AW(7)}) & LUT-RAM (7-series) / BRAM (ECP5), see 2.2 \\
\code{soc/sram_controller.sv} & the existing \code{SRAM_SKY130} two-port
FSM (single write port, dedicated registered-read port, 2-entry skid
buffer, \textbf{no flat array}) is reused through a file-local macro
\code{SRAM_CTRL_2P} (defined iff \code{SRAM_SKY130} or \code{FPGA};
undefined at end of file); \code{FPGA} instantiates
\code{fpga_sdp_ram_be} instead of the OpenRAM macro & BRAM (RAMB36 TDP),
instead of the 32,768-flop + 1024:1 mux flat array \\
\code{rv32i_clock_gate.sv} & \textbf{NO ARM} (section 6) & - \\
\end{longtable}
}

}

Why registered-read is timing-compatible: the Sky130 macro launches
\code{dout1} on the negedge after the address edge and the FSM samples
it at the next posedge; a registered-output block RAM presents the same
word over the same cycle, so \code{rd_pend_q} is high exactly when the
data is fresh. A same-cycle read+write of the same address is a race in
the Sky130 arm too; the FPGA arm returns the old word.

No arm contains an \code{initial}, a vendor primitive, an array reset,
or a runtime-indexed mux of flops. The \code{ma7} defect class
(\code{OPT_MUXTREE} on runtime-indexed muxes) is a Synlig issue; this
flow uses sv2v + yosys, and the coarse-grain check was not re-run for
the new files (nothing in them selects a flop array by a runtime index;
reasoning, not a tool run).

\subsection{Non-FPGA identity proof}\label{tm:dec:fpga-stage0-evaluation:4-non-fpga-identity-proof}\origsec{4}

\src{tools/fpga/stage0_sv2v_identity.sh} runs sv2v over the full SoC
(with GPU and with the tie-off; file list identical on both trees,
\src{rtl/fpga} deliberately excluded) under the five non-FPGA define
sets: none (cocotb default), \code{SIMULATION}, \code{__pnr__},
\texttt{\_\allowbreak{}\_\allowbreak{}pnr\_\allowbreak{}\_\allowbreak{}\ +\allowbreak{}\ SRAM\_\allowbreak{}SKY130}
(Sky130 SoC flow),
\texttt{USE\_\allowbreak{}ICG\_\allowbreak{}CELL\ +\allowbreak{}\ \_\allowbreak{}\_\allowbreak{}pnr\_\allowbreak{}\_\allowbreak{}\ +\allowbreak{}\ SRAM\_\allowbreak{}ASAP7}
(ASAP7 SoC flow), on a pristine
\texttt{git\ archive\ origin/\allowbreak{}main} export and on this
branch, and compares with \texttt{cmp\ -\allowbreak{}s}.

Result (measured): \textbf{10 of 10 outputs byte-identical after one
normalisation.} sv2v embeds
\texttt{<absolute\ path>:\allowbreak{}<line>:\allowbreak{}<col>} in
every elaboration-time \texttt{\$fatal} message; that is a diagnostic
source location, it differs between two checkouts of the same RTL and
shifts when a file gains lines above a \texttt{\$fatal}
(\code{sram_controller.sv} gained 10). The script strips exactly that
token
(\texttt{Fatal\ [.\allowbreak{}.\allowbreak{}]\ <path>:\allowbreak{}L:\allowbreak{}C\ }
-\textgreater{} \texttt{Fatal\ [.\allowbreak{}.\allowbreak{}]\ }) and
nothing else before comparing. Before normalisation all 10 differed by
those lines only (first run: 72 diff lines in \code{nogpu_sky130}, all
\texttt{\$display("Fatal\ .\allowbreak{}.\allowbreak{}.\allowbreak{}}
lines).

{\small\setlength{\tabcolsep}{4pt}

{
\begin{longtable}[]{@{}
  >{\raggedright\arraybackslash}p{\dimexpr(\linewidth-2\tabcolsep)*490/1000\relax}
  >{\raggedright\arraybackslash}p{\dimexpr(\linewidth-2\tabcolsep)*490/1000\relax}@{}}
\toprule\noalign{}
\begin{minipage}[b]{\linewidth}\raggedright
output
\end{minipage} & \begin{minipage}[b]{\linewidth}\raggedright
sha256 prefix (normalised)
\end{minipage} \\
\midrule\noalign{}
\endhead
\bottomrule\noalign{}
\endlastfoot
gpu\_sim / gpu\_\allowbreak{}simulation / gpu\_pnr /
gpu\_\allowbreak{}sky130 / gpu\_asap7 & 76abbdb2... / bed3b4d2... /
aaf5e54e... / 50ae3bdd... / d23ab691... \\
nogpu\_sim / nogpu\_\allowbreak{}simulation / nogpu\_pnr /
nogpu\_\allowbreak{}sky130 / nogpu\_\allowbreak{}asap7 & 4ac613f2... /
1361d449... / e064e71c... / d7f0bcd2... / 60f174dc... \\
\end{longtable}
}

}

Verilator \code{-Wall} lint of full \code{soc_top}
(\texttt{-\allowbreak{}Wno-\allowbreak{}IMPORTSTAR\ -\allowbreak{}Wno-\allowbreak{}UNUSEDSIGNAL\ -\allowbreak{}Wno-\allowbreak{}UNOPTFLAT\ -\allowbreak{}Wno-\allowbreak{}SYNCASYNCNET}),
measured: FPGA gpu 0 errors / 3 warnings, FPGA no-GPU 0 / 4, default gpu
0 / 4, default no-GPU 0 / 5; none of the warnings is in a file this
change touched (\code{rv32i_cache_pkg} unused parameters,
\code{DECLFILENAME} on the tie-off).
\code{check_struct_member_select.py}: exit 0.

\subsection{Regression with the FPGA define}\label{tm:dec:fpga-stage0-evaluation:5-regression-with-the-fpga-define}\origsec{5}

Command:
\texttt{nix\ develop\ -\allowbreak{}-\allowbreak{}command\ make\ -\allowbreak{}i\ -\allowbreak{}C\ tb/\allowbreak{}cocotb/\allowbreak{}soc\ FPGA=\allowbreak{}1\ SIM\_\allowbreak{}BUILD\_\allowbreak{}ROOT=\allowbreak{}/\allowbreak{}nobackup/\allowbreak{}fpga\_\allowbreak{}stage0/\allowbreak{}regr\ soc\_\allowbreak{}all\_\allowbreak{}ci}
(separate, content-keyed SIM\_BUILD root; \code{FPGA=1} adds
\texttt{+\allowbreak{}define+\allowbreak{}FPGA\ -\allowbreak{}y\ rtl/\allowbreak{}fpga},
a new opt-in hook in \src{tb/cocotb/soc/Makefile}). \code{-i} because
\code{soc_all_ci} is a chain of \texttt{\$(MAKE)} lines and stops at the
first failing suite; the first, un-\code{-i} run aborted at
\code{soc_boot} after 700 passed / 1 failed.

{\footnotesize\setlength{\tabcolsep}{3pt}

{
\begin{longtable}[]{@{}
  >{\raggedright\arraybackslash}p{\dimexpr(\linewidth-6\tabcolsep)*434/1000\relax}
  >{\raggedright\arraybackslash}p{\dimexpr(\linewidth-6\tabcolsep)*56/1000\relax}
  >{\raggedright\arraybackslash}p{\dimexpr(\linewidth-6\tabcolsep)*56/1000\relax}
  >{\raggedright\arraybackslash}p{\dimexpr(\linewidth-6\tabcolsep)*434/1000\relax}@{}}
\toprule\noalign{}
\begin{minipage}[b]{\linewidth}\raggedright
Run
\end{minipage} & \begin{minipage}[b]{\linewidth}\raggedright
Passed
\end{minipage} & \begin{minipage}[b]{\linewidth}\raggedright
Failed
\end{minipage} & \begin{minipage}[b]{\linewidth}\raggedright
Notes
\end{minipage} \\
\midrule\noalign{}
\endhead
\bottomrule\noalign{}
\endlastfoot
\code{soc_all_ci}, \code{FPGA=1}, no test edits (measured, CI-style grep
\texttt{cocotb.\allowbreak{}regression\ .\allowbreak{}*\ passed\$}),
\textbf{before the helper} & \textbf{742} & \textbf{20} & 72 suite
summaries; main is 762 \\
Diagnostic re-run of the 15 suites that contain the 20 failures, TB
handle path temporarily changed (edit reverted, never committed) & 31 &
\textbf{0} & the 20 pass; 742 + 20 = \textbf{762} = main \\
\code{soc_all_ci}, \code{FPGA=1}, \textbf{with the
\code{sram_backdoor.sram_mem} helper (\prref{282})}, same command,
\code{SIM_BUILD_ROOT=/nobackup/fpga_helper/regr} (measured, CI grep) &
\textbf{762} & \textbf{0} & 72 suite summaries, 17 skipped, 0
\texttt{contains\ no\ object\ named} lines, make exit 0; 2026-10-14
18:12 to 19:46 (cold build) \\
Default build, same commit, CI \texttt{cocotb\ soc\_\allowbreak{}all}
job (\prref{282}, run 38554821410) & \textbf{762} & \textbf{0} & log line
\texttt{cocotb:\allowbreak{}\ 762\ passed,\allowbreak{}\ 0\ failed};
floors untouched \\
\end{longtable}
}

}

\textbf{Resolved by \src{tb/cocotb/soc/sram_backdoor.py} (\prref{282}).}
\code{sram_mem(dut)} tries \code{dut.u_soc.u_sram.mem}, then
\code{dut.u_soc.u_sram.u_fpga_mem.mem}, and raises
\code{SramBackdoorError} naming every path tried if neither resolves;
all 31 backdoor sites in 18 SoC suites go through it (pure test
refactor, no assertion or RTL change;
\src{tb/tests/test_sram_backdoor.py} pins each path and the error). The
\code{SRAM_SKY130} macro arm (\code{u_sram.u_sram_macro}) is
deliberately \textbf{not} served: no SoC-level cocotb build defines
\code{SRAM_SKY130}, so a path for it could not be verified, and the
macro\textquotesingle s read/\allowbreak{}write timing differs from the
flat array; the helper fails loudly there. The statement below describes
the failure as it was before the helper.

\textbf{Root cause of all 20 failures (one cause, measured):} every one
is
\texttt{AttributeError:\allowbreak{}\ u\_\allowbreak{}sram\ contains\ no\ object\ named\ mem}
(20 of 20 tracebacks, \texttt{grep\ -\allowbreak{}c}). These tests
backdoor-read / backdoor-load the SoC SRAM through the hierarchical
handle \code{dut.u_soc.u_sram.mem} (about 20 sites across
\code{test_boot}, \code{test_periph_loopback},
\code{test_soc_coherency}, \code{test_cpu_gpu_irq},
\code{test_soc_gpio/pwm/wdt/trng/i2c/crypto/npu}, the multiclock and PMU
suites). In the \code{FPGA} arm (and equally in the \code{SRAM_SKY130}
arm) the array is not in \code{sram_controller}; it is
\code{u_sram.u_fpga_mem.mem}. It is a testbench hierarchy dependency,
not a functional RTL difference: with the path pointed at
\code{u_fpga_mem.mem} all 20 pass. It also means those 20 tests are the
ones that prove the \code{sram_controller} FPGA arm end to end through
the SoC; that evidence now exists from the unmodified-RTL regression
with the helper (762/762 above), not only from the diagnostic re-run.
The \code{sram_controller} unit suite (AXI-only) and the boot PC
scoreboard / 100-boot stability passed unmodified under \code{FPGA=1}.

\textbf{CI job: NOT added (user decision 2026-10-14: "path helper, no CI
job"), floors untouched} (\code{PASS_FLOOR} 759 / \code{CPU_PASS_FLOOR}
unchanged). The \code{FPGA=1} regression is a local-only check. Cost:
\textasciitilde35 min warm (measured: 06:57 to 07:24 for the cached
\code{-i} rerun), 94 min for the cold helper run above on this 15 GB
host. The testbench blocker that kept it red is removed by the helper.

\subsection{Clock gating: no arm written (decision needed)}\label{tm:dec:fpga-stage0-evaluation:6-clock-gating-no-arm-written-decision-needed}\origsec{6}

\code{rv32i_clock_gate} is a latch + AND on the clock. Synthesised as-is
(the numbers above assume this unchanged behavioural gate): one
\code{LDCE} per instance, and the AND output drives flop clocks from
fabric, not from a clock buffer. Vivado/\allowbreak{}nextpnr accept this
only with warnings; the gated net is not a global clock, so skew and
hold are untrustworthy. It gates \code{core_clk} into
\code{u_cpu}/\code{u_gpu} (PMU) and the DFT TCK path. \textbf{More than
one reasonable option exists, so nothing was committed:}

{\footnotesize\setlength{\tabcolsep}{3pt}

{
\begin{longtable}[]{@{}
  >{\raggedright\arraybackslash}p{\dimexpr(\linewidth-4\tabcolsep)*201/1000\relax}
  >{\raggedright\arraybackslash}p{\dimexpr(\linewidth-4\tabcolsep)*390/1000\relax}
  >{\raggedright\arraybackslash}p{\dimexpr(\linewidth-4\tabcolsep)*390/1000\relax}@{}}
\toprule\noalign{}
\begin{minipage}[b]{\linewidth}\raggedright
Option
\end{minipage} & \begin{minipage}[b]{\linewidth}\raggedright
What it is
\end{minipage} & \begin{minipage}[b]{\linewidth}\raggedright
Cost / risk
\end{minipage} \\
\midrule\noalign{}
\endhead
\bottomrule\noalign{}
\endlastfoot
A. Vendor primitive per family & \code{BUFGCE} (Xilinx) / \code{DCCA}
(ECP5) in the arm & Correct and glitch-free, but two vendor-specific
arms; unknown-module lint stubs; cannot be exercised by Verilator \\
B. Clock enable & Remove gating; add a CE to every flop in the CPU/GPU
domains & Closest to FPGA practice, but touches every flop (not a
\code{rv32i_clock_gate}-only edit) and changes the verified RTL
broadly \\
C. Pass-through & \texttt{gclk\ =\allowbreak{}\ clk}; the PMU still
sequences reset/\allowbreak{}isolation/\allowbreak{}enables & One-line
arm, no primitive; but the CPU/GPU keep toggling while "gated", so PMU
power-down tests no longer prove the gate; functional behaviour under
PMU sequencing is then only reset+isolation \\
D. Keep the latch+AND & Let the tool insert a BUFG on the fabric output
& Zero RTL change; tool-dependent, timing unsupported, one more
CDC-style hazard \\
\end{longtable}
}

}

Recommendation (not implemented): \textbf{C for Stage 1 bring-up}
(nothing to debug in the gate; PMU register interface still works), and
\textbf{A} only if a Stage 3 power-sequencing test needs real gating.

\subsection{What does not map to an FPGA (or maps only
conditionally)}\label{tm:dec:fpga-stage0-evaluation:7-what-does-not-map-to-an-fpga-or-maps-only-conditionally}\origsec{7}

{\footnotesize\setlength{\tabcolsep}{3pt}

{
\begin{longtable}[]{@{}
  >{\raggedright\arraybackslash}p{\dimexpr(\linewidth-2\tabcolsep)*490/1000\relax}
  >{\raggedright\arraybackslash}p{\dimexpr(\linewidth-2\tabcolsep)*490/1000\relax}@{}}
\toprule\noalign{}
\begin{minipage}[b]{\linewidth}\raggedright
Item
\end{minipage} & \begin{minipage}[b]{\linewidth}\raggedright
Result
\end{minipage} \\
\midrule\noalign{}
\endhead
\bottomrule\noalign{}
\endlastfoot
PLL (\code{pll_clkgen_stub}, \code{pll_subsystem}) & Stub is a pure
reference pass-through; synthesises to wires. A real prototype needs
MMCM/\allowbreak{}PLLE2 (7-series) or EHXPLLL (ECP5) at the top, not
written. \code{pll_apb_regs} just becomes registers with no effect. \\
DFT (\code{jtag_tap}, \code{dft_clk_switch}, \code{dft_clk_mux},
\code{dft_rst_mux}, TCK ICG) & Synthesises (latch + mux clock paths, see
2.4) but is dead weight on a prototype; the intentional-latch clock
switch maps to LDCE and creates fabric clock muxes the tools will warn
about. Options: tie scan/test off and accept, or strip under \code{FPGA}
(not done). \\
CDC primitives (\code{cdc_sync}, \code{cdc_reset_sync},
\code{cdc_gray_fifo}, \code{async_axi_fifo}, \code{apb_cdc_bridge}) &
Map: 2-FF chains become FDRE/\allowbreak{}FDCE, gray FIFO storage
becomes LUT-RAM with an async read. \textbf{Not} done: \code{ASYNC_REG}
attributes and the
\texttt{set\_\allowbreak{}max\_\allowbreak{}delay\ -\allowbreak{}datapath\_\allowbreak{}only}
constraints (the \#94 SDC is Tcl-validated only and is Synopsys-style,
not XDC/LPF). Without them a vendor tool will time the crossings as
synchronous. \\
Pad-less GPIO / I2C triplets (\code{*_o}, \code{*_oe_o}, \code{*_i}) &
No \code{inout} in the RTL. \code{-noiopad} leaves 292 port bits
unbuffered. A top-level wrapper with \code{IOBUF}/\code{TRELLIS_IO} is
needed at Stage 1; not written. \\
\code{boot_rom} & \textbf{Numbers understate it.} Under \code{__pnr__}
it is zero-initialised (1024 words, 4 KB), so yosys folds it to
constants. A real prototype loads firmware:
\code{MEM_INIT_FILE}/\code{initial} becomes BRAM init, about 1 RAMB36
(7-series). Not in the counts. \\
2 clock domains (\code{clk_i} 571 MHz class, \code{cpu_clk_i} 1282 MHz
class on ASIC) & Both exist as ports. Two outputs of \textbf{one} MMCM
are phase-related, so they would not exercise the async crossing the
ASIC has; see 9.4. \\
ASAP7/\allowbreak{}Sky130 hard macros & None remain: every macro site
has an \code{FPGA} arm or is behavioural. \\
\end{longtable}
}

}

\subsection{Device-class sizing (INFERENCE from vendor datasheet
figures as I remember them; verify before buying)}\label{tm:dec:fpga-stage0-evaluation:8-device-class-sizing-inference-from-vendor-datasheet-figures-as-i-remember-them-verify-before-buying}\origsec{8}

Target utilisation from the skill defaults: LUT \textless= 70 \%, BRAM
\textless= 80 \%, DSP \textless= 80 \%. yosys LUT counts are pessimistic
against vendor tools, so these margins are conservative.

{\footnotesize\setlength{\tabcolsep}{3pt}

{
\begin{longtable}[]{@{}
  >{\raggedright\arraybackslash}p{\dimexpr(\linewidth-10\tabcolsep)*252/1000\relax}
  >{\raggedright\arraybackslash}p{\dimexpr(\linewidth-10\tabcolsep)*93/1000\relax}
  >{\raggedright\arraybackslash}p{\dimexpr(\linewidth-10\tabcolsep)*88/1000\relax}
  >{\raggedright\arraybackslash}p{\dimexpr(\linewidth-10\tabcolsep)*55/1000\relax}
  >{\raggedright\arraybackslash}p{\dimexpr(\linewidth-10\tabcolsep)*153/1000\relax}
  >{\raggedright\arraybackslash}p{\dimexpr(\linewidth-10\tabcolsep)*339/1000\relax}@{}}
\toprule\noalign{}
\begin{minipage}[b]{\linewidth}\raggedright
Device (example board)
\end{minipage} & \begin{minipage}[b]{\linewidth}\raggedright
LUT
\end{minipage} & \begin{minipage}[b]{\linewidth}\raggedright
BRAM
\end{minipage} & \begin{minipage}[b]{\linewidth}\raggedright
DSP
\end{minipage} & \begin{minipage}[b]{\linewidth}\raggedright
no-GPU
\end{minipage} & \begin{minipage}[b]{\linewidth}\raggedright
with GPU
\end{minipage} \\
\midrule\noalign{}
\endhead
\bottomrule\noalign{}
\endlastfoot
XC7A35T / 50T (Arty A7-35/50) & 20.8 K / 32.6 K & 50 / 75 RAMB36 & 90 /
120 & \textbf{No} & \textbf{No} \\
XC7A75T & 47.2 K & 105 & 180 & Fits, \textasciitilde82 \% LUT: no margin
& No \\
XC7A100T (Arty A7-100, Nexys A7) & 63.4 K & 135 & 240 & \textbf{Fits,
\textasciitilde61 \% LUT} & No (\textasciitilde106 \% with LUT-RAM) \\
XC7A200T & 134.6 K & 365 & 740 & Fits (\textasciitilde29 \%) &
\textbf{Fits, \textasciitilde50 \%} \\
XC7K325T (Genesys 2, KC705) & 203.8 K & 445 & 840 & Fits
(\textasciitilde19 \%) & Fits (\textasciitilde33 \%) \\
ECP5 LFE5U-25F & 24 K & 56 DP16KD & 28 & \textbf{No} & \textbf{No} \\
ECP5 LFE5U-45F & 43.8 K LUT4 & 108 & 72 & Fits, \textasciitilde78 \%
LUT: no margin & \textbf{No} \\
ECP5 LFE5U-85F (ULX3S, Versa, OrangeCrab 85) & 83.6 K & 208 & 156 &
\textbf{Fits, \textasciitilde41 \%} & Marginal: 67.4 K LUT4 + LUT-RAM
(\textasciitilde8-13 K sites) is \textasciitilde90-95 \%, above the 70
\% target \\
\end{longtable}
}

}

BRAM and DSP are not the limit anywhere: 9 RAMB36-eq of 105+ (7-series),
17 / 49 DP16KD (ECP5), 19-45 multipliers. The GPU is LUT-bound, and
7-series mapped its 16 KB shared memory to LUT-RAM (\textasciitilde5 K
sites) where BRAM would be nearly free.

\subsection{Board options and go / no-go (options for the user, not a
decision)}\label{tm:dec:fpga-stage0-evaluation:9-board-options-and-go-no-go-options-for-the-user-not-a-decision}\origsec{9}

\begin{quote}
\textbf{Superseded for board choice by
\code{FPGA_BOARD_SELECTION.md} (\cref{tm:dec:fpga-board-selection})}
(2026-10-14): growth-adjusted sizing, Singapore
price/\allowbreak{}stock, recommendation (Nexys Video, XC7A200T),
second-clock add-on and fmax tool availability.
\end{quote}

{\footnotesize\setlength{\tabcolsep}{3pt}

{
\begin{longtable}[]{@{}
  >{\raggedright\arraybackslash}p{\dimexpr(\linewidth-6\tabcolsep)*267/1000\relax}
  >{\raggedright\arraybackslash}p{\dimexpr(\linewidth-6\tabcolsep)*99/1000\relax}
  >{\raggedright\arraybackslash}p{\dimexpr(\linewidth-6\tabcolsep)*307/1000\relax}
  >{\raggedright\arraybackslash}p{\dimexpr(\linewidth-6\tabcolsep)*307/1000\relax}@{}}
\toprule\noalign{}
\begin{minipage}[b]{\linewidth}\raggedright
Option
\end{minipage} & \begin{minipage}[b]{\linewidth}\raggedright
Fits
\end{minipage} & \begin{minipage}[b]{\linewidth}\raggedright
Toolchain
\end{minipage} & \begin{minipage}[b]{\linewidth}\raggedright
Note
\end{minipage} \\
\midrule\noalign{}
\endhead
\bottomrule\noalign{}
\endlastfoot
1. ECP5-85F board (ULX3S / Versa) & no-GPU with margin & Fully open:
yosys + nextpnr-ecp5 + prjtrellis + openFPGALoader
(nextpnr/\allowbreak{}prjtrellis not on host: flake change) & Only
option that needs no vendor tool; GPU does not fit with margin \\
2. Artix-7 100T board (Arty A7-100T / Nexys A7-100T) & no-GPU with
margin & Vivado (free WebPACK covers this part); nextpnr-xilinx is
experimental & Cheapest vendor-backed path; GPU does not fit \\
3. Artix-7 200T / Kintex-7 325T (Genesys 2, KC705) & both builds &
Vivado & Only way to prototype CPU + GPU together without shrinking the
GPU \\
\end{longtable}
}

}

Go / no-go as options: \textbf{no-GPU prototype: GO on option 1 or 2}
(measured fit by resource count; fmax and routing unproven).
\textbf{With-GPU prototype: NO-GO on options 1 and 2 as-is}; it needs
either option 3, or a smaller GPU, or moving the GPU shared memory to
BRAM and re-measuring (7-series: removes \textasciitilde5 K LUT sites,
still \textasciitilde62 K LUT, \textasciitilde98 \% of an A7-100T, so
still no).

\subsubsection{Open human decisions}\label{tm:dec:fpga-stage0-evaluation:open-human-decisions}

\begin{enumerate}
\def\labelenumi{\arabic{enumi}.}
\tightlist
\item
  \textbf{Board / vendor family} (options above).
\item
  \textbf{GPU does not fit} on the cheap boards: choose option 3, shrink
  the GPU (lane count, 16 KB shared memory), or ship the CPU-only
  prototype first.
\item
  \textbf{GPU shared memory style on FPGA} (\code{ram_style} BRAM vs
  LUT-RAM; vendor mappers disagree).
\item
  \textbf{Clock-gating strategy} (section 6; recommendation C then A).
\item
  \textbf{Dual-oscillator CDC plan} for later stages: two independent
  oscillators/\allowbreak{}PLLs (needed for the async-crossing soak and
  the 7 ns / 3 ns coprime ratio the ASIC regression used) vs one MMCM
  with two outputs (phase-related, does not stress the CDC). Also the
  XDC/LPF translation of the \#94 SDC and \code{ASYNC_REG}.
\item
  \textbf{Flake change} for
  yosys/\allowbreak{}sv2v/\allowbreak{}nextpnr/\allowbreak{}openFPGALoader
  (an \code{fpga} devShell) vs keeping absolute store paths.
\item
  \textbf{CI job for the FPGA-define regression}: proposal in section 5;
  not added.
\item
  Whether to strip DFT / TAP on the prototype and who owns the top-level
  IO wrapper and \code{boot_rom} image.
\end{enumerate}

\subsection{What was NOT verified}\label{tm:dec:fpga-stage0-evaluation:10-what-was-not-verified}\origsec{10}

No place-and-route, no timing (fmax \textbf{not measured}), no
bitstream, no hardware, no vendor-tool synthesis (so
Vivado/\allowbreak{}Quartus LUT counts will differ), no gate-level or
post-synthesis simulation of the yosys netlist, no
metastability/\allowbreak{}CDC behaviour on FPGA, no power. The
\code{boot_rom} content and IO buffers are outside the counts. Device
sizing is inference. The \code{ecp5} jobs initially aborted in
yosys\textquotesingle s cosmetic \code{autoname} step
(\code{std::bad_alloc} under the 10 GB cap, peak RSS \textasciitilde8.8
GB); the final runs use
\texttt{synth\_\allowbreak{}ecp5\ -\allowbreak{}run\ :\allowbreak{}check}
(every mapping pass runs; only \code{autoname} /
\texttt{hierarchy\ -\allowbreak{}check} / \code{stat} of the
\code{check} label are skipped, and \code{stat} +
\texttt{blackbox\ =\allowbreak{}A:\allowbreak{}whitebox} are run
explicitly). The 7-series runs used the full \code{synth_xilinx} script.
\end{relocated}

\chapter{High-level synthesis}\label{tm:ch:high-level-synthesis}

This chapter describes the high-level synthesis work and its results.

\begin{relocated}{docs/CPP_TO_RTL_HLS_EVALUATION.md}
% Converted from docs/CPP_TO_RTL_HLS_EVALUATION.md by tools/docs/md2tex.py; edit this file, not the original.
\section{C++ → SystemVerilog (HLS) Flow Evaluation \& RTL-Generation
Approach Decision}\label{tm:dec:cpp-to-rtl-hls-evaluation}

\textbf{Date:} 2026-07-14 \textbf{Status:} Decision recorded ·
\textbf{pilot RUN and COMPLETE --- Stage 1 + Stage 2, 2026-09-07} (see
"Stage 1 results" and "Stage 2 results --- the PPA table" at the end)
\textbf{Scope:} Should RTL be generated via
\texttt{NL\ →\ C+\allowbreak{}+\allowbreak{}\ (AI)\ →\ RTL\ (HLS\ tool)}
instead of the current
\texttt{NL\ →\ SystemVerilog\ (RTL-\allowbreak{}orchestrator\ agents)}
flow?

\begin{center}\rule{0.5\linewidth}{0.5pt}\end{center}

\subsection{Recommended Approach (decision)}\label{tm:dec:cpp-to-rtl-hls-evaluation:recommended-approach-decision}

\begin{quote}
\textbf{Use \texttt{NL\ →\ SystemVerilog} (RTL-orchestrator agents,
current flow) for the CPU, caches, AXI interconnect, and GPU control
logic.}

\textbf{Use
\texttt{NL\ →\ C+\allowbreak{}+\allowbreak{}\ →\ HLS\ →\ SystemVerilog}
only for \emph{future, datapath-dominant} accelerators} (e.g. the Phase
6 INT8 NPU, DSP/FFT, or crypto cores) where there is no tuned hand-RTL
to beat and dataflow HLS is strongest.
\end{quote}

Rationale summary (see Part 2 for the per-axis discussion and per-claim
citations): control-heavy, timing-critical, cycle-accurate
microarchitecture (hazard/\allowbreak{}forward unit, cache refill FSMs,
crossbar arbitration/\allowbreak{}handshake, SIMT divergence stack, warp
scheduler) is the regime where the surveyed literature and this
project\textquotesingle s own sign-off experience indicate hand-RTL
beats HLS on quality, PPA, and timing --- treat this as a well-supported
working hypothesis, not a measured result for \emph{this} RTL (that is
what the pilot below would settle). The two-stage flow also stacks two
error sources (LLM spec-drift in NL→C++ \emph{plus} HLS
QoR/\allowbreak{}subset limits) and forfeits the
cycle/\allowbreak{}timing control this project relied on for sign-off
(EX-stage retiming, forwarding-mux placement, per-path SDC
false/\allowbreak{}multicycle exemptions). Our current NL→SV path is
strictly shorter and keeps full control. HLS\textquotesingle s
documented wins are datapath results --- e.g. a memcached dataflow
pipeline reported \textasciitilde70\% register reduction at comparable
performance ({[}arXiv:1408.5387{]}, Xilinx Vivado HLS, FPGA target) ---
which is why HLS is reserved here for datapath accelerators.

\begin{center}\rule{0.5\linewidth}{0.5pt}\end{center}

\subsection{Part 1 --- Available open-source C++/HLS → (System)Verilog
tools}\label{tm:dec:cpp-to-rtl-hls-evaluation:part-1-available-open-source-chls-systemverilog-tools}

Survey as of \textbf{2026-07}; releases are \textbf{not pinned} below
(the pilot in the backlog section must pin exact
commits/\allowbreak{}tags before any measurement). "Front-end coverage"
= which input language/\allowbreak{}subset the tool accepts; "output
target" = HDL(s) it emits. Qualitative notes ("active", "broadest
front-end") are relative editorial judgments from the linked project
pages, not benchmarked scores --- verify against each project at pilot
time.

{\footnotesize\setlength{\tabcolsep}{3pt}

{
\begin{longtable}[]{@{}
  >{\raggedright\arraybackslash}p{\dimexpr(\linewidth-8\tabcolsep)*112/1000\relax}
  >{\raggedright\arraybackslash}p{\dimexpr(\linewidth-8\tabcolsep)*181/1000\relax}
  >{\raggedright\arraybackslash}p{\dimexpr(\linewidth-8\tabcolsep)*151/1000\relax}
  >{\raggedright\arraybackslash}p{\dimexpr(\linewidth-8\tabcolsep)*268/1000\relax}
  >{\raggedright\arraybackslash}p{\dimexpr(\linewidth-8\tabcolsep)*268/1000\relax}@{}}
\toprule\noalign{}
\begin{minipage}[b]{\linewidth}\raggedright
Tool
\end{minipage} & \begin{minipage}[b]{\linewidth}\raggedright
Front-end
\end{minipage} & \begin{minipage}[b]{\linewidth}\raggedright
Output target
\end{minipage} & \begin{minipage}[b]{\linewidth}\raggedright
License / activity (as of 2026-07)
\end{minipage} & \begin{minipage}[b]{\linewidth}\raggedright
Notes
\end{minipage} \\
\midrule\noalign{}
\endhead
\bottomrule\noalign{}
\endlastfoot
\textbf{Bambu (PandA)} & C / C++ (GCC-based subset) & Verilog & GPL;
repo actively maintained & General C/C++ HLS with the widest C/C++ input
coverage of this set; PPA competitive \emph{within the HLS class} per
its own literature. \\
\textbf{Google XLS} & DSLX (Rust-like); C++ via \code{xlscc} subset &
\textbf{Verilog + SystemVerilog} & Apache-2; actively developed &
"Mid-level" synthesis; explicit pipeline-stage / throughput control; one
source → SW model + RTL. Strongest for fixed-latency dataflow. \\
\textbf{SCCL} & SystemC (C++11 subset) & \textbf{SystemVerilog} / VHDL /
FIRRTL & Open source & Clang front-end; SystemC synthesis subset; emits
SystemVerilog directly. \\
\textbf{CIRCT + ScaleHLS / Calyx} & C/C++ → MLIR &
\textbf{SystemVerilog} (via ExportVerilog) & Apache-2; actively
developed & Infrastructure/\allowbreak{}frameworks, not a turnkey
C++→RTL button; strong future direction. \\
\textbf{LegUp} & C & Verilog & Open version dated; maintained path
proprietary (Microchip SmartHLS) & Historically better fuzz-reliability
than early Bambu (per the reliability study cited below). \\
\end{longtable}
}

}

\textbf{Takeaway:} viable open-source C++→RTL tools exist.
\textbf{Multiple} emit SystemVerilog --- \textbf{SCCL} (from SystemC),
\textbf{XLS} (from DSLX / \code{xlscc}), and \textbf{CIRCT} (via
ExportVerilog); \textbf{Bambu} and \textbf{LegUp} emit Verilog. So
SystemVerilog output is \emph{not} a differentiator; the real axes are
front-end (which C/C++ subset) and QoR. \textbf{Bambu} accepts the
widest plain C/C++ input; every tool imposes a \emph{synthesizable
subset} --- none ingests arbitrary modern C++.

\subsection{Part 2 --- Why not for this project\textquotesingle s core
(quality / PPA / timing)}\label{tm:dec:cpp-to-rtl-hls-evaluation:part-2-why-not-for-this-projects-core-quality-ppa-timing}

The claims below are drawn from the cited literature and this
project\textquotesingle s own sign-off data. Where a statement is a
general finding from the literature rather than a measurement of
\emph{this} RTL, it is framed as such --- the backlog pilot is what
would convert these from supported expectations into project-specific
numbers.

\begin{itemize}
\tightlist
\item
  \textbf{Quality:} The surveyed literature reports HLS excels at
  dataflow/\allowbreak{}statically-scheduled algorithms and is weaker
  than hand-RTL on control-intensive logic --- the regime that dominates
  this SoC (semiengineering.com "HLS for RISC-V"; BittWare RTL-vs-HLS
  study). Independent of QoR, the two-stage flow compounds LLM
  spec-drift with HLS subset/\allowbreak{}QoR limits, and the LLM must
  target each tool\textquotesingle s narrow subset (a less-trained
  target than "write SystemVerilog"). \emph{Hypothesis for this RTL,
  pending the pilot.}
\item
  \textbf{PPA:} Reliability/\allowbreak{}QoR studies place open HLS
  \textasciitilde on par with commercial HLS on datapath but behind
  hand-RTL for control blocks (johnwickerson fuzzing-HLS study;
  BittWare). Our \emph{measured} sign-off numbers (CPU 1418 MHz / 27.27
  mW ASAP7; SoC 571 MHz / 62.9 mW) come from hand micro-architecture and
  are the concrete anchor; the expectation that HLS control FSMs add
  registers/\allowbreak{}area and miss these fmax targets is a
  literature-based hypothesis the pilot would test.
\item
  \textbf{Timing:} HLS gives limited timing controllability
  (pragmas/\allowbreak{}target-period, not hand-placed critical paths).
  Full HLS RISC-V cores exist and work (HL5,
  arXiv/\allowbreak{}CICC\textquotesingle20) but are reported below
  tuned hand-RTL fmax. A regression on the paths that gate this design
  is expected but unproven for this RTL --- again, the
  pilot\textquotesingle s purpose.
\item
  \textbf{Process cost:} New toolchain + subset + pragma tuning +
  mandatory C↔RTL equivalence checking (HLS tools have documented
  reliability bugs --- johnwickerson study) --- a large new verification
  surface on top of the existing cocotb/\allowbreak{}pyuvm flow.
\end{itemize}

\subsubsection{References}\label{tm:dec:cpp-to-rtl-hls-evaluation:references}

\begin{itemize}
\tightlist
\item
  PandA-Bambu ---
  \href{https://github.com/ferrandi/PandA-bambu}{https:/\allowbreak{}/\allowbreak{}github.com/\allowbreak{}ferrandi/\allowbreak{}PandA-bambu}
\item
  Google XLS ---
  \href{https://github.com/google/xls}{https:/\allowbreak{}/\allowbreak{}github.com/\allowbreak{}google/\allowbreak{}xls}
\item
  SCCL: An open-source SystemC to RTL translator (IEEE) ---
  \href{https://ieeexplore.ieee.org/document/10171476/}{https:/\allowbreak{}/\allowbreak{}ieeexplore.ieee.org/\allowbreak{}document/\allowbreak{}10171476/}
\item
  An Empirical Study of the Reliability of HLS Tools ---
  \href{https://johnwickerson.github.io/papers/fuzzingHLS.pdf}{https:/\allowbreak{}/\allowbreak{}johnwickerson.github.io/\allowbreak{}papers/\allowbreak{}fuzzingHLS.pdf}
\item
  High-Level Synthesis for RISC-V (Semiconductor Engineering) ---
  \href{https://semiengineering.com/high-level-synthesis-for-risc-v/}{https:/\allowbreak{}/\allowbreak{}semiengineering.com/\allowbreak{}high-level-synthesis-for-risc-v/}
\item
  HL5: A 32-bit RISC-V Processor Designed with HLS ---
  \href{https://sld.cs.columbia.edu/pubs/mantovani_cicc20.pdf}{https:/\allowbreak{}/\allowbreak{}sld.cs.columbia.edu/\allowbreak{}pubs/\allowbreak{}mantovani\_\allowbreak{}cicc20.pdf}
\item
  Comparing FPGA RTL to HLS C/C++ (BittWare) ---
  \href{https://www.bittware.com/resources/comparing-rtl-to-hls/}{https:/\allowbreak{}/\allowbreak{}www.bittware.com/\allowbreak{}resources/\allowbreak{}comparing-rtl-to-hls/}
\item
  HLS Case Study: Memcached Server (\textasciitilde70\% register
  reduction claim) ---
  \href{https://arxiv.org/abs/1408.5387}{https:/\allowbreak{}/\allowbreak{}arxiv.org/\allowbreak{}abs/\allowbreak{}1408.5387}
\item
  CIRCT Verilog/\allowbreak{}SystemVerilog generation (ExportVerilog)
  ---
  \href{https://circt.llvm.org/docs/VerilogGeneration/}{https:/\allowbreak{}/\allowbreak{}circt.llvm.org/\allowbreak{}docs/\allowbreak{}VerilogGeneration/}
\end{itemize}

\begin{center}\rule{0.5\linewidth}{0.5pt}\end{center}

\subsection{Backlog item --- Empirical confirmation pilot}\label{tm:dec:cpp-to-rtl-hls-evaluation:backlog-item-empirical-confirmation-pilot}

\begin{quote}
\textbf{Tracked as GitHub issue
\href{https://github.com/chuanseng-ng/claude_verilog_test/issues/119}{\#119}.}
Filed on GitHub rather than beads (\code{bd}) because the remote beads
DB is at schema v32 while the available \code{bd} is v53; writing a
native \code{bd} issue would require migrating the shared schema (forces
every clone to re-bootstrap), a team-wide coordination decision left to
the maintainer. Convert to a \code{bd} issue later if desired.
\end{quote}

\textbf{Title:} Empirically confirm HLS-vs-hand-RTL PPA gap on one
datapath + one control block \textbf{Type:} task · \textbf{Priority:}
low (validation, not blocking) · \textbf{Phase:} 6+ exploration

\textbf{Goal:} Quantify (not just argue) the QoR gap between the two
flows on an already-signed-off block, to confirm the recommendation
above.

\textbf{Method (apples-to-apples on ASAP7):}

\begin{enumerate}
\def\labelenumi{\arabic{enumi}.}
\tightlist
\item
  Pick two poles from existing signed-off RTL:

  \begin{itemize}
  \tightlist
  \item
    \textbf{Memory coalescer} (\src{rtl/gpu/}, datapath-ish →
    HLS-favorable).
  \item
    A \textbf{cache refill FSM} (\src{rtl/mem/}, control-heavy →
    HLS-unfavorable).
  \end{itemize}
\item
  \textbf{Flow A (baseline):} existing hand-RTL through the current
  ASAP7 sign-off.
\item
  \textbf{Flow B:}
  \texttt{NL\ →\ C+\allowbreak{}+\allowbreak{}\ (LLM)\ →\ Bambu} (widest
  C/C++ FE) and/or \textbf{XLS} (SV output) → same
  \texttt{make\ librelane-\allowbreak{}asap7} sign-off.
\item
  \textbf{Verification --- two independent checks, both required:}

  \begin{itemize}
  \tightlist
  \item
    \textbf{Source→RTL equivalence:} C++ ↔ generated-RTL equivalence via
    the HLS tool\textquotesingle s own C/RTL co-simulation \emph{or}
    formal equivalence (e.g. the tool\textquotesingle s cosim harness).
    This is the check that actually validates the HLS step, and is the
    known gap flagged in the project\textquotesingle s existing
    source-to-RTL-equivalence backlog item.
  \item
    \textbf{Reference agreement:} generated RTL vs the Python reference
    model (\src{tb/models/}) as an independent oracle. Neither check
    substitutes for the other.
  \end{itemize}
\item
  Then compare \textbf{fmax / area / power / DRC} against Flow A.
\end{enumerate}

\textbf{Reproducibility --- pin before measuring (else the comparison
isn\textquotesingle t apples-to-apples):} Bambu/XLS release commit or
tag; LibreLane + ASAP7 PDK revision; the exact SDC constraints (clock
period, IO delays, false/\allowbreak{}multicycle paths) used for
\emph{both} flows; PVT/\allowbreak{}process corner (match the current
ASAP7 sign-off corner); clock/\allowbreak{}reset/\allowbreak{}IO wrapper
the generated block is placed in; power-activity assumptions (VCD-driven
vs default switching); and latency/\allowbreak{}throughput normalization
(HLS may pick a different cycle count --- normalize on throughput at
matched fmax, and report the II/\allowbreak{}latency of each version).

\textbf{Acceptance thresholds (measurable; tune before the run if
desired):}

\begin{itemize}
\tightlist
\item
  "Regresses noticeably" = Flow B is worse than Flow A by
  \textbf{\textgreater10 \% fmax} \emph{or} \textbf{\textgreater20 \%
  area} \emph{or} \textbf{\textgreater20 \% power} at matched function.
  \sout{\emph{or} introduces any DRC/\allowbreak{}antenna violation Flow
  A did not have.} --- \textbf{DRC clause struck, see below.}
\item
  "Competitive" = Flow B within \textbf{±5 \% fmax} and \textbf{±10 \%
  area/\allowbreak{}power} of Flow A. \sout{with 0 DRC.}
\item
  Any result between these bands is "inconclusive → widen the block set
  or re-tune pragmas."
\end{itemize}

\begin{quote}
\textbf{\warn{} The DRC/\allowbreak{}antenna axis is withdrawn for ASAP7
--- struck 2026-09-07 (Stage 1, bead \bead{cge}).} Bead \code{xy6}
(closed, root-caused) established that \textbf{no ASAP7 run in this
project\textquotesingle s history has ever completed detailed routing}:
\code{pnr/scripts/openroad/drt.tcl} wraps the whole
\code{detailed_route} call in a bare \texttt{catch\{\}}, so
DRT-0073/\allowbreak{}0074 pin-access failures abort routing and
\code{write_views} commits the pre-route placed netlist. 100 \% of
recoverable runs show zero ROUTED nets. Bead \code{ocm} (open) owns the
unsolved pin-access cause. Both arms of this pilot would hit the
identical wall, so a DRC column yields \textbf{no discriminating signal}
--- it would compare two vacuous zeros.
Timing/\allowbreak{}power/\allowbreak{}area survive as
GRT-estimate-based figures that are self-consistent between arms.
\textbf{Stage 2 must not promise a DRC column.} Revisit only if
\code{ocm} is solved.
\end{quote}

\textbf{Expected result (hypothesis):} HLS regresses per the thresholds
above on the FSM, is competitive on the coalescer --- which would
confirm "keep hand-RTL for the control core, reserve HLS for datapath
accelerators." A surprising result (HLS competitive on the FSM) would
itself be a useful signal to revisit the recommendation.

\textbf{Deliverable:} a short PPA comparison table appended to this doc.

\textbf{Note:} this pilot was optional. \textbf{It was run anyway ---
Stage 1 completed 2026-09-07; see "Stage 1 results" at the end of this
document.} Two of the three blocks returned categorical results rather
than PPA numbers, and one correction landed on this
section\textquotesingle s own framing (the "datapath-ish" pole is not a
datapath block). Stage 2 (P\&R + the PPA table) has not been run.

\begin{center}\rule{0.5\linewidth}{0.5pt}\end{center}

\subsection{Stage 1 results --- the pilot was run (2026-09-07)}\label{tm:dec:cpp-to-rtl-hls-evaluation:stage-1-results-the-pilot-was-run-2026-09-07}

\textbf{Status:} Stage 1 complete. \textbf{No PPA table yet} --- that is
Stage 2, which measures. Stage 1 installed the tool, authored the C,
generated the RTL, and proved functional equivalence.

The pilot was expected to return a PPA delta. It returned something more
decisive: on the three blocks attempted, \textbf{two of the three
results are categorical rather than numeric.} Details in
\code{hls/*/NOTES.md}; the reproducibility pin is
\src{hls/PROVENANCE.json}.

\subsection{What was pinned}\label{tm:dec:cpp-to-rtl-hls-evaluation:what-was-pinned}

Bambu (PandA) \textbf{2024.10}, revision
\code{c2ba6936ca2ed63137095fea0b630a1c66e20e63-main}. Upstream ships
only a prebuilt AppImage, so it is packaged as a nix devshell
(\src{flake.nix}, \code{devShells.hls}) via
\code{appimageTools.wrapType2} + \code{fetchurl} with the AppImage
sha256 enforced by nix. \textbf{The flake is the pin} --- no per-machine
install step, and nix refuses to build if the upstream artefact ever
changes.

Target \code{--device-name=asap7-TC}, \code{--clock-period=0.705} ns
(matching \src{pnr/asap7/template/config.json}). Bambu bundles
\code{asap7sc7p5t_SIMPLE_RVT_TT_nldm_201020.lib} --- the same cell
library and TT corner \src{pnr/asap7/cpu/config.json} already uses, so
Stage 2 can match the corner exactly. Co-simulation against the
flake\textquotesingle s Verilator 5.048.

Generated Verilog is \textbf{not committed}. Bambu stamps a timestamp
into a header comment of every \code{.v}, so the raw digest changes
every run (verified: three back-to-back runs of identical input gave
three digests, byte-identical except the
\texttt{-\allowbreak{}\ Date\ …} line). \src{tools/eda/wrap-bambu.sh}
therefore emits a \textbf{normalized} digest, which is stable and is
what \code{PROVENANCE.json} pins.

\subsection{Method: the reference RTL body was withheld}\label{tm:dec:cpp-to-rtl-hls-evaluation:method-the-reference-rtl-body-was-withheld}

Each C source was authored by a fresh agent given \textbf{only} the
block\textquotesingle s natural-language spec (the reference \code{.sv}
header comment) and its port list, explicitly barred from reading
anything under \src{rtl/}. Transliterating the SystemVerilog would have
measured transliteration and answered nothing. Each
block\textquotesingle s \code{ASSUMPTIONS.md} records every point where
the spec was silent and what was chosen --- those files are a measured
output of the experiment, not paperwork.

\subsection{Results}\label{tm:dec:cpp-to-rtl-hls-evaluation:results}

{\small\setlength{\tabcolsep}{4pt}

{
\begin{longtable}[]{@{}
  >{\raggedright\arraybackslash}p{\dimexpr(\linewidth-4\tabcolsep)*222/1000\relax}
  >{\raggedright\arraybackslash}p{\dimexpr(\linewidth-4\tabcolsep)*230/1000\relax}
  >{\raggedright\arraybackslash}p{\dimexpr(\linewidth-4\tabcolsep)*528/1000\relax}@{}}
\toprule\noalign{}
\begin{minipage}[b]{\linewidth}\raggedright
Block
\end{minipage} & \begin{minipage}[b]{\linewidth}\raggedright
Pole
\end{minipage} & \begin{minipage}[b]{\linewidth}\raggedright
Outcome
\end{minipage} \\
\midrule\noalign{}
\endhead
\bottomrule\noalign{}
\endlastfoot
\code{memory_coalescer} & AXI / serialiser & Equivalent on both legs.
One measured latency penalty. \\
\code{rv32i_hazard_unit} & control & \textbf{Form inexpressible}
(combinational → 6-9 cycle FSM) \textbf{and} functional divergence,
254/2000 vectors \\
\code{rv32i_cache_arbiter} & control / bus arbitration &
\textbf{Inexpressible --- no C authored} \\
\end{longtable}
}

}

Every block was checked on \textbf{two independent legs}:
Bambu\textquotesingle s own C/RTL co-simulation (which validates the
C→RTL step), and the \emph{same, unmodified} cocotb suite run against
both the hand-RTL and the HLS arm behind a wire-only shim (which
validates the C against the spec). Neither leg substitutes for the
other.

\subsubsection{\texorpdfstring{\code{memory_coalescer} ---
equivalent, with a 6× control-path latency
penalty}{--- equivalent, with a 6× control-path latency penalty}}\label{tm:dec:cpp-to-rtl-hls-evaluation:1-memory-coalescer-equivalent-with-a-6-control-path-latency-penalty}\origsec{1}

Both legs green: cosim 5/5 vectors; the existing 5-test cocotb suite
passes on both arms.

One test initially failed. Measured cause: on an empty lane mask the
hand-RTL asserts \code{done_o} after \textbf{1} rising edge
(\texttt{IDLE\ →\ DONE} directly), the HLS version after \textbf{6} ---
Bambu walks all eight lane predicates sequentially even when none fire.
The test\textquotesingle s 5-cycle bound was a testbench assumption
calibrated against the reference, not a specified requirement (the spec
sets no cycle budget), so it was rebased for \textbf{both} arms with the
AXI-silence assertions untouched.

\textbf{Correction to this document\textquotesingle s own framing:}
\code{memory_coalescer} was named here as the "datapath-ish →
HLS-favorable" pole. It is not a datapath block. It is 214 lines of
lane-walking FSM plus AXI handshakes with zero arithmetic, and
\texttt{GPU\_\allowbreak{}ENABLE\_\allowbreak{}COALESCE\ =\allowbreak{}\ 1'b0}
means nothing is actually coalesced --- it serialises one single-beat
transaction per active lane. It was kept because it is the block that
exercises Bambu\textquotesingle s \code{m_axi} path and carries the real
interface risk, but it never tested the datapath hypothesis.

\subsubsection{\texorpdfstring{\code{rv32i_hazard_unit} --- the
control pole, substituted and doubly
negative}{--- the control pole, substituted and doubly negative}}\label{tm:dec:cpp-to-rtl-hls-evaluation:2-rv32i-hazard-unit-the-control-pole-substituted-and-doubly-negative}\origsec{2}

Substituted for "a cache refill FSM", which turned out not to exist as a
standalone module --- the refill logic is smeared across
\code{rv32i_icache.sv} (619 lines, 5 SRAM macros) and
\code{rv32i_dcache.sv} (990 lines). Carving it out would have created a
\emph{new} baseline that was never signed off. The hazard unit is 427
lines, purely combinational, and is named in this
document\textquotesingle s own control-heavy list, so it tests the same
hypothesis against real signed-off RTL.

\textbf{Result 2a --- HLS cannot express the block\textquotesingle s
form.} The reference has no clock, no reset and \textbf{zero
flip-flops}. Bambu wraps every design in a start/\allowbreak{}done
handshake: \textbf{9 FSM states, 50 flip-flops, 6-9 cycles per
evaluation, data-dependent latency.} For a block whose job is resolving
hazards within one cycle, that is not a drop-in replacement at any clock
period. This is not an authoring artefact --- rebuilding the
stall/\allowbreak{}flush logic branch-free, with no priority chain at
all, produced an \emph{identical} schedule (area within 0.3 \%).

\textbf{Result 2b --- spec drift, precisely localised.} Hand-RTL 28/28
spec-derived directed tests; HLS 26/28; the 2000-vector cross-arm
differential reports \textbf{254 mismatches (12.7 \%)}. Root cause: the
failing vectors make the EX1b/EX1c producers \emph{loads}, whose data is
not ready, so those forwarding tiers must be excluded; the HLS version
gates only the EX2 tier on \code{mem_rd}.

The sharp part: \textbf{all 254 mismatches land on forwarding outputs.
Not one stall, flush or load-use output ever mismatched.} The 9-item
priority ladder --- which this document\textquotesingle s source spec
states explicitly and completely --- transferred perfectly. The drift is
confined entirely to what the spec left to inference.
\textbf{Spec-stated behaviour transferred; spec-implied behaviour did
not.}

The divergence was deliberately not fixed. It is the measurement. It was
also \emph{predicted in advance} by the C author in
\code{ASSUMPTIONS.md}, including the observation that it would "only
ever fail under random-vector equivalence, never under a trace" ---
which is exactly what happened.

\subsubsection{\texorpdfstring{\code{rv32i_cache_arbiter} --- outside
the expressible domain}{--- outside the expressible domain}}\label{tm:dec:cpp-to-rtl-hls-evaluation:3-rv32i-cache-arbiter-outside-the-expressible-domain}\origsec{3}

No C was authored, and none should be. Bambu emits AXI \textbf{masters}
only, and only as a side-effect of pointer dereference; asking for a
slave is rejected
(\texttt{error:\allowbreak{}\ Invalid\ HLS\ interface\ mode} for
\texttt{mode\ =\allowbreak{}\ s\_\allowbreak{}axilite}), the shipped
clang plugin contains no AXI-slave mode string at all, and
\code{--generate-interface} accepts only \code{MINIMAL} / \code{INFER} /
\code{WB4}. Under the first two every top module carries
\code{start_port}/\code{done_port} --- the C execution model made
structural. The arbiter is permanently reactive, with three incoming AXI
channel sets and grant-hold-until-\code{RLAST}; there is no invocation
to map onto. Full evidence, including a probe of the one slave-shaped
option (\code{WB4}, a memory-mapped \emph{control} interface in
Wishbone), is in \src{hls/mem/rv32i_cache_arbiter/INEXPRESSIBLE.md}.

Forcing it would have required a hand-written shim containing the grant
register, priority encoder, burst-hold logic and response routing ---
the entire design --- around an HLS core contributing nothing. That
comparison would be hand-RTL versus hand-RTL, reported as hand-RTL
versus HLS.

\subsection{What this does to the recommendation}\label{tm:dec:cpp-to-rtl-hls-evaluation:what-this-does-to-the-recommendation}

The recommendation at the top of this document \textbf{stands, and is
strengthened on its control-logic half.} For that half the claim can now
be stated more strongly than the literature supported:

\begin{quote}
Not "HLS produces worse control logic" but, for the two control blocks
attempted, \textbf{"HLS cannot express these blocks at all"} --- one
because a combinational block becomes a multi-cycle FSM, the other
because the interface has no representation. Neither is a scheduling or
pragma-tuning problem; both follow from the C execution model, so
neither is fixed by re-tuning.
\end{quote}

\textbf{The datapath half of the recommendation remains untested.}
Neither measured block is a datapath block, and the one named as such is
not. Nothing here supports or undermines reserving HLS for datapath
accelerators --- that hypothesis is exactly as evidenced as it was
before this pilot ran. If it matters, a genuine datapath pole
(\src{rtl/gpu/vector_alu.sv}, 8 × 32-bit multiply plus
shift/\allowbreak{}logic) would test it.

Two secondary findings worth carrying into any future HLS work:

\begin{itemize}
\tightlist
\item
  \textbf{Ordinary C factoring is not available.} Bambu refuses to
  inline \code{static} helpers --- it made a shared helper a 1-resource
  submodule with internal DISTRAM and \emph{failed the build}
  (\texttt{clock\ constraint\ too\ tight:\allowbreak{}\ BRAMs\ for\ this\ device\ cannot\ run\ so\ fast}).
  \code{__attribute__((always_inline))} did not change the decision;
  only preprocessor macros worked.
\item
  \textbf{The output is structurally explosive.} Bambu instantiates
  \textasciitilde550 modules for one small block. Verilator then emits a
  1.2 MB symbol-table constructor on which g++ spun at 99.9 \% CPU for
  32 minutes at \code{-O0}. \src{sim/Makefile} carries
  \texttt{-\allowbreak{}-\allowbreak{}output-\allowbreak{}split\ 500}
  and \code{NO_TRACE=1} on the HLS arms for this reason.
\end{itemize}

\subsection{Stage 2 --- what it may and may not claim}\label{tm:dec:cpp-to-rtl-hls-evaluation:stage-2-what-it-may-and-may-not-claim}

\begin{itemize}
\tightlist
\item
  \textbf{No DRC column.} See the struck threshold above (beads
  \code{xy6} / \code{ocm}).
\item
  Bambu\textquotesingle s own estimates (coalescer 820 FF / area 352
  744; hazard unit 50 FF / area 1 791 535) are \textbf{not comparable to
  hand-RTL} --- they cover the whole generated hierarchy including
  Bambu\textquotesingle s AXI master adapter, which the hand-RTL has no
  equivalent of. Stage 2 must synthesise both arms for real and state
  the counting boundary before quoting anything.
\item
  Both shims are wire-only (zero registers), so Flow B\textquotesingle s
  area needs no separate shim line.
\item
  Normalize on throughput per the pinning list above: the hand-RTL cycle
  counts are 1 (coalescer, empty mask) and \textbf{0} (hazard unit,
  combinational) against 6 and 6-9 respectively.
\end{itemize}

\begin{center}\rule{0.5\linewidth}{0.5pt}\end{center}

\subsection{Stage 2 results --- the PPA table (2026-09-07)}\label{tm:dec:cpp-to-rtl-hls-evaluation:stage-2-results-the-ppa-table-2026-09-07}

\textbf{Status:} Stage 2 complete (bead \bead{r8r}). This is the
comparison table this document names as the pilot\textquotesingle s
deliverable. Reproducibility pin: \src{hls/PROVENANCE.json}; configs
under
\texttt{pnr/\allowbreak{}asap7/\allowbreak{}\{coalescer,\allowbreak{}hazard\}\_\allowbreak{}\{rtl,\allowbreak{}hls\}/\allowbreak{}};
collector \src{tools/verif/collect_r8r_ppa.py}.

\subsection{Method}\label{tm:dec:cpp-to-rtl-hls-evaluation:method}

Four ASAP7 runs --- two hand-RTL, two Bambu-HLS --- \textbf{to
\code{OpenROAD.STAPrePNR} only}. That step is where area, timing
\emph{and} power all first appear, so placement, CTS and routing are
needed for none of them. All four use the same clock period (705 ps),
the same corner (\code{nom_tt_025C_0p7V}), the same synthesis strategy
and the same skip list; a diff of the four \code{config.json} differs
\textbf{only} in \code{DESIGN_NAME}, \code{VERILOG_FILES} and the clock
port. Both HLS arms are synthesised \emph{with} their wire-only shim, so
all four have identical port lists.

\subsection{The table}\label{tm:dec:cpp-to-rtl-hls-evaluation:the-table}

\begin{landscape}
\small\setlength{\tabcolsep}{4pt}

{
\begin{longtable}[]{@{}
  >{\raggedright\arraybackslash}p{\dimexpr(\linewidth-14\tabcolsep)*214/1000\relax}
  >{\raggedright\arraybackslash}p{\dimexpr(\linewidth-14\tabcolsep)*89/1000\relax}
  >{\raggedright\arraybackslash}p{\dimexpr(\linewidth-14\tabcolsep)*89/1000\relax}
  >{\raggedright\arraybackslash}p{\dimexpr(\linewidth-14\tabcolsep)*62/1000\relax}
  >{\raggedright\arraybackslash}p{\dimexpr(\linewidth-14\tabcolsep)*205/1000\relax}
  >{\raggedright\arraybackslash}p{\dimexpr(\linewidth-14\tabcolsep)*89/1000\relax}
  >{\raggedright\arraybackslash}p{\dimexpr(\linewidth-14\tabcolsep)*89/1000\relax}
  >{\raggedright\arraybackslash}p{\dimexpr(\linewidth-14\tabcolsep)*142/1000\relax}@{}}
\toprule\noalign{}
\begin{minipage}[b]{\linewidth}\raggedright
Block
\end{minipage} & \begin{minipage}[b]{\linewidth}\raggedright
Flow
\end{minipage} & \begin{minipage}[b]{\linewidth}\raggedright
Area µm²
\end{minipage} & \begin{minipage}[b]{\linewidth}\raggedright
Cells
\end{minipage} & \begin{minipage}[b]{\linewidth}\raggedright
Critical path ps
\end{minipage} & \begin{minipage}[b]{\linewidth}\raggedright
fmax MHz
\end{minipage} & \begin{minipage}[b]{\linewidth}\raggedright
Power mW
\end{minipage} & \begin{minipage}[b]{\linewidth}\raggedright
Undriven wires
\end{minipage} \\
\midrule\noalign{}
\endhead
\bottomrule\noalign{}
\endlastfoot
\code{memory_coalescer} & hand-RTL & 569.74 & 4219 & 1620.15 & 617.2 &
4.469 & \textbf{0} \\
\code{memory_coalescer} & \textbf{HLS} & 605.22 & 4814 &
\textbf{1059.09} & \textbf{944.2} & 4.845 & \textbf{9} \\
\code{rv32i_hazard_unit} & hand-RTL & 23.17 & 233 & \textbf{149.3}
(combinational) & n/a & \textbf{0.018} & \textbf{0} \\
\code{rv32i_hazard_unit} & \textbf{HLS} & 26.55 & 233 & 302.53 (per
cycle) & 3305.5 & \textbf{0.182} & \textbf{28} \\
\end{longtable}
}

\end{landscape}

\code{rv32i_hazard_unit}\textquotesingle s hand-RTL arm is purely
combinational --- no clock, no registers, zero register-to-register
paths --- so "fmax" is undefined for it. It is timed against a
\textbf{virtual clock} with a zero I/O budget, and its number is the
input-to-output path delay (705 − 555.7 slack).

\subsection{\texorpdfstring{Result 1 --- the coalescer: HLS \emph{won}
on speed, and that was not the
hypothesis}{Result 1 --- the coalescer: HLS won on speed, and that was not the hypothesis}}\label{tm:dec:cpp-to-rtl-hls-evaluation:result-1-the-coalescer-hls-won-on-speed-and-that-was-not-the-hypothesis}

Against this document\textquotesingle s own thresholds (\textgreater10
\% fmax or \textgreater20 \% area/\allowbreak{}power = "regresses"; ±5
\% fmax and ±10 \% area/\allowbreak{}power = "competitive"):

\begin{itemize}
\tightlist
\item
  \textbf{Critical path −34.6 \%} (1620 → 1059 ps), i.e. \textbf{fmax
  +53 \%}, in the HLS arm\textquotesingle s favour.
\item
  \textbf{Area +6.2 \%}, \textbf{power +8.4 \%} --- both inside the ±10
  \% "competitive" band.
\end{itemize}

So on raw PPA the HLS coalescer is \textbf{better than competitive}. The
reason is visible in the structure: Bambu split the lane walk across
more, shorter cycles, while the hand-RTL keeps a longer combinational
path through the 8-lane mux and its AXI FSM.

\textbf{This does not overturn the recommendation, and the reason is
throughput.} A shorter clock period bought with more cycles per
transaction is only a win if the cycle count holds, and it does not: on
the empty-mask case the hand-RTL completes in \textbf{1} rising edge and
the HLS arm in \textbf{6}. Per-transaction latency, not fmax, is what
this block\textquotesingle s consumer sees.

\emph{Honest gap:} a full 8-lane-transaction cycle count for the
hand-RTL arm was \textbf{not measured}, so the throughput comparison is
anchored only on the empty-mask case plus Bambu\textquotesingle s own
cosim average of 23 cycles/\allowbreak{}vector. A rigorous throughput
ratio needs that measurement; it is not claimed here.

\subsection{Result 2 --- the hazard unit: regression on every axis that
matters}\label{tm:dec:cpp-to-rtl-hls-evaluation:result-2-the-hazard-unit-regression-on-every-axis-that-matters}

\begin{itemize}
\tightlist
\item
  \textbf{Power ×10.1} (0.018 → 0.182 mW, +911 \%) --- far outside the
  \textgreater20 \% "regresses" threshold, and the single starkest
  number in this pilot.
\item
  \textbf{Area +14.6 \%} --- past the ±10 \% competitive band.
\item
  \textbf{Latency ×12.2 -- ×18.2}: 149.3 ps of combinational delay
  versus 302.53 ps × 6--9 cycles = 1815--2723 ps for the same
  evaluation.
\end{itemize}

\textbf{Verdict: regresses}, decisively. This is the block Stage 1
already showed to be \emph{inexpressible in form} --- a zero-flip-flop
combinational cloud becoming a 9-state, 33-flop FSM --- and the PPA
numbers now put a cost on that: an order of magnitude more power and
over an order of magnitude more latency, for logic the pipeline needs
resolved within a single cycle.

\subsection{Result 3 --- netlist quality: undriven wires, only in the
HLS arms}\label{tm:dec:cpp-to-rtl-hls-evaluation:result-3-netlist-quality-undriven-wires-only-in-the-hls-arms}

Yosys\textquotesingle{} own \code{CHECK} pass reports wires that are
used but never driven:

{
\begin{longtable}[]{@{}lll@{}}
\toprule\noalign{}
& hand-RTL & HLS \\
\midrule\noalign{}
\endhead
\bottomrule\noalign{}
\endlastfoot
\code{memory_coalescer} & 0 & \textbf{9} \\
\code{rv32i_hazard_unit} & 0 & \textbf{28} \\
\end{longtable}
}

They are \code{OUT_UNBOUNDED_*} functional-unit outputs plus an undriven
\code{s_start_port0}. Both hand-RTL arms are clean. The
flow\textquotesingle s \code{Checker.YosysSynthChecks} had to be skipped
\textbf{on all four} (to keep the flow identical) for the HLS arms to
synthesise at all; the skip is documented in \src{pnr/Makefile} and this
asymmetry is reported rather than hidden by it.

\subsection{What must not be read into this table}\label{tm:dec:cpp-to-rtl-hls-evaluation:what-must-not-be-read-into-this-table}

\begin{itemize}
\tightlist
\item
  \textbf{No DRC column, by construction.} Struck in Stage 1 (beads
  \code{xy6}/\code{ocm}): no ASAP7 run in this project has ever
  completed detailed routing, so both arms would compare two vacuous
  zeros. These are partial runs and never route.
\item
  \textbf{Pre-placement, zero-parasitic STA.} Absolute fmax is
  optimistic for all four. The \emph{comparison} holds because every arm
  is treated identically, but these are not sign-off numbers.
\item
  The 233-cell count matching across the two hazard arms is a
  coincidence, not an error --- the areas and flop counts differ (the
  HLS arm has 33 flip-flops; the hand-RTL arm has none).
\end{itemize}

\subsection{Effect on the recommendation}\label{tm:dec:cpp-to-rtl-hls-evaluation:effect-on-the-recommendation}

The recommendation at the top of this document \textbf{stands}, now with
numbers behind its control-logic half rather than only literature: on
the control block, HLS costs \textbf{10× the power and 12--18× the
latency}, on top of Stage 1\textquotesingle s finding that it cannot
reproduce the block\textquotesingle s combinational form at all.

The coalescer result is the honest complication, and it is worth stating
plainly: \textbf{on an AXI-serialiser workload, HLS produced a shorter
critical path at \textasciitilde6 \% area cost.} That is a point in
HLS\textquotesingle s favour on a block that is \emph{not}
control-dominated in the hazard-unit sense, and it is consistent with
reserving HLS for throughput-oriented dataflow rather than for
cycle-critical control.

\sout{\textbf{The datapath half of the recommendation remains untested}
--- neither measured block is a datapath block (see the Stage 1
correction above). Bead \texttt{gg8} tracks
\texttt{rtl/gpu/vector\_alu.sv} as the block that would actually test
it.}

\begin{quote}
\textbf{Superseded by Stage 3 (bead \bead{gg8}, below).} It was tested.
It does not hold.
\end{quote}

\begin{center}\rule{0.5\linewidth}{0.5pt}\end{center}

\subsection{\texorpdfstring{Stage 3 results --- the datapath pole
(2026-09-07, bead
\code{gg8})}{Stage 3 results --- the datapath pole (2026-09-07, bead )}}\label{tm:dec:cpp-to-rtl-hls-evaluation:stage-3-results-the-datapath-pole-2026-09-07-bead-gg8}

Stage 2 closed by admitting that the pilot\textquotesingle s
\emph{datapath} half rested on nothing: \code{memory_coalescer}, the
block chosen as the "datapath-ish → HLS-favorable" pole, is an AXI
transaction FSM. Stage 3 measures the real thing ---
\src{rtl/gpu/vector_alu.sv}, 8 lanes × 32 bits of pure combinational
arithmetic with \textbf{zero registers}.

Same protocol throughout: the C was authored from the natural-language
spec with the hand-RTL withheld
(\src{hls/gpu/vector_alu/ASSUMPTIONS.md}), then both arms taken through
identical ASAP7 runs to \code{OpenROAD.STAPrePNR} --- same 705 ps clock,
same \code{nom_tt_025C_0p7V} corner, same skip list. The HLS arm is
synthesised with its wire-only shim.

\subsection{The table}\label{tm:dec:cpp-to-rtl-hls-evaluation:the-table-1}

{
\begin{longtable}[]{@{}
  >{\raggedright\arraybackslash}p{\dimexpr(\linewidth-8\tabcolsep)*215/1000\relax}
  >{\raggedleft\arraybackslash}p{\dimexpr(\linewidth-8\tabcolsep)*280/1000\relax}
  >{\raggedleft\arraybackslash}p{\dimexpr(\linewidth-8\tabcolsep)*258/1000\relax}
  >{\raggedright\arraybackslash}p{\dimexpr(\linewidth-8\tabcolsep)*108/1000\relax}
  >{\raggedright\arraybackslash}p{\dimexpr(\linewidth-8\tabcolsep)*118/1000\relax}@{}}
\toprule\noalign{}
\begin{minipage}[b]{\linewidth}\raggedright
Metric
\end{minipage} & \begin{minipage}[b]{\linewidth}\raggedleft
hand-RTL \code{vector_alu}
\end{minipage} & \begin{minipage}[b]{\linewidth}\raggedleft
HLS \code{vector_alu_hls}
\end{minipage} & \begin{minipage}[b]{\linewidth}\raggedright
Δ
\end{minipage} & \begin{minipage}[b]{\linewidth}\raggedright
Band
\end{minipage} \\
\midrule\noalign{}
\endhead
\bottomrule\noalign{}
\endlastfoot
Area (µm²) & 4 655.15 & 7 234.00 & \textbf{+55.4 \%} & regresses \\
Cells & 54 634 & 70 895 & +29.8 \% & --- \\
Sequential cells & \textbf{0} & 9 412 & --- & --- \\
Power (mW) & 14.147 & 37.281 & \textbf{+163.5 \%} & regresses \\
Critical path & 759.29 ps (comb, in→out) & 4 142.52 ps (FF→FF) &
\textbf{×5.5} & regresses \\
Cycles / operation & 1 (combinational) & mean 14.35, range 8--38 &
×14--38 & regresses \\
Undriven wires & \textbf{0} & 16 & --- & --- \\
\end{longtable}
}

Every axis regresses, by margins several times this
document\textquotesingle s own thresholds (\textgreater10 \% fmax or
\textgreater20 \% area/\allowbreak{}power). Throughput --- the number
that actually matters for a datapath --- is 759.29 ps for one operation
against 4 142.52 ps × 14.35 = \textbf{59.4 ns} mean, 157.4 ns worst:
\textbf{≈78× mean, ≈207× worst}.

\subsection{Two caveats, both cutting in the HLS arm\textquotesingle s
favour}\label{tm:dec:cpp-to-rtl-hls-evaluation:two-caveats-both-cutting-in-the-hls-arms-favour}

\textbf{1. The 4 142.52 ps is partly a flow artifact.} 2 550.65 ps of it
--- \textbf{62.3 \%} --- is a \emph{single} arc: a minimum-drive
\code{DFFASRHQNx1} driving \textbf{fanout 1188 / 630.7 fF}, unbuffered
because
\texttt{-\allowbreak{}-\allowbreak{}to\ OpenROAD.\allowbreak{}STAPrePNR}
stops before the resizer. The other 12 arcs total 1 546.4 ps. Buffered,
the path would plausibly land near 1.7--1.8 ns. The hand-RTL arm has no
comparable net --- its worst arc is 60.6 ps, 8.0 \% of its path, spread
over 24 well-balanced arcs --- so skipping the resizer penalises the HLS
arm \textbf{asymmetrically}, even though both ran the identical flow.
Read the timing ratio as \textbf{×2.2--5.5}, not ×5.5 flat. Area and
power are resizer-independent and regress on their own, so the verdict
is unaffected.

This also refines Stage 2\textquotesingle s caveat that "comparison
valid because treatment is identical" --- true procedurally, but the
\emph{effect} is not symmetric. Measured share of critical path in one
unbuffered high-fanout arc, all six arms:

{
\begin{longtable}[]{@{}lrr@{}}
\toprule\noalign{}
Arm & share & worst-arc fanout \\
\midrule\noalign{}
\endhead
\bottomrule\noalign{}
\endlastfoot
\code{coalescer_rtl} & 71.5 \% & 532 \\
\code{coalescer_hls} & 50.6 \% & 267 \\
\code{hazard_rtl} & 34.2 \% & 2 \\
\code{hazard_hls} & 31.6 \% & 15 \\
\code{valu_rtl} & 8.0 \% & 6 \\
\code{valu_hls} & \textbf{62.3 \%} & \textbf{1188} \\
\end{longtable}
}

No Stage 2 verdict changes direction: the coalescer\textquotesingle s
\emph{hand-RTL} arm suffered more than its HLS arm, so that HLS win is
if anything understated, and both hazard arms are artifact-free. But
absolute timing ratios across this pilot are bounds, not exact figures.

\textbf{2. Bambu\textquotesingle s own
\texttt{CYCLES\ value=\allowbreak{}"7"} is not the latency.} Parsed
per-vector from \code{results.txt}, the real distribution is 8 cycles
×22, 22 ×9, 38 ×3 --- mean \textbf{14.35} over 488 cycles / 34 vectors.
The \code{7} is roughly half the true mean; it must not be quoted.

\subsection{The FSM is structural, not an authoring or clock artifact}\label{tm:dec:cpp-to-rtl-hls-evaluation:the-fsm-is-structural-not-an-authoring-or-clock-artifact}

Two ablations, neither requested, both run because "the C was written
badly" and "the clock was too tight" are the first two objections to a
result like this:

\begin{itemize}
\tightlist
\item
  \textbf{Relaxed clock (1.75 ns instead of 0.705 ns):} 62 control
  steps, 6 012 FF, \textbf{identical area}. A slower target does not buy
  a flatter datapath.
\item
  \textbf{"HLS-friendly" rewrite} --- the lane loop restructured into a
  ternary-select chain, the shape usually recommended to coax parallel
  hardware out of HLS: \textbf{3.3× area}, still sequential.
\end{itemize}

Bambu schedules a loop over lanes into a multi-cycle FSM over a shared
functional unit. It has no mode that emits 8 parallel combinational
slices. Whether it can be made to emit a single-control-step datapath
\emph{at all} is a genuinely open question, left as a follow-up.

\subsection{Effect on the recommendation}\label{tm:dec:cpp-to-rtl-hls-evaluation:effect-on-the-recommendation-1}

The recommendation\textquotesingle s control-logic half \textbf{stands}
and is now supported by two independent blocks.

Its datapath half \textbf{does not}. It was never evidence-backed ---
Stage 1 and Stage 2 both said so --- and now that it has been measured
on a genuine datapath block, it is contradicted: HLS is \textbf{not}
favourable here, it is decisively worse on area, power, latency and
throughput simultaneously. The useful reformulation is narrower than
"reserve HLS for datapath accelerators":

\begin{quote}
On this toolchain, HLS is worth considering for
\textbf{throughput-oriented transaction/\allowbreak{}dataflow logic}
(the coalescer result), and not for \textbf{combinational arithmetic
datapaths} (this result) or for \textbf{cycle-critical control} (the
hazard unit). The distinction that predicted the outcome was not
datapath-vs-control; it was whether the block\textquotesingle s natural
form is already a cycle-by-cycle state machine --- HLS reproduces those,
and rebuilds everything else as one.
\end{quote}

\subsection{Verification}\label{tm:dec:cpp-to-rtl-hls-evaluation:verification}

Both arms pass \src{tb/cocotb/gpu/test_vector_alu.py} \textbf{14/14}.
The suite had no handshake handling at all, so the HLS arm scored 1/14
until a runtime protocol branch (the \code{test_hazard_unit.py} pattern)
was added; expected values and assertions are unchanged, only the
sampling point moved into \code{ReadOnly()}. This directed suite does
\textbf{not} exercise \code{ASSUMPTIONS.md} items 8 (\code{result_o} for
branch opcodes) or 10 (\code{VMOV_*}), so those two spec-drift
candidates remain \textbf{unprobed, not cleared}.

\begin{center}\rule{0.5\linewidth}{0.5pt}\end{center}

\subsection{\texorpdfstring{GH \#119 \code{egt} follow-up --- the
combinational-candidate PPA attempt is INVALID
(2026-09-15)}{GH \#119  follow-up --- the combinational-candidate PPA attempt is INVALID (2026-09-15)}}\label{tm:dec:cpp-to-rtl-hls-evaluation:gh-119-egt-follow-up-the-combinational-candidate-ppa-attempt-is-invalid-2026-09-15}

\textbf{Question asked:} bead \bead{egt} established Bambu 2024.10
\emph{can} emit a functionally-correct, zero-register, zero-latency
\code{vector_alu}
(\texttt{-\allowbreak{}-\allowbreak{}speculative-\allowbreak{}sdc-\allowbreak{}scheduling\ -\allowbreak{}-\allowbreak{}clock-\allowbreak{}period=\allowbreak{}100},
cosim + \code{test_vector_alu.py} both 14/14 against the pre-synthesis
\code{.v}). The remaining step was to re-measure the Stage-2/3 PPA
verdict against this "comb" candidate. \textbf{Conclusion: do not quote
comb PPA numbers.} A real synthesis bug was found and fixed; a second,
deeper synthesis-correctness bug was found and is \textbf{not} fixed.
\code{rtl} and \code{hls} reproduce the Stage 2 (\code{r8r}) baseline
cleanly and remain trustworthy.

\subsection{Setup}\label{tm:dec:cpp-to-rtl-hls-evaluation:setup}

Three ASAP7 \code{OpenROAD.STAPrePNR}-only runs
(\texttt{pnr/\allowbreak{}asap7/\allowbreak{}valu\_\allowbreak{}\{rtl,\allowbreak{}hls,\allowbreak{}hls\_\allowbreak{}comb\}/\allowbreak{}},
same 705 ps SDC, same corner, driven by
\code{/nobackup/egt_ppa/run_all.sh} under memory-capped
\texttt{systemd\ -\allowbreak{}-\allowbreak{}user} scopes): \code{rtl} =
hand-RTL, \code{hls} = default-flags Bambu baseline (registered, FSM),
\code{hls_comb} = the \code{sdc_longclk} candidate re-synthesised at
leaf-block level.

\subsection{\texorpdfstring{Bug 1 (found and fixed):
\code{SYNTH_SHARE_RESOURCES} left
enabled}{Bug 1 (found and fixed):  left enabled}}\label{tm:dec:cpp-to-rtl-hls-evaluation:bug-1-found-and-fixed-synth-share-resources-left-enabled}

\src{pnr/asap7/valu_hls_comb/config.json} (like \code{valu_rtl/} and
\code{valu_hls/}) never set \code{SYNTH_SHARE_RESOURCES}, so
yosys\textquotesingle s SAT-based \code{share} resource-sharing pass ran
unchecked. Its log (\code{05-yosys-synthesis/yosys-synthesis.log}, run
\code{RUN_2026-09-15_14-10-19}) shows it declaring six per-lane
\texttt{\$shl} cells \textbf{"never active. Sharing is pointless, we
simply remove it"} --- wrong for a fully-parallel, zero-register 8-lane
datapath, where \code{active_mask_i} is a runtime input, not a constant,
so no lane\textquotesingle s shifter is ever provably dead. This is the
\emph{exact} failure class already documented and fixed in
\src{pnr/asap7/gpu/config.json}
(\texttt{SYNTH\_\allowbreak{}SHARE\_\allowbreak{}RESOURCES:\allowbreak{}\ false},
for the same "8-lane vector\_\allowbreak{}alu" shape --- see the
\code{project_yosys_share_pass_hang} memory note) but never propagated
to the leaf-block configs. \textbf{Fix applied:} added
\texttt{"SYNTH\_\allowbreak{}SHARE\_\allowbreak{}RESOURCES":\allowbreak{}\ false}
to \code{valu_hls_comb/config.json} only
(\code{valu_rtl}/\code{valu_hls} left untouched --- their numbers
already match the \code{r8r} baseline, so \code{share} is not corrupting
them). Re-running confirmed the
\texttt{SHARE\ pass\ (SAT-\allowbreak{}based\ resource\ sharing)} step
no longer appears in the log at all.

\subsection{\texorpdfstring{Bug 2 (found, NOT fixed): the post-fix
netlist still discards most of
\code{rs2_i}}{Bug 2 (found, NOT fixed): the post-fix netlist still discards most of }}\label{tm:dec:cpp-to-rtl-hls-evaluation:bug-2-found-not-fixed-the-post-fix-netlist-still-discards-most-of-rs2-i}

Re-running after the Bug 1 fix (\code{RUN_2026-09-15_14-25-30}) produced
\textbf{byte-for-byte identical} metrics to the broken run (1587
instances / 138.7 µm² / 0.085 mW / WS +449 ps / TNS 0) --- proving
\code{share} was never the dominant cause. Direct inspection of the
final gate netlist (\code{05-yosys-synthesis/vector_alu_hls.nl.v})
shows:

\begin{itemize}
\tightlist
\item
  \textbf{\code{rs2_i} collapses to exactly 8 live bits total} (bit 0 of
  every 32-bit lane: bits \code{0,32,64,96,128,160,192,224}) out of the
  256-bit port --- bits \code{[31:1]} of every lane\textquotesingle s
  second ALU operand are read by \textbf{nothing}. \code{rs1_i} by
  contrast uses all 256 bits correctly.
\item
  \textbf{Only 22 XOR/\allowbreak{}XNOR-family cells exist in the whole
  1587-cell netlist.} Binary addition structurally requires one XOR per
  bit for the sum term (a full/half adder cannot be built from AOI-only
  logic); 22 XOR-family gates cannot implement even one correct 32-bit
  adder, let alone eight adders, eight subtractors and eight 32×32
  multipliers.
\item
  \textbf{\code{ALUMACC} (\code{yosys-synthesis.log} step 67) converts
  only 1 of the 8 \texttt{\$mul} cells} present earlier in the flow
  (confirmed independently: a standalone
  \texttt{yosys\ -\allowbreak{}sv\ .\allowbreak{}.\allowbreak{}.\allowbreak{}\ opt\ -\allowbreak{}full}
  pass on the bare candidate module reproduces \textbf{8} genuine 32×32
  \texttt{\$mul} cells, all fed directly from \code{rs1_N}/\code{rs2_N},
  at the RTLIL level) --- the other 7 were already eliminated before
  reaching \code{ALUMACC}, and even the 1 survivor is gone by the final
  \code{techmap} (no \texttt{\$mul}/\texttt{\$macc}/\texttt{\$alu}
  mapping template is invoked for a multiply anywhere in the log).
\item
  Bambu\textquotesingle s own \textbf{pre-synthesis area estimate for
  this candidate is 11 067 363 units} --- 8.5× the registered
  baseline\textquotesingle s own estimate (1 298 853, from
  \code{ASSUMPTIONS.md}\textquotesingle s "Measured QoR" section) ---
  i.e. Bambu itself expects the fully-unrolled combinational form to be
  \emph{larger}, not 30--50× smaller than either baseline
  arm\textquotesingle s final mapped netlist.
\end{itemize}

Working hypothesis (not verified to completion): three
\code{read_verilog} frontend warnings in the candidate\textquotesingle s
generic Bambu shifter/\allowbreak{}comparator FU templates
(\texttt{Range\ [0:\allowbreak{}-\allowbreak{}1]\ select\ out\ of\ bounds\ on\ signal\ 'in2':\allowbreak{}\ Setting\ 1\ LSB\ bits\ to\ undef},
at the \emph{unparameterized} module body\textquotesingle s default
\code{BITSIZE_in2}/\allowbreak{}computed-bitsize mismatch) inject
spurious X-constants that a register-free, fully-combinational netlist
lets cascade through the whole design during
\code{opt}/\code{OPT_MUXTREE} (no clock-edge boundary contains the
damage the way it does in the \code{rtl}/\code{hls} arms) --- but this
was not confirmed by formal equivalence or gate-level simulation, and no
further fix was attempted this session (deep surgery on a 900 KB
machine-generated netlist without full re-verification was judged
higher-risk than reporting the gap). \textbf{Tracked next step: bead
\code{claude_verilog_test-gcd}
(\texttt{discovered-\allowbreak{}from:\allowbreak{}\ egt}).}

\textbf{Do not quote \code{hls_comb}/\code{comb} area, power, WS, TNS,
or any derived fmax as if it were the true implementation cost of a
combinational \code{vector_alu} --- the netlist those numbers were
measured on does not correctly implement the second ALU operand for most
opcodes.}

\subsection{\texorpdfstring{Comparison table (\code{rtl} and \code{hls}
only --- \code{comb} excluded per
above)}{Comparison table ( and  only ---  excluded per above)}}\label{tm:dec:cpp-to-rtl-hls-evaluation:comparison-table-rtl-and-hls-only-comb-excluded-per-above}

All four figures reproduced from fresh ASAP7 \code{OpenROAD.STAPrePNR}
runs, 2026-09-15, same SDC/\allowbreak{}corner as the original
\code{r8r} Stage 2 runs; \code{r8r}\textquotesingle s recorded numbers
(\src{docs/CPP_TO_RTL_HLS_EVALUATION.md} (\cref{tm:dec:cpp-to-rtl-hls-evaluation}), "Stage 2 results" above) are
reproduced alongside for drift-checking.

{\footnotesize\setlength{\tabcolsep}{3pt}

{
\begin{longtable}[]{@{}
  >{\raggedright\arraybackslash}p{\dimexpr(\linewidth-8\tabcolsep)*86/1000\relax}
  >{\raggedleft\arraybackslash}p{\dimexpr(\linewidth-8\tabcolsep)*92/1000\relax}
  >{\raggedleft\arraybackslash}p{\dimexpr(\linewidth-8\tabcolsep)*135/1000\relax}
  >{\raggedleft\arraybackslash}p{\dimexpr(\linewidth-8\tabcolsep)*334/1000\relax}
  >{\raggedleft\arraybackslash}p{\dimexpr(\linewidth-8\tabcolsep)*334/1000\relax}@{}}
\toprule\noalign{}
\begin{minipage}[b]{\linewidth}\raggedright
Metric
\end{minipage} & \begin{minipage}[b]{\linewidth}\raggedleft
\code{rtl} (this run)
\end{minipage} & \begin{minipage}[b]{\linewidth}\raggedleft
\code{rtl} (\code{r8r}, 2026-09-07)
\end{minipage} & \begin{minipage}[b]{\linewidth}\raggedleft
\code{hls} (this run)\footnote{This \code{hls} column was measured on
  the netlist later proven \textbf{corrupt} by gate-level simulation
  (bead \bead{b0t} part B, 2026-09-18) and superseded 2026-09-19 when
  the remedy (\code{USE_SYNLIG:false} + sv2v) was applied to
  \src{pnr/asap7/valu_hls/config.json}. The corrected, GLS-verified
  figures are \textbf{89 224 instances / 9188.40 µm² / 52.7 mW / Setup
  WNS −1507.04 ps / Setup TNS −7 942 937.40 ps / Hold WNS −11.95 ps} ---
  see the "REMEDY APPLIED" block under "Baseline FSM arm" below. This
  column is left unaltered for provenance; treat it as a measurement of
  a broken netlist, not a valid \code{hls}-arm PPA point.}
\end{minipage} & \begin{minipage}[b]{\linewidth}\raggedleft
\code{hls} (\code{r8r}, 2026-09-07)\footnote{This \code{hls} column was
  measured on the netlist later proven \textbf{corrupt} by gate-level
  simulation (bead \bead{b0t} part B, 2026-09-18) and superseded
  2026-09-19 when the remedy (\code{USE_SYNLIG:false} + sv2v) was
  applied to \src{pnr/asap7/valu_hls/config.json}. The corrected,
  GLS-verified figures are \textbf{89 224 instances / 9188.40 µm² / 52.7
  mW / Setup WNS −1507.04 ps / Setup TNS −7 942 937.40 ps / Hold WNS
  −11.95 ps} --- see the "REMEDY APPLIED" block under "Baseline FSM arm"
  below. This column is left unaltered for provenance; treat it as a
  measurement of a broken netlist, not a valid \code{hls}-arm PPA point.}
\end{minipage} \\
\midrule\noalign{}
\endhead
\bottomrule\noalign{}
\endlastfoot
Instances & 54 634 & 54 634 & 71 480 & 70 895 \\
Area (µm²) & 4655.15 & 4655.15 & 7254.37 & 7234.0 \\
Power (mW) & 14.147 & 14.147 & 37.667 & 37.281 \\
Setup WS (ps) & −54.29 & (n/a --- see below) & −996.18 & (n/a) \\
Setup TNS (ps) & −2465.20 & --- & −3 779 683.10 & --- \\
Hold WS (ps) & +8.77 & --- & −13.01 & --- \\
\end{longtable}
}

}

\code{rtl} reproduces exactly (0.00\% drift on all three shared metrics
--- same \texttt{705\ ps} combinational critical path,
\texttt{759.\allowbreak{}29\ ps} data arrival, matching to 2 decimal
places). \code{hls} drifts \textbf{≤1.1\%} (instances +0.83\%, area
+0.28\%, power +1.04\%) --- attributable to
environment/\allowbreak{}tool-version movement since the 2026-09-07
\code{r8r} run, not a structural change; both are within normal
run-to-run noise for this flow. \warn{} Both \code{hls} columns were
later proven to be measurements of a \textbf{corrupt} netlist (bead
\bead{b0t}, see footnote); the fixed, GLS-verified figure is the "REMEDY
APPLIED" block under "Baseline FSM arm" further down.

\subsubsection{Worst-arc decomposition (mandatory before quoting any
fmax ratio, per the}\label{tm:dec:cpp-to-rtl-hls-evaluation:worst-arc-decomposition-mandatory-before-quoting-any-fmax-ratio-per-the}

\code{asap7-to-openroad-staprepnr-comparisons-stopping-before-the}
memory note ---
\texttt{-\allowbreak{}-\allowbreak{}to\ OpenROAD.\allowbreak{}STAPrePNR}
never runs the resizer, so a single unbuffered high-fanout net can
dominate a path and the effect is asymmetric between designs)

{
\begin{longtable}[]{@{}
  >{\raggedright\arraybackslash}p{\dimexpr(\linewidth-8\tabcolsep)*56/1000\relax}
  >{\raggedleft\arraybackslash}p{\dimexpr(\linewidth-8\tabcolsep)*392/1000\relax}
  >{\raggedleft\arraybackslash}p{\dimexpr(\linewidth-8\tabcolsep)*317/1000\relax}
  >{\raggedleft\arraybackslash}p{\dimexpr(\linewidth-8\tabcolsep)*103/1000\relax}
  >{\raggedleft\arraybackslash}p{\dimexpr(\linewidth-8\tabcolsep)*112/1000\relax}@{}}
\toprule\noalign{}
\begin{minipage}[b]{\linewidth}\raggedright
Arm
\end{minipage} & \begin{minipage}[b]{\linewidth}\raggedleft
Critical path total
\end{minipage} & \begin{minipage}[b]{\linewidth}\raggedleft
Single worst arc
\end{minipage} & \begin{minipage}[b]{\linewidth}\raggedleft
Arc share
\end{minipage} & \begin{minipage}[b]{\linewidth}\raggedleft
Arc fanout
\end{minipage} \\
\midrule\noalign{}
\endhead
\bottomrule\noalign{}
\endlastfoot
\code{rtl} & 759.29 ps (combinational, virtual clock) & 60.57 ps
(\code{OAI211}, stage 12/23) & \textbf{8.0\%} & 6 \\
\code{hls} & 1667.84 ps (register-to-register) & 981.02 ps (DFF
\code{QN} → \code{NOR3}) & \textbf{58.8\%} & 469 \\
\end{longtable}
}

\code{rtl}\textquotesingle s critical path is not dominated by any
single unbuffered net --- its area/\allowbreak{}power numbers \emph{and}
its WS are reasonably representative of the design as-is.
\code{hls}\textquotesingle s critical path is \textbf{dominated} by one
register\textquotesingle s 469-fanout, unbuffered fan-out net (58.8\% of
the total path) --- its WS/\allowbreak{}TNS/\allowbreak{}derived fmax
are a lower bound inflated by the absence of a resizer pass, not a clean
measurement of the FSM\textquotesingle s intrinsic speed; only its area
and power (resizer-independent) should be quoted with confidence. Do not
compute an \code{rtl}-vs-\code{hls} fmax ratio from these WS figures
without carrying this caveat.

\subsection{Verdict}\label{tm:dec:cpp-to-rtl-hls-evaluation:verdict}

\begin{itemize}
\tightlist
\item
  \textbf{Area/\allowbreak{}power, \code{rtl} vs \code{hls}:} confirms
  the original Stage 2 (\code{r8r}) result --- HLS\textquotesingle s
  registered, multi-cycle-FSM arm costs \textbf{+55.8\% area} and
  \textbf{+166.3\% power} versus hand-RTL, consistent with (slightly
  worse than, within noise of) the \code{r8r}-recorded
  +55.4\%/\allowbreak{}+163.5\%.
\item
  \textbf{The combinational candidate remains unmeasured.} Bead
  \code{egt} answered "can Bambu emit a zero-register datapath" with
  \textbf{yes} (functionally, pre-synthesis); this follow-up shows that
  answer does not yet extend through physical-design synthesis --- the
  ASAP7 flow\textquotesingle s default optimization settings do not
  reliably preserve that datapath\textquotesingle s function for this
  design shape. The datapath half of the pilot\textquotesingle s
  original recommendation is therefore still untested at the PPA level,
  exactly as Stage 1/2 left it.
\end{itemize}

\begin{center}\rule{0.5\linewidth}{0.5pt}\end{center}

\subsection{\texorpdfstring{GH \#119 \code{gcd} follow-up --- comb
candidate root-caused (Synlig) and fixed
(2026-09-15)}{GH \#119  follow-up --- comb candidate root-caused (Synlig) and fixed (2026-09-15)}}\label{tm:dec:cpp-to-rtl-hls-evaluation:gh-119-gcd-follow-up-comb-candidate-root-caused-synlig-and-fixed-2026-09-15}

\textbf{Root cause: the Synlig/\allowbreak{}UHDM SystemVerilog frontend
(\code{USE_SYNLIG:true}, LibreLane\textquotesingle s default), not
LibreLane\textquotesingle s opt passes.} Two other hypotheses were
tested and empirically DISPROVED first:

\begin{enumerate}
\def\labelenumi{\arabic{enumi}.}
\tightlist
\item
  \texttt{opt\_\allowbreak{}expr\ -\allowbreak{}undriven} (gated by
  \code{SYNTH_TIE_UNDEFINED}, default \texttt{"low"}): setting it
  \code{null} reproduced the byte-identical broken 1587-instance
  netlist. Re-reading \code{librelane_synth()}
  (\code{librelane/scripts/pyosys/synthesize.py}) shows the offending
  \code{OPT_MUXTREE} call is inside the \emph{very first}
  \texttt{librelane\_\allowbreak{}opt(d,\allowbreak{}\ nodffe=\allowbreak{}True,\allowbreak{}\ nosdff=\allowbreak{}True)}
  call, which never receives \code{undriven=} at all ---
  \code{SYNTH_TIE_UNDEFINED} cannot affect it regardless of config.
\item
  \code{SYNTH_SHARE_RESOURCES} (already documented as a real, separate
  hazard for this design shape --- see the note in
  \src{pnr/asap7/valu_hls_comb/config.json}): real, but re-running with
  only that fix applied reproduced byte-identical broken metrics,
  proving it was not the dominant cause.
\end{enumerate}

\textbf{Proof of divergence.} A SAT-based miter (yosys 0.46 --- the
exact binary LibreLane\textquotesingle s flow uses --- with real ASAP7
functional cell models parsed directly from the same liberty files
LibreLane loads, via
\texttt{read\_\allowbreak{}liberty\ -\allowbreak{}ignore\_\allowbreak{}miss\_\allowbreak{}func})
between the pristine candidate and the broken netlist was
unconstrained-SAT (a mismatch exists, but includes
don\textquotesingle t-care states under reset/\allowbreak{}idle).
\textbf{Constrained to the valid handshake window} (\code{rst_n=1},
\code{start_i=1} fixed) it is \textbf{still SAT}:

\begin{quote}
Counterexample: \code{opcode_i=VMUL(3)},
\texttt{active\_\allowbreak{}mask\_\allowbreak{}i=\allowbreak{}8'b00000001}
(lane 0 only active), \code{rs1_i}/\code{rs2_i} nonzero only in
\textbf{inactive} lanes 2 and 7. The pristine RTL gives
\code{result_o=0} (matches the already-passing \code{test_partial_mask}
spec: "inactive lanes produce 0"); the broken netlist gives
\texttt{result\_\allowbreak{}o[lane2]\ =\allowbreak{}\ rs1\_\allowbreak{}i[lane2]}
(\code{0x40000000}) leaked straight through, unmasked.
\code{done_o}/\code{branch_taken_o} match on both sides.
\end{quote}

\textbf{Localisation.} The corrupted construct is \code{ui_cond_expr_FU}
(a Bambu/\allowbreak{}PANDA generic IP-library module,
\texttt{assign\ out1\ =\allowbreak{}\ in1\ !=\allowbreak{}\ 0\ ?\ in2\ :\allowbreak{}\ in3;},
\code{vector_alu_synth.v:330-347}), instantiated per-call-site with its
own parameterization to implement "output 0 when
\code{active_mask_i[lane]=0}".
\code{05-yosys-synthesis.log}\textquotesingle s \code{OPT_MUXTREE} step
deletes 271 of its \texttt{\$ternary\$} cells\textquotesingle{} mux
ports in the broken run (only 10, legitimately dead, once fixed) ---
this is exactly where the masking branch gets pruned. Full writeup, the
exposure audit of the baseline FSM arm, and the CPU/GPU hand-RTL
exposure verdict are on bead \code{claude_verilog_test-b0t} (opened for
the general \code{USE_SYNLIG:true} audit) and bead
\code{claude_verilog_test-gcd}.

\textbf{Fix.} Bypass Synlig for this leaf: pre-process \code{gpu_pkg.sv}
+ the shim through \code{sv2v} (the project\textquotesingle s existing
Synlig fallback, already used for the SoC ASAP7 flow --- see
\code{sky130-soc-sv2v} in \src{pnr/Makefile}) into plain Verilog-2005,
point \code{VERILOG_FILES} at the sv2v output instead of the raw
\code{.sv} + \code{gpu_pkg.sv}, and set \code{USE_SYNLIG:false}. Made
reproducible via
\texttt{make\ -\allowbreak{}C\ pnr\ egt-\allowbreak{}valu-\allowbreak{}hls-\allowbreak{}comb-\allowbreak{}sv2v}
(auto-run by \code{librelane-asap7-valu-hls-comb}) --- see
\src{pnr/asap7/valu_hls_comb/config.json}\textquotesingle s
\code{USE_SYNLIG_NOTE}.

\textbf{Fix verification.} A full-width SAT re-proof (fixed netlist vs.
RTL candidate) did not converge within the 3G/\textasciitilde10min
budget --- 32×32 multiplier equivalence at full width is classically
SAT-hard (238K cells / 1.1M variables after \code{flatten}). A full
14-vector Verilator/\allowbreak{}cocotb regression against the fixed
netlist\textquotesingle s functional-liberty-modeled Verilog also could
not be completed: \code{flatten} exploded the 47,791-cell netlist to
\textasciitilde238K primitive gates before \code{write_verilog}, and the
resulting Verilator build did not finish in over 50 minutes wall-clock.
The fix was first verified (same session) by direct \code{eval} (yosys,
real ASAP7 functional cell models, no SAT) on \textbf{5 discriminating
vectors}: the original counterexample plus
VADD/\allowbreak{}VSLL/\allowbreak{}VAND with garbage data in inactive
lanes and VMUL with all 8 lanes active and distinct per-lane operands
--- all 5 matched the RTL candidate exactly, byte-for-byte, on
\code{done_o}/\code{result_o}/\code{branch_taken_o}.

\textbf{2026-09-15, follow-up (bead \bead{egt}, coordinator-directed):
verification widened to 252 vectors.} The 5-vector check above is
superseded, not retracted, by a much larger independent cross-check
against the \emph{golden model}, not the RTL candidate:
\src{tools/verif/gls/gen_vector_alu_vectors.py} generates 252 vectors
(random + corners) computed from the Python expected-value formulas
embedded directly in \src{tb/cocotb/gpu/test_vector_alu.py} --- every
non-branch opcode the suite exercises
(VADD/\allowbreak{}VSUB/\allowbreak{}VMUL/\allowbreak{}VAND/\allowbreak{}VOR/\allowbreak{}VXOR/\allowbreak{}VSLL/\allowbreak{}VSRL/\allowbreak{}VSRA),
VADDI/\allowbreak{}VANDI immediate sweeps, VBEQ/VBLT branch
corners+random, and a dedicated 12-mask-pattern × 8-op
partial-active-mask sweep with garbage forced into every inactive lane
--- the exact corruption class the original SAT-miter counterexample
exposed. \src{tools/verif/gls/make_eval_ys.py} batches one yosys
\code{eval} call per vector
(\texttt{read\_\allowbreak{}liberty\ -\allowbreak{}ignore\_\allowbreak{}miss\_\allowbreak{}func}
over the 5 ASAP7 \code{asap7sc7p5t_SIMPLE} libs, \code{read_verilog},
\code{hierarchy}, \code{proc}, \code{flatten},
\texttt{opt\ -\allowbreak{}full}, then 252×
\texttt{eval\ -\allowbreak{}set\ .\allowbreak{}.\allowbreak{}.\allowbreak{}\ -\allowbreak{}show\ done\_\allowbreak{}o\ -\allowbreak{}show\ result\_\allowbreak{}o\ -\allowbreak{}show\ branch\_\allowbreak{}taken\_\allowbreak{}o})
against the netlist actually referenced by the \textbf{tracked}
\src{pnr/asap7/valu_hls_comb/config.json} --- run
\code{RUN_2026-09-15_15-06-29}
(\code{05-yosys-synthesis/vector_alu_hls.nl.v}, \code{USE_SYNLIG:false}
+ sv2v, 47 791 instances; \emph{not} the later
\code{RUN_2026-09-15_16-29-18}, a leftover \code{USE_SYNLIG:true}
diagnostic re-run of the already-known-broken 1587-instance netlist).
All 252 evals completed in one batched yosys invocation
(\textasciitilde6 minutes wall clock, \textasciitilde1 GB peak, no SAT,
no flatten-to-primitives explosion since only
\texttt{opt\ -\allowbreak{}full} runs once and is amortized across all
252 evals). \src{tools/verif/gls/parse_eval_log.py} compared every
vector: \textbf{252/252 PASS --- 0 mismatches, 0 \code{X} in any
output}, on both \code{done_o} and the full
\code{result_o}/\code{branch_taken_o} vectors, including every
garbage-in-inactive-lane partial-mask case. Combined with the original
SAT-miter proof (which showed \emph{what} the pre-fix netlist got wrong)
and the 5-vector direct match to the RTL candidate, this closes the
fix-verification gap the original 5-vector check left open. \textbf{The
\code{comb} (fixed) row of the comparison table below is now quotable
without the earlier caveat} --- see the next section. The RTL-candidate
side of this specific 252-vector run could not be repeated
(yosys\textquotesingle s own, non-Synlig Verilog frontend does not
accept the port-adapting shim\textquotesingle s
\texttt{module\ vector\_\allowbreak{}alu\_\allowbreak{}hls\ import\ gpu\_\allowbreak{}pkg:\allowbreak{}:\allowbreak{}*;\ (.\allowbreak{}.\allowbreak{}.\allowbreak{})}
import-in-header SystemVerilog syntax); that side still rests on the
original 5-vector + 14-vector Verilator/\allowbreak{}cocotb evidence
above, which is unaffected.

\subsection{\texorpdfstring{Corrected comb metrics
(\code{RUN_2026-09-15_15-06-29}, \code{USE_SYNLIG:false} +
sv2v)}{Corrected comb metrics (,  + sv2v)}}\label{tm:dec:cpp-to-rtl-hls-evaluation:corrected-comb-metrics-run-2026-09-15-15-06-29-use-synligfalse-sv2v}

\textbf{These numbers are now quotable without caveat} (2026-09-15
update, bead \bead{egt}): the 252-vector yosys-\code{eval} cross-check
above confirms this exact run\textquotesingle s netlist is functionally
correct, so \code{comb} (fixed) = \textbf{47 791 instances / 4139.77 µm²
/ 14.060 mW}, vs. hand-RTL \code{rtl} = \textbf{54 634 / 4655.15 /
14.147}, is a valid area/\allowbreak{}power comparison --- \code{comb}
is smaller and lower-power than hand-RTL. The timing figures
(Setup/\allowbreak{}Hold WS, TNS, and any derived fmax) still carry the
worst-arc and no-resizer caveats below; only area and power are
resizer-independent and safe to quote unconditionally.

\begin{quote}
\warn{} \textbf{Update 2026-09-19 (bead \bead{b0t}):} the "\code{hls}
remains functionally unresolved" sentence below is superseded --- the
baseline FSM arm was proven \textbf{corrupt} by gate-level simulation
(bead \bead{b0t} part B) and the tracked config has since been fixed
(\code{USE_SYNLIG:false} + sv2v). The \code{hls} column in the table
immediately below is the \textbf{pre-fix, corrupt-netlist} measurement;
see the footnote on that column and the "Baseline FSM arm" §"REMEDY
APPLIED" block further down for the corrected, functionally-verified
figures. The baseline FSM arm (\code{hls}, \code{USE_SYNLIG:true})
remains functionally unresolved --- see "Baseline FSM arm ... ---
verdict: INCONCLUSIVE" further down and bead
\code{claude_verilog_test-b0t} --- so do not compare \code{comb} against
\code{hls} as if \code{hls}\textquotesingle s own netlist correctness
were settled; the \code{hls} row below is included for reference only,
carried over unchanged from the Stage-2 baseline.
\end{quote}

{
\begin{longtable}[]{@{}
  >{\raggedright\arraybackslash}p{\dimexpr(\linewidth-6\tabcolsep)*152/1000\relax}
  >{\raggedleft\arraybackslash}p{\dimexpr(\linewidth-6\tabcolsep)*143/1000\relax}
  >{\raggedleft\arraybackslash}p{\dimexpr(\linewidth-6\tabcolsep)*95/1000\relax}
  >{\raggedleft\arraybackslash}p{\dimexpr(\linewidth-6\tabcolsep)*590/1000\relax}@{}}
\toprule\noalign{}
\begin{minipage}[b]{\linewidth}\raggedright
Metric
\end{minipage} & \begin{minipage}[b]{\linewidth}\raggedleft
\code{comb} (fixed)
\end{minipage} & \begin{minipage}[b]{\linewidth}\raggedleft
\code{rtl}
\end{minipage} & \begin{minipage}[b]{\linewidth}\raggedleft
\code{hls}\footnote{This \code{hls} column was measured on the netlist
  later proven \textbf{corrupt} by gate-level simulation (bead
  \bead{b0t} part B, 2026-09-18) and superseded 2026-09-19 when the
  remedy (\code{USE_SYNLIG:false} + sv2v) was applied to
  \src{pnr/asap7/valu_hls/config.json}. The corrected, GLS-verified
  figures are \textbf{89 224 instances / 9188.40 µm² / 52.7 mW / Setup
  WNS −1507.04 ps / Setup TNS −7 942 937.40 ps / Hold WNS −11.95 ps} ---
  see the "REMEDY APPLIED" block under "Baseline FSM arm" below. This
  column is left unaltered for provenance; treat it as a measurement of
  a broken netlist, not a valid \code{hls}-arm PPA point.}
\end{minipage} \\
\midrule\noalign{}
\endhead
\bottomrule\noalign{}
\endlastfoot
Instances & 47 791 & 54 634 & 71 480 \\
Area (µm²) & 4139.77 & 4655.15 & 7254.37 \\
Power (mW) & 14.060 & 14.147 & 37.667 \\
Setup WS (ps) & −190.50 & −54.29 & −996.18 \\
Setup TNS (ps) & −14 657.70 & −2465.20 & −3 779 683.10 \\
Hold WS (ps) & −1.01 & +8.77 & −13.01 \\
\end{longtable}
}

\code{comb} is the \textbf{smallest and lowest-power} of the three arms
(−11.1\% area, −0.6\% power vs. \code{rtl}; −42.9\% area, −62.7\% power
vs. \code{hls}) --- consistent with a zero-register, zero-clock-tree
combinational datapath. Its setup WS is worse than
\code{rtl}\textquotesingle s (both are pure combinational virtual-clock
paths, so directly comparable): \code{comb}\textquotesingle s critical
path is a naive one-shot Bambu-generated 8-lane datapath with no hand
RTL retiming/\allowbreak{}balancing, \code{rtl}\textquotesingle s is
hand-optimized.

\subsubsection{\texorpdfstring{Worst-arc decomposition (\code{comb},
mandatory per the}{Worst-arc decomposition (, mandatory per the}}\label{tm:dec:cpp-to-rtl-hls-evaluation:worst-arc-decomposition-comb-mandatory-per-the}

\code{asap7-to-openroad-staprepnr-comparisons-stopping-before-the}
memory note)

{\small\setlength{\tabcolsep}{4pt}

{
\begin{longtable}[]{@{}
  >{\raggedright\arraybackslash}p{\dimexpr(\linewidth-8\tabcolsep)*118/1000\relax}
  >{\raggedleft\arraybackslash}p{\dimexpr(\linewidth-8\tabcolsep)*486/1000\relax}
  >{\raggedleft\arraybackslash}p{\dimexpr(\linewidth-8\tabcolsep)*196/1000\relax}
  >{\raggedleft\arraybackslash}p{\dimexpr(\linewidth-8\tabcolsep)*86/1000\relax}
  >{\raggedleft\arraybackslash}p{\dimexpr(\linewidth-8\tabcolsep)*94/1000\relax}@{}}
\toprule\noalign{}
\begin{minipage}[b]{\linewidth}\raggedright
Arm
\end{minipage} & \begin{minipage}[b]{\linewidth}\raggedleft
Critical path total
\end{minipage} & \begin{minipage}[b]{\linewidth}\raggedleft
Single worst arc
\end{minipage} & \begin{minipage}[b]{\linewidth}\raggedleft
Arc share
\end{minipage} & \begin{minipage}[b]{\linewidth}\raggedleft
Arc fanout
\end{minipage} \\
\midrule\noalign{}
\endhead
\bottomrule\noalign{}
\endlastfoot
\code{comb} (fixed) & 885.50 ps (combinational, virtual clock,
\code{rs2_i[226]}→\code{result_o[251]}) & 60.68 ps (\code{AOI32xp33}) &
\textbf{6.9\%} & 4 \\
\end{longtable}
}

}

Unlike \code{hls}\textquotesingle s FSM arm (58.8\% of its critical path
on one 469-fanout net --- see above), \code{comb}\textquotesingle s
worst single arc is a modest 6.9\% share at fanout 4: no single
unbuffered net dominates, so its WNS/TNS are a reasonably representative
(if unoptimized) measurement, not a resizer-absence artifact. Do not
compute a \code{comb}-vs-\code{rtl}/\code{hls} fmax ratio without
carrying the "no resizer at
\texttt{-\allowbreak{}-\allowbreak{}to\ OpenROAD.\allowbreak{}STAPrePNR}"
caveat regardless.

\subsection{\texorpdfstring{Baseline FSM arm (\code{valu_hls},
\code{USE_SYNLIG:true}) --- \warn{} RESOLVED 2026-09-18:
\textbf{CONFIRMED
CORRUPT}}{Baseline FSM arm (, ) ---  RESOLVED 2026-09-18: CONFIRMED CORRUPT}}\label{tm:dec:cpp-to-rtl-hls-evaluation:baseline-fsm-arm-valu-hls-use-synligtrue-warning-resolved-2026-09-18-confirmed-corrupt}

\begin{quote}
\warn{} \textbf{The "INCONCLUSIVE" verdict below is SUPERSEDED (bead
\bead{b0t}, 2026-09-18).} A gate-level simulation of both netlists
settled it: the tracked \textbf{\code{USE_SYNLIG:true} netlist is
CORRUPT} and the \code{USE_SYNLIG:false} + sv2v candidate is
\textbf{CLEAN}. Three independent golden vectors (VMUL partial-mask,
VADD full-mask, VAND alternating-mask) were driven through both 43-cycle
simulations: HLS-SYNLIG produced garbage on \textbf{all three} (no match
in either lane order), HLS-SV2V matched the golden model
\textbf{bit-exact on all three, all 8 lanes}. Same
Synlig/\allowbreak{}UHDM defect class already root-caused for the
\code{comb} arm --- it corrupts the FSM/\allowbreak{}registered arm too.

\textbf{What broke the three prior sessions\textquotesingle{} tooling
wall:} \code{flatten} \emph{before} yosys \texttt{sim\ -\allowbreak{}r}
(not \texttt{opt\ -\allowbreak{}full}), using the committed hand
SEQ-cell models \src{tools/verif/gls/asap7_seq_cell_models.v}. Both full
43-cycle simulations then complete in \textbf{under 2 minutes each at
\textasciitilde1 GB peak}, inside a 4 GB cap --- versus the earlier
\texttt{sat\ -\allowbreak{}seq} 3 GB OOM on a single unrolled timestep
and the 560 s+ wire-matching timeout on the unflattened netlist. Driver:
\src{tools/verif/gls/run_hls_seq_check.py}.

\textbf{Consequence for the numbers in this document:} every
\code{valu_hls} PPA figure quoted here (\textbf{71 480 instances /
7254.37 µm² / 37.667 mW}) was measured on the \textbf{corrupt} netlist
and is therefore not a valid PPA point for a working design. The honest
figure for a functionally correct \code{valu_hls} is the sv2v
candidate\textquotesingle s \textbf{89 224 instances / 9188.40 µm² /
52.7 mW}. Any \code{rtl}-vs-\code{hls} comparison in this document that
uses the 7254 µm² / 37.667 mW column \textbf{understates HLS area and
power by \textasciitilde25--40 \%}. The historical numbers are left
unaltered above for provenance; treat them as measurements of a broken
netlist.

\textbf{Remedy NOT applied, pending approval:} switching
\src{pnr/asap7/valu_hls/config.json} to \code{USE_SYNLIG:false} + the
sv2v shim (the same fix already applied to \code{valu_hls_comb}) would
make the block correct at the cost of materially worse headline PPA,
because the current sign-off numbers reflect the smaller, broken
netlist. \code{valu_hls} is an isolated HLS experiment block and is
\textbf{not} part of the signed-off SoC.

Evidence: netlists
\code{/nobackup/asap7_valu_hls_runs/RUN_2026-09-15_14-05-40/final/nl/vector_alu_hls.nl.v}
(corrupt) and
\code{/nobackup/pnr_ab_valu_hls/runs/RUN_2026-09-15_16-34-19/final/nl/vector_alu_hls.nl.v}
(clean); logs/VCDs under \code{/nobackup/b0t_partB/}. Full detail on
bead \code{claude_verilog_test-b0t}.
\end{quote}

\begin{quote}
\ok{} \textbf{REMEDY APPLIED 2026-09-19 (user-approved, bead
\bead{b0t}).} \src{pnr/asap7/valu_hls/config.json} is now
\code{USE_SYNLIG:false} + the sv2v-preprocessed shim (mirroring
\code{valu_hls_comb/config.json}\textquotesingle s
\code{USE_SYNLIG_NOTE}/fix exactly), reproducible via the new
\texttt{make\ -\allowbreak{}C\ pnr\ r8r-\allowbreak{}valu-\allowbreak{}hls-\allowbreak{}sv2v}
target (auto-run by \code{librelane-asap7-valu-hls}). The sv2v output
was regenerated and confirmed byte-for-byte identical (\code{sha256sum},
not a filtered \code{diff} --- see the RTK \code{diff}-exit-code caveat
elsewhere in this repo) to the file already used by the 2026-09-15 A/B
run.

\textbf{Reused, not re-run:} the tracked config now differs from the A/B
candidate run (\code{RUN_2026-09-15_16-34-19}) in \textbf{zero}
functional keys --- a programmatic key-by-key diff of both
\code{config.json}s (ignoring \code{//}-comment keys) found only the two
intentional changes (\code{VERILOG_FILES}, \code{USE_SYNLIG}) and no
others (same SDC, same \code{CLOCK_PERIOD}, same
\code{FP_CORE_UTIL}/density, same everything else). Per the
coordinator\textquotesingle s explicit reuse condition, that
run\textquotesingle s numbers are adopted directly rather than
re-running LibreLane:

{
\begin{longtable}[]{@{}lrr@{}}
\toprule\noalign{}
Metric & Corrupt (tracked pre-fix) & \textbf{Fixed (now tracked)} \\
\midrule\noalign{}
\endhead
\bottomrule\noalign{}
\endlastfoot
Instances & 71 480 & \textbf{89 224} \\
Area (µm²) & 7254.37 & \textbf{9188.40} \\
Power (mW) & 37.667 & \textbf{52.7} \\
Setup WNS (ps) & −996.18 & \textbf{−1507.04} \\
Setup TNS (ps) & −3 779 683.10 & \textbf{−7 942 937.40} \\
Hold WNS (ps) & −13.01 & \textbf{−11.95} \\
\end{longtable}
}

\textbf{Functional verification of the certified netlist:}
\src{tools/verif/gls/run_hls_seq_check.py} was re-run (fresh, this
session, not reusing the earlier session\textquotesingle s ad-hoc logs)
against \code{RUN_2026-09-15_16-34-19/final/nl/vector_alu_hls.nl.v} ---
the exact netlist whose PPA is now adopted above --- for all 3 golden
vectors used to originally confirm corruption:

{
\begin{longtable}[]{@{}
  >{\raggedright\arraybackslash}p{\dimexpr(\linewidth-4\tabcolsep)*91/1000\relax}
  >{\raggedright\arraybackslash}p{\dimexpr(\linewidth-4\tabcolsep)*467/1000\relax}
  >{\raggedright\arraybackslash}p{\dimexpr(\linewidth-4\tabcolsep)*422/1000\relax}@{}}
\toprule\noalign{}
\begin{minipage}[b]{\linewidth}\raggedright
Vector
\end{minipage} & \begin{minipage}[b]{\linewidth}\raggedright
Tag
\end{minipage} & \begin{minipage}[b]{\linewidth}\raggedright
Result
\end{minipage} \\
\midrule\noalign{}
\endhead
\bottomrule\noalign{}
\endlastfoot
109 & \code{VADD_mask01_0xff} (full mask) & \textbf{PASS}, bit-exact,
forward lane order \\
134 & \code{VMUL_mask02_0x01} (partial mask) & \textbf{PASS}, bit-exact,
forward lane order \\
148 & \code{VAND_mask04_0x55} (alternating mask) & \textbf{PASS},
bit-exact, forward lane order \\
\end{longtable}
}

3/3 PASS. Commit applying this remedy: \code{f2445ad} ("{[}PD{]} b0t:
apply approved Synlig fix to valu\_hls baseline FSM arm" ---
\src{pnr/Makefile} + \src{pnr/asap7/valu_hls/config.json}). The
historical "corrupt" figures above are left unaltered for provenance.
\end{quote}

\subsubsection{Original (superseded) verdict: INCONCLUSIVE, not
re-fixed}\label{tm:dec:cpp-to-rtl-hls-evaluation:original-superseded-verdict-inconclusive-not-re-fixed}

An A/B synthesis (temp config under \code{/nobackup/pnr_ab_valu_hls/},
tracked \src{pnr/asap7/valu_hls/config.json} left unmodified) reran the
SAME baseline candidate through \code{USE_SYNLIG:false} + an
sv2v-preprocessed shim: \textbf{89 224 instances / 9188.40 µm² / 52.7
mW}, vs. the tracked Synlig run\textquotesingle s \textbf{71 480 /
7254.37 µm² / 37.667 mW} --- a \textbf{\textasciitilde25--40\% gap},
present proportionally from the very first synthesis-log cell-count
checkpoint onward (43 659 vs. 55 084), not the \textasciitilde1--2\%
that would indicate no corruption. \textbf{This does not by itself prove
which side (if either) is correct.} A clocked functional cross-check of
the baseline (needed because it is a registered, multi-cycle FSM, unlike
the zero-register \code{comb} candidate) was attempted and blocked by
two independent tool/\allowbreak{}resource walls: yosys
0.46\textquotesingle s \texttt{sim\ -\allowbreak{}set} throws a syntax
error despite being documented in \texttt{help\ sim} (order-independent,
reproduced minimally); a full Verilator/\allowbreak{}cocotb cosim of
either baseline gate netlist OOMs at the mandated 3G cap even without
\code{flatten} (Verilator\textquotesingle s own front-end, not the
compile step, on a \textasciitilde71--89K-cell liberty-functional-model
netlist). \textbf{Do not apply the \code{comb} arm\textquotesingle s fix
to the tracked \code{valu_hls} config} without a functional proof of
which size is correct --- blind reuse could make it worse, not better.
Tracked on bead \code{claude_verilog_test-b0t} for follow-up
(candidates: get \texttt{sim\ -\allowbreak{}set} working, possibly via
the yosys 0.62 build already on this host, or budget a \textgreater3G
Verilator cosim).

\subsection{CPU/GPU hand-RTL exposure verdict: NOT exposed to this
specific construct}\label{tm:dec:cpp-to-rtl-hls-evaluation:cpugpu-hand-rtl-exposure-verdict-not-exposed-to-this-specific-construct}

\begin{quote}
\bad{} \textbf{UPDATE 2026-09-19 (bead \bead{u99} → \code{ma7}): the
"NOT exposed" verdict below is WRONG for the CPU macro.} A differential
--- the \emph{same} testbench, ROM and cycle sequencing run against the
Synlig gate netlist (\code{pnr/asap7/soc/macro/rv32i_cpu_top.nl.v.gz},
sha256-verified) and against RTL simulation --- showed the gate netlist
returns \textbf{0x30} for an APB debug read of \code{DBG_GPR[4]} where
RTL returns the correct \textbf{0x0c}. Since only the Synlig-built
netlist diverges, harness error, timing races and
debug-register-semantics explanations are all excluded. Mechanism is the
same class as \code{gcd}/\code{b0t} part B: \code{OPT_MUXTREE}
mis-elaborating a \textbf{runtime-indexed mux}.

The code-shape argument failed because it searched for
Bambu\textquotesingle s \emph{specific} \code{ui_cond_expr_FU} template.
The actual trigger is more general --- a plain runtime-indexed read.
\src{rtl/cpu/core/rv32i_regfile.sv} uses that idiom on \textbf{all
three} read ports:
\texttt{rd\_\allowbreak{}data1\ =\allowbreak{}\ regs[rd\_\allowbreak{}addr1]}
and
\texttt{rd\_\allowbreak{}data2\ =\allowbreak{}\ regs[rd\_\allowbreak{}addr2]}
(both feed ID/\allowbreak{}execute) and
\texttt{dbg\_\allowbreak{}rd\_\allowbreak{}data\ =\allowbreak{}\ regs[dbg\_\allowbreak{}rd\_\allowbreak{}addr]}
(proven corrupt). \src{rtl/gpu/warp_scheduler.sv:152}
(\code{div_depth[exec_warp_id_i]}) is the same class, untested.

\warn{} \textbf{Do not read this as "only the debug path is affected."}
The \code{u99} run verified instruction words and PC retire \emph{order}
over a 6-instruction straight-line program --- it did \textbf{not}
verify architectural register \emph{values}, and corrupted main read
ports would not necessarily perturb PC order on straight-line code. The
execute read ports are \textbf{UNVERIFIED, not proven clean}. Bounding
that exposure, and the remedy (switch the CPU/GPU configs to
\code{USE_SYNLIG:false} + sv2v, which forces a re-harden of both macros
and a SoC re-run, and will likely worsen PPA the way \code{valu_hls}
did) are tracked on bead \textbf{\code{ma7}} (P1). Until then, every
ASAP7 CPU/GPU macro PPA figure rests on a netlist from a frontend with a
\textbf{proven} miscompile in this repo.

\bad{} \textbf{UPDATE 2026-09-19 (bead \bead{ma7} step 1): CONFIRMED ---
the corruption reaches the execute datapath, not just the debug port.}
Method: extended the same \src{tools/verif/gls/tb_cpu_macro_check.v}
two-arm harness (arm 1 = Synlig gate netlist, arm 3 =
RTL/\allowbreak{}Verilator, same testbench file, same netlist ---
decompressed sha256 \code{a416d23f0380...}, exact match to
\code{u99}\textquotesingle s originally-verified copy) with a new
BRANCH-DEPENDENT program
(\src{tools/verif/gls/gen_cpu_check_rom_branch.py}, verified bit-exact
against \src{tb/models/rv32i_model.py} first):
\texttt{x1=\allowbreak{}5;\ x2=\allowbreak{}9;\ x3=\allowbreak{}9;\ x4=\allowbreak{}x1+\allowbreak{}x2;\ BEQ\ x2,\allowbreak{}x3\ -\allowbreak{}>\ PASS\ if\ taken\ (x2=\allowbreak{}=\allowbreak{}x3,\allowbreak{}\ must\ be\ TAKEN)\ else\ WRONG\ marker;\ .\allowbreak{}.\allowbreak{}.\allowbreak{}}.
This makes an ID-stage regfile read corruption
(\code{rd_data1}/\code{rd_data2}, which feed EX --- \textbf{not} the
known-bad debug port) observable through
\code{commit_pc_o}/\code{commit_insn_o}, the channel \code{u99} already
showed is reliable.

Result: \textbf{both arms retire an identical trace through the BEQ
instruction\textquotesingle s own fetch}
(\code{0x00,0x04,0x08,0x0c,0x10}, same insn words) and \textbf{diverge
exactly at that instruction\textquotesingle s commit}. Arm 3 (RTL)
continues \code{0x10,0x20,0x24,0x28} and parks at PASS --- BEQ correctly
TAKEN, matching the Python reference model bit-exact. Arm 1 (the same
Synlig netlist \code{u99} used) continues \code{0x10,0x14,0x18} and
parks at the dead WRONG marker --- BEQ \textbf{incorrectly evaluated
NOT-TAKEN}. 0 yosys ERRORs in the run. Since fetch/\allowbreak{}decode
are proven identical up to and including the branch instruction itself,
the divergence is isolated to the branch comparator\textquotesingle s
consumption of the ID-stage read-port values for x2/x3 --- a genuinely
different signal than \code{u99}\textquotesingle s \code{dbg_rd_data}.

\textbf{Exact coverage claim (not a general proof):} proven corrupt for
register indices x2 (value 9) and/or x3 (value 9) as read by the branch
comparator\textquotesingle s rs1/rs2 ports in this specific access
pattern; cannot distinguish which port (or both) is wrong. The preceding
\texttt{ADD\ x4,\allowbreak{}x1,\allowbreak{}x2} fetched identically in
both arms but x4 is a dead value in this program (never read back), so
that instruction\textquotesingle s own operand-read correctness is
\textbf{unproven} by this experiment. Not proven: other register
indices, other instruction types
(loads/\allowbreak{}stores/\allowbreak{}JALR), or
repeated/\allowbreak{}back-to-back-branch access patterns.

\textbf{Read-port cone-sharing (RTL source, not confirmed in netlist):}
\code{rv32i_regfile.sv} has three independent \code{always_comb} blocks
(\code{rd_data1}, \code{rd_data2}, \code{dbg_rd_data}), each its own
runtime-indexed mux keyed by an independent address signal (ID-stage
rs1/rs2 vs. APB-derived \code{dbg_rd_addr}) --- no shared select logic
at RTL, only the underlying storage flops are shared. The post-synthesis
netlist is fully flattened with anonymous net names (grep for
\code{rd_data1}/\code{rd_data2}/\code{dbg_rd_addr}/\code{regfile}
returns zero hits), so cone-sharing could not be confirmed by netlist
inspection without a hierarchy-preserving resynthesis (not attempted
here). This is consistent with independent per-port instances of the
same \code{OPT_MUXTREE} defect class, not one bug already fully
characterized by \code{u99}. \code{ma7} stays OPEN --- GPU untested,
remedy not applied, both still need a human decision. Full write-up:
bead \bead{ma7} notes.
\end{quote}

\begin{quote}
\bad{} \textbf{UPDATE 2026-09-19 (bead \bead{ma7} remedy, PARTIALLY
APPLIED --- CPU re-harden BLOCKED).} The \code{USE_SYNLIG:false} + sv2v
remedy (same fix already applied to \code{valu_hls}) was applied to
\src{pnr/asap7/cpu/config.json}/\code{config_3014.json}, but the CPU
macro\textquotesingle s own re-harden PD run has \textbf{not completed
cleanly}: it reaches \code{OpenROAD.DetailedRouting}
(synth/\allowbreak{}floorplan/\allowbreak{}CTS/\allowbreak{}GRT/ resizer
all pass) and then hits a new, unresolved routing-congestion regression
--- 108716 DRC violations vs. this block\textquotesingle s own best-ever
prior routed baseline of 2045, \code{DRT-0255} failures on SRAM macro
pins, zero parasitics extracted, \code{STAPostPNR} failure.
\texttt{pnr/\allowbreak{}asap7/\allowbreak{}soc/\allowbreak{}macro/\allowbreak{}\ rv32i\_\allowbreak{}cpu\_\allowbreak{}top.\allowbreak{}*}
therefore remain the OLD, Synlig-built, proven-corrupt views --- nothing
was regenerated. GPU config was switched (\code{USE_SYNLIG:false}) for
consistency only; the GPU re-harden itself remains fully deferred
(host-gated). Full attempt history, ruled-out causes, and the open
decision point: bead \bead{ma7} notes, bead
\code{claude_verilog_test-lxv},
\texttt{Doc:\allowbreak{}\ cr:\allowbreak{}run:\allowbreak{}asap7}\textquotesingle s
2026-09-19 addendum.
\end{quote}

\begin{quote}
\ok{} \textbf{UPDATE 2026-09-18 (bead \bead{b0t} part A)}: the
code-shape argument below has since been backed by a real frontend
\textbf{differential} --- replaying LibreLane\textquotesingle s exact
\texttt{librelane\_\allowbreak{}opt(nodffe=\allowbreak{}True,\allowbreak{}\ nosdff=\allowbreak{}True)}
sequence through two independent frontends on the same RTL and comparing
mux-port removal. \textbf{CPU macro (\code{rv32i_cpu_top}): clean} ---
the real Synlig run removes \emph{fewer} mux ports than an independent
slang differential (800 vs 1358), matching 1:1 by register name on
sampled branches: the opposite signature from the corruption proven
above. \textbf{GPU macro (\code{gpu_top}): clean on every point
investigated} --- \code{vector_alu} matches exactly (16/16) across
frontends, and \code{shared_memory}\textquotesingle s large raw
divergence (1001 vs 1) resolved to \textbf{0 disproven points} under a
direct EQY, with all 256 unproven points being SRAM-blackbox read-data
bits. Still unproven: no full gate-level check of either complete macro,
and three smaller GPU divergences (\code{gpu_compute_unit},
\code{memory_coalescer}, \code{gpu_top} glue) are differential-only.
Scripts: \src{tools/verif/synlig_audit/}.
\end{quote}

\texttt{grep\ -\allowbreak{}rl\ "ui\_\allowbreak{}cond\_\allowbreak{}expr\_\allowbreak{}FU|\_\allowbreak{}bambu\_\allowbreak{}artificial\_\allowbreak{}|BAMBU/\allowbreak{}PANDA"\ rtl/\allowbreak{}}
returns \textbf{zero matches} anywhere in \src{rtl/}. The hand-written
RTL (\src{rtl/gpu/vector_alu.sv}, \code{rtl/cpu/*}) implements the
identical "inactive lane → 0" masking semantics via a plain
\texttt{if\ (active\_\allowbreak{}mask\_\allowbreak{}i[lane])} inside a
generate/\allowbreak{}for-loop --- a completely different code shape
from Bambu\textquotesingle s massively per-instance-parameterized
generic FU template instantiation pattern that Synlig mis-elaborates
here. The \emph{specific} failure mode proven in this section does not
structurally apply to the CPU/GPU macros. This is a
structural/\allowbreak{}code-shape argument, not a functional
re-verification of the CPU/GPU macros --- per explicit instruction this
session did not re-run those flows, so their \code{USE_SYNLIG:true}
configs remain functionally unaudited (tracked on bead
\code{claude_verilog_test-b0t}).
\end{relocated}

\chapter{Analog PLL}\label{tm:ch:analog-pll}

This chapter describes the analog PLL flow. It covers both process kits.
